% Generated mechanically from frozen input inputs/p02/ja/tex/derived.tex.
% Do not edit this generated file; edit the bound locale source instead.

% OK, start here.
%

\setcounter{chapter}{12}
\chapter{導来圏}



\phantomsection
\label{derived-section-phantom}


\section{序論}
\label{derived-section-introduction}

\noindent
まず、三角圏および三角圏における局所化について論じる。次に、加法圏の
複体のホモトピー圏が三角圏であることを証明する。その上で、アーベル圏の
導来圏を、擬同型に関するホモトピー圏の局所化として定義する。
よい参考文献として Verdier の学位論文 \cite{Verdier} がある。



\section{三角圏}
\label{derived-section-triangulated-categories}

\noindent
三角圏は、アーベル圏の導来圏に内在する構造を記述するための便利な道具である。
参考文献として \cite{Verdier}、\cite{KS}、\cite{Neeman} がある。




\section{三角圏の定義}
\label{derived-section-triangulated-definitions}

\noindent
本節では、三角圏および前三角圏に関する定義の大部分をまとめる。

\begin{definition}
\label{derived-definition-triangle}
$\mathcal{D}$ を加法圏とする。また
$[1] : \mathcal{D} \to \mathcal{D}$、$E \mapsto E[1]$
を $\mathcal{D}$ の加法的自己同値関手とする。
\begin{enumerate}
\item {\it 三角}とは、$X, Y, Z \in \Ob(\mathcal{D})$ であり、
$f : X \to Y$、$g : Y \to Z$、$h : Z \to X[1]$ が
$\mathcal{D}$ の射であるような六つ組
$(X, Y, Z, f, g, h)$ のことである。
\item {\it 三角の射}
$(X, Y, Z, f, g, h) \to (X', Y', Z', f', g', h')$
とは、$\mathcal{D}$ の射 $a : X \to X'$、$b : Y \to Y'$、
$c : Z \to Z'$ であって、
$b \circ f = f' \circ a$、$c  \circ g = g' \circ b$、
$a[1] \circ h = h' \circ c$ を満たすもののことである。
\end{enumerate}
\end{definition}

\noindent
三角の射は、次の可換図式によって表される。
$$
\xymatrix{
X \ar[r] \ar[d]^a &
Y \ar[r] \ar[d]^b &
Z \ar[r] \ar[d]^c &
X[1] \ar[d]^{a[1]} \\
X' \ar[r] &
Y' \ar[r] &
Z' \ar[r] &
X'[1]
}
$$
定義 \ref{derived-definition-triangle} の設定において、$[0] = \text{id}$ と書く。
$n > 0$ に対しては $[1]$ の $n$ 回の合成を $[n]$ と書き、$[1]$ の擬逆
$[-1]$ を一つ選び、$[-1]$ の $n$ 回の合成を $[-n]$ と定める。このとき
$\{[n]\}_{n \in \mathbf{Z}}$ は $n \in \mathbf{Z}$ を添字とする
$\mathcal{D}$ の加法的自己同値関手の族であり、関手の同型
$[n] \circ [m] \cong [n + m]$ が与えられている。

\medskip\noindent
以下は Verdier の学位論文で与えられた三角圏の定義である。

\begin{definition}
\label{derived-definition-triangulated-category}
{\it 三角圏}とは、三つ組
$(\mathcal{D}, \{[n]\}_{n\in \mathbf{Z}}, \mathcal{T})$
であって、次を満たすものをいう。
\begin{enumerate}
\item $\mathcal{D}$ は加法圏である。
\item $[1] : \mathcal{D} \to \mathcal{D}$、$E \mapsto E[1]$
は加法的自己同値関手であり、$n \in \mathbf{Z}$ に対する $[n]$ は
上で述べたとおりである。
\item $\mathcal{T}$ は三角（定義 \ref{derived-definition-triangle}）の集合であり、
その元を {\it 完全三角} と呼ぶ。
\end{enumerate}
さらに、次の条件を課す。
\begin{enumerate}
\item[TR1] 完全三角と同型な三角はすべて完全三角である。
$(X, X, 0, \text{id}, 0, 0)$ の形の三角はすべて完全三角である。
$\mathcal{D}$ の任意の射 $f : X \to Y$ に対して、
$(X, Y, Z, f, g, h)$ の形の完全三角が存在する。
\item[TR2] 三角 $(X, Y, Z, f, g, h)$ が完全三角であるための必要十分条件は、
三角 $(Y, Z, X[1], g, h, -f[1])$ が完全三角であることである。
\item[TR3] 実線で描かれた図式
$$
\xymatrix{
X \ar[r]^f \ar[d]^a &
Y \ar[r]^g \ar[d]^b &
Z \ar[r]^h \ar@{-->}[d] &
X[1] \ar[d]^{a[1]} \\
X' \ar[r]^{f'} &
Y' \ar[r]^{g'} &
Z' \ar[r]^{h'} &
X'[1]
}
$$
が与えられ、その各行が完全三角であり、$b \circ f = f' \circ a$ を
満たすとする。このとき、$(a, b, c)$ が三角の射となるような射
$c : Z \to Z'$ が存在する。
\item[TR4] $\mathcal{D}$ の対象 $X$、$Y$、$Z$、射
$f : X \to Y$、$g : Y \to Z$、および完全三角
$(X, Y, Q_1, f, p_1, d_1)$、
$(X, Z, Q_2, g \circ f, p_2, d_2)$、
$(Y, Z, Q_3, g, p_3, d_3)$
が与えられたとする。このとき、次を満たす射
$a : Q_1 \to Q_2$ および $b : Q_2 \to Q_3$ が存在する。
\begin{enumerate}
\item $(Q_1, Q_2, Q_3, a, b, p_1[1] \circ d_3)$ は完全三角である。
\item 三つ組 $(\text{id}_X, g, a)$ は三角の射
$(X, Y, Q_1, f, p_1, d_1) \to (X, Z, Q_2, g \circ f, p_2, d_2)$
である。
\item 三つ組 $(f, \text{id}_Z, b)$ は三角の射
$(X, Z, Q_2, g \circ f, p_2, d_2) \to (Y, Z, Q_3, g, p_3, d_3)$
である。
\end{enumerate}
\end{enumerate}
TR1、TR2、TR3 が成り立つとき、$(\mathcal{D}, [\ ], \mathcal{T})$ を
{\it 前三角圏} と呼ぶ。
\footnote{$[\ ]$ は族 $\{[n]\}_{n\in \mathbf{Z}}$ の略記として用いる。}
\end{definition}

\noindent
TR4 は次のように説明できる。$Q_1$ を $Y/X$、$Q_2$ を $Z/X$、
$Q_3$ を $Z/Y$ と考えれば、TR4(a) は同型
$(Z/X)/(Y/X) \cong Z/Y$ を表し、TR4(b) と TR4(c) は、三角の射によって
$X \to Y \to Q_1 \to X[1]$ などの三角を比較できることを表している。
この考えをより精密に述べ直したものについては、
補題 \ref{derived-lemma-two-split-injections} の証明を参照されたい。

\medskip\noindent
TR2 の符号は、$(X, Y, Z, f, g, h)$ が完全三角ならば、長い列
\begin{equation}
\label{derived-equation-rotate}
\ldots \to
Z[-1] \xrightarrow{-h[-1]}
X \xrightarrow{f}
Y \xrightarrow{g}
Z \xrightarrow{h}
X[1] \xrightarrow{-f[1]}
Y[1] \xrightarrow{-g[1]}
Z[1] \to \ldots
\end{equation}
の各連続する四項が完全三角を与えることを意味する。

\medskip\noindent
通常どおり記号を濫用し、付加的なデータの記号を明示的に導入せず、単に
（前）三角圏 $\mathcal{D}$ と呼ぶことにする。前三角圏という概念は、TR4 と
同値な命題を見いだす際に有用である。

\medskip\noindent
三角関手を次のように定義する。

\begin{definition}
\label{derived-definition-exact-functor-triangulated-categories}
$\mathcal{D}$、$\mathcal{D}'$ を前三角圏とする。$\mathcal{D}$ から
$\mathcal{D}'$ への {\it 完全関手}、または {\it 三角関手} とは、関手
$F : \mathcal{D} \to \mathcal{D}'$ と、関手的同型
$\xi_X : F(X[1]) \to F(X)[1]$ の組であって、$\mathcal{D}$ の任意の完全三角
$(X, Y, Z, f, g, h)$ に対し、三角
$(F(X), F(Y), F(Z), F(f), F(g), \xi_X \circ F(h))$
が $\mathcal{D}'$ の完全三角となるものをいう。
\end{definition}

\noindent
完全関手は加法的である。補題 \ref{derived-lemma-exact-functor-additive} を参照されたい。
二つの三角圏が同値であるというとき、それらが三角圏の $2$-圏において
同値であることを意味する。この $2$-圏における $2$-射
$a : (F, \xi) \to (F', \xi')$ とは、単に $\xi$ と $\xi'$ に整合する
関手の自然変換 $a : F \to F'$、すなわち図式
$$
\xymatrix{
F \circ [1] \ar[r]_\xi \ar[d]_{a \star 1} & [1] \circ F \ar[d]^{1 \star a} \\
F' \circ [1] \ar[r]^{\xi'} & [1] \circ F'
}
$$
が可換となるようなもののことである。

\begin{definition}
\label{derived-definition-triangulated-subcategory}
$(\mathcal{D}, [\ ], \mathcal{T})$ を前三角圏とする。
{\it 前三角部分圏}\footnote{この定義は標準的でない可能性がある。
$\mathcal{D}'$ が充満部分圏ならば、$\mathcal{T}'$ は
$\mathcal{D}'$ における三角の集合と $\mathcal{T}$ との共通部分である。
補題 \ref{derived-lemma-triangulated-subcategory} を参照されたい。この場合、記号から
$\mathcal{T}'$ を省略する。}
とは、次を満たす組 $(\mathcal{D}', \mathcal{T}')$ のことである。
\begin{enumerate}
\item $\mathcal{D}'$ は $\mathcal{D}$ の加法部分圏であり、$[1]$ によって
保たれ、かつ $[1] : \mathcal{D}' \to \mathcal{D}'$ は自己同値である。
\item $\mathcal{T}' \subset \mathcal{T}$ は、任意の
$(X, Y, Z, f, g, h) \in \mathcal{T}'$ に対して
$X, Y, Z \in \Ob(\mathcal{D}')$ および
$f, g, h \in \text{Arrows}(\mathcal{D}')$ が成り立つような部分集合である。
\item $(\mathcal{D}', [\ ], \mathcal{T}')$ は前三角圏である。
\end{enumerate}
$\mathcal{D}$ が三角圏であるとする。このとき、
$(\mathcal{D}', \mathcal{T}')$ が前三角部分圏であり、かつ
$(\mathcal{D}', [\ ], \mathcal{T}')$ が三角圏であるならば、
この組を {\it 三角部分圏} と呼ぶ。
\end{definition}

\noindent
この状況では、包含関手 $\mathcal{D}' \to \mathcal{D}$ は完全関手であり、
$\xi_X : X[1] \to X[1]$ は $X[1]$ 上の恒等射によって与えられる。

\medskip\noindent
補題 \ref{derived-lemma-composition-zero} において、前三角圏の完全三角
$(X, Y, Z, f, g, h)$ に対して合成 $g \circ f : X \to Z$ が零であることを見る。
したがって、列 (\ref{derived-equation-rotate}) は複体である。ホモロジー的関手とは、
この複体を長完全列に移す関手である。

\begin{definition}
\label{derived-definition-homological}
$\mathcal{D}$ を前三角圏、$\mathcal{A}$ をアーベル圏とする。
加法的関手 $H : \mathcal{D} \to \mathcal{A}$ が
{\it ホモロジー的} であるとは、任意の完全三角
$(X, Y, Z, f, g, h)$ に対して列
$$
H(X) \to H(Y) \to H(Z)
$$
がアーベル圏 $\mathcal{A}$ において完全であることをいう。加法的関手
$H : \mathcal{D}^{opp} \to \mathcal{A}$ が {\it コホモロジー的} であるとは、
対応する関手 $\mathcal{D} \to \mathcal{A}^{opp}$ がホモロジー的であることをいう。
\end{definition}

\noindent
$H : \mathcal{D} \to \mathcal{A}$ がホモロジー的関手であるとき、しばしば
$H^n(X) = H(X[n])$ と書き、$H(X) = H^0(X)$ とする。上での TR2 の議論から、
完全三角 $(X, Y, Z, f, g, h)$ は長完全列
\begin{equation}
\label{derived-equation-long-exact-cohomology-sequence}
\xymatrix@C=3pc{
H^{-1}(Z) \ar[r]^{H(h[-1])} &
H^0(X) \ar[r]^{H(f)} &
H^0(Y) \ar[r]^{H(g)} &
H^0(Z) \ar[r]^{H(h)} &
H^1(X)
}
\end{equation}
を定める。これを、その完全三角とホモロジー的関手に付随する
{\it 長完全列} と呼ぶ。表示したとおり、長完全列の射には符号を付けない。
そのため、完全三角を回転して得られる三角（TR2）に付随する長完全列の写像は、
上の列の写像といくつかの符号だけ異なることになる。

\begin{definition}
\label{derived-definition-delta-functor}
$\mathcal{A}$ をアーベル圏、$\mathcal{D}$ を三角圏とする。
{\it $\mathcal{A}$ から $\mathcal{D}$ への $\delta$-関手}とは、関手
$G : \mathcal{A} \to \mathcal{D}$ と、任意の短完全列
$$
0 \to A \xrightarrow{a} B \xrightarrow{b} C \to 0
$$
に射 $\delta = \delta_{A \to B \to C} : G(C) \to G(A)[1]$ を対応させる
規則であって、次を満たすものをいう。
\begin{enumerate}
\item 上のような任意の短完全列に対して、三角
$(G(A), G(B), G(C), G(a), G(b), \delta_{A \to B \to C})$
は $\mathcal{D}$ の完全三角である。
\item 短完全列の任意の射
$(A \to B \to C) \to (A' \to B' \to C')$
に対して、図式
$$
\xymatrix{
G(C) \ar[d] \ar[rr]_{\delta_{A \to B \to C}} & &
G(A)[1] \ar[d] \\
G(C') \ar[rr]^{\delta_{A' \to B' \to C'}} & &
G(A')[1]
}
$$
は可換である。
\end{enumerate}
この状況で、
$(G(A), G(B), G(C), G(a), G(b), \delta_{A \to B \to C})$
を、与えられた $\delta$-関手による短完全列の {\it 像} と呼ぶ。
\end{definition}

\noindent
$\delta$-関手には付加的な構造が備わっていることに注意する。厳密にいえば、
与えられた関手 $\mathcal{A} \to \mathcal{D}$ が $\delta$-関手であるというのは
意味をなさないが、それでもしばしばそのように述べる。


\section{三角圏に関する初等的結果}
\label{derived-section-elementary-results}

\noindent
本節の結果の大部分は前三角圏に対して証明されるので、特に任意の三角圏においても
成り立つ。

\begin{lemma}
\label{derived-lemma-composition-zero}
$\mathcal{D}$ を前三角圏とする。
$(X, Y, Z, f, g, h)$ を完全三角とする。このとき
$g \circ f = 0$、$h \circ g = 0$、$f[1] \circ h = 0$ である。
\end{lemma}

\begin{proof}
TR1 により、$(X, X, 0, 1, 0, 0)$ は完全三角である。図式
$$
\xymatrix{
X \ar[r] \ar[d]^1 &
X \ar[r] \ar[d]^f &
0 \ar[r] \ar@{-->}[d] &
X[1] \ar[d]^{1[1]} \\
X \ar[r]^f &
Y \ar[r]^g &
Z \ar[r]^h &
X[1]
}
$$
に TR3 を適用する。もちろん点線の矢印は零写像である。したがって図式の可換性から
$g \circ f = 0$ が従う。他の場合については、三角を回転する、すなわち TR2 を
適用すればよい。
\end{proof}

\begin{lemma}
\label{derived-lemma-representable-homological}
$\mathcal{D}$ を前三角圏とする。$\mathcal{D}$ の任意の対象 $W$ に対して、関手
$\Hom_\mathcal{D}(W, -)$ はホモロジー的であり、関手
$\Hom_\mathcal{D}(-, W)$ はコホモロジー的である。
\end{lemma}

\begin{proof}
完全三角 $(X, Y, Z, f, g, h)$ を考える。すでに
$g \circ f = 0$ を見た。補題 \ref{derived-lemma-composition-zero} を参照されたい。
$g \circ a = 0$ を満たす射 $a : W \to Y$ が与えられたとする。このとき可換図式
$$
\xymatrix{
W \ar[r]_1 \ar@{..>}[d]^b &
W \ar[r] \ar[d]^a &
0 \ar[r] \ar[d]^0 &
W[1] \ar@{..>}[d]^{b[1]} \\
X \ar[r] & Y \ar[r] & Z \ar[r] & X[1]
}
$$
を得る。両方の行は完全三角である（上の行には TR1 を用いる）。したがって、
点線の矢印 $b$ を補うことができる（まず TR2 により回転し、次に TR3 を適用して、
最後に元へ回転する）。これで補題が証明された。
\end{proof}

\begin{lemma}
\label{derived-lemma-third-isomorphism-triangle}
$\mathcal{D}$ を前三角圏とする。
$$
(a, b, c) : (X, Y, Z, f, g, h) \to (X', Y', Z', f', g', h')
$$
を完全三角の射とする。$a, b, c$ のうち二つが同型ならば、残る一つも同型である。
\end{lemma}

\begin{proof}
$a$ と $c$ が同型であると仮定する。$\mathcal{D}$ の任意の対象 $W$ に対して
$H_W( - ) = \Hom_\mathcal{D}(W, -)$
と書く。このとき、アーベル群の可換図式
$$
\xymatrix{
H_W(Z[-1]) \ar[r] \ar[d] &
H_W(X) \ar[r] \ar[d] &
H_W(Y) \ar[r] \ar[d] &
H_W(Z) \ar[r] \ar[d] &
H_W(X[1]) \ar[d] \\
H_W(Z'[-1]) \ar[r] &
H_W(X') \ar[r] &
H_W(Y') \ar[r] &
H_W(Z') \ar[r] &
H_W(X'[1])
}
$$
を得る。仮定により、右側の二本および左側の二本の縦の矢印は全単射である。
補題 \ref{derived-lemma-representable-homological} により $H_W$ はホモロジー的であるから、
五項補題（Homology, Lemma \ref{homology-lemma-five-lemma}）によって中央の縦の矢印は
同型である。したがって、米田の補題
（Categories, Lemma \ref{categories-lemma-yoneda}）により、$b$ が同型であることが
分かる。他の場合は、TR2 を用いて回転することにより従う。
\end{proof}

\begin{remark}
\label{derived-remark-special-triangles}
$\mathcal{D}$ を加法圏とし、定義 \ref{derived-definition-triangle} のような
シフト関手 $[n]$ を備えているとする。三角
$(X, Y, Z, f, g, h)$ が {\it 特殊}\footnote{これは標準的な用語ではない。}
であるとは、$\mathcal{D}$ の任意の対象 $W$ に対してアーベル群の長い列
$$
\ldots \to
\Hom_\mathcal{D}(W, X) \to
\Hom_\mathcal{D}(W, Y) \to
\Hom_\mathcal{D}(W, Z) \to
\Hom_\mathcal{D}(W, X[1]) \to \ldots
$$
が完全であることとする。補題 \ref{derived-lemma-third-isomorphism-triangle} の証明から、
$$
(a, b, c) : (X, Y, Z, f, g, h) \to (X', Y', Z', f', g', h')
$$
が特殊三角の射であり、$a, b, c$ のうち二つが同型ならば、残る一つも同型である。
{\it 余特殊} 三角、すなわち関手 $\Hom_\mathcal{D}(-, W)$ を適用すると
長完全列になる三角については、双対な主張が成り立つ。したがって完全三角は
特殊かつ余特殊であるが、一般には（余）特殊三角は完全三角よりはるかに多い。
\end{remark}
\begin{lemma}
\label{derived-lemma-third-map-square-zero}
$\mathcal{D}$ を前三角圏とする。
$$
(0, b, 0), (0, b', 0) : (X, Y, Z, f, g, h) \to (X, Y, Z, f, g, h)
$$
を、ある完全三角の自己射とする。このとき $bb' = 0$ である。
\end{lemma}

\begin{proof}
図式
$$
\xymatrix{
X \ar[r] \ar[d]^0 &
Y \ar[r] \ar[d]^{b, b'} \ar@{..>}[ld]^\alpha &
Z \ar[r] \ar[d]^0 \ar@{..>}[ld]^\beta &
X[1] \ar[d]^0 \\
X \ar[r] & Y \ar[r] & Z \ar[r] & X[1]
}
$$
を考える。補題 \ref{derived-lemma-representable-homological} を適用すると、
$b' = f \circ \alpha$ および $b = \beta \circ g$ を満たす点線の矢印
$\alpha$ と $\beta$ が得られる。ゆえに
$bb' = \beta \circ g \circ f \circ \alpha = 0$ である。実際、
補題 \ref{derived-lemma-composition-zero} により $g \circ f = 0$ である。
\end{proof}

\begin{lemma}
\label{derived-lemma-third-map-idempotent}
$\mathcal{D}$ を前三角圏とし、$(X, Y, Z, f, g, h)$ を完全三角とする。
図式
$$
\xymatrix{
Z \ar[r]_h \ar[d]_c & X[1] \ar[d]^{a[1]} \\
Z \ar[r]^h & X[1]
}
$$
が可換で、$a^2 = a$ および $c^2 = c$ が成り立つとする。このとき、
$b^2 = b$ を満たす射 $b : Y \to Y$ であって、$(a, b, c)$ が三角
$(X, Y, Z, f, g, h)$ の自己射となるものが存在する。
\end{lemma}

\begin{proof}
TR3 により、$(a, b', c)$ が $(X, Y, Z, f, g, h)$ の自己射となるような射
$b'$ が存在する。このとき $(0, (b')^2 - b', 0)$ も自己射である。
補題 \ref{derived-lemma-third-map-square-zero} により、$(b')^2 - b'$ の二乗は零である。
$b = b' - (2b' - 1)((b')^2 - b') = 3(b')^2 - 2(b')^3$ と置く。計算により
$(a, b, c)$ が自己射であり、
$b^2 - b = (4(b')^2 - 4b' - 3)((b')^2 - b')^2 = 0$ であることが分かる。
\end{proof}

\begin{lemma}
\label{derived-lemma-cone-triangle-unique-isomorphism}
$\mathcal{D}$ を前三角圏とし、$f : X \to Y$ を $\mathcal{D}$ の射とする。
完全三角 $(X, Y, Z, f, g, h)$ が存在し、これは三角の（一意とは限らない）
同型を除いて一意である。より正確には、もう一つのそのような完全三角
$(X, Y, Z', f, g', h')$ が与えられると、同型
$$
(1, 1, c) : (X, Y, Z, f, g, h) \longrightarrow (X, Y, Z', f, g', h')
$$
が存在する。
\end{lemma}

\begin{proof}
存在は TR1 による。同型を除く一意性は TR3 と
補題 \ref{derived-lemma-third-isomorphism-triangle} による。
\end{proof}

\begin{lemma}
\label{derived-lemma-uniqueness-third-arrow}
$\mathcal{D}$ を前三角圏とする。
$$
(a, b, c) : (X, Y, Z, f, g, h) \to (X', Y', Z', f', g', h')
$$
を完全三角の射とする。次の条件のいずれかが成り立つとする。
\begin{enumerate}
\item $\Hom(Y, X') = 0$,
\item $\Hom(Z, Y') = 0$,
\item $\Hom(X, X') = \Hom(Z, X') = 0$,
\item $\Hom(Z, X') = \Hom(Z, Z') = 0$, or
\item $\Hom(X[1], Z') = \Hom(Z, X') = 0$
\end{enumerate}
このとき、$(a, b, c)$ が三角の射となるような $Y \to Y'$ の射として
$b$ は一意である。
\end{lemma}

\begin{proof}
もう一つの三角の射 $(a, b', c)$ があるとすれば、$(0, b - b', 0)$ も
三角の射である。したがって示すべきことは、$X \to Y \to Y'$ と
$Y \to Y' \to Z'$ がともに零となるような射 $b : Y \to Y'$ は零射だけだ、
ということである。以下、補題 \ref{derived-lemma-representable-homological} を
断りなく用いる。特に、条件 (3) は (1) を含意する。条件 (1) を仮定する。
合成 $g' \circ b : Y \to Y' \to Z'$ が零ならば、$b$ は射 $Y \to X'$ に
持ち上がり、この射は $0$ でなければならない。これで (1) が証明された。

\medskip\noindent
(2) と (4) の証明はこの議論の双対である。

\medskip\noindent
(5) を仮定する。図式
$$
\xymatrix{
X \ar[r]_f \ar[d]^0 &
Y \ar[r]_g \ar[d]^b &
Z \ar[r]_h \ar[d]^0 \ar@{..>}[ld]^\epsilon &
X[1] \ar[d]^0 \\
X' \ar[r]^{f'} &
Y' \ar[r]^{g'} &
Z' \ar[r]^{h'} &
X'[1]
}
$$
を考える。$b = \epsilon \circ g$ となるように $\epsilon$ を選べる。
このとき $g' \circ \epsilon \circ g = 0$ なので、ある
$\delta \in \Hom(X[1], Z')$ に対して
$g' \circ \epsilon = \delta \circ h$ となる。
$\Hom(X[1], Z') = 0$ であるから $g' \circ \epsilon = 0$ と結論できる。
したがって、ある $\gamma \in \Hom(Z, X')$ に対して
$\epsilon = f' \circ \gamma$ となる。$\Hom(Z, X') = 0$ であるから
$\epsilon = 0$、したがって望みどおり $b = 0$ である。
\end{proof}

\begin{lemma}
\label{derived-lemma-third-object-zero}
$\mathcal{D}$ を前三角圏とし、$f : X \to Y$ を $\mathcal{D}$ の射とする。
次は同値である。
\begin{enumerate}
\item $f$ は同型である。
\item $(X, Y, 0, f, 0, 0)$ は完全三角である。
\item 任意の完全三角 $(X, Y, Z, f, g, h)$ に対して $Z = 0$ である。
\end{enumerate}
\end{lemma}

\begin{proof}
TR1 により、三角 $(X, X, 0, 1, 0, 0)$ は完全である。
$(X, Y, Z, f, g, h)$ を完全三角とする。TR3 により、完全三角の射
$(1, f, 0) : (X, X, 0) \to (X, Y, Z)$ が存在する。
$f$ が同型ならば、補題 \ref{derived-lemma-third-isomorphism-triangle} により
$(1, f, 0)$ は三角の同型であり、$Z = 0$ である。逆に、$Z = 0$ ならば
$(1, f, 0)$ も三角の同型であり、したがって $f$ は同型である。
\end{proof}

\begin{lemma}
\label{derived-lemma-direct-sum-triangles}
$\mathcal{D}$ を前三角圏とする。
$(X, Y, Z, f, g, h)$ および $(X', Y', Z', f', g', h')$ を三角とする。
次は同値である。
\begin{enumerate}
\item $(X \oplus X', Y \oplus Y', Z \oplus Z',
f \oplus f', g \oplus g', h \oplus h')$
は完全三角である。
\item $(X, Y, Z, f, g, h)$ と $(X', Y', Z', f', g', h')$ はともに
完全三角である。
\end{enumerate}
\end{lemma}

\begin{proof}
(2) を仮定する。TR1 により、完全三角
$(X \oplus X', Y \oplus Y', Q, f \oplus f', g'', h'')$
を選べる。TR3 により、完全三角の射
$(X, Y, Z, f, g, h) \to
(X \oplus X', Y \oplus Y', Q, f \oplus f', g'', h'')$
および
$(X', Y', Z', f', g', h') \to
(X \oplus X', Y \oplus Y', Q, f \oplus f', g'', h'')$
を見いだすことができる。これらの射の直和を取ると、三角の射
$$
\xymatrix{
(X \oplus X', Y \oplus Y', Z \oplus Z',
f \oplus f', g \oplus g', h \oplus h')
\ar[d]^{(1, 1, c)} \\
(X \oplus X', Y \oplus Y', Q, f \oplus f', g'', h'').
}
$$
を得る。注意 \ref{derived-remark-special-triangles} の用語では、これは特殊三角の射である
（特殊三角の直和は特殊である）。したがって $c$ は同型であり、(1) が成り立つ。

\medskip\noindent
(1) を仮定する。$(X, Y, Z, f, g, h)$ が完全三角であることを示す。
まず、$(X, Y, Z, f, g, h)$ は、完全三角、したがって特殊三角である
$(X \oplus X', Y \oplus Y', Z \oplus Z',
f \oplus f', g \oplus g', h \oplus h')$
の直和因子なので、特殊三角である
（注意 \ref{derived-remark-special-triangles} の用語を用いた）。
TR1 により、完全三角 $(X, Y, Q, f, g'', h'')$ を取る。
TR3 により、完全三角の射
 $(X \oplus X', Y \oplus Y', Z \oplus Z',
f \oplus f', g \oplus g', h \oplus h') \to (X, Y, Q, f, g'', h'')$
が存在する。これを包含写像と合成すると、三角の射
$$
(1, 1, c) :
(X, Y, Z, f, g, h)
\longrightarrow
(X, Y, Q, f, g'', h'')
$$
を得る。注意 \ref{derived-remark-special-triangles} により $c$ は同型であり、(2) が
成り立つ。
\end{proof}
\begin{lemma}
\label{derived-lemma-split}
$\mathcal{D}$ を前三角圏とし、$(X, Y, Z, f, g, h)$ を完全三角とする。
\begin{enumerate}
\item $h = 0$ ならば、$g$ の右逆射 $s : Z \to Y$ が存在する。
\item $g$ の任意の右逆射 $s : Z \to Y$ に対して、写像
$f \oplus s : X \oplus Z \to Y$ は同型である。
\item $\mathcal{D}$ の任意の対象 $X', Z'$ に対して、三角
$(X', X' \oplus Z', Z', (1, 0), (0, 1), 0)$ は完全である。
\end{enumerate}
\end{lemma}

\begin{proof}
(1) を見るには、補題 \ref{derived-lemma-representable-homological} により
$\Hom_\mathcal{D}(Z, Y) \to \Hom_\mathcal{D}(Z, Z) \to
\Hom_\mathcal{D}(Z, X[1])$
が完全であることを用いる。同様に、$s$ が (2) のような射ならば $h = 0$ であり、
列
$$
0 \to \Hom_\mathcal{D}(W, X) \to \Hom_\mathcal{D}(W, Y)
\to \Hom_\mathcal{D}(W, Z) \to 0
$$
は分裂完全である（$s : Z \to Y$ によって分裂する）。したがって、米田の補題により
$X \oplus Z \to Y$ は同型である。最後の主張は TR1 と
補題 \ref{derived-lemma-direct-sum-triangles} から従う。
\end{proof}

\begin{lemma}
\label{derived-lemma-when-split}
$\mathcal{D}$ を前三角圏とし、$f : X \to Y$ を $\mathcal{D}$ の射とする。
次は同値である。
\begin{enumerate}
\item $f$ は核をもつ。
\item $f$ は余核をもつ。
% P02-E0019: frozen source has the grammatical typo “is the isomorphic”.
\item $\mathcal{D}$ のある対象 $K, Z, Q$ に対して、$f$ は射影と余射影の合成
$K \oplus Z \to Z \to Z \oplus Q$ と同型である。
\end{enumerate}
\end{lemma}

\begin{proof}
$X' \oplus Z \to Z \oplus Y'$ の形の写像と同型な任意の射は、核と余核の両方を
もつ。したがって (3) $\Rightarrow$ (1), (2) である。次に
(1) $\Rightarrow$ (3) を証明する。まず、$f : X \to Y$ が単射、すなわち
その核が零であるとする。TR1 により、完全三角 $(X, Y, Z, f, g, h)$ が存在する。
補題 \ref{derived-lemma-composition-zero} により、合成 $f \circ h[-1] = 0$ である。
$f$ は単射なので $h[-1] = 0$、したがって $h = 0$ である。すると
補題 \ref{derived-lemma-split} から $Y = X \oplus Z$ となり、(3) が成り立つ。
次に、$f$ が核 $K$ をもつと仮定する。$K \to X$ は単射なので、
$X = K \oplus X'$ であり、$f|_{X'} : X' \to Y$ は単射である。したがって
$Y = X' \oplus Y'$ となり、主張を得る。(2) $\Rightarrow$ (3) はこれの双対である。
\end{proof}

\begin{lemma}
\label{derived-lemma-products-sums-shifts-triangles}
$\mathcal{D}$ を前三角圏とし、$I$ を集合とする。
\begin{enumerate}
\item $X_i$、$i \in I$ を $\mathcal{D}$ の対象の族とする。
\begin{enumerate}
\item $\prod X_i$ が存在すれば、$(\prod X_i)[1] = \prod X_i[1]$ である。
\item $\bigoplus X_i$ が存在すれば、$(\bigoplus X_i)[1] = \bigoplus X_i[1]$ である。
\end{enumerate}
\item $X_i \to Y_i \to Z_i \to X_i[1]$ を $\mathcal{D}$ の完全三角の族とする。
\begin{enumerate}
\item $\prod X_i$、$\prod Y_i$、$\prod Z_i$ が存在すれば、
$\prod X_i \to \prod Y_i \to \prod Z_i \to \prod X_i[1]$
は完全三角である。
\item $\bigoplus X_i$、$\bigoplus Y_i$、$\bigoplus Z_i$ が存在すれば、
$\bigoplus X_i \to \bigoplus Y_i \to \bigoplus Z_i \to \bigoplus X_i[1]$
は完全三角である。
\end{enumerate}
\end{enumerate}
\end{lemma}

\begin{proof}
(1) は、$[1]$ が $\mathcal{D}$ の自己同値であり、直和と積が圏の構造によって
定義されるため成り立つ。(2)(a) を証明しよう。完全三角
$\prod X_i \to \prod Y_i \to Z \to \prod X_i[1]$ を選ぶ。
各 $j$ に対して TR3 を用い、射影写像 $\prod X_i \to X_j$ および
$\prod Y_i \to Y_j$ とともに完全三角の射をなす射 $p_j : Z \to Z_j$ を選べる。
積の定義を用いると、完全三角から積で作られた三角への三角の射に組み込まれる写像
$\prod p_i : Z \to \prod Z_i$ が得られる。アーベル群の完全列の積は完全なので、
「積」の三角
$\prod X_i \to \prod Y_i \to \prod Z_i \to \prod X_i[1]$
は注意 \ref{derived-remark-special-triangles} の用語で特殊である。したがって
注意 \ref{derived-remark-special-triangles} により、この三角の射は同型であり、
TR1 から結論が従う。(2)(b) の証明は双対である。
\end{proof}

\begin{lemma}
\label{derived-lemma-projectors-have-images-triangulated}
$\mathcal{D}$ を前三角圏とする。$\mathcal{D}$ が可算積をもてば、
$\mathcal{D}$ はカルービ圏である。$\mathcal{D}$ が可算余積をもつ場合も、
$\mathcal{D}$ はカルービ圏である。
\end{lemma}

\begin{proof}
$\mathcal{D}$ が可算積をもつと仮定する。Homology, Lemma
\ref{homology-lemma-projectors-have-images} により、右逆射をもつ射が核をもつことを
確かめれば十分である。右逆射をもつ任意の射は全射であり、したがって
補題 \ref{derived-lemma-when-split} により核をもつ。二番目の主張は一番目の双対である。
\end{proof}

\noindent
次の補題を用いると、前三角圏が三角圏であることを証明するのが少し容易になる。

\begin{lemma}
\label{derived-lemma-easier-axiom-four}
$\mathcal{D}$ を前三角圏とする。TR4 を証明するためには、任意の合成可能な射の組
$f : X \to Y$ および $g : Y \to Z$ が与えられたとき、次が存在することを
示せば十分である。
\begin{enumerate}
\item 同型 $i : X' \to X$、$j : Y' \to Y$、$k : Z' \to Z$ が存在し、さらに
$f' = j^{-1}fi : X' \to Y'$ および
$g' = k^{-1}gj : Y' \to Z'$ と置くと、
\item 完全三角
$(X', Y', Q_1, f', p_1, d_1)$、
$(X', Z', Q_2, g' \circ f', p_2, d_2)$、
$(Y', Z', Q_3, g', p_3, d_3)$
が存在して、TR4 の主張が成り立つ。
\end{enumerate}
\end{lemma}

\begin{proof}
$X, Y, Z$ を $X', Y', Z'$ で置き換えることは、完全三角およびその同型の定義から
問題にならない。この補題は、補題 \ref{derived-lemma-cone-triangle-unique-isomorphism}
により完全三角
$(X', Y', Q_1, f', p_1, d_1)$、
$(X', Z', Q_2, g' \circ f', p_2, d_2)$、
$(Y', Z', Q_3, g', p_3, d_3)$
が同型を除いて一意であることから従う。
\end{proof}

\begin{lemma}
\label{derived-lemma-triangulated-subcategory}
$\mathcal{D}$ を前三角圏とし、$\mathcal{D}'$ を $\mathcal{D}$ の加法的充満部分圏と
する。次は同値である。
\begin{enumerate}
\item $(\mathcal{D}', \mathcal{T}')$ が $\mathcal{D}$ の前三角部分圏となるような
三角の集合 $\mathcal{T}'$ が存在する。
\item $\mathcal{D}'$ は $[1]$ によって保たれ、
$[1] : \mathcal{D}' \to \mathcal{D}'$ は自己同値であり、かつ
$\mathcal{D}'$ の任意の射 $f : X \to Y$ に対して、$\mathcal{D}$ の完全三角
$(X, Y, Z, f, g, h)$ であって、$Z$ が $\mathcal{D}'$ の対象と同型となるものが
存在する。
\end{enumerate}
この場合、(1) の $\mathcal{T}'$ は、$X, Y, Z \in \Ob(\mathcal{D}')$ を満たす
$\mathcal{D}$ の完全三角 $(X, Y, Z, f, g, h)$ の集合である。最後に、
$\mathcal{D}$ が三角圏ならば、(1) と (2) は次とも同値である。
\begin{enumerate}
\item[(3)] $\mathcal{D}'$ は三角部分圏である。
\end{enumerate}
\end{lemma}

\begin{proof}
省略する。
\end{proof}

\begin{lemma}
\label{derived-lemma-exact-functor-additive}
前三角圏の完全関手は加法的である。
\end{lemma}

\begin{proof}
$F : \mathcal{D} \to \mathcal{D}'$ を前三角圏の完全関手とする。
$(0, 0, 0, 1_0, 1_0, 0)$ は $\mathcal{D}$ の完全三角なので、三角
$$
(F(0), F(0), F(0), 1_{F(0)}, 1_{F(0)}, F(0))
$$
は $\mathcal{D}'$ において完全である。これは
$1_{F(0)} \circ 1_{F(0)}$ が零であることを意味する。補題
\ref{derived-lemma-composition-zero} を参照されたい。したがって $F(0)$ は
$\mathcal{D}'$ の零対象である。このことから、任意の零射に $F$ を適用したものも
零である（加法圏の射が零であるための必要十分条件は、それが零対象を経由して
分解することである）。次に、補題 \ref{derived-lemma-split} の (3) により
$(X, X \oplus Y, Y, (1, 0), (0, 1), 0)$ が完全三角であることを用いると、
$(F(X), F(X \oplus Y), F(Y), F(1, 0), F(0, 1), 0)$ も完全三角である。
補題 \ref{derived-lemma-split} の (2) により、写像
$F(X) \oplus F(Y) \to F(X \oplus Y)$ は同型である。最後に、
Homology, Lemma \ref{homology-lemma-additive-functor} を適用すればよい。
\end{proof}

\begin{lemma}
\label{derived-lemma-exact-equivalence}
$F : \mathcal{D} \to \mathcal{D}'$ を前三角圏の充満忠実な完全関手とする。
このとき、$\mathcal{D}$ の三角 $(X, Y, Z, f, g, h)$ が完全であるための
必要十分条件は、
$(F(X), F(Y), F(Z), F(f), F(g), F(h))$ が $\mathcal{D}'$ において
完全であることである。
\end{lemma}

\begin{proof}
必要性は明らかである。$(F(X), F(Y), F(Z))$ が $\mathcal{D}'$ において
完全であると仮定する。$\mathcal{D}$ の完全三角
$(X, Y, Z', f, g', h')$ を選ぶ。補題
\ref{derived-lemma-cone-triangle-unique-isomorphism} により、三角の同型
$$
(1, 1, c') : (F(X), F(Y), F(Z)) \longrightarrow (F(X), F(Y), F(Z')).
$$
が存在する。$F$ は充満忠実なので、$F(c) = c'$ となる射 $c : Z \to Z'$ が
存在する。このとき $(1, 1, c)$ は $(X, Y, Z)$ と $(X, Y, Z')$ の間の同型である。
ゆえに TR1 により $(X, Y, Z)$ は完全である。
\end{proof}

\begin{lemma}
\label{derived-lemma-composition-exact}
$\mathcal{D}, \mathcal{D}', \mathcal{D}''$ を前三角圏とする。
$F : \mathcal{D} \to \mathcal{D}'$ および
$F' : \mathcal{D}' \to \mathcal{D}''$ を完全関手とする。
このとき $F' \circ F$ は完全関手である。
\end{lemma}

\begin{proof}
省略する。
\end{proof}

\begin{lemma}
\label{derived-lemma-exact-compose-homological-functor}
$\mathcal{D}$ を前三角圏、$\mathcal{A}$ をアーベル圏とし、
$H : \mathcal{D} \to \mathcal{A}$ をホモロジー的関手とする。
\begin{enumerate}
\item $\mathcal{D}'$ を前三角圏とし、
$F : \mathcal{D}' \to \mathcal{D}$ を完全関手とする。
このとき、合成 $H \circ F$ もホモロジー的関手である。
\item $\mathcal{A}'$ をアーベル圏とし、
$G : \mathcal{A} \to \mathcal{A}'$ を完全関手とする。
このとき $G \circ H$ もホモロジー的関手である。
\end{enumerate}
\end{lemma}

\begin{proof}
省略する。
\end{proof}

\begin{lemma}
\label{derived-lemma-exact-compose-delta-functor}
$\mathcal{D}$ を三角圏、$\mathcal{A}$ をアーベル圏とし、
$G : \mathcal{A} \to \mathcal{D}$ を $\delta$-関手とする。
\begin{enumerate}
\item $\mathcal{D}'$ を三角圏とし、
$F : \mathcal{D} \to \mathcal{D}'$ を完全関手とする。
このとき、合成 $F \circ G$ も $\delta$-関手である。
\item $\mathcal{A}'$ をアーベル圏とし、
$H : \mathcal{A}' \to \mathcal{A}$ を完全関手とする。
このとき $G \circ H$ も $\delta$-関手である。
\end{enumerate}
\end{lemma}

\begin{proof}
省略する。
\end{proof}

\begin{lemma}
\label{derived-lemma-compose-delta-functor-homological}
$\mathcal{D}$ を三角圏、$\mathcal{A}$ および $\mathcal{B}$ をアーベル圏とする。
$G : \mathcal{A} \to \mathcal{D}$ を $\delta$-関手、
$H : \mathcal{D} \to \mathcal{B}$ をホモロジー的関手とする。
$\mathcal{A}$ のすべての $A$ に対して $H^{-1}(G(A)) = 0$ であると仮定する。
このとき族
$$
\{H^n \circ G, H^n(\delta_{A \to B \to C})\}_{n \geq 0}
$$
は $\mathcal{A} \to \mathcal{B}$ の $\delta$-関手である。
Homology, Definition \ref{homology-definition-cohomological-delta-functor} を
参照されたい。
\end{lemma}

\begin{proof}
この記号は次を意味する。
$0 \to A \xrightarrow{a} B \xrightarrow{b} C \to 0$ が
$\mathcal{A}$ の短完全列ならば、
$$
\delta = \delta_{A \to B \to C} : G(C) \to G(A)[1]
$$
は $\mathcal{D}$ の射であって、
$(G(A), G(B), G(C), a, b, \delta)$ は完全三角である。
定義 \ref{derived-definition-delta-functor} を参照されたい。このとき
$H^n(\delta) : H^n(G(C)) \to H^n(G(A)[1]) = H^{n + 1}(G(A))$
が短完全列に関して関手的であることは明らかである。最後に、長完全コホモロジー列
(\ref{derived-equation-long-exact-cohomology-sequence}) と $H^{-1}(G(C))$ の消滅を
組み合わせると、望みどおり $\mathcal{B}$ の長完全列
$$
0 \to H^0(G(A)) \to H^0(G(B)) \to H^0(G(C))
\xrightarrow{H^0(\delta)} H^1(G(A)) \to \ldots
$$
が得られる。
\end{proof}

\noindent
次の結果の証明では TR4 を用いる。

\begin{proposition}
\label{derived-proposition-9}
$\mathcal{D}$ を三角圏とする。任意の可換図式
$$
\xymatrix{
X \ar[r] \ar[d] & Y \ar[d] \\
X' \ar[r] & Y'
}
$$
は、図式
$$
\xymatrix{
X \ar[r] \ar[d] & Y \ar[r] \ar[d] & Z \ar[r] \ar[d] & X[1] \ar[d] \\
X' \ar[r] \ar[d] & Y' \ar[r] \ar[d] & Z' \ar[r] \ar[d] & X'[1] \ar[d] \\
X'' \ar[r] \ar[d] & Y'' \ar[r] \ar[d] & Z'' \ar[r] \ar[d] & X''[1] \ar[d] \\
X[1] \ar[r] & Y[1] \ar[r] & Z[1] \ar[r] & X[2]
}
$$
へ拡張できる。ここですべての正方形は可換であるが、右下の正方形だけは
反可換である。さらに、各行と各列は完全三角である。最後に、下の行
（それぞれ右の列）の射は、上の行（それぞれ左の列）の射に $[1]$ を適用して
得られる。
\end{proposition}

\begin{proof}
証明を読みやすくするため、この証明では矢印を書かない。完全三角
$(X, Y, Z)$、$(X', Y', Z')$、$(X, X', X'')$、$(Y, Y', Y'')$、
$(X, Y', A)$ を選ぶ。射 $X \to Y'$ は、合成 $X \to Y \to Y'$ と
合成 $X \to X' \to Y'$ の両方に等しいことに注意する。したがって TR4 の
とおり、次の射を見いだせる。
\begin{enumerate}
\item $a : Z \to A$ および $b : A \to Y''$、
\item $a' : X'' \to A$ および $b' : A \to Z'$。
\end{enumerate}
$c : Y'' \to Z[1]$ を合成 $Y'' \to Y[1] \to Z[1]$ とし、
$c' : Z' \to X''[1]$ を合成 $Z' \to X'[1] \to X''[1]$ とする。
この TR4 の適用から次が得られる。
\begin{enumerate}
\item $(Z, A, Y'', a, b, c)$ および $(X'', A, Z', a', b', c')$
は完全三角である。
\item $(X, Y, Z) \to (X, Y', A)$、
$(X, Y', A) \to (Y, Y', Y'')$、
$(X, X', X'') \to (X, Y', A)$、
$(X, Y', A) \to (X', Y', Z')$
は三角の射である。
\end{enumerate}
まず、$(X, X', X'') \to (X, Y', A)$ と
$(X, Y', A) \to (Y, Y', Y'')$ が三角の射であることから、次の二図式の
うち一番目が可換であると分かる。
$$
\vcenter{
\xymatrix{
X' \ar[r] \ar[d] & Y' \ar[d] \\
X'' \ar[r]^{b \circ a'} \ar[d] & Y'' \ar[d] \\
X[1] \ar[r] & Y[1]
}
}
\quad\text{and}\quad
\vcenter{
\xymatrix{
Y \ar[r] \ar[d] & Z \ar[d]^{b' \circ a} \ar[r] & X[1] \ar[d] \\
Y' \ar[r] & Z' \ar[r] & X'[1]
}
}
$$
$(X, Y, Z) \to (X, Y', A)$ と $(X, Y', A) \to (X', Y', Z')$ が
三角の射であることから、二番目も可換である。この時点で、写像
$b \circ a' : X'' \to Y''$ から始まる完全三角 $(X'', Y'' , Z'')$ を選ぶ。

\medskip\noindent
次に、射 $X'' \to A \to Y''$ と三角
$(X'', A, Z', a', b', c')$、
$(X'', Y'', Z'')$、
$(A, Y'', Z[1], b, c , -a[1])$ に TR4 をもう一度適用し、射
$a'' : Z' \to Z''$ および $b'' : Z'' \to Z[1]$ を得る。
すると $(Z', Z'', Z[1], a'', b'', - b'[1] \circ a[1])$ は完全三角であり、
したがって $(Z, Z', Z'', -b' \circ a, a'', -b'')$ も完全三角であり、
したがって $(Z, Z', Z'', b' \circ a, a'', b'')$ も完全三角である。
さらに、$(X'', A, Z') \to (X'', Y'', Z'')$ および
$(X'', Y'', Z'') \to (A, Y'', Z[1], b, c , -a[1])$
は三角の射である。この時点ですべての完全三角とすべての射を定義し終えたので、
残るのはいくつかの可換性関係を確かめることだけである。

\medskip\noindent
図式の中央の正方形が可換であることを見る。射 $Y' \to Z'$ は
$Y' \to A \to Z'$ と分解する。なぜなら
$(X, Y', A) \to (X', Y', Z')$ が三角の射だからである。同様に、射
$Y' \to Y''$ は $Y' \to A \to Y''$ と分解する。なぜなら
$(X, Y', A) \to (Y, Y', Y'')$ が三角の射だからである。したがって、
辺が $(A, Z', Z'', Y'')$ である正方形は可換である。実際、TR4 により
$(X'', A, Z') \to (X'', Y'', Z'')$ は三角の射である。
辺が $(Y'', Z'', Y[1], Z[1])$ である正方形は可換である。なぜなら
$(X'', Y'', Z'') \to (A, Y'', Z[1], b, c , -a[1])$
は三角の射であり、$c : Y'' \to Z[1]$ は合成
$Y'' \to Y[1] \to Z[1]$ だからである。
辺が $(Z', X'[1], X''[1], Z'')$ である正方形は可換である。なぜなら
$(X'', A, Z') \to (X'', Y'', Z'')$ は三角の射であり、
$c' : Z' \to X''[1]$ は合成 $Z' \to X'[1] \to X''[1]$ だからである。
最後に、辺が $(Z'', X''[1], Z[1], X[2])$ である正方形が反可換であることを
示さなければならない。これは
$(X'', Y'', Z'') \to (A, Y'', Z[1], b, c , -a[1])$
が三角の射であることから従い、証明が完了する。
\end{proof}


\section{三角圏の局所化}
\label{derived-section-localization}

\noindent
複体のホモトピー圏から導来圏を構成するために、局所化の手続きを用いる。

\begin{definition}
\label{derived-definition-localization}
$\mathcal{D}$ を前三角圏とする。乗法系 $S$ が
{\it 三角構造と両立する} とは、次の二条件が成り立つことをいう。
\begin{enumerate}
\item[MS5] $\mathcal{D}$ の射 $f$ に対して
$f \in S \Leftrightarrow f[1] \in S$\footnote{注意 \ref{derived-remark-MS5} を参照。}
が成り立つ。
\item[MS6] 実線部分が可換な図式
$$
\xymatrix{
X \ar[r] \ar[d]^s &
Y \ar[r] \ar[d]^{s'} &
Z \ar[r] \ar@{-->}[d] &
X[1] \ar[d]^{s[1]} \\
X' \ar[r] &
Y' \ar[r] &
Z' \ar[r] &
X'[1]
}
$$
が与えられ、その各行が完全三角で $s, s' \in S$ を満たすならば、
$(s, s', s'')$ が三角の射となるような $S$ の射 $s'' : Z \to Z'$ が存在する。
\end{enumerate}
\end{definition}

\noindent
これらの公理は、乗法系を定義する公理から独立ではないことが分かる。

\begin{lemma}
\label{derived-lemma-localization-conditions}
$\mathcal{D}$ を前三角圏とし、$S \subset \text{Arrows}(\mathcal{D})$ とする。
\begin{enumerate}
\item $S$ がすべての恒等射を含み、MS6
（定義 \ref{derived-definition-localization}）が成り立つならば、
$\mathcal{D}$ のすべての同型が $S$ に属する。
\item MS1、MS5（Categories, Definition
\ref{categories-definition-multiplicative-system}）、MS6 が成り立つならば、
MS2 も成り立つ。
\end{enumerate}
\end{lemma}

\begin{proof}
$S$ がすべての恒等射を含み、MS6 が成り立つと仮定する。
$f : X \to Y$ を $\mathcal{D}$ の同型とする。図式
$$
\xymatrix{
0 \ar[r] \ar[d]^1 &
X \ar[r]_1 \ar[d]^1 &
X \ar[r] \ar@{-->}[d] &
0[1] \ar[d]^{1[1]} \\
0 \ar[r] &
X \ar[r]^f &
Y \ar[r] &
0[1]
}
$$
を考える。補題 \ref{derived-lemma-third-object-zero} により各行は完全三角である。
MS6 により点線の矢印が存在して $S$ に属するので、$f$ は $S$ に属する。

\medskip\noindent
MS1、MS5、MS6 を仮定する。$f : X \to Y$ を $\mathcal{D}$ の射とし、
$t : X \to X'$ を $S$ の元とする。完全三角 $(X, Y, Z, f, g, h)$ を選ぶ。
次に完全三角 $(X', Y', Z, f', g', t[1] \circ h)$ を選ぶ
（ここでは TR1 と TR2 を用いる）。MS5、MS6（および回転のための TR2）により、
可換図式
$$
\xymatrix{
X \ar[r] \ar[d]^t &
Y \ar[r] \ar@{..>}[d]^{s'} &
Z \ar[r] \ar[d]^1 &
X[1] \ar[d]^{t[1]} \\
X' \ar[r] &
Y' \ar[r] &
Z \ar[r] &
X'[1]
}
$$
の点線の矢印を、さらに $s' \in S$ となるように見いだせる。これは LMS2 を
証明する。RMS2 の証明は双対である。
\end{proof}

\begin{remark}
\label{derived-remark-MS5}
MS1 と MS6 のもとで、条件 MS5 は、すべての $s \in S$ と
$n \in \mathbf{Z}$ に対して $s[n] \in S$ であるという条件と同値である。
例えば、MS5 が成り立ち、$s \in S$ であり、$s[-1] \in S$ を示したいとする。
$s[-1][1]$ は $s$ と等しいのではなく、$\mathcal{D}$ の矢印として $s$ と
同型なだけなので、これは直ちには分からない。それでも、ある同型 $f$, $g$ に
対して $s[-1][1] = f \circ s \circ g$ となることは従う。補題
\ref{derived-lemma-localization-conditions} (1) により $f, g \in S$ であるから、
MS1 により $s[-1][1] \in S$、さらに MS5 により $s[-1] \in S$ となる。
完全な証明は演習として読者に委ねる。
\end{remark}

\begin{lemma}
\label{derived-lemma-triangle-functor-localize}
$F : \mathcal{D} \to \mathcal{D}'$ を前三角圏の完全関手とする。
$$
S = \{f \in \text{Arrows}(\mathcal{D})
\mid F(f)\text{ is an isomorphism}\}
$$
と置く。このとき $S$ は飽和しており（Categories,
Definition \ref{categories-definition-saturated-multiplicative-system} を参照）、
$\mathcal{D}$ の三角構造と両立する乗法系である。
\end{lemma}

\begin{proof}
公理 MS1 -- MS6 を証明しなければならない。Categories, Definitions
\ref{categories-definition-multiplicative-system}、
\ref{categories-definition-saturated-multiplicative-system}、および
定義 \ref{derived-definition-localization} を参照されたい。
MS1、MS4、MS5 は定義から直ちに従う。MS6 は TR3 と
補題 \ref{derived-lemma-third-isomorphism-triangle} から従う。補題
\ref{derived-lemma-localization-conditions} により MS2 が成り立つ。証明を終えるには
MS3 を示さなければならない。このため、$f, g : X \to Y$ を
$\mathcal{D}$ の射とし、$t : Z \to X$ を
$f \circ t = g \circ t$ を満たす $S$ の元とする。$\mathcal{D}$ は加法的なので、
これは $a = f - g$ と置いたとき $a \circ t = 0$ であることにほかならない。
TR1 により完全三角 $(Z, X, Q, t, d, h)$ を選ぶ。
$a \circ t = 0$ なので、補題 \ref{derived-lemma-representable-homological} により、
$i \circ d = a$ を満たす射 $i : Q \to Y$ が存在する。最後に、TR1 を再び用いて
三角 $(Q, Y, W, i, j, k)$ を選べる。図示すると
$$
\xymatrix{
Z \ar[r]_t & X \ar[r]_d \ar[d]^1 & Q \ar[r] \ar[d]^i & Z[1] \\
& X \ar[r]_a & Y \ar[d]^j \\
& & W
}
$$
となる。さて、この図式に関手 $F$ を適用する。$t \in S$ なので
$F(Q) = 0$ である。補題 \ref{derived-lemma-third-object-zero} を参照されたい。
したがって同じ補題により $F(j)$ は同型、すなわち $j \in S$ である。
最後に、$j \circ i = 0$ なので
$j \circ a = j \circ i \circ d = 0$ である。したがって
$j \circ f = j \circ g$ となり、LMS3 が成り立つ。RMS3 の証明は双対である。
\end{proof}

\begin{lemma}
\label{derived-lemma-homological-functor-localize}
$H : \mathcal{D} \to \mathcal{A}$ を、前三角圏からアーベル圏への
ホモロジー的関手とする。
$$
S = \{f \in \text{Arrows}(\mathcal{D})
\mid H^i(f)\text{ is an isomorphism for all }i \in \mathbf{Z}\}
$$
と置く。このとき $S$ は飽和しており（Categories,
Definition \ref{categories-definition-saturated-multiplicative-system} を参照）、
$\mathcal{D}$ の三角構造と両立する乗法系である。
\end{lemma}

\begin{proof}
公理 MS1 -- MS6 を証明しなければならない。Categories, Definitions
\ref{categories-definition-multiplicative-system}、
\ref{categories-definition-saturated-multiplicative-system}、および
定義 \ref{derived-definition-localization} を参照されたい。
MS1、MS4、MS5 は定義から直ちに従う。MS6 は TR3 と長完全コホモロジー列
(\ref{derived-equation-long-exact-cohomology-sequence}) から従う。補題
\ref{derived-lemma-localization-conditions} により MS2 が成り立つ。証明を終えるには
MS3 を示さなければならない。このため、$f, g : X \to Y$ を
$\mathcal{D}$ の射とし、$t : Z \to X$ を
$f \circ t = g \circ t$ を満たす $S$ の元とする。$\mathcal{D}$ は加法的なので、
これは $a = f - g$ と置いたとき $a \circ t = 0$ であることにほかならない。
TR1 と TR2 により完全三角 $(Z, X, Q, t, u, v)$ を選ぶ。
$a \circ t = 0$ なので、補題 \ref{derived-lemma-representable-homological} により、
% P02-E0020: frozen source writes g where the factorization requires u.
$i \circ g = a$ を満たす射 $i : Q \to Y$ が存在する。最後に、TR1 を再び用いて
三角 $(Q, Y, W, i, j, k)$ を選べる。図示すると
$$
\xymatrix{
Z \ar[r]_t & X \ar[r]_u \ar[d]^1 & Q \ar[r]_v \ar[d]^i & Z[1] \\
& X \ar[r]_a & Y \ar[d]^j \\
& & W
}
$$
となる。さて、この図式に関手 $H^i$ を適用する。$t \in S$ なので、
長完全コホモロジー列 (\ref{derived-equation-long-exact-cohomology-sequence}) により
$H^i(Q) = 0$ である。同じ議論から、すべての $i$ に対して $H^i(j)$ は同型、
すなわち $j \in S$ である。最後に、$j \circ i = 0$ なので
$j \circ a = j \circ i \circ u = 0$ である。したがって
$j \circ f = j \circ g$ となり、LMS3 が成り立つ。RMS3 の証明は双対である。
\end{proof}

\begin{proposition}
\label{derived-proposition-construct-localization}
$\mathcal{D}$ を前三角圏とし、$S$ を三角構造と両立する乗法系とする。
このとき、$[1] \circ Q = Q \circ [1]$ であり、局所化関手
$Q : \mathcal{D} \to S^{-1}\mathcal{D}$ が完全となるような
$S^{-1}\mathcal{D}$ 上の前三角圏構造が一意に存在する。さらに、
$\mathcal{D}$ が三角圏ならば $S^{-1}\mathcal{D}$ も三角圏である。
\end{proposition}

\begin{proof}
Homology, Lemma \ref{homology-lemma-localization-additive} において、
$S^{-1}\mathcal{D}$ が加法圏であり、局所化関手 $Q$ が加法的であることを見た。
MS5 から、$Q \circ [1] = [1] \circ Q$（関手としての等式）を満たす
加法的自己同値
$[1] : S^{-1}\mathcal{D} \to S^{-1}\mathcal{D}$
が一意に存在することが従う。詳細は省略する。$S^{-1}\mathcal{D}$ の三角が
完全であるとは、それが局所化関手 $Q$ による完全三角の像と同型であることと定める。

\medskip\noindent
TR1 の証明。ここで証明すべき唯一のことは、
$a : Q(X) \to Q(Y)$ が $S^{-1}\mathcal{D}$ の射ならば、$a$ が完全三角に
組み込まれることである。$S$ のある $s : Y \to Y'$ と射
$f : X \to Y'$ によって $a = Q(s)^{-1} \circ Q(f)$ と書く。
$\mathcal{D}$ の完全三角 $(X, Y', Z, f, g, h)$ を選ぶ。このとき
$(Q(X), Q(Y), Q(Z), a, Q(g) \circ Q(s), Q(h))$ は
$S^{-1}\mathcal{D}$ の完全三角である。

\medskip\noindent
TR2 の証明。定義から直ちに従う。

\medskip\noindent
TR3 の証明。存在を要求される点線の矢印は、二つの三角を同型な三角で置き換えた後に
構成してよいことに注意する。したがって、$\mathcal{D}$ の完全三角
$(X, Y, Z, f, g, h)$、$(X', Y', Z', f', g', h')$ と、
$S^{-1}\mathcal{D}$ における可換図式
$$
\xymatrix{
Q(X) \ar[r]_{Q(f)} \ar[d]_a & Q(Y) \ar[d]^b \\
Q(X') \ar[r]^{Q(f')} & Q(Y')
}
$$
が与えられていると仮定してよい。Categories, Lemma
\ref{categories-lemma-left-localization-diagram} を適用すると、
$\mathcal{D}$ の射 $f'' : X'' \to Y''$ と可換図式
$$
\xymatrix{
X \ar[d]_f \ar[r]_k & X'' \ar[d]^{f''} & X' \ar[d]^{f'} \ar[l]^s \\
Y \ar[r]^l & Y'' & Y' \ar[l]_t
}
$$
であって、$s, t \in S$ かつ $a = s^{-1}k$、$b = t^{-1}l$ を満たすものが
得られる。この時点で $\mathcal{D}$ に対する TR3 と MS6 を用い、可換図式
% P02-E0021: frozen source labels the shifted vertical arrow g[1], not k[1].
$$
\xymatrix{
X \ar[r] \ar[d]^k &
Y \ar[r] \ar[d]^l &
Z \ar[r] \ar[d]^m &
X[1] \ar[d]^{g[1]} \\
X'' \ar[r] &
Y'' \ar[r] &
Z'' \ar[r] &
X''[1] \\
X' \ar[r] \ar[u]_s &
Y' \ar[r] \ar[u]_t &
Z' \ar[r] \ar[u]_r &
X'[1] \ar[u]_{s[1]}
}
$$
であって $r \in S$ となるものを得る。したがって
$c = Q(r)^{-1}Q(m)$ と置けば、求める三角の射
$$
\xymatrix{
(Q(X), Q(Y), Q(Z), Q(f), Q(g), Q(h))
\ar[d]^{(a, b, c)} \\
(Q(X'), Q(Y'), Q(Z'), Q(f'), Q(g'), Q(h'))
}
$$
が得られる。

\medskip\noindent
以上で最初の主張が証明された。さらに $\mathcal{D}$ が三角圏であるとき、
$S^{-1}\mathcal{D}$ も三角圏であることを示すには、なお TR4 を証明しなければ
ならない。このため、補題 \ref{derived-lemma-easier-axiom-four} によって次の主張に
帰着する。合成可能な射 $a : Q(X) \to Q(Y)$ と $b : Q(Y) \to Q(Z)$ が
与えられたとき、$Q(X), Q(Y), Q(Z)$ を必要に応じて同型な対象で置き換えた後、
八面体を構成しなければならない。まず、$\mathcal{D}$ のある射
$f : X \to Y$ によって $a = Q(f)$ となるような対象で $Y$ を置き換えてよい。
（より正確には、$S$ の $s : Y \to Y'$ と射 $f : X \to Y'$ によって
$a = s^{-1}f$ と書き、$Y$ を $Y'$ で置き換える。）その後、同様に
$\mathcal{D}$ のある射 $g : Y \to Z$ によって $b = Q(g)$ となるような対象で
$Z$ を置き換える。ここで $\mathcal{D}$ の完全三角
$(X, Y, Q_1, f, p_1, d_1)$、
$(X, Z, Q_2, g \circ f, p_2, d_2)$、
$(Y, Z, Q_3, g, p_3, d_3)$
を TR1 により取り、TR4 のとおりの射 $a : Q_1 \to Q_2$ および
$b : Q_2 \to Q_3$ を取る。これらすべてのデータに関手 $Q$ を適用すれば、
$S^{-1}\mathcal{D}$ に対応する構造が得られることは直ちに確かめられる。
\end{proof}

\noindent
三角圏の局所化の普遍性は次のとおりである
（前三角圏について定式化するので、特に三角圏についても成り立つ）。

\begin{lemma}
\label{derived-lemma-universal-property-localization}
$\mathcal{D}$ を前三角圏とし、$S$ を三角構造と両立する乗法系とする。
$Q : \mathcal{D} \to S^{-1}\mathcal{D}$ を局所化関手とする。
命題 \ref{derived-proposition-construct-localization} を参照されたい。
\begin{enumerate}
\item $H : \mathcal{D} \to \mathcal{A}$ をアーベル圏 $\mathcal{A}$ への
ホモロジー的関手とし、すべての $s \in S$ に対して $H(s)$ が同型であるとする。
このとき、$H = H' \circ Q$ を満たす一意な分解
$H' : S^{-1}\mathcal{D} \to \mathcal{A}$
（Categories, Lemma \ref{categories-lemma-properties-left-localization} を参照）
もホモロジー的関手である。
\item $F : \mathcal{D} \to \mathcal{D}'$ を前三角圏 $\mathcal{D}'$ への
完全関手とし、すべての $s \in S$ に対して $F(s)$ が同型であるとする。
このとき、$F = F' \circ Q$ を満たす一意な分解
$F' : S^{-1}\mathcal{D} \to \mathcal{D}'$
（Categories, Lemma \ref{categories-lemma-properties-left-localization} を参照）
も完全関手である。
\end{enumerate}
\end{lemma}

\begin{proof}
この補題はほとんど自明である。詳細は省略する。
\end{proof}

\begin{lemma}
\label{derived-lemma-localization-subcategory}
$\mathcal{D}$ を前三角圏とし、
$\mathcal{D}' \subset \mathcal{D}$ を充満な前三角部分圏とする。
$S$ を $\mathcal{D}$ の、三角構造と両立する飽和乗法系とする。
$\mathcal{D}$ の各 $X$ に対して、$X'$ が $\mathcal{D}'$ の対象であるような
$S$ の元 $s : X' \to X$ が存在すると仮定する。このとき
$S' = S \cap \text{Arrows}(\mathcal{D}')$ は三角構造と両立する飽和乗法系であり、
関手
$$
(S')^{-1}\mathcal{D}' \longrightarrow S^{-1}\mathcal{D}
$$
は前三角圏の同値である。
\end{lemma}

\begin{proof}
命題 \ref{derived-proposition-construct-localization} の商関手
$Q : \mathcal{D} \to S^{-1}\mathcal{D}$ を考える。$S$ は飽和しているので、
射 $f : X \to Y$ が $S$ に属するための必要十分条件は $Q(f)$ が可逆である
ことである。Categories, Lemma \ref{categories-lemma-what-gets-inverted} を
参照されたい。したがって $S'$ は、合成
$\mathcal{D}' \to \mathcal{D} \to S^{-1}\mathcal{D}$
によって同型に移される矢印の集合である。ゆえに、
% P02-E0022: frozen source duplicates the verb in “is is”.
補題 \ref{derived-lemma-triangle-functor-localize} により $S'$ は三角構造と両立する
飽和乗法系である。補題 \ref{derived-lemma-universal-property-localization} により、
前三角圏の完全関手 $(S')^{-1}\mathcal{D}' \to S^{-1}\mathcal{D}$ を得る。
仮定により、この関手は本質的全射である。$X', Y'$ を $\mathcal{D}'$ の対象とする。
Categories, Remark \ref{categories-remark-right-localization-morphisms-colimit} により
$$
\Mor_{S^{-1}\mathcal{D}}(X', Y') =
\colim_{s : X \to X'\text{ in }S} \Mor_\mathcal{D}(X, Y')
$$
である。仮定から、任意の $S$ の元 $s : X \to X'$ に対して、
$X''$ が $\mathcal{D}'$ に属するような $S$ の射 $s' : X'' \to X$ が
存在する。このとき $s \circ s' : X'' \to X'$ は $S'$ に属する。
したがって上の余極限は
$$
\colim_{s'' : X'' \to X'\text{ in }S'} \Mor_{\mathcal{D}'}(X'', Y') =
\Mor_{(S')^{-1}\mathcal{D}'}(X', Y')
$$
に等しい。これにより関手が充満忠実でもあることが示され、証明が完了する。
\end{proof}

\noindent
次の補題は、局所化関手の核
（定義 \ref{derived-definition-kernel-category} を参照）を記述する。

\begin{lemma}
\label{derived-lemma-kernel-localization}
$\mathcal{D}$ を前三角圏、$S$ を三角構造と両立する乗法系とし、
$Z$ を $\mathcal{D}$ の対象とする。次は同値である。
\begin{enumerate}
\item $S^{-1}\mathcal{D}$ において $Q(Z) = 0$ である。
\item $0 : Z \to Z'$ が $S$ の元となるような
$Z' \in \Ob(\mathcal{D})$ が存在する。
\item $0 : Z' \to Z$ が $S$ の元となるような
$Z' \in \Ob(\mathcal{D})$ が存在する。
\item ある対象 $Z'$ と完全三角 $(X, Y, Z \oplus Z', f, g, h)$ であって
$f \in S$ を満たすものが存在する。
\end{enumerate}
$S$ が飽和していれば、これらは次とも同値である。
\begin{enumerate}
\item[(5)] 射 $0 \to Z$ は $S$ の元である。
\item[(6)] 射 $Z \to 0$ は $S$ の元である。
\item[(7)] $f \in S$ を満たす完全三角 $(X, Y, Z, f, g, h)$ が存在する。
\end{enumerate}
\end{lemma}

\begin{proof}
(1)、(2)、(3) の同値性は Homology, Lemma
\ref{homology-lemma-kernel-localization} である。(2) が成り立つならば、
$(Z'[-1], Z'[-1] \oplus Z, Z, (1, 0), (0, 1), 0)$ は
「$0 \in S$」をもつ完全三角である
（補題 \ref{derived-lemma-split} を参照）。回転することにより (4) が従う。
$(X, Y, Z \oplus Z', f, g, h)$ が $f \in S$ を満たす完全三角ならば、
$Q(f)$ は同型であり、したがって $Q(Z \oplus Z') = 0$、さらに
$Q(Z) = 0$ である。よって (1) -- (4) はすべて同値である。

\medskip\noindent
次に $S$ が飽和していると仮定する。(5)、(6)、(7) はいずれも、同値な条件
(1) -- (4) の一つを含意することに注意する。$Q(Z) = 0$ と仮定する。
すると $0 \to Z$ は、$S^{-1}\mathcal{D}$ において同型となる
$\mathcal{D}$ の射である。Categories, Lemma
\ref{categories-lemma-what-gets-inverted} によれば、$S$ が飽和していることから
$0 \to Z$ は $S$ に属する。したがって (1) $\Rightarrow$ (5) である。
双対的に (1) $\Rightarrow$ (6) である。最後に、$0 \to Z$ が $S$ に
属するならば、三角 $(0, Z, Z, 0, \text{id}_Z, 0)$ は TR1 と TR2 により
完全であり、(7) のような三角である。
\end{proof}

\begin{lemma}
\label{derived-lemma-limit-triangles}
$\mathcal{D}$ を前三角圏とする。$S$ を $\mathcal{D}$ の三角構造と両立する
飽和乗法系とし、$(X, Y, Z, f, g, h)$ を $\mathcal{D}$ の完全三角とする。
三角の射からなる圏
$$
\mathcal{I} =
\{(s, s', s'') : (X, Y, Z, f, g, h) \to (X', Y', Z', f', g', h')
\mid s, s', s'' \in S\}
$$
を考える。このとき $\mathcal{I}$ はフィルター圏であり、関手
$\mathcal{I} \to X/S$、$\mathcal{I} \to Y/S$、$\mathcal{I} \to Z/S$
は余終である。
\end{lemma}

\begin{proof}
読者には、この補題の証明を読み飛ばし、代わりに紙片の上で自ら導くことを強く勧める。

\medskip\noindent
最初に、完全三角の回転（TR2）を用いると、$\mathcal{I}$ と、完全三角
$(Y, Z, X[1], g, h, -f[1])$ に対応する圏との圏同値が得られることに注意する。
これを用いると、例えば関手 $\mathcal{I} \to X/S$ が余終であることを証明すれば、
関手 $\mathcal{I} \to Y/S$ と $\mathcal{I} \to Z/S$ についても同じことが
成り立つと分かる。

\medskip\noindent
$s : X \to X'$ を $S$ の射とする。MS2 により、$f' \circ s = s' \circ f$ を
満たす $s' : Y \to Y'$ と $f' : X' \to Y'$ を見いだし、次いで MS6 により
これを $\mathcal{I}$ の対象へ完成できる。したがって関手
$\mathcal{I} \to X/S$ は対象上全射である。上のように回転することで、
関手 $\mathcal{I} \to Y/S$ と $\mathcal{I} \to Z/S$ についても同じことが
成り立つ。

\medskip\noindent
$X/S$ の対象 $s_1 : X \to X_1 $、$s_2 : X \to X_2$ と、$X/S$ の射
$a : X_1 \to X_2$ が与えられたとする。$S$ は飽和しているので
$a \in S$ である。Categories, Lemma
\ref{categories-lemma-what-gets-inverted} を参照されたい。前段落の議論により、
$s_1 : X \to X_1$ を $\mathcal{I}$ の対象
$(s_1, s'_1, s''_1) : (X, Y, Z, f, g, h) \to (X_1, Y_1, Z_1, f_1, g_1, h_1)$
へ完成できる。次にこれを繰り返し、与えられた $a : X_1 \to X_2$ を完成する
$(a, b, c) : (X_1, Y_1, Z_1, f_1, g_1, h_1) \to (X_2, Y_2, Z_2, f_2, g_2, h_2)$
であって $a, b, c \in S$ を満たすものを見いだせる。このとき
$(a, b, c)$ は $\mathcal{I}$ の射である。このようにして、関手
$\mathcal{I} \to X/S$ は射上でも全射である。上のように回転すれば、
関手 $\mathcal{I} \to Y/S$ と $\mathcal{I} \to Z/S$ についても同じことが
成り立つ。

\medskip\noindent
恒等射が対象を与えるので、圏 $\mathcal{I}$ は空でない。これでフィルター圏の
定義の条件 (1) が証明された。Categories, Definition
\ref{categories-definition-directed} を参照されたい。

\medskip\noindent
圏 $\mathcal{I}$ について Categories, Definition
\ref{categories-definition-directed} の条件 (2) を確かめる。$\mathcal{I}$ の対象
$(s_1, s'_1, s''_1) : (X, Y, Z, f, g, h) \to (X_1, Y_1, Z_1, f_1, g_1, h_1)$
および
$(s_2, s'_2, s''_2) : (X, Y, Z, f, g, h) \to (X_2, Y_2, Z_2, f_2, g_2, h_2)$
が与えられたとする。$(X_1, Y_1, Z_1, f_1, g_1, h_1)$ と
$(X_2, Y_2, Z_2, f_2, g_2, h_2)$ の両方から射を受ける $\mathcal{I}$ の対象を
見いだしたい。Categories, Remark
\ref{categories-remark-left-localization-morphisms-colimit} により、圏
$X/S$、$Y/S$、$Z/S$ はフィルター圏である。したがって、$X/S$ の射
$X \to X_3$、$s : X_2 \to X_3$、$a : X_1 \to X_3$ を見いだせる。
上の議論により、$s', s'' \in S$ を満たす三角の射
$(s, s', s'') : (X_2, Y_2, Z_2, f_2, g_2, h_2) \to
(X_3, Y_3, Z_3, f_3, g_3, h_3)$
を見いだせる。$(X_2, Y_2, Z_2)$ を $(X_3, Y_3, Z_3)$ で置き換えれば、
$X/S$ の射 $a : X_1 \to X_2$ が存在すると仮定してよい。$Y$ と $Z$ について
同じ議論を繰り返すことにより（上のように回転する）、$X/S$ の射
$a : X_1 \to X_2$、$Y/S$ の射 $b : Y_1 \to Y_2$、$Z/S$ の射
$c : Z_1 \to Z_2$ が存在すると仮定してよい。しかし、これらの射が必ずしも
完全三角の射を与えるとは限らない。一方、必要な図式は
$S^{-1}\mathcal{D}$ において可換である。したがって、例えば
$s'_2 \circ f_2 \circ a = s'_2 \circ b \circ f_1$
を満たす $S$ の射 $s'_2 : Y_2 \to Y_3$ が存在することが分かる。
上のように $(X_2, Y_2, Z_2)$ をもう一度置き換えれば、
$f_2 \circ a = b \circ f_1$ となる状況に帰着する。さらに二度回転して同じ議論を
適用すると、$(a, b, c)$ が三角の射であると仮定してよい。これで条件 (2) が
証明された。

\medskip\noindent
次に Categories, Definition \ref{categories-definition-directed} の条件 (3) を
確かめる。$(s_1, s_1', s_1'') : (X, Y, Z) \to (X_1, Y_1, Z_1)$ と
$(s_2, s_2', s_2'') : (X, Y, Z) \to (X_2, Y_2, Z_2)$
を $\mathcal{I}$ の対象とし、$(a, b, c), (a', b', c')$ をそれらの間の二つの
射とする。$a \circ s_1 = a' \circ s_1$ なので、
$s_3 \circ a = s_3 \circ a'$ を満たす射 $s_3 : X_2 \to X_3$ が存在する。
上の全射性の主張を用いると、これを $s_3, s_3', s_3'' \in S$ を満たす
三角の射
$(s_3, s_3', s_3'') : (X_2, Y_2, Z_2) \to (X_3, Y_3, Z_3)$
へ完成できる。したがって
$(s_3 \circ s_2, s_3' \circ s_2', s_3'' \circ s_2'') :
(X, Y, Z) \to (X_3, Y_3, Z_3)$
も $\mathcal{I}$ の対象であり、写像 $(a, b, c), (a', b', c')$ を
$(s_3, s_3', s_3'')$ と合成すると $a = a'$ を得る。回転することにより、
同じ操作を行って $b = b'$ と $c = c'$ も得られる。

\medskip\noindent
最後に、$\mathcal{I} \to X/S$ が余終であることを確かめる。Categories,
Definition \ref{categories-definition-cofinal} を参照されたい。この関手は
全射なので、一番目の条件は成り立つ。$X/S$ の対象 $s : X \to X'$ と、
$\mathcal{I}$ の二対象
$(s_1, s'_1, s''_1) : (X, Y, Z, f, g, h) \to (X_1, Y_1, Z_1, f_1, g_1, h_1)$
および
$(s_2, s'_2, s''_2) : (X, Y, Z, f, g, h) \to (X_2, Y_2, Z_2, f_2, g_2, h_2)$、
さらに $X/S$ の射 $t_1 : X' \to X_1$、$t_2 : X' \to X_2$ が
与えられたとする。上で証明した $\mathcal{I}$ の性質 (2) により、射
$(s_3, s'_3, s''_3) : (X_1, Y_1, Z_1, f_1, g_1, h_1) \to
(X_3, Y_3, Z_3, f_3, g_3, h_3)$
および
$(s_4, s'_4, s''_4) : (X_2, Y_2, Z_2, f_2, g_2, h_2) \to
(X_3, Y_3, Z_3, f_3, g_3, h_3)$
を $\mathcal{I}$ において見いだせる。合成
$X' \to X_1 \to X_3$ と $X' \to X_2 \to X_3$ が等しければ完了する
（Categories, Definition \ref{categories-definition-cofinal} の表示式を参照）。
等しくない場合、$X/S$ はフィルター圏なので、合成
$X' \to X_1 \to X_3 \to X_4$ と $X' \to X_2 \to X_3 \to X_4$ が
等しくなるような $X/S$ の射 $X_3 \to X_4$ を選べる。最後に
$X_3 \to X_4$ を $\mathcal{I}$ の射
$(X_3, Y_3, Z_3) \to (X_4, Y_4, Z_4)$
へ完成し、その射と合成すれば所望の結果を得る。
\end{proof}
\section{三角圏の商}
\label{derived-section-quotients}

\noindent
三角圏とその三角部分圏が与えられると、「商」を取ることによって別の三角圏を
構成できる。この構成には局所化を用いる。これはアーベル圏をセール部分圏で
割った商と類似している。Homology, Section
\ref{homology-section-serre-subcategories} を参照されたい。実際の構成を行う前に、
完全関手の核について簡単に論じる。

\begin{definition}
\label{derived-definition-saturated}
$\mathcal{D}$ を前三角圏とする。$\mathcal{D}$ の充満前三角部分圏
$\mathcal{D}'$ が {\it 飽和している} とは、$X \oplus Y$ が
$\mathcal{D}'$ の対象と同型ならば、$X$ と $Y$ の両方が
$\mathcal{D}'$ の対象と同型であることをいう。
\end{definition}

\noindent
飽和三角部分圏は {\it 厚い三角部分圏} と呼ばれることもある。一部の文献では、
この語は厳密充満三角部分圏に対してだけ用いられる（また、定義そのものが
厳密充満性を含意するように書かれることもある）。これとは別に
{\it \'epaisse 三角部分圏} という概念もある。その定義は次のとおりである。
可換図式
$$
\xymatrix{
& S \ar[rd] \\
X \ar[ru] \ar[rr] & & Y \ar[r] & T \ar[r] & X[1]
}
$$
が与えられ、第二の行が完全三角であり、$S$ と $T$ が $\mathcal{D}'$ の対象と
同型ならば、$X$ と $Y$ も $\mathcal{D}'$ の対象と同型であると要求する。
これは飽和性と同値であることが分かる
（これは初等的であり、\cite{Rickard-derived} に見いだせる）。
実際に扱うには、飽和圏という概念の方が容易である。

\begin{lemma}
\label{derived-lemma-triangle-functor-kernel}
$F : \mathcal{D} \to \mathcal{D}'$ を前三角圏の完全関手とする。
$\mathcal{D}''$ を、対象が
$$
\Ob(\mathcal{D}'') =
\{X \in \Ob(\mathcal{D}) \mid F(X) = 0\}
$$
である $\mathcal{D}$ の充満部分圏とする。このとき $\mathcal{D}''$ は
$\mathcal{D}$ の厳密充満な飽和前三角部分圏である。
$\mathcal{D}$ が三角圏ならば、$\mathcal{D}''$ は三角部分圏である。
\end{lemma}

\begin{proof}
$\mathcal{D}''$ が $[1]$ と $[-1]$ によって保たれることは明らかである。
$(X, Y, Z, f, g, h)$ を $\mathcal{D}$ の完全三角とし、
$F(X) = F(Y) = 0$ とする。このとき
$(F(X), F(Y), F(Z), F(f), F(g), F(h))$ は完全なので、$F(Z) = 0$ でもある。
したがって補題 \ref{derived-lemma-triangulated-subcategory} を適用し、
$\mathcal{D}''$ が前三角部分圏
（それぞれ、$\mathcal{D}$ が三角圏ならば三角部分圏）であることが分かる。
最後の飽和性の主張は
$F(X) \oplus F(Y) = 0 \Rightarrow F(X) = F(Y) = 0$
から従う。
\end{proof}

\begin{lemma}
\label{derived-lemma-homological-functor-kernel}
$H : \mathcal{D} \to \mathcal{A}$ を、前三角圏からアーベル圏への
ホモロジー的関手とする。$\mathcal{D}'$ を、対象が
$$
\Ob(\mathcal{D}') =
\{X \in \Ob(\mathcal{D}) \mid
H(X[n]) = 0\text{ for all }n \in \mathbf{Z}\}
$$
である $\mathcal{D}$ の充満部分圏とする。このとき $\mathcal{D}'$ は
$\mathcal{D}$ の厳密充満な飽和前三角部分圏である。
$\mathcal{D}$ が三角圏ならば、$\mathcal{D}'$ は三角部分圏である。
\end{lemma}

\begin{proof}
$\mathcal{D}'$ が $[1]$ と $[-1]$ によって保たれることは明らかである。
$(X, Y, Z, f, g, h)$ を $\mathcal{D}$ の完全三角とし、すべての $n$ に対して
$H(X[n]) = H(Y[n]) = 0$ とする。このとき長完全列
(\ref{derived-equation-long-exact-cohomology-sequence}) により、すべての $n$ に対して
$H(Z[n]) = 0$ でもある。したがって補題
\ref{derived-lemma-triangulated-subcategory} を適用し、$\mathcal{D}'$ が前三角部分圏
（それぞれ、$\mathcal{D}$ が三角圏ならば三角部分圏）であることが分かる。
飽和性の主張は、すべての $n \in \mathbf{Z}$ に対する
\begin{align*}
H((X \oplus Y)[n]) = 0 & \Rightarrow H(X[n] \oplus Y[n]) = 0 \\
& \Rightarrow H(X[n]) \oplus H(Y[n]) = 0 \\
& \Rightarrow H(X[n]) = H(Y[n]) = 0
\end{align*}
から従う。
\end{proof}

\begin{lemma}
\label{derived-lemma-homological-functor-bounded}
$H : \mathcal{D} \to \mathcal{A}$ を、前三角圏からアーベル圏への
ホモロジー的関手とする。
$\mathcal{D}_H^{+}, \mathcal{D}_H^{-}, \mathcal{D}_H^b$
を、それぞれ対象が
$$
\begin{matrix}
\Ob(\mathcal{D}_H^{+}) =
\{X \in \Ob(\mathcal{D}) \mid
H(X[n]) = 0\text{ for all }n \ll 0\} \\
\Ob(\mathcal{D}_H^{-}) =
\{X \in \Ob(\mathcal{D}) \mid
H(X[n]) = 0\text{ for all }n \gg 0\} \\
\Ob(\mathcal{D}_H^b) =
\{X \in \Ob(\mathcal{D}) \mid
H(X[n]) = 0\text{ for all }|n| \gg 0\}
\end{matrix}
$$
である $\mathcal{D}$ の充満部分圏とする。これらはすべて
$\mathcal{D}$ の厳密充満な飽和前三角部分圏である。
$\mathcal{D}$ が三角圏ならば、これらはすべて三角部分圏である。
\end{lemma}

\begin{proof}
$\mathcal{D}_H^{+}$ について証明しよう。これが $[1]$ と $[-1]$ によって
保たれることは明らかである。$(X, Y, Z, f, g, h)$ を $\mathcal{D}$ の
完全三角とし、すべての $n \ll 0$ に対して $H(X[n]) = H(Y[n]) = 0$ とする。
このとき長完全列 (\ref{derived-equation-long-exact-cohomology-sequence}) により、
すべての $n \ll 0$ に対して $H(Z[n]) = 0$ でもある。したがって補題
\ref{derived-lemma-triangulated-subcategory} を適用し、$\mathcal{D}_H^{+}$ が
前三角部分圏（それぞれ、$\mathcal{D}$ が三角圏ならば三角部分圏）であることが
分かる。飽和性の主張は、すべての $n \in \mathbf{Z}$ に対する
\begin{align*}
H((X \oplus Y)[n]) = 0 & \Rightarrow H(X[n] \oplus Y[n]) = 0 \\
& \Rightarrow H(X[n]) \oplus H(Y[n]) = 0 \\
& \Rightarrow H(X[n]) = H(Y[n]) = 0
\end{align*}
から従う。
\end{proof}

\begin{definition}
\label{derived-definition-kernel-category}
$\mathcal{D}$ を（前）三角圏とする。
\begin{enumerate}
\item $F : \mathcal{D} \to \mathcal{D}'$ を完全関手とする。
{\it $F$ の核} とは、補題 \ref{derived-lemma-triangle-functor-kernel} で記述した
厳密充満な飽和（前）三角部分圏のことである。
\item $H : \mathcal{D} \to \mathcal{A}$ をホモロジー的関手とする。
{\it $H$ の核} とは、補題 \ref{derived-lemma-homological-functor-kernel} で記述した
厳密充満な飽和（前）三角部分圏のことである。
\end{enumerate}
これらを $\Ker(F)$ または $\Ker(H)$ と書くこともある。
\end{definition}

\noindent
次の補題の証明では TR4 を用いる。

\begin{lemma}
\label{derived-lemma-construct-multiplicative-system}
$\mathcal{D}$ を三角圏とし、
$\mathcal{D}' \subset \mathcal{D}$ を充満三角部分圏とする。
\begin{equation}
\label{derived-equation-multiplicative-system}
S =
\left\{
\begin{matrix}
f \in \text{Arrows}(\mathcal{D})
\text{ such that there exists a distinguished triangle }\\
(X, Y, Z, f, g, h) \text{ of }\mathcal{D}\text{ with }
Z\text{ isomorphic to an object of }\mathcal{D}'
\end{matrix}
\right\}
\end{equation}
と置く。このとき $S$ は $\mathcal{D}$ の三角構造と両立する乗法系である。
この状況で、次は同値である。
\begin{enumerate}
\item $S$ は飽和乗法系である。
\item $\mathcal{D}'$ は飽和三角部分圏である。
\end{enumerate}
\end{lemma}

\begin{proof}
最初の主張を証明するには、MS1、MS2、MS3、MS5、MS6 が成り立つことを
示さなければならない。

\medskip\noindent
MS1 の証明。$\mathcal{D}$ の各対象 $X$ に対して
$(X, X, 0, 1, 0, 0)$ は完全であり、$0$ は $\mathcal{D}'$ の対象なので、
恒等射が $S$ に属することは明らかである。$f : X \to Y$ と $g : Y \to Z$ を、
$S$ に属する合成可能な射とする。完全三角
$(X, Y, Q_1, f, p_1, d_1)$、
$(X, Z, Q_2, g \circ f, p_2, d_2)$、
$(Y, Z, Q_3, g, p_3, d_3)$
を選ぶ。仮定により、$Q_1$ と $Q_3$ は $\mathcal{D}'$ の対象と同型である。
TR4 により完全三角 $(Q_1, Q_2, Q_3, a, b, c)$ が存在する。
$\mathcal{D}'$ は三角部分圏なので、$Q_2$ は $\mathcal{D}'$ の対象と同型である。
したがって $g \circ f \in S$ である。

\medskip\noindent
MS3 の証明。$a : X \to Y$ を射とし、$t : Z \to X$ を
$a \circ t = 0$ を満たす $S$ の元とする。LMS3 を証明するには、
$s \circ a = 0$ を満たす $s \in S$ を見いだせば十分である。
補題 \ref{derived-lemma-triangle-functor-localize} の証明と比較されたい。
TR1 と TR2 により完全三角 $(Z, X, Q, t, g, h)$ を選ぶ。
$a \circ t = 0$ なので、補題 \ref{derived-lemma-representable-homological} により、
$i \circ g = a$ を満たす射 $i : Q \to Y$ が存在する。最後に、TR1 を再び用いて
三角 $(Q, Y, W, i, s, k)$ を選べる。図示すると
$$
\xymatrix{
Z \ar[r]_t & X \ar[r]_g \ar[d]^1 & Q \ar[r] \ar[d]^i & Z[1] \\
& X \ar[r]_a & Y \ar[d]^s \\
& & W
}
$$
となる。$t \in S$ なので、$Q$ は $\mathcal{D}'$ の対象と同型である。
したがって $s \in S$ である。最後に、補題
\ref{derived-lemma-composition-zero} により $s \circ i = 0$ なので、
$s \circ a = s \circ i \circ g = 0$ である。ゆえに LMS3 が成り立つ。
RMS3 の証明は双対である。

\medskip\noindent
MS5 の証明。完全三角および $\mathcal{D}'$ がシフトによって保たれることから
従う。

\medskip\noindent
MS6 の証明。$s, s' \in S$ を満たす可換図式
$$
\xymatrix{
X \ar[r] \ar[d]^s &
Y \ar[d]^{s'} \\
X' \ar[r] &
Y'
}
$$
が与えられたとする。命題 \ref{derived-proposition-9} により、これを九つの正方形から
なる図式へ拡張できる。$s, s'$ は $S$ の元なので、$X'', Y''$ は
$\mathcal{D}'$ の対象と同型である。$\mathcal{D}'$ は充満三角部分圏なので、
$Z''$ も $\mathcal{D}'$ の対象と同型である。したがって射 $Z \to Z'$ は
$S$ の元である。これで MS6 が証明された。

\medskip\noindent
MS2 は MS1、MS5、MS6 の形式的帰結である。補題
\ref{derived-lemma-localization-conditions} を参照されたい。
これで補題の最初の主張の証明が完了する。

\medskip\noindent
$S$ が飽和していると仮定しよう
（以下、完全三角の回転を断りなく用いる）。$X \oplus Y$ が
$\mathcal{D}'$ の対象と同型であるとする。射 $f : 0 \to X$ を考える。
$(0, X \oplus Y, X \oplus Y, 0, 1, 0)$ は完全三角なので、合成
$0 \to X \to X \oplus Y$ は $S$ の元である。また、三角
$(X, X \oplus Y, Y, (1, 0), (0, 1), 0)$ は完全なので、合成
$Y[-1] \to 0 \to X$ は $S$ の元である。補題 \ref{derived-lemma-split} を
参照されたい。したがって、$S$ が飽和していることから $0 \to X$ は
$S$ の元である。ゆえに、望みどおり $X$ は $\mathcal{D}'$ の対象と同型である。

\medskip\noindent
最後に、$\mathcal{D}'$ が飽和三角部分圏であると仮定する。
$$
W \xrightarrow{h}
X \xrightarrow{g}
Y \xrightarrow{f} Z
$$
を、$fg, gh \in S$ を満たす $\mathcal{D}$ の合成可能な射とする。
次の図式にある対象を順次構成する。
$$
\xymatrix{
 & &
Q_{12} \ar[rd] & &
Q_{23} \ar[rd] \\
 &
Q_1 \ar[ld]_{\! + \! 1} \ar[ru] & &
Q_2 \ar[ld]_{\! + \! 1} \ar[ll]_{\! + \! 1} \ar[ru] & &
Q_3 \ar[ld]_{\! + \! 1} \ar[ll]_{\! + \! 1} \\
W \ar[rr] & &
X \ar[lu] \ar[rr] & &
Y \ar[lu] \ar[rr] & &
Z \ar[lu]
}
$$
まず、完全三角
$(W, X, Q_1)$、$(X, Y, Q_2)$、$(Y, Z, Q_3)$ $(W, Y, Q_{12})$、
$(X, Z, Q_{23})$ を選ぶ。$s : Q_2 \to Q_1[1]$ を合成
$Q_2 \to X[1] \to Q_1[1]$ とし、$t : Q_3 \to Q_2[1]$ を合成
$Q_3 \to Y[1] \to Q_2[1]$ とする。合成 $W \to X \to Y$ および
$X \to Y \to Z$ に TR4 を適用すると、射 $s$ と $t$ を用いる完全三角
$(Q_1, Q_{12}, Q_2)$ と $(Q_2, Q_{23}, Q_3)$ が存在する。
$W \to Y$ と $X \to Z$ は $S$ に属すると仮定したので、$Q_{12}$ と
$Q_{23}$ は $\mathcal{D}'$ の対象と同型である。したがって $S$ は合成と
シフトで閉じているので、$s[1]t$ も $S$ の元である。
$Y[1] \to Q_2[1] \to X[2]$ は零なので、$s[1]t = 0$ であることに注意する。
補題 \ref{derived-lemma-composition-zero} を参照されたい。したがって補題
\ref{derived-lemma-split} により、$Q_3[1] \oplus Q_1[2]$ は
$\mathcal{D}'$ の対象と同型である。$\mathcal{D}'$ に関する仮定から、
$Q_3$ と $Q_1$ は $\mathcal{D}'$ の対象と同型である。完全三角
$(Q_1, Q_{12}, Q_2)$ を見ると、$Q_2$ も $\mathcal{D}'$ の対象と同型である。
完全三角 $(X, Y, Q_2)$ を見て、最後に $g \in S$ と結論する。
% P02-E0023: frozen source has the ungrammatical phrase “It is also follows”.
（$h, f \in S$ も従うが、これは必要ない。）
\end{proof}

\begin{definition}
\label{derived-definition-quotient-category}
$\mathcal{D}$ を三角圏、$\mathcal{B}$ を充満三角部分圏とする。
{\it 商圏 $\mathcal{D}/\mathcal{B}$} を
$\mathcal{D}/\mathcal{B} = S^{-1}\mathcal{D}$ によって定義する。
ここで $S$ は、補題 \ref{derived-lemma-construct-multiplicative-system} により
$\mathcal{B}$ に付随する $\mathcal{D}$ の乗法系である。この場合、
局所化関手 $Q : \mathcal{D} \to \mathcal{D}/\mathcal{B}$ を
{\it 商関手} と呼ぶ。
\end{definition}

\noindent
商関手 $Q : \mathcal{D} \to \mathcal{D}/\mathcal{B}$ は三角圏の
完全関手であることに注意する。命題 \ref{derived-proposition-construct-localization} を
参照されたい。この構成の普遍性は次のとおりである。

\begin{lemma}
\label{derived-lemma-universal-property-quotient}
\begin{slogan}
Verdier 商の普遍性。
\end{slogan}
$\mathcal{D}$ を三角圏、$\mathcal{B}$ を $\mathcal{D}$ の充満三角部分圏とし、
$Q : \mathcal{D} \to \mathcal{D}/\mathcal{B}$ を商関手とする。
\begin{enumerate}
\item $H : \mathcal{D} \to \mathcal{A}$ をアーベル圏 $\mathcal{A}$ への
ホモロジー的関手とし、$\mathcal{B} \subset \Ker(H)$ を満たすとする。
このとき、$H = H' \circ Q$ を満たす一意な分解
$H' : \mathcal{D}/\mathcal{B} \to \mathcal{A}$ が存在し、
$H'$ もホモロジー的関手である。
\item $F : \mathcal{D} \to \mathcal{D}'$ を前三角圏 $\mathcal{D}'$ への
完全関手とし、$\mathcal{B} \subset \Ker(F)$ を満たすとする。
このとき、$F = F' \circ Q$ を満たす一意な分解
$F' : \mathcal{D}/\mathcal{B} \to \mathcal{D}'$ が存在し、
$F'$ も完全関手である。
\end{enumerate}
\end{lemma}

\begin{proof}
この補題は補題 \ref{derived-lemma-universal-property-localization} から従う。
実際、$f : X \to Y$ を $\mathcal{D}$ の射とし、ある完全三角
$(X, Y, Z, f, g, h)$ に対して $Z$ が $\mathcal{B}$ の対象と同型であるとする。
このとき、(1) の仮定のもとでは $H(f)$ が、(2) の仮定のもとでは $F(f)$ が
それぞれ同型である。詳細は省略する。
\end{proof}

\noindent
商関手の核は次のように記述できる。

\begin{lemma}
\label{derived-lemma-kernel-quotient}
$\mathcal{D}$ を三角圏、$\mathcal{B}$ を充満三角部分圏とする。
商関手 $Q : \mathcal{D} \to \mathcal{D}/\mathcal{B}$ の核は、
対象が
$$
\Ob(\Ker(Q)) =
\left\{
\begin{matrix}
Z \in \Ob(\mathcal{D})
\text{ such that there exists a }Z' \in \Ob(\mathcal{D}) \\
\text{ such that }Z \oplus Z'\text{ is isomorphic to an object of }\mathcal{B}
\end{matrix}
\right\}
$$
である $\mathcal{D}$ の厳密充満部分圏である。言い換えれば、これは
$\mathcal{B}$ を含む $\mathcal{D}$ の最小の厳密充満な飽和三角部分圏である。
\end{lemma}

\begin{proof}
まず、核は自動的に、各対象の直和因子を含む厳密充満三角部分圏である。
補題 \ref{derived-lemma-triangle-functor-kernel} を参照されたい。その対象の記述は、
定義と補題 \ref{derived-lemma-kernel-localization} の (4) から従う。
\end{proof}

\noindent
$\mathcal{D}$ を三角圏とする。この時点で、次の二つの半順序集合の間に
順序保存写像を誘導する構成が得られている。
\begin{enumerate}
\item $\mathcal{D}$ の三角構造と両立する乗法系 $S$ の半順序集合。
\item 充満三角部分圏 $\mathcal{B} \subset \mathcal{D}$ の半順序集合。
\end{enumerate}
具体的には、構成は
$S \mapsto \mathcal{B}(S) = \Ker(Q : \mathcal{D} \to S^{-1}\mathcal{D})$
および $\mathcal{B} \mapsto S(\mathcal{B})$ によって与えられる。
ここで $S(\mathcal{B})$ は (\ref{derived-equation-multiplicative-system}) の乗法系、
すなわち
$$
S(\mathcal{B}) =
\left\{
\begin{matrix}
f \in \text{Arrows}(\mathcal{D})
\text{ such that there exists a distinguished triangle }\\
(X, Y, Z, f, g, h) \text{ of }\mathcal{D}\text{ with }
Z\text{ isomorphic to an object of }\mathcal{B}
\end{matrix}
\right\}
$$
である。これらの操作は互いに逆ではないことに注意する。

\begin{lemma}
\label{derived-lemma-operations}
$\mathcal{D}$ を三角圏とする。上で述べた操作は次の性質をもつ。
\begin{enumerate}
\item $S(\mathcal{B}(S))$ は $S$ の「飽和」、すなわち $S$ を含む
$\mathcal{D}$ の最小の飽和乗法系である。
\item $\mathcal{B}(S(\mathcal{B}))$ は $\mathcal{B}$ の「飽和」、
すなわち $\mathcal{B}$ を含む $\mathcal{D}$ の最小の厳密充満な
飽和三角部分圏である。
\end{enumerate}
特に、これらの構成は、$\mathcal{D}$ の飽和乗法系であって
$\mathcal{D}$ の三角構造と両立するものの
（半順序）集合と、$\mathcal{D}$ の厳密充満な飽和三角部分圏の（半順序）集合との
間の互いに逆な写像を定める。
\end{lemma}

\begin{proof}
まず、充満三角部分圏 $\mathcal{B}$ から始める。このとき
$\mathcal{B}(S(\mathcal{B})) =
\Ker(Q : \mathcal{D} \to \mathcal{D}/\mathcal{B})$
であり、したがって (2) は補題 \ref{derived-lemma-kernel-quotient} の内容である。

\medskip\noindent
次に、$S$ を $\mathcal{D}$ の乗法系であって、$\mathcal{D}$ の
三角構造と両立するものとする。このとき
$\mathcal{B}(S) = \Ker(Q : \mathcal{D} \to S^{-1}\mathcal{D})$ である。
したがって、局所化された圏において補題 \ref{derived-lemma-third-object-zero} を用いると、
\begin{align*}
S(\mathcal{B}(S))
& =
\left\{
\begin{matrix}
f \in \text{Arrows}(\mathcal{D})
\text{ such that there exists a distinguished}\\
\text{triangle }(X, Y, Z, f, g, h) \text{ of }\mathcal{D}\text{ with }Q(Z) = 0
\end{matrix}
\right\}.
\\
& =
\{f \in \text{Arrows}(\mathcal{D}) \mid Q(f)\text{ is an isomorphism}\} \\
& = \hat S = S'
\end{align*}
を得る。Categories, Lemma \ref{categories-lemma-what-gets-inverted} の
記号を用いた。その補題の最後の主張によって証明が完了する。
\end{proof}

\begin{lemma}
\label{derived-lemma-acyclic-general}
$H : \mathcal{D} \to \mathcal{A}$ を、三角圏 $\mathcal{D}$ から
アーベル圏 $\mathcal{A}$ へのホモロジー的関手とする。定義
\ref{derived-definition-homological} を参照されたい。$\mathcal{D}$ の部分圏
$\Ker(H)$ は $\mathcal{D}$ の厳密充満な飽和三角部分圏であり、それに対応する
飽和乗法系（補題 \ref{derived-lemma-operations} を参照）は
$$
S = \{f \in \text{Arrows}(\mathcal{D}) \mid
H^i(f)\text{ is an isomorphism for all }i \in \mathbf{Z}\}.
$$
である。関手 $H$ は商関手
$Q : \mathcal{D} \to \mathcal{D}/\Ker(H)$ を経由して分解する。
\end{lemma}

\begin{proof}
補題 \ref{derived-lemma-homological-functor-kernel} により、圏 $\Ker(H)$ は
$\mathcal{D}$ の厳密充満な飽和三角部分圏である。補題
\ref{derived-lemma-homological-functor-localize} により、集合 $S$ は
三角構造と両立する飽和乗法系である。$\Ker(H)$ に対応する乗法系は集合
$$
\left\{
\begin{matrix}
f \in \text{Arrows}(\mathcal{D})
\text{ such that there exists a distinguished triangle }\\
(X, Y, Z, f, g, h)\text{ with } H^i(Z) = 0 \text{ for all }i
\end{matrix}
\right\}.
$$
であった。長完全コホモロジー列
(\ref{derived-equation-long-exact-cohomology-sequence}) により、$f$ がこの集合の元で
あるための必要十分条件は $f$ が $S$ の元であることだと明らかに分かる。
最後に、$H$ が $Q$ を経由して分解することは
補題 \ref{derived-lemma-universal-property-quotient} の帰結である。
\end{proof}
\section{完全関手の随伴}
\label{derived-section-adjoints}

\noindent
三角圏の間の随伴関手に関する結果を述べる。

\begin{lemma}
\label{derived-lemma-adjoint-is-exact}
$F : \mathcal{D} \to \mathcal{D}'$ を三角圏の間の完全関手とする。
$F$ が右随伴 $G: \mathcal{D'} \to \mathcal{D}$ をもつならば、
$G$ も完全関手である。
\end{lemma}

\begin{proof}
$\xi_X : F(X[1]) \to F(X)[1]$ を定義
\ref{derived-definition-exact-functor-triangulated-categories} のものとする。
$\epsilon_A : F(G(A)) \to A$ を随伴射とする。合成
$$
F(G(A)[1]) \xrightarrow{\xi_{G(A)}} F(G(A))[1] \xrightarrow{\epsilon_A[1]} A[1]
$$
を考える。この写像に随伴する写像 $G(A)[1] \to G(A[1])$ は同型であると
主張する。これを見るには、米田の補題により、$\mathcal{D}$ の各対象 $X$ に対して
$\Mor_\mathcal{D}(X[1], -)$ を適用すると全単射が得られることを示せば十分である。
実際、任意の射 $f : X[1] \to G(A)[1]$ は一意な
$g : X \to G(A)$ によって $g[1]$ と書かれ、任意の $g$ は一意な
$h : F(X) \to A$ に随伴し、さらに $h[1] \circ \xi_X$ は一意な
$i : X[1] \to G(A[1])$ に随伴する。$f$ を $i$ に送る規則は
$\Mor_\mathcal{D}(X[1], G(A)[1])$ と
$\Mor_\mathcal{D}(X[1], G(A[1]))$ の間の全単射を与える。
この全単射が上の射によって与えられることは、Categories, Section
\ref{categories-section-adjoint} の議論と次の計算から従う。
\begin{align*}
\epsilon_A[1] \circ \xi_{G(A)} \circ F(f)
& =
\epsilon_A[1] \circ \xi_{G(A)} \circ F(g[1]) \\
& =
\epsilon_A[1] \circ F(g)[1] \circ \xi_X \\
& = 
(\epsilon_A \circ F(g))[1] \circ \xi_X \\
& =
h[1] \circ \xi_X
\end{align*}
詳細の一部は省略する。同型 $G(A)[1] \to G(A[1])$ の逆を用いて、
$G$ が完全関手であることを示す。これらの同型は $A$ に関して関手的である
（詳細は省略する）。

\medskip\noindent
$A \to B \to C \to A[1]$ を $\mathcal{D}'$ の完全三角とする。
$\mathcal{D}$ の完全三角
$$
G(A) \to G(B) \to X \to G(A)[1]
$$
を選ぶ。このとき
$F(G(A)) \to F(G(B)) \to F(X) \to F(G(A))[1]$
は $\mathcal{D}'$ の完全三角である。TR3 により、完全三角の射
$$
\xymatrix{
F(G(A)) \ar[r] \ar[d]^{\epsilon_A} &
F(G(B)) \ar[r] \ar[d]^{\epsilon_B} &
F(X) \ar[r] \ar[d] &
F(G(A))[1] \ar[d]^{\epsilon_A[1]} \\
A \ar[r] & B \ar[r] & C \ar[r] & A[1]
}
$$
を選べる。射 $F(X) \to F(G(A))[1]$ は、$F(X) \to F(G(A)[1])$ と
$\xi_{G(A)}$ の合成であることを思い出そう。したがって、$G$ が $F$ の
右随伴であることと上の議論から、図式
$$
\xymatrix{
G(A) \ar[r] \ar[d] & G(B) \ar[r] \ar[d] & X \ar[r] \ar[d] & G(A)[1] \ar[d] \\
G(A) \ar[r] & G(B) \ar[r] & G(C) \ar[r] & G(A[1])
}
$$
が可換となるような射 $X \to G(C)$ が存在し、右端の縦の射は上で構成した
射である。
% P02-E0024: frozen source uses D' for W and both Hom groups, although X and G(C) lie in D.
対象 $W$ が $\mathcal{D}'$ に属するとしてホモロジー的関手
$\Hom_{\mathcal{D}'}(W, -)$ を適用すると、$5$ 項補題から
$$
\Hom_{\mathcal{D}'}(W, X) \to \Hom_{\mathcal{D}'}(W, G(C))
$$
が全単射であることが従う。米田の補題をもう一度用いて、$X \to G(C)$ が
同型であると結論する。したがって
$G(A) \to G(B) \to G(C) \to G(A)[1]$ は完全三角であり、これが示すべきことで
あった。
\end{proof}

\begin{lemma}
\label{derived-lemma-fully-faithful-adjoint-kernel-zero}
$\mathcal{D}$、$\mathcal{D}'$ を三角圏とし、
$F : \mathcal{D} \to \mathcal{D}'$ および
$G : \mathcal{D}' \to \mathcal{D}$ を関手とする。次を仮定する。
\begin{enumerate}
\item $F$ と $G$ は完全関手である。
\item $F$ は充満忠実である。
\item $G$ は $F$ の右随伴である。
\item $G$ の核は零である。
\end{enumerate}
このとき $F$ は圏同値である。
\end{lemma}

\begin{proof}
$F$ は充満忠実なので、随伴射 $\text{id} \to G \circ F$ は同型である
（Categories, Lemma \ref{categories-lemma-adjoint-fully-faithful}）。
$X$ を $\mathcal{D}'$ の対象とする。$\mathcal{D}'$ の完全三角
$$
F(G(X)) \to X \to Y \to F(G(X))[1]
$$
を選ぶ。$G$ を適用し、$G(F(G(X))) = G(X)$ を用いると、完全三角
$$
G(X) \to G(X) \to G(Y) \to G(X)[1]
$$
を得る。したがって $G(Y) = 0$ であり、ゆえに $Y = 0$ である。
したがって $F(G(X)) \to X$ は同型である。
\end{proof}


\section{ホモトピー圏}
\label{derived-section-homotopy}

\noindent
$\mathcal{A}$ を加法圏とする。$\mathcal{A}$ のホモトピー圏
$K(\mathcal{A})$ とは、$\mathcal{A}$ の複体を対象とし、複体の射の
ホモトピー類を射とする圏である。形式的な定義は次のとおりである。

\begin{definition}
\label{derived-definition-complexes-notation}
$\mathcal{A}$ を加法圏とする。
\begin{enumerate}
\item $\text{Comp}(\mathcal{A}) = \text{CoCh}(\mathcal{A})$ を
{\it （余鎖）複体の圏} とする。
\item 複体 $K^\bullet$ が {\it 下に有界} であるとは、すべての $n \ll 0$ に
対して $K^n = 0$ であることをいう。
\item 複体 $K^\bullet$ が {\it 上に有界} であるとは、すべての $n \gg 0$ に
対して $K^n = 0$ であることをいう。
\item 複体 $K^\bullet$ が {\it 有界} であるとは、すべての $|n| \gg 0$ に
対して $K^n = 0$ であることをいう。
\item $\text{Comp}^{+}(\mathcal{A})$、$\text{Comp}^{-}(\mathcal{A})$、
それぞれ $\text{Comp}^b(\mathcal{A})$ を、
$\text{Comp}(\mathcal{A})$ の充満部分圏で、対象がそれぞれ下に有界、上に有界、
有界な複体であるものとする。
\item $K(\mathcal{A})$ を、$\text{Comp}(\mathcal{A})$ と同じ対象をもち、
複体の写像のホモトピー類を射とする圏とする
（Homology, Lemma \ref{homology-lemma-compose-homotopy-cochain} を参照）。
\item $K^{+}(\mathcal{A})$、$K^{-}(\mathcal{A})$、
それぞれ $K^b(\mathcal{A})$ を、$K(\mathcal{A})$ の充満部分圏で、
対象がそれぞれ $\mathcal{A}$ の下に有界、上に有界、有界な複体であるものとする。
\end{enumerate}
\end{definition}

\noindent
圏 $K(\mathcal{A})$、$K^{+}(\mathcal{A})$、$K^{-}(\mathcal{A})$、
$K^b(\mathcal{A})$ は三角圏になることが分かる。その証明のため、まず錐と
分裂完全列に関する道具立てを整える。

\section{錐と項ごとに分裂する列}
\label{derived-section-cones}

\noindent
$\mathcal{A}$ を加法圏とし、$K(\mathcal{A})$ を $\mathcal{A}$ の複体を対象、
ホモトピーを法とする複体の射を射とする圏とする。複体上のシフト関手
$[n]$（Homology, Definition \ref{homology-definition-shift-cochain} を参照）は、
$[n] \circ [m] = [n + m]$ かつ $[0] = \text{id}$ を満たす関手
$[n] : K(\mathcal{A}) \to K(\mathcal{A})$ を誘導することに注意する。

\begin{definition}
\label{derived-definition-cone}
$\mathcal{A}$ を加法圏とする。$f : K^\bullet \to L^\bullet$ を
$\mathcal{A}$ の複体の射とする。$f$ の {\it 錐} とは、
$C(f)^n = L^n \oplus K^{n + 1}$ であり、微分が
$$
d_{C(f)}^n =
\left(
\begin{matrix}
d^n_L & f^{n + 1} \\
0 & -d_K^{n + 1}
\end{matrix}
\right)
$$
で与えられる複体 $C(f)^\bullet$ のことである。これには、自明な写像
$L^n \to C(f)^n \to K^{n + 1}$ から誘導される標準的な複体の射
$i : L^\bullet \to C(f)^\bullet$ および
$p : C(f)^\bullet \to K^\bullet[1]$ が備わっている。
\end{definition}

\noindent
言い換えれば、$(K, L, C(f), f, i, p)$ は三角
$$
K^\bullet \to L^\bullet \to C(f)^\bullet \to K^\bullet[1]
$$
をなす。この三角の構成は、次の意味で関手的である。

\begin{lemma}
\label{derived-lemma-functorial-cone}
$$
\xymatrix{
K_1^\bullet \ar[r]_{f_1} \ar[d]_a & L_1^\bullet \ar[d]^b \\
K_2^\bullet \ar[r]^{f_2} & L_2^\bullet
}
$$
が、ホモトピーの意味で可換な複体の射の図式であるとする。このとき、
$K(\mathcal{A})$ における三角の射
$$
(a, b, c) : (K_1^\bullet, L_1^\bullet, C(f_1)^\bullet, f_1, i_1, p_1)
\to
(K_2^\bullet, L_2^\bullet, C(f_2)^\bullet, f_2, i_2, p_2)
$$
を与える射 $c : C(f_1)^\bullet \to C(f_2)^\bullet$ が存在する。
\end{lemma}

\begin{proof}
$b \circ f_1 - f_2 \circ a= d \circ h + h \circ d$ を満たす射の族
$h^n : K_1^n \to L_2^{n - 1}$ を取る。$c^n$ を行列
$$
c^n =
\left(
\begin{matrix}
b^n & h^{n + 1} \\
0 & a^{n + 1}
\end{matrix}
\right) :
L_1^n \oplus K_1^{n + 1} \to L_2^n \oplus K_2^{n + 1}
$$
によって定める。
% P02-E0025: frozen source has “A matrix computation show”.
行列計算により $c$ は複体の射である。$c \circ i_1 = i_2 \circ b$ は自明であり、
$p_2 \circ c = a \circ p_1$ も直ちに確かめられる。
\end{proof}

\noindent
補題 \ref{derived-lemma-functorial-cone} の証明で構成した射
$c : C(f_1)^\bullet \to C(f_2)^\bullet$ は、一般には
$f_2 \circ a$ と $b \circ f_1$ の間に選んだホモトピー $h$ に依存する。

\begin{lemma}
\label{derived-lemma-map-from-cone}
$f: K^\bullet \to L^\bullet$ および $g : L^\bullet \to M^\bullet$ を、
$g \circ f$ が零とホモトピックであるような複体の射とする。このとき
\begin{enumerate}
\item $g$ は射 $C(f)^\bullet \to M^\bullet$ を経由して分解し、
\item $f$ は射 $K^\bullet \to C(g)^\bullet[-1]$ を経由して分解する。
\end{enumerate}
\end{lemma}

\begin{proof}
仮定は、図式
$$
\xymatrix{
K^\bullet \ar[r]_f \ar[d] & L^\bullet \ar[d]^g \\
0 \ar[r] & M^\bullet
}
$$
がホモトピーの意味で可換であることをいう。$0 \to M^\bullet$ の錐は
$M^\bullet$ なので、補題 \ref{derived-lemma-functorial-cone} の写像
$C(f)^\bullet \to C(0 \to M^\bullet) = M^\bullet$ が (1) の写像である。
$K^\bullet \to 0$ の錐は $K^\bullet[1]$ であり、補題 \ref{derived-lemma-functorial-cone} を適用すると
$K^\bullet[1] \to C(g)^\bullet$ が得られる。$[-1]$ を適用して (2) の写像を得る。
\end{proof}

\noindent
補題 \ref{derived-lemma-map-from-cone} の証明で構成した射
$C(f)^\bullet \to M^\bullet$ および $K^\bullet \to C(g)^\bullet[-1]$ は、
一般には選んだホモトピーに依存することに注意する。

\begin{definition}
\label{derived-definition-termwise-split-map}
$\mathcal{A}$ を加法圏とする。{\it 項ごとに分裂する単射
$\alpha : A^\bullet \to B^\bullet$} とは、各 $A^n \to B^n$ が
直和因子の包含と同型であるような複体の射のことである。
{\it 項ごとに分裂する全射 $\beta : B^\bullet \to C^\bullet$} とは、
各 $B^n \to C^n$ が直和因子への射影と同型であるような複体の射のことである。
\end{definition}

\begin{lemma}
\label{derived-lemma-make-commute-map}
$\mathcal{A}$ を加法圏とする。
$$
\xymatrix{
A^\bullet \ar[r]_f \ar[d]_a & B^\bullet \ar[d]^b \\
C^\bullet \ar[r]^g & D^\bullet
}
$$
を、ホモトピーの意味で可換な複体の射の図式とする。$f$ が項ごとに分裂する
単射ならば、$b$ はこの図式を可換にする射とホモトピックである。$g$ が
項ごとに分裂する全射ならば、$a$ はこの図式を可換にする射とホモトピックである。
\end{lemma}

\begin{proof}
$bf - ga = dh + hd$ を満たす射の族 $h^n : A^n \to D^{n - 1}$ を取る。
$\pi^n : B^n \to A^n$ を $f^n$ の分裂射とする。
$b' = b - dh\pi - h\pi d$ と置く。$s^n : D^n \to C^n$ を
$g^n : C^n \to D^n$ の分裂射とする。$a' = a + dsh + shd$ と置く。
計算は省略する。
\end{proof}

\noindent
次の補題を用いると、複体の射を、各次数で直和因子の包含となる射に
置き換えることができる。

\begin{lemma}
\label{derived-lemma-make-injective}
$\mathcal{A}$ を加法圏とする。$\alpha : K^\bullet \to L^\bullet$ を
$\mathcal{A}$ の複体の射とする。このとき、分解
$$
\xymatrix{
K^\bullet \ar[r]^{\tilde \alpha} \ar@/_1pc/[rr]_\alpha &
\tilde L^\bullet \ar[r]^\pi &
L^\bullet
}
$$
であって、次を満たすものが存在する。
\begin{enumerate}
\item $\tilde \alpha$ は項ごとに分裂する単射である
（Definition \ref{derived-definition-termwise-split-map} を参照）。
\item $\pi \circ s = \text{id}_{L^\bullet}$ を満たし、かつ
$s \circ \pi$ が $\text{id}_{\tilde L^\bullet}$ とホモトピックであるような
複体の射 $s : L^\bullet \to \tilde L^\bullet$ が存在する。
\end{enumerate}
さらに、$K^\bullet$ と $L^\bullet$ がともに $K^{+}(\mathcal{A})$、
$K^{-}(\mathcal{A})$、または $K^b(\mathcal{A})$ に属するならば、
$\tilde L^\bullet$ も同じ圏に属する。
\end{lemma}

\begin{proof}
$$
\tilde L^n = L^n \oplus K^n \oplus K^{n + 1}
$$
と置き、
$$
d^n_{\tilde L} =
\left(
\begin{matrix}
d^n_L & 0 & 0 \\
0 & d^n_K & \text{id}_{K^{n + 1}} \\
0 & 0 & -d^{n + 1}_K
\end{matrix}
\right)
$$
と定める。言い換えれば、
$\tilde L^\bullet = L^\bullet \oplus C(1_{K^\bullet})$ である。さらに、
$$
\tilde \alpha =
\left(
\begin{matrix}
\alpha \\
\text{id}_{K^n} \\
0
\end{matrix}
\right)
$$
と置く。これは明らかに分裂単射であり、複体の射を定めることも明らかである。
また、
$$
\pi =
\left(
\begin{matrix}
\text{id}_{L^n} &
0 &
0
\end{matrix}
\right)
$$
と定めれば、明らかに $\pi \circ \tilde \alpha = \alpha$ である。さらに、
$$
s =
\left(
\begin{matrix}
\text{id}_{L^n} \\
0 \\
0
\end{matrix}
\right)
$$
と置けば、$\pi \circ s = \text{id}_{L^\bullet}$ である。最後に、
$h^n : \tilde L^n \to \tilde L^{n - 1}$ を、$\tilde L^n$ の直和因子 $K^n$ を
恒等射によって $\tilde L^{n - 1}$ の直和因子 $K^n$ に写す射とする。このとき
$$
\text{id}_{\tilde L^\bullet} - s \circ \pi
=
d \circ h + h \circ d
$$
であることは直ちに確かめられる（下の remark の計算を参照）。これで補題の証明が完了する。
\end{proof}

\begin{remark}
\label{derived-remark-compute-modules}
上の証明の最後に表示した等式は、元を用いて次のように確かめられる。
$s\pi(l, k, k^{+}) = (l, 0, 0)$ である。したがって、左辺の射は
$(l, k, k^{+})$ を $(0, k, k^{+})$ に写す。一方、
$h(l, k, k^{+}) = (0, 0, k)$ および
$d(l, k, k^{+}) = (dl, dk + k^{+}, -dk^{+})$ である。したがって、望みどおり
$$(dh + hd)(l, k, k^{+}) =
d(0, 0, k) + h(dl, dk + k^{+}, -dk^{+}) =
(0, k, -dk) + (0, 0, dk + k^{+}) = (0, k, k^{+})$$
を得る。
\end{remark}

\begin{lemma}
\label{derived-lemma-make-surjective}
$\mathcal{A}$ を加法圏とする。$\alpha : K^\bullet \to L^\bullet$ を
$\mathcal{A}$ の複体の射とする。このとき、分解
$$
\xymatrix{
K^\bullet \ar[r]^i \ar@/_1pc/[rr]_\alpha &
\tilde K^\bullet \ar[r]^{\tilde \alpha} &
L^\bullet
}
$$
であって、次を満たすものが存在する。
\begin{enumerate}
\item $\tilde \alpha$ は項ごとに分裂する全射である
（Definition \ref{derived-definition-termwise-split-map} を参照）。
\item $s \circ i = \text{id}_{K^\bullet}$ を満たし、かつ
$i \circ s$ が $\text{id}_{\tilde K^\bullet}$ とホモトピックであるような
複体の射 $s : \tilde K^\bullet \to K^\bullet$ が存在する。
\end{enumerate}
さらに、$K^\bullet$ と $L^\bullet$ がともに $K^{+}(\mathcal{A})$、
$K^{-}(\mathcal{A})$、または $K^b(\mathcal{A})$ に属するならば、
$\tilde K^\bullet$ も同じ圏に属する。
\end{lemma}

\begin{proof}
補題 \ref{derived-lemma-make-injective} の双対である。
$$
\tilde K^n = K^n \oplus L^{n - 1} \oplus L^n
$$
と置き、
$$
d^n_{\tilde K} =
\left(
\begin{matrix}
d^n_K & 0 & 0 \\
0 & - d^{n - 1}_L & \text{id}_{L^n} \\
0 & 0 & d^n_L
\end{matrix}
\right)
$$
と定める。言い換えれば、
$\tilde K^\bullet = K^\bullet \oplus C(1_{L^\bullet[-1]})$ である。さらに、
$$
\tilde \alpha =
\left(
\begin{matrix}
\alpha &
0 &
\text{id}_{L^n}
\end{matrix}
\right)
$$
と置く。これは明らかに分裂全射であり、複体の射を定めることも明らかである。
また、
$$
i =
\left(
\begin{matrix}
\text{id}_{K^n} \\
0 \\
0
\end{matrix}
\right)
$$
と定めれば、明らかに $\tilde \alpha \circ i = \alpha$ である。さらに、
$$
s =
\left(
\begin{matrix}
\text{id}_{K^n} &
0 &
0
\end{matrix}
\right)
$$
と置けば、$s \circ i = \text{id}_{K^\bullet}$ である。最後に、
$h^n : \tilde K^n \to \tilde K^{n - 1}$ を、$\tilde K^n$ の直和因子
$L^{n - 1}$ を恒等射によって $\tilde K^{n - 1}$ の直和因子
$L^{n - 1}$ に写す射とする。このとき
$$
\text{id}_{\tilde K^\bullet} - i \circ s
=
d \circ h + h \circ d
$$
であることは直ちに確かめられる。これで補題の証明が完了する。
\end{proof}

\begin{definition}
\label{derived-definition-split-ses}
$\mathcal{A}$ を加法圏とする。$\mathcal{A}$ の複体の
{\it 項ごとに分裂する完全列}とは、複体を項とする複体
$$
0 \to
A^\bullet \xrightarrow{\alpha}
B^\bullet \xrightarrow{\beta}
C^\bullet \to 0
$$
と、$\alpha^n$ および $\beta^n$ と両立する指定された直和分解
$B^n = A^n \oplus C^n$ との組のことである。直和分解から誘導される写像を、
しばしば $s^n : C^n \to B^n$ および $\pi^n : B^n \to A^n$ と書く。
Homology, Lemma \ref{homology-lemma-ses-termwise-split-cochain} により、
対応する複体の射
$$
\delta : C^\bullet \longrightarrow A^\bullet[1]
$$
が得られ、その次数 $n$ における写像は
$\pi^{n + 1} \circ d_B^n \circ s^n$ である。言い換えれば、
$(A^\bullet, B^\bullet, C^\bullet, \alpha, \beta, \delta)$ は三角
$$
A^\bullet \to B^\bullet \to C^\bullet \to A^\bullet[1]
$$
をなす。これを {\it 項ごとに分裂する複体の列に付随する三角} と呼ぶ。
\end{definition}

\begin{lemma}
\label{derived-lemma-triangle-independent-splittings}
$\mathcal{A}$ を加法圏とする。
$0 \to A^\bullet \to B^\bullet \to C^\bullet \to 0$ を、
Definition \ref{derived-definition-split-ses} の意味で項ごとに分裂する完全列とする。
$(\pi')^n$, $(s')^n$ を第二の分裂射の族とし、この第二の分裂の組に付随する
射を $\delta' : C^\bullet \longrightarrow A^\bullet[1]$ と書く。このとき
$$
(1, 1, 1) :
(A^\bullet, B^\bullet, C^\bullet, \alpha, \beta, \delta)
\longrightarrow
(A^\bullet, B^\bullet, C^\bullet, \alpha, \beta, \delta')
$$
は $K(\mathcal{A})$ における三角の同型である。
\end{lemma}

\begin{proof}
主張は単に、$\delta$ と $\delta'$ がホモトピックな複体の写像であることをいう。
これは Homology, Lemma
\ref{homology-lemma-ses-termwise-split-homotopy-cochain} である。
\end{proof}

\begin{remark}
\label{derived-remark-make-commute}
$\mathcal{A}$ を加法圏とする。
$0 \to A_i^\bullet \to B_i^\bullet \to C_i^\bullet \to 0$、$i = 1, 2$ を、
項ごとに分裂する完全列とする。また、$a : A_1^\bullet \to A_2^\bullet$、
$b : B_1^\bullet \to B_2^\bullet$、$c : C_1^\bullet \to C_2^\bullet$ を、
図式
$$
\xymatrix{
A_1^\bullet \ar[d]_a \ar[r] &
B_1^\bullet \ar[r] \ar[d]_b &
C_1^\bullet \ar[d]_c \\
A_2^\bullet \ar[r] & B_2^\bullet \ar[r] & C_2^\bullet
}
$$
が $K(\mathcal{A})$ で可換となるような複体の射とする。一般には、上の図式を
複体の圏で可換にする、$b$ とホモトピックな射
$b' : B_1^\bullet \to B_2^\bullet$ は {\bf 存在しない}。
実際、Examples, Equation (\ref{examples-equation-commutes-up-to-homotopy})
を考える。そこで中央の写像を、図式を可換にするホモトピックな写像で
置き換えられたとすれば、成り立たないはずの跡の加法性が成り立ってしまう。
\end{remark}

\begin{lemma}
\label{derived-lemma-nilpotent}
$\mathcal{A}$ を加法圏とする。
$0 \to A_i^\bullet \to B_i^\bullet \to C_i^\bullet \to 0$、
$i = 1, 2, 3$ を、項ごとに分裂する複体の完全列とする。
$b : B_1^\bullet \to B_2^\bullet$ および
$b' : B_2^\bullet \to B_3^\bullet$ を、図式
$$
\vcenter{
\xymatrix{
A_1^\bullet \ar[d]_0 \ar[r] &
B_1^\bullet \ar[r] \ar[d]_b &
C_1^\bullet \ar[d]_0 \\
A_2^\bullet \ar[r] & B_2^\bullet \ar[r] & C_2^\bullet
}
}
\quad\text{and}\quad
\vcenter{
\xymatrix{
A_2^\bullet \ar[d]^0 \ar[r] &
B_2^\bullet \ar[r] \ar[d]^{b'} &
C_2^\bullet \ar[d]^0 \\
A_3^\bullet \ar[r] & B_3^\bullet \ar[r] & C_3^\bullet
}
}
$$
が $K(\mathcal{A})$ で可換となるような複体の射とする。このとき
$b' \circ b = 0$ が $K(\mathcal{A})$ で成り立つ。
\end{lemma}

\begin{proof}
補題 \ref{derived-lemma-make-commute-map} により、$b$ と $b'$ をホモトピックな写像で
置き換えて、左の図式の右の正方形と右の図式の左の正方形を可換にできる。
言い換えれば、
$\Im(b^n) \subset \Im(A_2^n \to B_2^n)$ および
$\Ker((b')^n) \supset \Im(A_2^n \to B_2^n)$ である。
したがって、複体の写像として $b' \circ b = 0$ である。
\end{proof}

\begin{lemma}
\label{derived-lemma-third-isomorphism}
$\mathcal{A}$ を加法圏とする。
$f_1 : K_1^\bullet \to L_1^\bullet$ および
$f_2 : K_2^\bullet \to L_2^\bullet$ を複体の射とする。
$$
(a, b, c) :
(K_1^\bullet, L_1^\bullet, C(f_1)^\bullet, f_1, i_1, p_1)
\longrightarrow
(K_2^\bullet, L_2^\bullet, C(f_2)^\bullet, f_2, i_2, p_2)
$$
を $K(\mathcal{A})$ の三角の任意の射とする。$a$ と $b$ が
ホモトピー同値ならば、$c$ もホモトピー同値である。
\end{lemma}

\begin{proof}
$a^{-1} : K_2^\bullet \to K_1^\bullet$ を、$K(\mathcal{A})$ で $a$ の逆となる
複体の射とする。同様に、$b^{-1} : L_2^\bullet \to L_1^\bullet$ を、
$K(\mathcal{A})$ で $b$ の逆となる複体の射とする。
$c' : C(f_2)^\bullet \to C(f_1)^\bullet$ を、等式
$f_1 \circ a^{-1} = b^{-1} \circ f_2$ に補題
\ref{derived-lemma-functorial-cone} を適用して得られる射とする。
$c \circ c'$ と $c' \circ c$ が $K(\mathcal{A})$ の同型であることを示せばよい。
したがって、次を示せば十分である。$K(\mathcal{A})$ に三角の射
$$
(1, 1, c) : (K^\bullet, L^\bullet, C(f)^\bullet, f, i, p)
$$
が与えられれば、$c$ は $K(\mathcal{A})$ の同型である。仮定により、図式
$$
\xymatrix{
L^\bullet \ar[r] \ar[d]_1 &
C(f)^\bullet \ar[r] \ar[d]_c &
K^\bullet[1] \ar[d]_1 \\
L^\bullet \ar[r] &
C(f)^\bullet \ar[r] &
K^\bullet[1]
}
$$
の二つの正方形はホモトピーの意味で可換である。$C(f)^\bullet$ の構成により、
各行は項ごとに分裂する複体の列をなす。したがって補題
\ref{derived-lemma-nilpotent} により、$K(\mathcal{A})$ で
$(c - 1)^2 = 0$ である。ゆえに、$c$ は $2 - c$ を逆射とする
$K(\mathcal{A})$ の同型である。
\end{proof}

\noindent
したがって、$a$ と $b$ がホモトピー同値ならば、それらから得られる三角の射は
$K(\mathcal{A})$ における三角の同型である。実際、$K(\mathcal{A})$ において
錐から得られる三角の族と、$K(\mathcal{A})$ において項ごとに分裂する
複体の列から得られる三角の族とは、
少なくとも符号を除けば、同型を除いて一致する。

\begin{lemma}
\label{derived-lemma-the-same-up-to-isomorphisms}
$\mathcal{A}$ を加法圏とする。
\begin{enumerate}
\item 項ごとに分裂する複体の列
$(\alpha : A^\bullet \to B^\bullet,
\beta : B^\bullet \to C^\bullet, s^n, \pi^n)$
が与えられたとする。このとき、図式
$$
\xymatrix{
A^\bullet \ar[r] \ar[d] & B^\bullet \ar[d] \ar[r] &
C(\alpha)^\bullet \ar[r]_{-p} \ar[d] & A^\bullet[1] \ar[d] \\
A^\bullet \ar[r] & B^\bullet \ar[r] &
C^\bullet \ar[r]^\delta & A^\bullet[1]
}
$$
が $K(\mathcal{A})$ における三角の同型を定めるようなホモトピー同値
$C(\alpha)^\bullet \to C^\bullet$ が存在する。
\item 複体の射 $f : K^\bullet \to L^\bullet$ が与えられたとする。このとき
三角の同型
$$
\xymatrix{
K^\bullet \ar[r] \ar[d] & \tilde L^\bullet \ar[d] \ar[r] &
M^\bullet \ar[r]_{\delta} \ar[d] & K^\bullet[1] \ar[d] \\
K^\bullet \ar[r] & L^\bullet \ar[r] &
C(f)^\bullet \ar[r]^{-p} & K^\bullet[1]
}
$$
が存在する。ここで上の三角は、項ごとに分裂する完全列
$K^\bullet \to \tilde L^\bullet \to M^\bullet$ に付随する三角である。
\end{enumerate}
\end{lemma}

\begin{proof}
(1) の証明。$C(\alpha)^n = B^n \oplus A^{n + 1}$ である。
$C(\alpha)^n \to C^n$ を、$B^n$ への射影とそれに続く $\beta^n$ によって
定める。合成 $A^{n + 1} \to B^{n + 1} \to C^{n + 1}$ が零なので、これは
複体の射を定める。ホモトピー逆として、次数 $n$ で
$(s^n , -\delta^n)$ によって与えられる
$C^\bullet \to C(\alpha)^\bullet$ を取る。これは複体の射である。実際、
$\delta^n$ は
$$
d \circ s^n - s^{n + 1} \circ d = \alpha \circ \delta^n
$$
を満たす一意な射 $C^n \to A^{n + 1}$ として特徴づけられる。Homology,
Lemma \ref{homology-lemma-ses-termwise-split-cochain} の証明を参照せよ。
合成 $C^\bullet \to C(\alpha)^\bullet \to C^\bullet$ は恒等射である。
逆向きの合成 $C(\alpha)^\bullet \to C^\bullet \to C(\alpha)^\bullet$ は射
$$
\left(
\begin{matrix}
s^n \circ \beta^n & 0 \\
-\delta^n \circ \beta^n & 0
\end{matrix}
\right)
$$
に等しい。これが恒等射とホモトピックであることを見るには、行列
$$
\left(
\begin{matrix}
0 & 0 \\
\pi^n & 0
\end{matrix}
\right) : C(\alpha)^n = B^n \oplus A^{n + 1} \to
B^{n - 1} \oplus A^n = C(\alpha)^{n - 1}
$$
で与えられるホモトピー
$h^n : C(\alpha)^n \to C(\alpha)^{n - 1}$ を用いる。実際、
$$
\left(
\begin{matrix}
1 & 0 \\
0 & 1
\end{matrix}
\right)
-
\left(
\begin{matrix}
s^n \\
-\delta^n
\end{matrix}
\right)
\left(
\begin{matrix}
\beta^n & 0
\end{matrix}
\right)
=
\left(
\begin{matrix}
d & \alpha^n \\
0 & -d
\end{matrix}
\right)
\left(
\begin{matrix}
0 & 0 \\
\pi^n & 0
\end{matrix}
\right)
+
\left(
\begin{matrix}
0 & 0 \\
\pi^{n + 1} & 0
\end{matrix}
\right)
\left(
\begin{matrix}
d & \alpha^{n + 1} \\
0 & -d
\end{matrix}
\right)
$$
は直ちに確かめられる。(1) の証明を終えるには、射
$-p : C(\alpha)^\bullet \to A^\bullet[1]$（Definition
\ref{derived-definition-cone} を参照）と、合成
$C(\alpha)^\bullet \to C^\bullet \to A^\bullet[1]$ とがホモトピーを除いて
一致することを示す必要がある。これは上の計算から明らかである。すなわち、
ホモトピー逆 $(s, -\delta) : C^\bullet \to C(\alpha)^\bullet$ を用いて、代わりに
二つの写像 $C^\bullet \to A^\bullet[1]$ が一致することを確かめればよい。
実際、望みどおり $p \circ (s, -\delta) = -\delta$ である。

\medskip\noindent
(2) の証明。$\tilde f : K^\bullet \to \tilde L^\bullet$、
$s : L^\bullet \to \tilde L^\bullet$、および
$\pi : \tilde L^\bullet \to L^\bullet$ を補題
\ref{derived-lemma-make-injective} のものとする。補題
\ref{derived-lemma-functorial-cone} および補題 \ref{derived-lemma-third-isomorphism} により、三角
$(K^\bullet, L^\bullet, C(f), i, p)$ と
$(K^\bullet, \tilde L^\bullet, C(\tilde f), \tilde i, \tilde p)$ は同型である。
三角の同型は合成できることに注意する。したがって、$L^\bullet$ を
$\tilde L^\bullet$ で、$f$ を $\tilde f$ で置き換えてよい。言い換えれば、
$f$ は項ごとに分裂する単射であると仮定してよい。この場合、結果は (1) から従う。
\end{proof}

\begin{lemma}
\label{derived-lemma-sequence-maps-split}
$\mathcal{A}$ を加法圏とする。
$A_1^\bullet \to A_2^\bullet \to \ldots \to A_n^\bullet$ を、
合成可能な複体の射の列とする。このとき、可換図式
$$
\xymatrix{
A_1^\bullet \ar[r] &
A_2^\bullet \ar[r] &
\ldots \ar[r] &
A_n^\bullet \\
B_1^\bullet \ar[r] \ar[u] &
B_2^\bullet \ar[r] \ar[u] &
\ldots \ar[r] &
B_n^\bullet \ar[u]
}
$$
であって、各射 $B_i^\bullet \to B_{i + 1}^\bullet$ が項ごとに分裂する単射、
各射 $B_i^\bullet \to A_i^\bullet$ がホモトピー同値となるものが存在する。
さらに、すべての $A_i^\bullet$ が $K^{+}(\mathcal{A})$、
$K^{-}(\mathcal{A})$、または $K^b(\mathcal{A})$ に属するならば、
$B_i^\bullet$ もすべて同じ圏に属する。
\end{lemma}

\begin{proof}
$n = 1$ の場合は主張すべきことがない。補題 \ref{derived-lemma-make-injective} が
$n = 2$ の場合である。$B_n^\bullet$ を除く図式を構成できたと仮定する。
合成 $B_{n - 1}^\bullet \to A_{n - 1}^\bullet \to A_n^\bullet$ に補題
\ref{derived-lemma-make-injective} を適用する。すると、望む分解
$B_{n - 1}^\bullet \to B_n^\bullet \to A_n^\bullet$ が得られる。
\end{proof}

\begin{lemma}
\label{derived-lemma-rotate-triangle}
$\mathcal{A}$ を加法圏とする。
$(\alpha : A^\bullet \to B^\bullet, \beta : B^\bullet \to C^\bullet, s^n,
\pi^n)$ を項ごとに分裂する複体の列とし、
$(A^\bullet, B^\bullet, C^\bullet, \alpha, \beta, \delta)$ を
それに付随する三角とする。このとき、三角
$(C^\bullet[-1], A^\bullet, B^\bullet, \delta[-1], \alpha, \beta)$ は、三角
$(C^\bullet[-1], A^\bullet, C(\delta[-1])^\bullet, \delta[-1], i, p)$ と同型である。
\end{lemma}

\begin{proof}
$B^n = A^n \oplus C^n$ と書き、$\alpha^n$ と $\beta^n$ をそれぞれ自然な
包含写像と射影写像と同一視する。$\delta$ の構成により、
$$
d_B^n =
\left(
\begin{matrix}
d_A^n & \delta^n \\
0 & d_C^n
\end{matrix}
\right)
$$
である。一方、$\delta[-1] : C^\bullet[-1] \to A^\bullet$ の錐は
$C(\delta[-1])^n = A^n \oplus C^n$ で与えられ、その微分は上の行列と同一である。
これで補題が従う。
\end{proof}

\begin{lemma}
\label{derived-lemma-rotate-cone}
$\mathcal{A}$ を加法圏とする。$f : K^\bullet \to L^\bullet$ を複体の射とする。
三角 $(L^\bullet, C(f)^\bullet, K^\bullet[1], i, p, f[1])$ は、$f$ の錐の
定義から得られる項ごとに分裂する列
$$
0 \to L^\bullet \to C(f)^\bullet \to K^\bullet[1] \to 0
$$
に付随する三角である。
\end{lemma}

\begin{proof}
定義から直ちに従う。
\end{proof}

\section{ホモトピー圏における完全三角}
\label{derived-section-homotopy-triangulated}

\noindent
コホモロジー長完全列の境界写像が符号を伴わず蛇の補題の写像によって
与えられるようにしたいので、ホモトピー圏の完全三角を次のように定義する。

\begin{definition}
\label{derived-definition-distinguished-triangle}
$\mathcal{A}$ を加法圏とする。$K(\mathcal{A})$ の三角
$(X, Y, Z, f, g, h)$ が、項ごとに分裂する複体の完全列に付随する三角と
同型であるとき、これを {\it $K(\mathcal{A})$ の完全三角} と呼ぶ。
Definition \ref{derived-definition-split-ses} を参照せよ。
$K^{+}(\mathcal{A})$、$K^{-}(\mathcal{A})$、$K^b(\mathcal{A})$ についても
同様に定義する。
\end{definition}

\noindent
補題 \ref{derived-lemma-the-same-up-to-isomorphisms} により、形
$(K^\bullet, L^\bullet, C(f)^\bullet, f, i, -p)$ の三角は完全三角である。
実際、この定義から三角圏が得られる。命題
\ref{derived-proposition-homotopy-category-triangulated} を参照せよ。この命題を
証明する前に、TR4 を証明するための補題をもう一つ用意する。

\begin{lemma}
\label{derived-lemma-two-split-injections}
$\mathcal{A}$ を加法圏とする。$\alpha : A^\bullet \to B^\bullet$ および
$\beta : B^\bullet \to C^\bullet$ を複体の分裂単射とする。このとき、TR4 が
成り立つような完全三角
$(A^\bullet, B^\bullet, Q_1^\bullet, \alpha, p_1, d_1)$、
$(A^\bullet, C^\bullet, Q_2^\bullet, \beta \circ \alpha, p_2, d_2)$、および
$(B^\bullet, C^\bullet, Q_3^\bullet, \beta, p_3, d_3)$ が存在する。
\end{lemma}

\begin{proof}
$\pi_1^n : B^n \to A^n$ および $\pi_3^n : C^n \to B^n$ を分裂射とする。
すると $A^\bullet \to C^\bullet$ も分裂射
$\pi_2^n = \pi_1^n \circ \pi_3^n$ を持つ分裂単射である。
$Q_1^\bullet$、$Q_2^\bullet$、$Q_3^\bullet$ を「商」複体と書く。言い換えれば、
$Q_1^n = \Ker(\pi_1^n)$、$Q_3^n = \Ker(\pi_3^n)$、および
$Q_2^n = \Ker(\pi_2^n)$ である。これらの核は存在することに注意する。このとき、
$B^n = A^n \oplus Q_1^n$ および $C^n = B^n \oplus Q_3^n$ である。ここでは
$A^n$ を $B^n$ の部分対象とみなし、以下も同様とする。したがって
$C^n = A^n \oplus Q_1^n \oplus Q_3^n$ である。また、
$\pi_2^n = \pi_1^n \circ \pi_3^n$ は $Q_1^n$ と $Q_3^n$ の双方で零なので、
$Q_2^n = Q_1^n \oplus Q_3^n$ である。可換図式
$$
\begin{matrix}
0 & \to & A^\bullet & \to & B^\bullet & \to & Q_1^\bullet & \to & 0 \\
  &     & \downarrow &     & \downarrow &     & \downarrow  & \\
0 & \to & A^\bullet & \to & C^\bullet & \to & Q_2^\bullet & \to & 0 \\
  &     & \downarrow &     & \downarrow &     & \downarrow  & \\
0 & \to & B^\bullet & \to & C^\bullet & \to & Q_3^\bullet & \to & 0
\end{matrix}
$$
を考える。この図式の各行は項ごとに分裂する完全列なので、定義により
完全三角を定める。さらに、上の図式の下向きの矢印は選んだ分裂と両立するため、
三角の射
$$
(A^\bullet \to B^\bullet \to Q_1^\bullet \to A^\bullet[1])
\longrightarrow
(A^\bullet \to C^\bullet \to Q_2^\bullet \to A^\bullet[1])
$$
および
$$
(A^\bullet \to C^\bullet \to Q_2^\bullet \to A^\bullet[1])
\longrightarrow
(B^\bullet \to C^\bullet \to Q_3^\bullet \to B^\bullet[1]).
$$
を定める。上の図式の最下段の分裂列の分裂射 $Q_3^n \to C^n$ を
$C^n \to Q_2^n$ と合成すると、分裂列
$0 \to Q_1^\bullet \to Q_2^\bullet \to Q_3^\bullet \to 0$ の分裂射が得られる。
これから、対応する完全三角
$$
(Q_1^\bullet \to Q_2^\bullet \to Q_3^\bullet \to Q_1^\bullet[1])
$$
の射 $\delta : Q_3^\bullet \to Q_1^\bullet[1]$ は、合成
$Q_3^\bullet \to B^\bullet[1] \to Q_1^\bullet[1]$ に等しいことが容易に従う。
したがって、公理 TR4 の結論にある構造が得られる。
\end{proof}

\begin{proposition}
\label{derived-proposition-homotopy-category-triangulated}
$\mathcal{A}$ を加法圏とする。ホモトピーを法とする複体の圏
$K(\mathcal{A})$ は、その自然なシフト関手と上で定義した完全三角によって
三角圏となる。
\end{proposition}

\begin{proof}
TR1 の証明。定義により、完全三角と同型な三角はすべて完全である。また、任意の三角
$(A^\bullet, A^\bullet, 0, 1, 0, 0)$ は完全である。実際、
$0 \to A^\bullet \to A^\bullet \to 0 \to 0$ は項ごとに分裂する複体の列である。
最後に、複体の任意の射 $f : K^\bullet \to L^\bullet$ に対し、三角
$(K, L, C(f), f, i, -p)$ は補題 \ref{derived-lemma-the-same-up-to-isomorphisms} により完全である。

\medskip\noindent
TR2 の証明。$(X, Y, Z, f, g, h)$ を三角とし、
$(Y, Z, X[1], g, h, -f[1])$ が完全であると仮定する。このとき、項ごとに
分裂する複体の列 $A^\bullet \to B^\bullet \to C^\bullet$ であって、それに
付随する三角 $(A^\bullet, B^\bullet, C^\bullet, \alpha, \beta, \delta)$ が
$(Y, Z, X[1], g, h, -f[1])$ と同型となるものが存在する。逆向きに回転すると、
$(X, Y, Z, f, g, h)$ は
$(C^\bullet[-1], A^\bullet, B^\bullet, -\delta[-1], \alpha, \beta)$ と同型である。
補題 \ref{derived-lemma-rotate-triangle} により、三角
$(C^\bullet[-1], A^\bullet, B^\bullet, \delta[-1], \alpha, \beta)$ は
$(C^\bullet[-1], A^\bullet, C(\delta[-1])^\bullet, \delta[-1], i, p)$ と同型である。
前の三角の同型を $Y$ 上の $-1$ と前合成すると、
$(X, Y, Z, f, g, h)$ は
$(C^\bullet[-1], A^\bullet, C(\delta[-1])^\bullet, \delta[-1], i, -p)$ と同型になる。
したがって補題 \ref{derived-lemma-the-same-up-to-isomorphisms} により完全である。
逆に、$(X, Y, Z, f, g, h)$ が完全であると仮定する。補題
\ref{derived-lemma-the-same-up-to-isomorphisms} により、これはある複体の射 $f$ に対する
形 $(K^\bullet, L^\bullet, C(f), f, i, -p)$ の三角と同型である。すると回転した
三角 $(Y, Z, X[1], g, h, -f[1])$ は
$(L^\bullet, C(f), K^\bullet[1], i, -p, -f[1])$ と同型であり、これはさらに
三角 $(L^\bullet, C(f), K^\bullet[1], i, p, f[1])$ と同型である。
補題 \ref{derived-lemma-rotate-cone} により、この三角は完全である。ゆえに、望みどおり
$(Y, Z, X[1], g, h, -f[1])$ は完全である。

\medskip\noindent
TR3 の証明。$(X, Y, Z, f, g, h)$ および
$(X', Y', Z', f', g', h')$ を $K(\mathcal{A})$ の完全三角とし、
$a : X \to X'$ および $b : Y \to Y'$ を
$f' \circ a = b \circ f$ を満たす射とする。補題
\ref{derived-lemma-the-same-up-to-isomorphisms} により、
$(X, Y, Z, f, g, h) = (X, Y, C(f), f, i, -p)$ および
$(X', Y', Z', f', g', h') = (X', Y', C(f'), f', i', -p')$ と仮定してよい。
ここで、$f, f', a, b$ が与える可換図式に補題
\ref{derived-lemma-functorial-cone} を適用すればよい。

\medskip\noindent
TR4 の証明。ここまでで $K(\mathcal{A})$ が前三角圏であることは分かっている。
したがって補題 \ref{derived-lemma-easier-axiom-four} を用いることができる。
$A^\bullet \to B^\bullet$ および $B^\bullet \to C^\bullet$ を
$K(\mathcal{A})$ の合成可能な射とする。補題 \ref{derived-lemma-sequence-maps-split} により、
$A^\bullet \to B^\bullet$ と $B^\bullet \to C^\bullet$ は分裂単射であると
仮定してよい。この場合、結果は補題 \ref{derived-lemma-two-split-injections} から従う。
\end{proof}

\begin{remark}
\label{derived-remark-boundedness-conditions-triangulated}
$\mathcal{A}$ を加法圏とする。命題
\ref{derived-proposition-homotopy-category-triangulated} とまったく同じ証明により、圏
$K^{+}(\mathcal{A})$、$K^{-}(\mathcal{A})$、$K^b(\mathcal{A})$ も三角圏である。
実際、上または下に有界な複体の間の射の錐も、それぞれ上または下に有界である。
ただし、下でこれらが $K(\mathcal{A})$ の三角部分圏であることを証明するので、
それによって別の証明も得られる。
\end{remark}

\begin{lemma}
\label{derived-lemma-bounded-triangulated-subcategories}
$\mathcal{A}$ を加法圏とする。圏 $K^{+}(\mathcal{A})$、
$K^{-}(\mathcal{A})$、$K^b(\mathcal{A})$ は $K(\mathcal{A})$ の
充満三角部分圏である。
\end{lemma}

\begin{proof}
挙げた各圏は充満加法部分圏である。これらが三角部分圏であることを示すため、
補題 \ref{derived-lemma-triangulated-subcategory} の判定条件を用いる。各圏
$K^{+}(\mathcal{A})$、$K^{-}(\mathcal{A})$、$K^b(\mathcal{A})$ が
シフト関手 $[1], [-1]$ で保たれることは明らかである。最後に、
$f : A^\bullet \to B^\bullet$ を $K^{+}(\mathcal{A})$、
$K^{-}(\mathcal{A})$、または $K^b(\mathcal{A})$ の射とする。このとき
$(A^\bullet, B^\bullet, C(f)^\bullet, f, i, -p)$ は $K(\mathcal{A})$ の完全三角であり、
錐の構成から明らかなように $C(f)^\bullet \in K^{+}(\mathcal{A})$、
$K^{-}(\mathcal{A})$、または $K^b(\mathcal{A})$ である。これで補題が示された。
（別の方法として、
% P02-E0026: frozen source has “an termwise split injection”.
$K^\bullet \to L^\bullet$ は $K^{+}(\mathcal{A})$、$K^{-}(\mathcal{A})$、または
$K^b(\mathcal{A})$ において、項ごとに分裂する複体の単射と同型である
（補題 \ref{derived-lemma-make-injective} を参照）から、それに付随する完全三角を直接取ってもよい。）
\end{proof}

\begin{lemma}
\label{derived-lemma-additive-exact-homotopy-category}
$\mathcal{A}$、$\mathcal{B}$ を加法圏とし、
$F : \mathcal{A} \to \mathcal{B}$ を加法関手とする。誘導される関手
$$
\begin{matrix}
F : K(\mathcal{A}) \longrightarrow K(\mathcal{B}) \\
F : K^{+}(\mathcal{A}) \longrightarrow K^{+}(\mathcal{B}) \\
F : K^{-}(\mathcal{A}) \longrightarrow K^{-}(\mathcal{B}) \\
F : K^b(\mathcal{A}) \longrightarrow K^b(\mathcal{B})
\end{matrix}
$$
は三角圏の完全関手である。
\end{lemma}

\begin{proof}
$A^\bullet \to B^\bullet \to C^\bullet$ を、分裂
$(s^n, \pi^n)$ と付随する射 $\delta : C^\bullet \to A^\bullet[1]$ を持つ
$\mathcal{A}$ の複体の項ごとに分裂する列とする。Definition
\ref{derived-definition-split-ses} を参照せよ。このとき
$F(A^\bullet) \to F(B^\bullet) \to F(C^\bullet)$ は、分裂
$(F(s^n), F(\pi^n))$ と付随する射
$F(\delta) : F(C^\bullet) \to F(A^\bullet)[1]$ を持つ、項ごとに分裂する
複体の列である。したがって $F$ は完全三角を完全三角に移す。
\end{proof}

\begin{lemma}
\label{derived-lemma-improve-distinguished-triangle-homotopy}
$\mathcal{A}$ を加法圏とし、
$(A^\bullet, B^\bullet, C^\bullet, a, b, c)$ を $K(\mathcal{A})$ の完全三角とする。
このとき、すべての $n$ に対して
$0 \to A^n \to (B')^n \to C^n \to 0$ が分裂短完全列となるような、これと
同型な完全三角 $(A^\bullet, (B')^\bullet, C^\bullet, a', b', c)$ が存在する。
\end{lemma}

\begin{proof}
命題 \ref{derived-proposition-homotopy-category-triangulated} により
$K(\mathcal{A})$ が三角圏であることを用いる。
$W^\bullet$ を $c : C^\bullet \to A^\bullet[1]$ の錐とし、その写像を
$i : A^\bullet[1] \to W^\bullet$ および $p : W^\bullet \to C^\bullet[1]$ とする。
補題 \ref{derived-lemma-the-same-up-to-isomorphisms} により、
$(C^\bullet, A^\bullet[1], W^\bullet, c, i, -p)$ は完全三角である。
二回逆向きに回転すると、
$(A^\bullet, W^\bullet[-1], C^\bullet, -i[-1], p[-1], c)$ が完全三角であると分かる。
TR3 により、完全三角の射
$$
(\text{id}, \beta, \text{id}) : (A^\bullet, B^\bullet, C^\bullet, a, b, c) \to
(A^\bullet, W^\bullet[-1], C^\bullet, -i[-1], p[-1], c)
$$
が存在し、補題 \ref{derived-lemma-third-isomorphism-triangle} によりこれは同型でなければならない。
錐の構成そのものにより
$0 \to A^\bullet \to W^\bullet[-1] \to C^\bullet \to 0$ は項ごとに分裂する
複体の短完全列だから、これで証明が完了する。Section \ref{derived-section-cones} を参照せよ。
\end{proof}

\begin{remark}
\label{derived-remark-homotopy-double}
$\mathcal{A}$ を可算直和を持つ加法圏とする。
$\text{DoubleComp}(\mathcal{A})$ を $\mathcal{A}$ における二重複体の圏とする。
Homology, Section \ref{homology-section-double-complexes} を参照せよ。この圏を
用いて二つの三角圏を構成できる。
\begin{enumerate}
\item $\text{DoubleComp}(\mathcal{A})$ の対象 $A^{\bullet, \bullet}$ を、次のように
複体を項とする複体とみなす。
$$
\ldots \to A^{\bullet, -1} \to A^{\bullet, 0} \to A^{\bullet, 1} \to \ldots
$$
そして、命題 \ref{derived-proposition-homotopy-category-triangulated} による対応する三角構造を
持つホモトピー圏 $K_{first}(\text{DoubleComp}(\mathcal{A}))$ を取る。
Homology, Remark
\ref{homology-remark-double-complex-complex-of-complexes-first} により、関手
$$
\text{Tot} :
K_{first}(\text{DoubleComp}(\mathcal{A}))
\longrightarrow
K(\mathcal{A})
$$
は三角圏の完全関手である。
\item $\text{DoubleComp}(\mathcal{A})$ の対象 $A^{\bullet, \bullet}$ を、次のように
複体を項とする複体とみなす。
$$
\ldots \to A^{-1, \bullet} \to A^{0, \bullet} \to A^{1, \bullet} \to \ldots
$$
そして、命題 \ref{derived-proposition-homotopy-category-triangulated} による対応する三角構造を
持つホモトピー圏 $K_{second}(\text{DoubleComp}(\mathcal{A}))$ を取る。
Homology, Remark
\ref{homology-remark-double-complex-complex-of-complexes-second} により、関手
$$
\text{Tot} :
K_{second}(\text{DoubleComp}(\mathcal{A}))
\longrightarrow
K(\mathcal{A})
$$
は三角圏の完全関手である。
\end{enumerate}
\end{remark}

\begin{remark}
\label{derived-remark-double-complex-as-tensor-product-of}
$\mathcal{A}$、$\mathcal{B}$、$\mathcal{C}$ を加法圏とし、$\mathcal{C}$ は
可算直和を持つと仮定する。射について双線形な関手
$$
\otimes : \mathcal{A} \times \mathcal{B} \longrightarrow \mathcal{C},
\quad
(X, Y) \longmapsto X \otimes Y
$$
が与えられたとする。これは関手
$$
\text{Comp}(\mathcal{A}) \times \text{Comp}(\mathcal{B})
\longrightarrow
\text{DoubleComp}(\mathcal{C}), \quad
(X^\bullet, Y^\bullet)
\longmapsto
X^\bullet \otimes Y^\bullet
$$
を定める。Homology, Example
\ref{homology-example-double-complex-as-tensor-product-of} を参照せよ。
\begin{enumerate}
\item $\text{Comp}(\mathcal{A})$ の対象 $X^\bullet$ を固定すると、関手
$$
K(\mathcal{B}) \longrightarrow K(\mathcal{C}), \quad
Y^\bullet \longmapsto \text{Tot}(X^\bullet \otimes Y^\bullet)
$$
は三角圏の完全関手である。
\item $\text{Comp}(\mathcal{B})$ の対象 $Y^\bullet$ を固定すると、関手
$$
K(\mathcal{A}) \longrightarrow K(\mathcal{C}), \quad
X^\bullet \longmapsto \text{Tot}(X^\bullet \otimes Y^\bullet)
$$
は三角圏の完全関手である。
\end{enumerate}
これは remark \ref{derived-remark-homotopy-double} から従う。実際、関手
$\text{Comp}(\mathcal{A}) \to \text{DoubleComp}(\mathcal{C})$、
$Y^\bullet \mapsto X^\bullet \otimes Y^\bullet$ および
$\text{Comp}(\mathcal{B}) \to \text{DoubleComp}(\mathcal{C})$、
$X^\bullet \mapsto X^\bullet \otimes Y^\bullet$ はホモトピーおよび
項ごとに分裂する短完全列と両立することが直ちに分かるため、三角圏の完全関手
$$
K(\mathcal{B}) \to K_{first}(\text{DoubleComp}(\mathcal{C}))
\quad\text{and}\quad
K(\mathcal{A}) \to K_{second}(\text{DoubleComp}(\mathcal{C}))
$$
を誘導する。二つのうち最初のものについては、同型
$$
\text{Tot}(X^\bullet \otimes Y^\bullet[1]) \cong
\text{Tot}(X^\bullet \otimes Y^\bullet)[1]
$$
が符号を伴うことに注意する（これは Homology, Remark
\ref{homology-remark-shift-double-complex} で選んだ符号に由来する）。
\end{remark}

\section{導来圏}
\label{derived-section-derived-categories}

\noindent
この節では、アーベル圏 $\mathcal{A}$ の導来圏を、$K(\mathcal{A})$ において
擬同型を可逆にすることによって構成する。その前に、関手
$H^i : \text{Comp}(\mathcal{A}) \to \mathcal{A}$ は $K(\mathcal{A})$ を経由して
分解することを思い出そう。Homology, Lemma
\ref{homology-lemma-map-cohomology-homotopy-cochain} を参照せよ。さらに、
Homology, Definition \ref{homology-definition-cohomology-shift} では同一視
$H^i(K^\bullet[n]) = H^{i + n}(K^\bullet)$ を定義した。ここでは混乱と起こり得る
符号の誤りを避けるため、
$$
H^i(K^\bullet) = H^0(K^\bullet[i])
$$
と改めて定義することにする。

\begin{lemma}
\label{derived-lemma-cohomology-homological}
$\mathcal{A}$ をアーベル圏とする。関手
$$
H^0 : K(\mathcal{A}) \longrightarrow \mathcal{A}
$$
はホモロジー的である。
\end{lemma}

\begin{proof}
$H^0$ は関手なので、完全三角の定義により、項ごとに分裂する複体の短完全列
$0 \to A^\bullet \to B^\bullet \to C^\bullet \to 0$ が与えられたとき、列
$H^0(A^\bullet) \to H^0(B^\bullet) \to H^0(C^\bullet)$ が完全であることを
示せば十分である。これは Homology, Lemma
\ref{homology-lemma-long-exact-sequence-cochain} から従う。
\end{proof}

\noindent
特に、この補題により、$K(\mathcal{A})$ の完全三角 $(X, Y, Z, f, g, h)$ から
コホモロジー長完全列
\begin{equation}
\label{derived-equation-long-exact-cohomology-sequence-D}
\xymatrix{
\ldots \ar[r] &
H^i(X) \ar[r]^{H^i(f)} &
H^i(Y) \ar[r]^{H^i(g)} &
H^i(Z) \ar[r]^{H^i(h)} &
H^{i + 1}(X) \ar[r] & \ldots
}
\end{equation}
が得られる。式 (\ref{derived-equation-long-exact-cohomology-sequence}) を参照せよ。
さらに、これは複体の短完全列に付随するコホモロジー長完全列と両立する
% P02-E0027: frozen source contains the editorial placeholder “insert future reference here”.
（将来の参照をここに挿入）。例えば、
$(A^\bullet, B^\bullet, C^\bullet, \alpha, \beta, \delta)$ が、項ごとに
分裂する複体の完全列に付随する完全三角ならば（Definition
\ref{derived-definition-split-ses} を参照）、上のコホモロジー列は蛇の補題を用いて
定義した列と一致する。Homology, Lemma
\ref{homology-lemma-long-exact-sequence-cochain}、また列の一致については
Homology, Lemma \ref{homology-lemma-ses-termwise-split-long-cochain} を参照せよ。

\medskip\noindent
複体 $K^\bullet$ は、すべての $i \in \mathbf{Z}$ に対して
$H^i(K^\bullet) = 0$ であるとき {\it 非輪状} であることを思い出そう。
また、複体の射 $f : K^\bullet \to L^\bullet$ は、すべての $i$ に対して
$H^i(f)$ が同型であるとき、かつそのときに限り {\it 擬同型} である。
Homology, Definition \ref{homology-definition-quasi-isomorphism-cochain} を参照せよ。

\begin{lemma}
\label{derived-lemma-acyclic}
$\mathcal{A}$ をアーベル圏とする。非輪状複体からなる $K(\mathcal{A})$ の充満部分圏
$\text{Ac}(\mathcal{A})$ は、$K(\mathcal{A})$ の厳密充満飽和三角部分圏である。
これに対応する $K(\mathcal{A})$ の飽和乗法系（補題 \ref{derived-lemma-operations} を参照）は、
擬同型の集合 $\text{Qis}(\mathcal{A})$ である。特に、局所化関手
$Q : K(\mathcal{A}) \to \text{Qis}(\mathcal{A})^{-1}K(\mathcal{A})$ の核は
$\text{Ac}(\mathcal{A})$ であり、関手 $H^0$ は $Q$ を経由して分解する。
\end{lemma}

\begin{proof}
補題 \ref{derived-lemma-cohomology-homological} により $H^0$ がホモロジー的関手であることは
分かっている。したがって、この補題は補題 \ref{derived-lemma-acyclic-general} の特別な場合である。
\end{proof}

\begin{definition}
\label{derived-definition-unbounded-derived-category}
$\mathcal{A}$ をアーベル圏とし、$\text{Ac}(\mathcal{A})$ および
$\text{Qis}(\mathcal{A})$ を補題 \ref{derived-lemma-acyclic} のものとする。
{\it $\mathcal{A}$ の導来圏}とは、三角圏
$$
D(\mathcal{A}) =
K(\mathcal{A})/\text{Ac}(\mathcal{A}) =
\text{Qis}(\mathcal{A})^{-1} K(\mathcal{A}).
$$
のことである。商関手との合成が上で定義した関手 $H^0$ に戻る一意な関手を
$H^0 : D(\mathcal{A}) \to \mathcal{A}$ と書く。補題
\ref{derived-lemma-homological-functor-bounded} を用いて、対象の集合が
$$
\begin{matrix}
\Ob(D^{+}(\mathcal{A})) =
\{X \in \Ob(D(\mathcal{A})) \mid
H^n(X) = 0\text{ for all }n \ll 0\} \\
\Ob(D^{-}(\mathcal{A})) =
\{X \in \Ob(D(\mathcal{A})) \mid
H^n(X) = 0\text{ for all }n \gg 0\} \\
\Ob(D^b(\mathcal{A})) =
\{X \in \Ob(D(\mathcal{A})) \mid
H^n(X) = 0\text{ for all }|n| \gg 0\}
\end{matrix}
$$
である厳密充満飽和三角部分圏
$D^{+}(\mathcal{A}), D^{-}(\mathcal{A}), D^b(\mathcal{A})$ を導入する。
圏 $D^b(\mathcal{A})$ を $\mathcal{A}$ の {\it 有界導来圏} と呼ぶ。
\end{definition}

\noindent
$K^\bullet$ と $L^\bullet$ が $\mathcal{A}$ の複体であるとき、
$K^\bullet$ と $L^\bullet$ が $D(\mathcal{A})$ の対象として同型であることを、
「$K^\bullet$ は $L^\bullet$ と {\it 擬同型である}」と言うことがある。

\begin{remark}
\label{derived-remark-existence-derived}
この章では一貫して「小さい」アーベル圏を扱う（Stacks project の慣例である）。
「大きい」アーベル圏 $\mathcal{A}$ については、導来圏の射が集合であるかどうか
明らかでないため、導来圏 $D(\mathcal{A})$ が存在するかどうかも明らかでない。
実際、一般には射は集合でない。Examples, Lemma
\ref{examples-lemma-big-abelian-category} を参照せよ。しかし、$\mathcal{A}$ が
Grothendieck アーベル圏で、$K^\bullet, L^\bullet$ が $K(\mathcal{A})$ に属するならば、
Injectives, Theorem \ref{injectives-theorem-K-injective-embedding-grothendieck} により、
K-入射的複体 $I^\bullet$ への擬同型 $L^\bullet \to I^\bullet$ が存在し、
補題 \ref{derived-lemma-K-injective} により
$$
\Hom_{D(\mathcal{A})}(K^\bullet, L^\bullet) =
\Hom_{K(\mathcal{A})}(K^\bullet, I^\bullet)
$$
となる。これは集合である。Grothendieck アーベル圏の例には、環上の加群の圏や、
より一般に環付きサイト上の加群の層の圏がある。
\end{remark}

\noindent
各変種 $D^{+}(\mathcal{A}), D^{-}(\mathcal{A}), D^b(\mathcal{A})$ は、対応する
ホモトピー圏の局所化として構成できる。これは次の簡単な補題に基づく。

\begin{lemma}
\label{derived-lemma-complex-cohomology-bounded}
$\mathcal{A}$ をアーベル圏とし、$K^\bullet$ を複体とする。
\begin{enumerate}
\item すべての $n \ll 0$ に対して $H^n(K^\bullet) = 0$ ならば、
$L^\bullet$ が下に有界であるような擬同型
$K^\bullet \to L^\bullet$ が存在する。
\item すべての $n \gg 0$ に対して $H^n(K^\bullet) = 0$ ならば、
$M^\bullet$ が上に有界であるような擬同型
$M^\bullet \to K^\bullet$ が存在する。
\item すべての $|n| \gg 0$ に対して $H^n(K^\bullet) = 0$ ならば、
複体の射の可換図式
$$
\xymatrix{
K^\bullet \ar[r] & L^\bullet \\
M^\bullet \ar[u] \ar[r] & N^\bullet \ar[u]
}
$$
であって、すべての矢印が擬同型、$L^\bullet$ が下に有界、$M^\bullet$ が
上に有界、$N^\bullet$ が有界複体となるものが存在する。
\end{enumerate}
\end{lemma}

\begin{proof}
$a \ll 0 \ll b$ を取り、$M^\bullet = \tau_{\leq b}K^\bullet$、
$L^\bullet = \tau_{\geq a}K^\bullet$、および
$N^\bullet = \tau_{\leq b}L^\bullet = \tau_{\geq a}M^\bullet$ と置く。
切断関手については Homology, Section \ref{homology-section-truncations} を参照せよ。
\end{proof}

\noindent
次の補題を述べるため、$\text{Ac}^{+}(\mathcal{A})$、
$\text{Ac}^{-}(\mathcal{A})$、および $\text{Ac}^b(\mathcal{A})$ を、それぞれ
$K^{+}(\mathcal{A})$、$K^{-}(\mathcal{A})$、および $K^b(\mathcal{A})$ と
$\text{Ac}(\mathcal{A})$ との共通部分とする。同様に、
$\text{Qis}^{+}(\mathcal{A})$、$\text{Qis}^{-}(\mathcal{A})$、および
$\text{Qis}^b(\mathcal{A})$ を、それぞれ $K^{+}(\mathcal{A})$、
$K^{-}(\mathcal{A})$、および $K^b(\mathcal{A})$ と
$\text{Qis}(\mathcal{A})$ との共通部分とする。

\begin{lemma}
\label{derived-lemma-bounded-derived}
$\mathcal{A}$ をアーベル圏とする。部分圏 $\text{Ac}^{+}(\mathcal{A})$、
$\text{Ac}^{-}(\mathcal{A})$、および $\text{Ac}^b(\mathcal{A})$ は、それぞれ
$K^{+}(\mathcal{A})$、$K^{-}(\mathcal{A})$、および $K^b(\mathcal{A})$ の
厳密充満飽和三角部分圏である。対応する飽和乗法系（補題
\ref{derived-lemma-operations} を参照）は、それぞれ集合
$\text{Qis}^{+}(\mathcal{A})$、$\text{Qis}^{-}(\mathcal{A})$、および
$\text{Qis}^b(\mathcal{A})$ である。
\begin{enumerate}
\item 関手 $K^{+}(\mathcal{A}) \to D^{+}(\mathcal{A})$ の核は
$\text{Ac}^{+}(\mathcal{A})$ であり、これは三角圏の同値
$$
K^{+}(\mathcal{A})/\text{Ac}^{+}(\mathcal{A}) =
\text{Qis}^{+}(\mathcal{A})^{-1}K^{+}(\mathcal{A})
\longrightarrow
D^{+}(\mathcal{A})
$$
を誘導する。
\item 関手 $K^{-}(\mathcal{A}) \to D^{-}(\mathcal{A})$ の核は
$\text{Ac}^{-}(\mathcal{A})$ であり、これは三角圏の同値
$$
K^{-}(\mathcal{A})/\text{Ac}^{-}(\mathcal{A}) =
\text{Qis}^{-}(\mathcal{A})^{-1}K^{-}(\mathcal{A})
\longrightarrow
D^{-}(\mathcal{A})
$$
を誘導する。
\item 関手 $K^b(\mathcal{A}) \to D^b(\mathcal{A})$ の核は
$\text{Ac}^b(\mathcal{A})$ であり、これは三角圏の同値
$$
K^b(\mathcal{A})/\text{Ac}^b(\mathcal{A}) =
\text{Qis}^b(\mathcal{A})^{-1}K^b(\mathcal{A})
\longrightarrow
D^b(\mathcal{A})
$$
を誘導する。
\end{enumerate}
\end{lemma}

\begin{proof}
最初の主張は、ホモロジー的関手 $H^0$ の制限を考え、補題
\ref{derived-lemma-acyclic-general} を適用すれば従う。(1)、(2)、(3) の核に関する主張は、
それぞれの場合の定義から従う。補題 \ref{derived-lemma-complex-cohomology-bounded} により、
各関手は本質的全射である。証明を終えるには、各関手が充満忠実であることを
示す必要がある。まず下に有界な場合を示す。

\medskip\noindent
$K^\bullet, L^\bullet$ を下に有界な複体とする。これらの間の
$D(\mathcal{A})$ における射は、$s$ が擬同型であるような組
$f : K^\bullet \to (L')^\bullet$, $s : L^\bullet \to (L')^\bullet$ に対する
$s^{-1}f$ の形である。このことから $(L')^\bullet$ のコホモロジーは下に有界である。
したがって補題 \ref{derived-lemma-complex-cohomology-bounded} により、
$(L'')^\bullet$ が下に有界となる擬同型
$s' : (L')^\bullet \to (L'')^\bullet$ を選べる。このとき組
$(s' \circ f, s' \circ s)$ は
$\text{Qis}^{+}(\mathcal{A})^{-1}K^{+}(\mathcal{A})$ の射を定める。
したがって関手は「充満」である。最後に、組
$f : K^\bullet \to (L')^\bullet$, $s : L^\bullet \to (L')^\bullet$ が
$\text{Qis}^{+}(\mathcal{A})^{-1}K^{+}(\mathcal{A})$ の射を定め、その像が
$D(\mathcal{A})$ で零であると仮定する。これは、$s' \circ f = 0$ を満たす
擬同型 $s' : (L')^\bullet \to (L'')^\bullet$ が存在することを意味する。
補題 \ref{derived-lemma-complex-cohomology-bounded} をもう一度用いると、
$(L''')^\bullet$ が下に有界となる擬同型
$s'' : (L'')^\bullet \to (L''')^\bullet$ が得られる。したがって
$s'' \circ s' \circ f = 0$ であり、これは $s^{-1}f$ が
$\text{Qis}^{+}(\mathcal{A})^{-1}K^{+}(\mathcal{A})$ で零であることを意味する。
これで (1) の関手が同値であることが示された。

\medskip\noindent
(2) の証明は (1) の証明の双対である。(3) を証明するために (2) の結果を
用いてよい。したがって、関手
$\text{Qis}^b(\mathcal{A})^{-1}K^b(\mathcal{A})
\to \text{Qis}^{-}(\mathcal{A})^{-1}K^{-}(\mathcal{A})$
が充満忠実であることを示せば十分である。前段落の議論は、初めから上に有界な
複体だけを一貫して用いれば、この場合にもそのまま適用できる。
\end{proof}

\section{\texorpdfstring{標準 \ensuremath{\delta}-関手}{標準 delta-関手}}
\label{derived-section-canonical-delta-functor}

\noindent
導来圏は普遍コホモロジー関手の受け皿であるべきである。その結果を述べるため、
アーベル圏から三角圏への $\delta$-関手という概念を用いる。Definition
\ref{derived-definition-delta-functor} を参照せよ。

\medskip\noindent
関手 $\text{Comp}(\mathcal{A}) \to K(\mathcal{A})$ を考える。この関手は一般には
{\bf $\delta$-関手でない}。これを見る最も簡単な方法は、$\mathcal{A}$ の対象の
分裂しない短完全列 $0 \to A \to B \to C \to 0$ を考えることである。
$\Hom_{K(\mathcal{A})}(C[0], A[1]) = 0$ なので、この短完全列から生じる
完全三角は $(A[0], B[0], C[0], a, b, 0)$ の形でなければならない。しかし、
$K(\mathcal{A})$ にそのような完全三角が存在すれば、この拡大は分裂することになる。
これは矛盾である。

\medskip\noindent
一方、関手 $\text{Comp}(\mathcal{A}) \to D(\mathcal{A})$ は $\delta$-関手である。
これを見るため、アーベル圏 $\mathcal{A}$ の複体の短完全列
$$
0 \to A^\bullet \xrightarrow{a} B^\bullet \xrightarrow{b} C^\bullet \to 0
$$
に付随する射 $\delta$ を定義する必要がある。射 $a$ の錐 $C(a)^\bullet$ を考える。
$C(a)^n = B^n \oplus A^{n + 1}$ であり、$q^n : C(a)^n \to C^n$ を
$B^n$ への射影とそれに続く $b^n$ によって定める。こうして複体の射
$$
q : C(a)^\bullet \longrightarrow C^\bullet.
$$
を得る。ここで $i$ は Definition \ref{derived-definition-cone} のものとすると、
$q \circ i = b$ であることは明らかである。$a^\bullet$ は各次数で単射なので、
$q$ の核は $\text{id}_{A^\bullet}$ の錐と同一視され、これは非輪状である。
したがって $q$ は擬同型である。補題 \ref{derived-lemma-the-same-up-to-isomorphisms} により、三角
$$
(A, B, C(a), a, i, -p)
$$
は $K(\mathcal{A})$ の完全三角である。局所化関手
$K(\mathcal{A}) \to D(\mathcal{A})$ は完全なので、
$(A, B, C(a), a, i, -p)$ は $D(\mathcal{A})$ の完全三角でもある。
$q$ は擬同型だから、$q$ は $D(\mathcal{A})$ の同型である。したがって
$$
(A, B, C, a, b, -p \circ q^{-1})
$$
は $D(\mathcal{A})$ の完全三角である。これは次の補題を示唆する。

\begin{lemma}
\label{derived-lemma-derived-canonical-delta-functor}
$\mathcal{A}$ をアーベル圏とする。上で定義した関手
$\text{Comp}(\mathcal{A}) \to D(\mathcal{A})$ は、
$$
\delta_{A^\bullet \to B^\bullet \to C^\bullet} = - p \circ q^{-1}
$$
によって自然な $\delta$-関手の構造を持つ。ここで $p$ と $q$ は上で説明した
ものである。同じ構成により、関手
$\text{Comp}^{+}(\mathcal{A}) \to D^{+}(\mathcal{A})$、
$\text{Comp}^{-}(\mathcal{A}) \to D^{-}(\mathcal{A})$、および
$\text{Comp}^b(\mathcal{A}) \to D^b(\mathcal{A})$ も $\delta$-関手となる。
\end{lemma}

\begin{proof}
この選択により、複体の任意の短完全列から完全三角が得られることはすでに見た。
可換図式
$$
\xymatrix{
0 \ar[r] &
A^\bullet \ar[r]_a \ar[d]_f &
B^\bullet \ar[r]_b \ar[d]_g &
C^\bullet \ar[r] \ar[d]_h &
0 \\
0 \ar[r] &
(A')^\bullet \ar[r]^{a'} &
(B')^\bullet \ar[r]^{b'} &
(C')^\bullet \ar[r] &
0
}
$$
が与えられたとき、Definition \ref{derived-definition-delta-functor} (2) の望ましい
可換図式が得られることを示す必要がある。補題
\ref{derived-lemma-functorial-cone} により、組 $(f, g)$ は標準射
$c : C(a)^\bullet \to C(a')^\bullet$ を誘導する。簡単な計算により
$q' \circ c = h \circ q$ および $f[1] \circ p = p' \circ c$ が分かる。
これから結果が直ちに従う。
\end{proof}

\begin{lemma}
\label{derived-lemma-derived-compare-triangles-ses}
$\mathcal{A}$ をアーベル圏とする。
$$
\xymatrix{
0 \ar[r] &
A^\bullet \ar[r] \ar[d] &
B^\bullet \ar[r] \ar[d] &
C^\bullet \ar[r] \ar[d] &
0 \\
0 \ar[r] &
D^\bullet \ar[r] &
E^\bullet \ar[r] &
F^\bullet \ar[r] &
0
}
$$
を、各行が複体の短完全列、各縦矢印が擬同型であるような複体の射の可換図式とする。
上の補題 \ref{derived-lemma-derived-canonical-delta-functor} の $\delta$-関手は、短完全列
$0 \to A^\bullet \to B^\bullet \to C^\bullet \to 0$ および
$0 \to D^\bullet \to E^\bullet \to F^\bullet \to 0$ を、互いに同型な
完全三角へ移す。
\end{lemma}

\begin{proof}
$K(\mathcal{A}) \to D(\mathcal{A})$ が擬同型を同型に移し、付随する完全三角が
関手的であることから直ちに従う。
\end{proof}

\begin{lemma}
\label{derived-lemma-derived-compare-triangles-split-case}
$\mathcal{A}$ をアーベル圏とする。
$$
\xymatrix{
0 \ar[r] &
A^\bullet \ar[r] &
B^\bullet \ar[r] &
C^\bullet \ar[r] &
0
}
$$
を複体の短完全列とし、この列は項ごとに分裂すると仮定する。
% P02-E0028: frozen source has “a short exact sequences”.
$(A^\bullet, B^\bullet, C^\bullet, \alpha, \beta, \delta)$ を、この列に付随する
$K(\mathcal{A})$ の完全三角とする。上の補題
\ref{derived-lemma-derived-canonical-delta-functor} の $\delta$-関手は短完全列
$0 \to A^\bullet \to B^\bullet \to C^\bullet \to 0$ を、完全三角
$$
(A^\bullet, B^\bullet, C^\bullet, \alpha, \beta, \delta).
$$
と同型な三角へ移す。
\end{lemma}

\begin{proof}
補題 \ref{derived-lemma-the-same-up-to-isomorphisms} から従う。
\end{proof}

\begin{remark}
\label{derived-remark-truncation-distinguished-triangle}
$\mathcal{A}$ をアーベル圏とし、$K^\bullet$ を $\mathcal{A}$ の複体、
$a \in \mathbf{Z}$ とする。このとき $D(\mathcal{A})$ に標準的な完全三角
$$
\tau_{\leq a}K^\bullet \to K^\bullet \to \tau_{\geq a + 1}K^\bullet \to
(\tau_{\leq a}K^\bullet)[1]
$$
が存在すると主張する。ここでは Homology, Section
\ref{homology-section-truncations} の標準切断関手 $\tau$ を用いた。実際、まず
短完全列
$$
0 \to \tau_{\leq a}K^\bullet \to K^\bullet \to
K^\bullet/\tau_{\leq a}K^\bullet \to 0
$$
に、われわれの $\delta$-関手（補題
\ref{derived-lemma-derived-canonical-delta-functor}）が付随させる完全三角を取る。次に、
Homology, Section \ref{homology-section-truncations} のコホモロジー群の記述により、
写像 $K^\bullet \to \tau_{\geq a + 1}K^\bullet$ は擬同型
$K^\bullet/\tau_{\leq a}K^\bullet \to \tau_{\geq a + 1}K^\bullet$ を
経由して分解することを用いる。同様にして、標準的な完全三角
$$
\tau_{\leq a}K^\bullet \to \tau_{\leq a + 1}K^\bullet \to
H^{a + 1}(K^\bullet)[-a-1] \to (\tau_{\leq a}K^\bullet)[1]
$$
および
$$
H^a(K^\bullet)[-a] \to \tau_{\geq a}K^\bullet \to \tau_{\geq a + 1}K^\bullet
\to H^a(K^\bullet)[-a + 1]
$$
を得る。
\end{remark}

\begin{lemma}
\label{derived-lemma-trick-vanishing-composition}
$\mathcal{A}$ をアーベル圏とする。
$$
K_0^\bullet \to K_1^\bullet \to \ldots \to K_n^\bullet
$$
を、次を満たす複体の写像とする。
\begin{enumerate}
\item $i > 0$ に対して $H^i(K_0^\bullet) = 0$ である。
\item $H^{-j}(K_j^\bullet) \to H^{-j}(K_{j + 1}^\bullet)$ は零である。
\end{enumerate}
このとき、合成 $K_0^\bullet \to K_n^\bullet$ は $D(\mathcal{A})$ において
$\tau_{\leq -n}K_n^\bullet \to K_n^\bullet$ を経由して分解する。双対的に、
$$
K_n^\bullet \to K_{n - 1}^\bullet \to \ldots \to K_0^\bullet
$$
を、次を満たす複体の写像とする。
\begin{enumerate}
\item $i < 0$ に対して $H^i(K_0^\bullet) = 0$ である。
\item $H^j(K_{j + 1}^\bullet) \to H^j(K_j^\bullet)$ は零である。
\end{enumerate}
このとき、合成 $K_n^\bullet \to K_0^\bullet$ は $D(\mathcal{A})$ において
$K_n^\bullet \to \tau_{\geq n}K_n^\bullet$ を経由して分解する。
\end{lemma}

\begin{proof}
$n = 1$ の場合。$D(\mathcal{A})$ では
$\tau_{\leq 0}K_0^\bullet = K_0^\bullet$ なので、$K_0^\bullet$ を
$\tau_{\leq 0}K_0^\bullet$ で、$K_1^\bullet$ を $\tau_{\leq 0}K_1^\bullet$ で
置き換えてよい。完全三角
$$
\tau_{\leq -1}K_1^\bullet \to K_1^\bullet \to
H^0(K_1^\bullet)[0] \to (\tau_{\leq -1}K_1^\bullet)[1]
$$
を考える（remark \ref{derived-remark-truncation-distinguished-triangle}）。合成
$K_0^\bullet \to K_1^\bullet \to H^0(K_1^\bullet)[0]$ は、合成
$K_0^\bullet \to H^0(K_0^\bullet)[0] \to H^0(K_1^\bullet)[0]$ に等しく、
仮定により後者は零なので、零である。$\Hom_{D(\mathcal{A})}(K_0^\bullet, -)$ が
ホモロジー的関手であること（補題 \ref{derived-lemma-representable-homological}）から、
望む分解が得られる。$n = 2$ の場合は、$n = 1$ の場合から分解
$K_0^\bullet \to \tau_{\leq -1}K_1^\bullet$ を得て、複体の写像
$\tau_{\leq -1}K_1^\bullet \to \tau_{\leq -1}K_2^\bullet$ に再び
$n = 1$ の場合を適用し、分解
$\tau_{\leq -1}K_1^\bullet \to \tau_{\leq -2}K_2^\bullet$ を得る。
一般の場合もまったく同様に証明される。
\end{proof}

\section{フィルトレーション付き導来圏}
\label{derived-section-filtered-derived-category}

\noindent
この節の参考文献は \cite[I, Chapter V]{cotangent} である。
$\mathcal{A}$ をアーベル圏とする。この節では $\mathcal{A}$ の
フィルトレーション付き導来圏 $DF(\mathcal{A})$ を定義する。簡潔に言えば、
$\mathcal{A}$ の有限フィルトレーションを備えた対象からなる完全圏の導来圏として
定義する。（したがって、われわれの構成は完全圏の導来圏というより一般的な構成の
特別な場合である。例えば \cite{Buhler}、\cite{Keller} を参照せよ。）
Illusie のフィルトレーション付き導来圏は、フィルトレーションが有限である対象から
なる、われわれの圏の充満部分圏である。（われわれの圏でもフィルトレーションは
各次数で有限だが、一様に有界とは限らない。）この選択をする理由は、構成が難しく
なるわけではなく、補題 \ref{derived-lemma-two-ss-complex-functor} のスペクトル系列に議論を
適用できるからである。remark \ref{derived-remark-functorial-ss} も参照せよ。

\medskip\noindent
Homology, Section \ref{homology-section-filtrations} で導入した
フィルトレーション付き対象の記法を用いる。$\mathcal{A}$ のフィルトレーション付き
対象の圏を $\text{Fil}(\mathcal{A})$ と書く。すべてのフィルトレーションは減少であると定める。

\begin{definition}
\label{derived-definition-finite-filtered}
$\mathcal{A}$ をアーベル圏とする。{\it $\mathcal{A}$ の有限フィルトレーション付き
対象の圏}とは、フィルトレーション $F$ が有限である $\mathcal{A}$ の
フィルトレーション付き対象 $(A, F)$ の圏である。これを
$\text{Fil}^f(\mathcal{A})$ と書く。
\end{definition}

\noindent
したがって $\text{Fil}^f(\mathcal{A})$ は $\text{Fil}(\mathcal{A})$ の充満部分圏である。
各 $p \in \mathbf{Z}$ に対して関手
$\text{gr}^p : \text{Fil}^f(\mathcal{A}) \to \mathcal{A}$ がある。また、関手
$$
\text{gr} = \bigoplus\nolimits_{p \in \mathbf{Z}} \text{gr}^p :
\text{Fil}^f(\mathcal{A}) \to \text{Gr}(\mathcal{A})
$$
がある。ここで $\text{Gr}(\mathcal{A})$ は $\mathcal{A}$ の次数付き対象の圏である。
Homology, Definition \ref{homology-definition-graded} を参照せよ。最後に、
フィルトレーション付き対象 $(A, F)$ に $\mathcal{A}$ の台となる対象を対応させる関手
$$
(\text{forget }F) : \text{Fil}^f(\mathcal{A}) \longrightarrow \mathcal{A}
$$
がある。圏 $\text{Fil}^f(\mathcal{A})$ は加法圏だが、一般にはアーベル圏でない。
Homology, Example \ref{homology-example-not-abelian} を参照せよ。

\medskip\noindent
関手 $\text{gr}^p$、$\text{gr}$、$(\text{forget }F)$ は加法的なので、補題
\ref{derived-lemma-additive-exact-homotopy-category} により三角圏の完全関手
$$
\text{gr}^p, (\text{forget }F) :
K(\text{Fil}^f(\mathcal{A}))
\to
K(\mathcal{A})
\quad\text{and}\quad
\text{gr} :
K(\text{Fil}^f(\mathcal{A}))
\to
K(\text{Gr}(\mathcal{A}))
$$
を誘導する。アーベル圏のホモトピー圏の場合との類似により、次を定義する。

\begin{definition}
\label{derived-definition-filtered-acyclic}
$\mathcal{A}$ をアーベル圏とする。
\begin{enumerate}
\item $\alpha : K^\bullet \to L^\bullet$ を
$K(\text{Fil}^f(\mathcal{A}))$ の射とする。射 $\text{gr}(\alpha)$ が擬同型であるとき、
$\alpha$ を {\it フィルトレーション付き擬同型} と呼ぶ。
\item $K^\bullet$ を $K(\text{Fil}^f(\mathcal{A}))$ の対象とする。複体
$\text{gr}(K^\bullet)$ が非輪状であるとき、$K^\bullet$ を
{\it フィルトレーション付き非輪状} と呼ぶ。
\end{enumerate}
\end{definition}

\noindent
$\alpha : K^\bullet \to L^\bullet$ がフィルトレーション付き擬同型であるための
必要十分条件は、各 $\text{gr}^p(\alpha)$ が擬同型であることである。同様に、複体
$K^\bullet$ がフィルトレーション付き非輪状であるための必要十分条件は、各
$\text{gr}^p(K^\bullet)$ が非輪状であることである。

\begin{lemma}
\label{derived-lemma-filtered-cohomology-homological}
$\mathcal{A}$ をアーベル圏とする。
\begin{enumerate}
\item 関手
$K(\text{Fil}^f(\mathcal{A})) \longrightarrow \text{Gr}(\mathcal{A})$、
$K^\bullet \longmapsto H^0(\text{gr}(K^\bullet))$ はホモロジー的である。
\item 関手
$K(\text{Fil}^f(\mathcal{A})) \rightarrow \mathcal{A}$、
$K^\bullet \longmapsto H^0(\text{gr}^p(K^\bullet))$ はホモロジー的である。
\item 関手
$K(\text{Fil}^f(\mathcal{A})) \longrightarrow \mathcal{A}$、
$K^\bullet \longmapsto H^0((\text{forget }F)K^\bullet)$ はホモロジー的である。
\end{enumerate}
\end{lemma}

\begin{proof}
$H^0 : K(\mathcal{A}) \to \mathcal{A}$ がホモロジー的であること（補題
\ref{derived-lemma-cohomology-homological}）と、関手
$\text{gr}, \text{gr}^p, (\text{forget }F)$ が三角圏の完全関手であることから従う。
補題 \ref{derived-lemma-exact-compose-homological-functor} を参照せよ。
\end{proof}

\begin{lemma}
\label{derived-lemma-filtered-acyclic}
$\mathcal{A}$ をアーベル圏とする。フィルトレーション付き非輪状複体からなる
$K(\text{Fil}^f(\mathcal{A}))$ の充満部分圏 $\text{FAc}(\mathcal{A})$ は、
$K(\text{Fil}^f(\mathcal{A}))$ の厳密充満飽和三角部分圏である。対応する
$K(\text{Fil}^f(\mathcal{A}))$ の飽和乗法系（補題 \ref{derived-lemma-operations} を参照）は、
フィルトレーション付き擬同型の集合 $\text{FQis}(\mathcal{A})$ である。特に、局所化関手
$$
Q :
K(\text{Fil}^f(\mathcal{A}))
\longrightarrow
\text{FQis}(\mathcal{A})^{-1}K(\text{Fil}^f(\mathcal{A}))
$$
の核は $\text{FAc}(\mathcal{A})$ であり、関手 $H^0 \circ \text{gr}$ は $Q$ を
経由して分解する。
\end{lemma}

\begin{proof}
補題 \ref{derived-lemma-filtered-cohomology-homological} により
$H^0 \circ \text{gr}$ はホモロジー的関手である。したがって、この補題は補題
\ref{derived-lemma-acyclic-general} の特別な場合である。
\end{proof}

\begin{definition}
\label{derived-definition-filtered-derived}
$\mathcal{A}$ をアーベル圏とし、$\text{FAc}(\mathcal{A})$ および
$\text{FQis}(\mathcal{A})$ を補題 \ref{derived-lemma-filtered-acyclic} のものとする。
{\it $\mathcal{A}$ のフィルトレーション付き導来圏}とは、三角圏
$$
DF(\mathcal{A}) =
K(\text{Fil}^f(\mathcal{A}))/\text{FAc}(\mathcal{A}) =
\text{FQis}(\mathcal{A})^{-1} K(\text{Fil}^f(\mathcal{A})).
$$
のことである。
\end{definition}

\begin{lemma}
\label{derived-lemma-filtered-derived-functors}
関手 $\text{gr}^p, \text{gr}, (\text{forget }F)$ は、局所化関手と可換する標準的な
完全関手
$$
\text{gr}^p, (\text{forget }F):
DF(\mathcal{A})
\longrightarrow
D(\mathcal{A})
$$
および
$$
\text{gr}:
DF(\mathcal{A})
\longrightarrow
D(\text{Gr}(\mathcal{A}))
$$
を誘導する。
\end{lemma}

\begin{proof}
各関手 $\text{gr}^p, \text{gr}, (\text{forget }F)$ がフィルトレーション付き擬同型を
擬同型に移すことを示せば、局所化の普遍性（補題
\ref{derived-lemma-universal-property-localization}）から従う。最初の二つについては定義から
明らかである。最後のものについては少し議論が必要である。
$f : K^\bullet \to L^\bullet$ を $K(\text{Fil}^f(\mathcal{A}))$ の
フィルトレーション付き擬同型とする。$f$ を含む完全三角
$(K^\bullet, L^\bullet, M^\bullet, f, g, h)$ を選ぶ。このとき、補題
\ref{derived-lemma-filtered-acyclic} により $M^\bullet$ はフィルトレーション付き非輪状である。
したがって、$K(\mathcal{A})$ に対する対応する補題により、フィルトレーションを
忘れたときフィルトレーション付き非輪状複体が非輪状複体になることを示せば十分である。
これは Homology, Lemma \ref{homology-lemma-filtered-acyclic} から従う。
\end{proof}

\begin{definition}
\label{derived-definition-filtered-derived-bounded}
$\mathcal{A}$ をアーベル圏とする。{\it 有界フィルトレーション付き導来圏}
$DF^b(\mathcal{A})$ とは、$\text{gr}(X) \in D^b(\mathcal{A})$ を満たす対象 $X$ から
なる $DF(\mathcal{A})$ の充満部分圏である。下に有界なフィルトレーション付き導来圏
$DF^{+}(\mathcal{A})$ および上に有界なフィルトレーション付き導来圏
$DF^{-}(\mathcal{A})$ も同様に定義する。
\end{definition}

\begin{lemma}
\label{derived-lemma-filtered-complex-cohomology-bounded}
$\mathcal{A}$ をアーベル圏とし、
$K^\bullet \in K(\text{Fil}^f(\mathcal{A}))$ とする。
\begin{enumerate}
\item すべての $n < a$ に対して $H^n(\text{gr}(K^\bullet)) = 0$ ならば、
すべての $n < a$ に対して $L^n = 0$ となるフィルトレーション付き擬同型
$K^\bullet \to L^\bullet$ が存在する。
\item すべての $n > b$ に対して $H^n(\text{gr}(K^\bullet)) = 0$ ならば、
すべての $n > b$ に対して $M^n = 0$ となるフィルトレーション付き擬同型
$M^\bullet \to K^\bullet$ が存在する。
\item すべての $|n| \gg 0$ に対して $H^n(\text{gr}(K^\bullet)) = 0$ ならば、
複体の射の可換図式
$$
\xymatrix{
K^\bullet \ar[r] & L^\bullet \\
M^\bullet \ar[u] \ar[r] & N^\bullet \ar[u]
}
$$
であって、すべての矢印がフィルトレーション付き擬同型、$L^\bullet$ が下に有界、
$M^\bullet$ が上に有界、$N^\bullet$ が有界複体となるものが存在する。
\end{enumerate}
\end{lemma}

\begin{proof}
すべての $n < a$ に対して $H^n(\text{gr}(K^\bullet)) = 0$ であると仮定する。
Homology, Lemma \ref{homology-lemma-filtered-acyclic} により、列
$$
K^{a - 2} \xrightarrow{d^{a - 2}} K^{a - 1} \xrightarrow{d^{a - 1}} K^a
$$
は $\mathcal{A}$ の対象の完全列であり、射 $d^{a - 2}$ と $d^{a - 1}$ は厳密である。
したがって $\text{Fil}^f(\mathcal{A})$ において
$\Coim(d^{a - 1}) = \Im(d^{a - 1})$ であり、写像
$\text{gr}(\Im(d^{a - 1})) \to \text{gr}(K^a)$ は単射で、その像は
$\text{gr}(K^{a - 1}) \to \text{gr}(K^a)$ の像に等しい。Homology, Lemma
\ref{homology-lemma-characterize-strict} を参照せよ。これは、切断への写像
$K^\bullet \to \tau_{\geq a}K^\bullet$、ただし
$$
\tau_{\geq a}K^\bullet =
(\ldots \to 0 \to K^a/\Im(d^{a - 1}) \to K^{a + 1} \to \ldots)
$$
がフィルトレーション付き擬同型であることを意味する。これで (1) が示された。
(2) の証明は (1) の証明の双対である。(3) は (1) と (2) から形式的に従う。
\end{proof}

\noindent
次の補題を述べるため、$\text{FAc}^{+}(\mathcal{A})$、
$\text{FAc}^{-}(\mathcal{A})$、および $\text{FAc}^b(\mathcal{A})$ を、それぞれ
$K^{+}(\text{Fil}^f\mathcal{A})$、$K^{-}(\text{Fil}^f\mathcal{A})$、および
$K^b(\text{Fil}^f\mathcal{A})$ と $\text{FAc}(\mathcal{A})$ との共通部分とする。
同様に、$\text{FQis}^{+}(\mathcal{A})$、$\text{FQis}^{-}(\mathcal{A})$、および
$\text{FQis}^b(\mathcal{A})$ を、それぞれ
$K^{+}(\text{Fil}^f\mathcal{A})$、$K^{-}(\text{Fil}^f\mathcal{A})$、および
$K^b(\text{Fil}^f\mathcal{A})$ と $\text{FQis}(\mathcal{A})$ との共通部分とする。

\begin{lemma}
\label{derived-lemma-filtered-bounded-derived}
$\mathcal{A}$ をアーベル圏とする。部分圏 $\text{FAc}^{+}(\mathcal{A})$、
$\text{FAc}^{-}(\mathcal{A})$、および $\text{FAc}^b(\mathcal{A})$ は、それぞれ
$K^{+}(\text{Fil}^f\mathcal{A})$、$K^{-}(\text{Fil}^f\mathcal{A})$、および
$K^b(\text{Fil}^f\mathcal{A})$ の厳密充満飽和三角部分圏である。対応する飽和乗法系
（補題 \ref{derived-lemma-operations} を参照）は、それぞれ集合
$\text{FQis}^{+}(\mathcal{A})$、$\text{FQis}^{-}(\mathcal{A})$、および
$\text{FQis}^b(\mathcal{A})$ である。
\begin{enumerate}
\item 関手 $K^{+}(\text{Fil}^f\mathcal{A}) \to DF^{+}(\mathcal{A})$ の核は
$\text{FAc}^{+}(\mathcal{A})$ であり、これは三角圏の同値
$$
K^{+}(\text{Fil}^f\mathcal{A})/\text{FAc}^{+}(\mathcal{A}) =
\text{FQis}^{+}(\mathcal{A})^{-1}K^{+}(\text{Fil}^f\mathcal{A})
\longrightarrow
DF^{+}(\mathcal{A})
$$
を誘導する。
\item 関手 $K^{-}(\text{Fil}^f\mathcal{A}) \to DF^{-}(\mathcal{A})$ の核は
$\text{FAc}^{-}(\mathcal{A})$ であり、これは三角圏の同値
$$
K^{-}(\text{Fil}^f\mathcal{A})/\text{FAc}^{-}(\mathcal{A}) =
\text{FQis}^{-}(\mathcal{A})^{-1}K^{-}(\text{Fil}^f\mathcal{A})
\longrightarrow
DF^{-}(\mathcal{A})
$$
を誘導する。
\item 関手 $K^b(\text{Fil}^f\mathcal{A}) \to DF^b(\mathcal{A})$ の核は
$\text{FAc}^b(\mathcal{A})$ であり、これは三角圏の同値
$$
K^b(\text{Fil}^f\mathcal{A})/\text{FAc}^b(\mathcal{A}) =
\text{FQis}^b(\mathcal{A})^{-1}K^b(\text{Fil}^f\mathcal{A})
\longrightarrow
DF^b(\mathcal{A})
$$
を誘導する。
\end{enumerate}
\end{lemma}

\begin{proof}
上の結果、特に補題 \ref{derived-lemma-filtered-complex-cohomology-bounded} と、補題
\ref{derived-lemma-bounded-derived} の証明で用いたものとまったく同じ議論から従う。
\end{proof}

\section{一般の導来関手}
\label{derived-section-derived-functors}

\noindent
この節の参考文献は \cite{SGA4} に収められた Deligne の expos\'e XVII である。
三角圏の間の完全関手と始域圏の乗法系が与えられ、完全関手を局所化された圏へ
「正しく」延長したいときの、右導来関手および左導来関手の非常に一般的な概念がある。

\begin{situation}
\label{derived-situation-derived-functor}
ここで $F : \mathcal{D} \to \mathcal{D}'$ は三角圏の完全関手であり、$S$ は
$\mathcal{D}$ の三角圏構造と両立する $\mathcal{D}$ の飽和乗法系である。
\end{situation}

\noindent
$X \in \Ob(\mathcal{D})$ とする。Categories, Remark
\ref{categories-remark-left-localization-morphisms-colimit} で、始域が $X$ である
$S$ の矢印 $s : X \to X'$ からなるフィルター圏 $X/S$ を定義したことを思い出そう。
双対的に、Categories, Remark
\ref{categories-remark-right-localization-morphisms-colimit} では、終域が $X$ である
$S$ の矢印 $s : X' \to X$ からなる余フィルター圏 $S/X$ を定義した。

\begin{definition}
\label{derived-definition-right-derived-functor-defined}
Situation \ref{derived-situation-derived-functor} の仮定と記法を用い、
$X \in \Ob(\mathcal{D})$ とする。
\begin{enumerate}
\item ind-対象
$$
(X/S) \longrightarrow \mathcal{D}', \quad
(s : X \to X') \longmapsto F(X')
$$
が本質的に定値であるとき、{\it 右導来関手 $RF$ は $X$ で定義される}という
\footnote{圏の ind-対象または pro-対象が本質的に定値となる条件については、
Categories, Section \ref{categories-section-essentially-constant} を参照せよ。}。
この場合、$\mathcal{D}'$ における値 $Y$ を {\it $X$ における $RF$ の値}と呼ぶ。
\item pro-対象
$$
(S/X) \longrightarrow \mathcal{D}', \quad
(s: X' \to X) \longmapsto F(X')
$$
が本質的に定値であるとき、{\it 左導来関手 $LF$ は $X$ で定義される}という。
この場合、$\mathcal{D}'$ における値 $Y$ を {\it $X$ における $LF$ の値}と呼ぶ。
\end{enumerate}
記法を濫用して、これらの値を単に $RF(X)$ または $LF(X)$ と書くことが多い。
\end{definition}

\noindent
$RF$ が定義される対象からなる $\mathcal{D}$ の充満部分圏は三角部分圏となり、
$RF$ はこの部分圏上で $S$ の射を同型に移す関手を定めることが後で分かる。

\begin{lemma}
\label{derived-lemma-derived-functor}
Situation \ref{derived-situation-derived-functor} の仮定と記法を用いる。
$f : X \to Y$ を $\mathcal{D}$ の射とする。
\begin{enumerate}
\item $RF$ が $X$ と $Y$ で定義されるならば、値の間に一意な射
$RF(f) : RF(X) \to RF(Y)$ が存在し、$s, s' \in S$ を持つ任意の可換図式
$$
\xymatrix{
X \ar[d]_f \ar[r]_s & X' \ar[d]^{f'} \\
Y \ar[r]^{s'} & Y'
}
$$
に対して図式
$$
\xymatrix{
F(X) \ar[d] \ar[r] & F(X') \ar[d] \ar[r] & RF(X) \ar[d] \\
F(Y) \ar[r] & F(Y') \ar[r] & RF(Y)
}
$$
が可換となる。
\item $LF$ が $X$ と $Y$ で定義されるならば、値の間に一意な射
$LF(f) : LF(X) \to LF(Y)$ が存在し、$s, s'$ が $S$ に属する任意の可換図式
$$
\xymatrix{
X' \ar[d]_{f'} \ar[r]_s & X \ar[d]^f \\
Y' \ar[r]^{s'} & Y
}
$$
に対して図式
$$
\xymatrix{
LF(X) \ar[d] \ar[r] & F(X') \ar[d] \ar[r] & F(X) \ar[d] \\
LF(Y) \ar[r] & F(Y') \ar[r] & F(Y)
}
$$
が可換となる。
\end{enumerate}
\end{lemma}

\begin{proof}
(1) は、余極限
$$
RF(X) = \colim_{s : X \to X'} F(X')
\quad\text{and}\quad
RF(Y) = \colim_{s' : Y \to Y'} F(Y')
$$
が存在することだけを仮定しても成り立つ。実際、余極限の間の射
$RF(X) \to RF(Y)$ を与えることは、$\Ob(X/S)$ の各 $s : X \to X'$ に対し、
圏 $X/S$ の射と両立する射 $F(X') \to RF(Y)$ を与えることと同じである。
その射を得るには、MS2 により可能な、$s, s'$ が $S$ に属する可換図式
$$
\xymatrix{
X \ar[d]_f \ar[r]_s & X' \ar[d]^{f'} \\
Y \ar[r]^{s'} & Y'
}
$$
を選び、$F(X') \to RF(Y)$ を合成 $F(X') \to F(Y') \to RF(Y)$ に等しいと置く。
これが上の図式の選択に依存しないことは MS3 を用いて分かる。詳細は省略する。
(2) の証明は双対である。
\end{proof}

\begin{lemma}
\label{derived-lemma-derived-inverts}
Situation \ref{derived-situation-derived-functor} の仮定と記法を用いる。
$s : X \to Y$ を $S$ の元とする。
\begin{enumerate}
\item $RF$ が $X$ で定義されるための必要十分条件は $Y$ で定義されることである。
この場合、値の間の写像 $RF(s) : RF(X) \to RF(Y)$ は同型である。
\item $LF$ が $X$ で定義されるための必要十分条件は $Y$ で定義されることである。
この場合、値の間の写像 $LF(s) : LF(X) \to LF(Y)$ は同型である。
\end{enumerate}
\end{lemma}

\begin{proof}
省略する。
\end{proof}

\begin{lemma}
\label{derived-lemma-derived-shift}
\begin{slogan}
導来関手はシフトと両立する
\end{slogan}
Situation \ref{derived-situation-derived-functor} の仮定と記法を用いる。
$X$ を $\mathcal{D}$ の対象、$n \in \mathbf{Z}$ とする。
\begin{enumerate}
\item $RF$ が $X$ で定義されるための必要十分条件は $X[n]$ で定義されることである。
この場合、値の間に標準同型 $RF(X)[n]= RF(X[n])$ がある。
\item $LF$ が $X$ で定義されるための必要十分条件は $X[n]$ で定義されることである。
この場合、値の間に標準同型 $LF(X)[n] \to LF(X[n])$ がある。
\end{enumerate}
\end{lemma}

\begin{proof}
省略する。
\end{proof}

\begin{lemma}
\label{derived-lemma-2-out-of-3-defined}
Situation \ref{derived-situation-derived-functor} の仮定と記法を用いる。
$(X, Y, Z, f, g, h)$ を $\mathcal{D}$ の完全三角とする。$RF$ が $X, Y, Z$ の
三つのうち二つで定義されるならば、残る一つでも定義される。さらに、この場合
$$
(RF(X), RF(Y), RF(Z), RF(f), RF(g), RF(h))
$$
は $\mathcal{D}'$ の完全三角である。$LF$ についても同様である。
\end{lemma}

\begin{proof}
$RF$ が $X, Y$ で定義され、その値が $A, B$ であるとする。補題
\ref{derived-lemma-derived-functor} の誘導射を $RF(f) : A \to B$ とする。
$\mathcal{D}'$ の完全三角 $(A, B, C, RF(f), b, c)$ を選べる。
$C$ が $Z$ における $RF$ の値であると主張する。

\medskip\noindent
これを見るため、Categories, Definition
\ref{categories-definition-essentially-constant-diagram} のような射
$\alpha : A \to F(X')$ が存在する $S$ の元 $s : X \to X'$ を選ぶ。MS2 により、
$s' \in S$ を持つ可換図式
$$
\xymatrix{
X \ar[d]_f \ar[r]_s & X' \ar[d]^{f'} \\
Y \ar[r]^{s'} & Y'
}
$$
を選べる。$Y/S$ はフィルター圏なので、$s'$ を $S$ のある
$s'' : Y \to Y''$ で置き換えた後、Categories, Definition
\ref{categories-definition-essentially-constant-diagram} のような射
$\beta : B \to F(Y')$ が存在すると仮定できる。図示すると
$$
\xymatrix{
A \ar[d]_{RF(f)} \ar[r]_-\alpha &
F(X') \ar[r] \ar[d]^{F(f')} &
A \ar[d]^{RF(f)} \\
B \ar[r]^-\beta & F(Y') \ar[r] & B
}
$$
である。左の正方形は可換とは限らないが、外側の正方形と右の正方形は可換である。
% P02-E0029: frozen source orients the ind-system index as s': Y' -> Y rather than Y -> Y'.
ind-対象 $\{F(Y')\}_{s' : Y' \to Y}$ が本質的に定値であるという仮定は、
$S$ の元 $s'' : Y \to Y''$ と、$s'' = h \circ s'$ を満たす射
$h : Y' \to Y''$ が存在し、$F(h)$ が合成
$F(Y') \to B \to F(Y') \to F(Y'')$ に等しいことを意味する。したがって、
$Y'$ を $Y''$ で、$\beta$ を $F(h) \circ \beta$ で置き換えると、この図式は
可換になる（矢印を用いた直接計算による）。

\medskip\noindent
MS6 を用いて、$s'' \in S$ を持つ三角の射
$$
(s, s', s'') : (X, Y, Z, f, g, h) \longrightarrow (X', Y', Z', f', g', h')
$$
を選ぶ。TR3 により、三角の射
$$
(\alpha, \beta, \gamma) :
(A, B, C, RF(f), b, c)
\longrightarrow
(F(X'), F(Y'), F(Z'), F(f'), F(g'), F(h'))
$$
を選ぶ。

\medskip\noindent
補題 \ref{derived-lemma-derived-inverts} により、$RF(Z')$ が定義され、矢印
$\gamma : C \to F(Z')$ が同型 $C \to RF(Z')$ を誘導することを示せば十分である。
実際、このとき三角の同型
$$
(A, B, C, RF(f), b, c)
\longrightarrow
(RF(X'), RF(Y'), RF(Z'), RF(f'), RF(g'), RF(h'))
$$
が得られ、TR1 により終域は完全三角である。補題 \ref{derived-lemma-limit-triangles} の
三角からなる圏 $\mathcal{I}$
$$
\mathcal{I} =
\{(t, t', t'') : (X', Y', Z', f', g', h') \to (X'', Y'', Z'', f'', g'', h'')
\mid (t, t', t'') \in S\}
$$
を考える。系 $F(Z'')$ が圏 $Z'/S$ 上で本質的に定値であることを示すことは、
$\mathcal{I} \to Z'/S$ が余終であるため、$F(Z'')$ の系が $\mathcal{I}$ 上で
本質的に定値であることを示すことと同値である。Categories, Lemma
\ref{categories-lemma-cofinal-essentially-constant} を参照せよ（余終性は補題
\ref{derived-lemma-limit-triangles} で証明されている）。$\mathcal{D}'$ の任意の対象 $W$ に対して図式
$$
\xymatrix{
\colim_\mathcal{I} \Mor_{\mathcal{D}'}(W, F(X'')) \ar[d] &
\Mor_{\mathcal{D}'}(W, A) \ar[l] \ar[d] \\
\colim_\mathcal{I} \Mor_{\mathcal{D}'}(W, F(Y'')) \ar[d] &
\Mor_{\mathcal{D}'}(W, B) \ar[d] \ar[l] \\
\colim_\mathcal{I} \Mor_{\mathcal{D}'}(W, F(Z'')) \ar[d] &
\Mor_{\mathcal{D}'}(W, C) \ar[d] \ar[l] \\
\colim_\mathcal{I} \Mor_{\mathcal{D}'}(W, F(X''[1])) \ar[d] &
\Mor_{\mathcal{D}'}(W, A[1]) \ar[d] \ar[l] \\
\colim_\mathcal{I} \Mor_{\mathcal{D}'}(W, F(Y''[1])) &
\Mor_{\mathcal{D}'}(W, B[1]) \ar[l]
}
$$
を考える。横の矢印は $(\alpha, \beta, \gamma)$ との合成によって与えられる。
フィルター余極限は完全なので（Algebra, Lemma
\ref{algebra-lemma-directed-colimit-exact}）、左の列は完全である。したがって $5$ 項補題
（Homology, Lemma \ref{homology-lemma-five-lemma}）により写像
$$
\colim_\mathcal{I} \Mor_{\mathcal{D}'}(W, F(Z''))
\longrightarrow
\Mor_{\mathcal{D}'}(W, C)
$$
は全単射である。Categories, Lemma
\ref{categories-lemma-characterize-essentially-constant-ind} (4) により、
$F(Z'')$ は値 $C$ を持つ $\mathcal{I}$ 上の本質的に定値な系であると結論する。
\end{proof}

\begin{lemma}
\label{derived-lemma-direct-sum-defined}
Situation \ref{derived-situation-derived-functor} の仮定と記法を用いる。
$X, Y$ を $\mathcal{D}$ の対象とする。
\begin{enumerate}
\item $RF$ が $X$ と $Y$ で定義されるならば、$RF$ は $X \oplus Y$ で定義される。
\item $\mathcal{D}'$ がカルービ圏で、$RF$ が $X \oplus Y$ で定義されるならば、
$RF$ は $X$ と $Y$ の双方で定義される。
\end{enumerate}
いずれの場合も $RF(X \oplus Y) = RF(X) \oplus RF(Y)$ である。
$LF$ についても同様である。
\end{lemma}

\begin{proof}
$RF$ が $X$ と $Y$ で定義されるならば、完全三角
$X \to X \oplus Y \to Y \to X[1]$（補題 \ref{derived-lemma-split}）と補題
\ref{derived-lemma-2-out-of-3-defined} により、$RF$ は $X \oplus Y$ で定義され、完全三角
$RF(X) \to RF(X \oplus Y) \to RF(Y) \to RF(X)[1]$ が得られる。
これに補題 \ref{derived-lemma-split} をもう一度適用すると、
$RF(X \oplus Y) = RF(X) \oplus RF(Y)$ を得る。これで (1) と最後の主張が示された。

\medskip\noindent
逆に、$RF$ が $X \oplus Y$ で定義され、$\mathcal{D}'$ がカルービ圏であると仮定する。
$S$ は飽和系なので、$S$ は加法的局所化関手
$Q : \mathcal{D} \to S^{-1}\mathcal{D}$ によって可逆になる矢印の集合である。
Categories, Lemma \ref{categories-lemma-what-gets-inverted} を参照せよ。したがって、
$S$ の任意の $s : X \to X'$ および $s' : Y \to Y'$ に対し、射
$s \oplus s' : X \oplus Y \to X' \oplus Y'$ は $S$ の元である。これにより関手
$$
X/S \times Y/S \longrightarrow (X \oplus Y)/S
$$
を得る。圏 $X/S, Y/S, (X \oplus Y)/S$ はフィルター圏であることを思い出そう
（Categories, Remark
\ref{categories-remark-left-localization-morphisms-colimit}）。Categories, Lemma
\ref{categories-lemma-essentially-constant-over-product} により、$X/S \times Y/S$ は
フィルター圏であり、$F|_{X/S} : X/S \to \mathcal{D}'$（それぞれ
% P02-E0030: frozen source uses the undefined G in both Y/S restrictions instead of F.
$G|_{Y/S} : Y/S \to \mathcal{D}'$）が本質的に定値であるための必要十分条件は、
$F|_{X/S} \circ \text{pr}_1 : X/S \times Y/S \to \mathcal{D}'$
（それぞれ $G|_{Y/S} \circ \text{pr}_2 : X/S \times Y/S \to \mathcal{D}'$）が
本質的に定値であることである。下で上の表示した関手が余終であることを示す。
したがって Categories, Lemma \ref{categories-lemma-cofinal-essentially-constant} により、
$F|_{(X \oplus Y)/S}$ が本質的に定値ならば、
$F|_{X/S} \circ \text{pr}_1 \oplus F|_{Y/S} \circ \text{pr}_2 :
X/S \times Y/S \to \mathcal{D}'$ が本質的に定値である。
Homology, Lemma
\ref{homology-lemma-direct-sum-from-product-essentially-constant} により
（ここで $\mathcal{D}'$ がカルービ圏であることを用いる）、
$F|_{X/S} \circ \text{pr}_1 \oplus F|_{Y/S} \circ \text{pr}_2$ が
本質的に定値ならば、$F|_{X/S} \circ \text{pr}_1$ と
$F|_{Y/S} \circ \text{pr}_2$ は本質的に定値である。したがって $RF$ は $X$ と
$Y$ で定義される。

\medskip\noindent
表示した関手が余終であることを証明する。任意の $S$ の元
$t : X \oplus Y \to Z$ を取る。MS2 を用いると、図式
$$
\xymatrix{
X \ar[d]^s & X \oplus Y \ar[d] \ar[l] \ar[r] & Y \ar[d]_{s'} \\
X' & Z \ar[l] \ar[r] & Y'
}
$$
が可換となるような射 $Z \to X'$、$Z \to Y'$ と $S$ の元
$s : X \to X'$、$s' : Y \to Y'$ を見いだせる。したがって
$(X \oplus Y)/S$ に写像 $Z \to X' \oplus Y'$ が存在する。すなわち、
Categories, Definition \ref{categories-definition-cofinal} の (1) が成り立つ。
(2) を証明するには、$t : X \oplus Y \to Z$ と
$(X \oplus Y)/S$ の射
$s_i \oplus s'_i : Z \to X'_i \oplus Y'_i$、$i = 1, 2$ が与えられたとき、
$a \circ s_1 = b \circ s_2$ および $c \circ s'_1 = d \circ s'_2$ を満たす
$S$ の射 $a : X'_1 \to X'$、$b : X'_2 \to X'$、
$c : Y'_1 \to Y'$、$d : Y'_2 \to Y'$ を見いだせることを示せば十分である。
そのため、まず任意の $X'$、$Y'$ と $S$ の写像 $a, b, c, d$ を選ぶ。
$X/S$ と $Y/S$ はフィルター圏なので、これは可能である。このとき二つの写像
$a \circ s_1, b \circ s_2 : Z \to X'$ は $S^{-1}\mathcal{D}$ で等しくなる。
したがって、これらを等しくする $S$ の射 $X' \to X''$ を見いだせる。同様に、
$c \circ s'_1$ と $d \circ s'_2$ を等しくする $S$ の射 $Y' \to Y''$ を
見いだせる。$X'$ を $X''$ で、$Y'$ を $Y''$ で置き換えれば、
$a \circ s_1 = b \circ s_2$ および $c \circ s'_1 = d \circ s'_2$ を得る。

\medskip\noindent
% P02-E0031: frozen source has singular “The proof” with plural “are dual”.
$LF$ に対する対応する主張の証明は双対である。
\end{proof}

\begin{proposition}
\label{derived-proposition-derived-functor}
Situation \ref{derived-situation-derived-functor} の仮定と記法を用いる。
\begin{enumerate}
\item $RF$ が定義される対象からなる $\mathcal{D}$ の充満部分圏
$\mathcal{E}$ は、$\mathcal{D}$ の厳密充満三角部分圏である。
\item 三角圏の完全関手 $RF : \mathcal{E} \longrightarrow \mathcal{D}'$ を得る。
\item 始域または終域が $\mathcal{E}$ に属する $S$ の元は、$\mathcal{E}$ の射である。
\item $S_\mathcal{E} = \text{Arrows}(\mathcal{E}) \cap S$ の任意の元は
$RF$ によって同型に移される。
\item 集合 $S_\mathcal{E}$ は三角構造と両立する $\mathcal{E}$ の飽和乗法系である。
\item 関手 $S_\mathcal{E}^{-1}\mathcal{E} \to S^{-1}\mathcal{D}$ は
三角圏の充満忠実な完全関手である。
\item 完全関手
$$
RF : S_\mathcal{E}^{-1}\mathcal{E} \longrightarrow \mathcal{D}'.
$$
を得る。
\item $\mathcal{D}'$ がカルービ圏ならば、$\mathcal{E}$ は
$\mathcal{D}$ の飽和三角部分圏である。
\end{enumerate}
$LF$ についても同様の結果が成り立つ。
\end{proposition}

\begin{proof}
$S$ は飽和なので、すべての同型を含む（Categories, Definition
\ref{categories-definition-saturated-multiplicative-system} に続く remark を参照）。
したがって (1) は補題 \ref{derived-lemma-derived-inverts}、
\ref{derived-lemma-2-out-of-3-defined}、\ref{derived-lemma-derived-shift} から従う。
(2) は補題 \ref{derived-lemma-derived-functor}、\ref{derived-lemma-derived-shift}、
\ref{derived-lemma-2-out-of-3-defined} から得られる。(3) は補題
\ref{derived-lemma-derived-inverts} から従う。(4) も補題
\ref{derived-lemma-derived-inverts} から従う。(5) は定義と (3) から従う。
(6) の充満忠実性は (3) と定義から従う。関手
$S_\mathcal{E}^{-1}\mathcal{E} \to S^{-1}\mathcal{D}$ が完全であることは、
$S_\mathcal{E}^{-1}\mathcal{E}$ の三角が完全であるための必要十分条件が、
$\mathcal{E}$ の完全三角の像と同型であることだから従う。命題
\ref{derived-proposition-construct-localization} の証明を参照せよ。
$RF : \mathcal{E} \to \mathcal{D}'$ が完全関手
$S_\mathcal{E}^{-1}\mathcal{E} \to \mathcal{D}'$ を経由して分解することは、補題
\ref{derived-lemma-universal-property-localization} から従う。最後に、(8) は補題
\ref{derived-lemma-direct-sum-defined} から従う。
\end{proof}

\noindent
命題 \ref{derived-proposition-derived-functor} により、$RF$ は
$S^{-1}\mathcal{D}$ の極大な厳密充満三角部分圏上にあり、その三角圏上の完全関手である。
図示すると
$$
\xymatrix{
\mathcal{D} \ar[d]_Q \ar[rrr]_F & & & \mathcal{D}' \\
S^{-1}\mathcal{D} & &
S_\mathcal{E}^{-1}\mathcal{E}
\ar[ll]_{\text{fully faithful}}^{\text{exact}} \ar[ur]_{RF}
}
$$
である。

\begin{definition}
\label{derived-definition-everywhere-defined}
Situation \ref{derived-situation-derived-functor} において、$RF$ が $\mathcal{D}$ のすべての
対象で定義されるとき、$F$ は {\it 右導来可能}、または $RF$ は
{\it 至る所で定義される}という。$LF$ が $\mathcal{D}$ のすべての対象で
定義されるとき、$F$ は {\it 左導来可能}、または $LF$ は
{\it 至る所で定義される}という。
\end{definition}

\noindent
この場合、右（それぞれ左）導来関手
\begin{equation}
\label{derived-equation-everywhere}
RF : S^{-1}\mathcal{D} \longrightarrow \mathcal{D}',
\quad\text{(resp. }
LF : S^{-1}\mathcal{D} \longrightarrow \mathcal{D}'),
\end{equation}
を得る。命題 \ref{derived-proposition-derived-functor} を参照せよ。興味深い状況の大半では、
$RF \circ Q$ は $F$ と等しくない。実際、標準写像 $F(X) \to RF(X)$ が決して
同型にならないことさえあり得る。実際の応用ではこれは起こらない。というのも、
実際には、三角圏 $\mathcal{D}$ の適切に選んだ対象で $F$ を評価することによって
$RF$ を計算できると示すことでしか、$F$ が右導来可能であることを証明できないからである。
これは次の定義の根拠となる。

\begin{definition}
\label{derived-definition-computes}
Situation \ref{derived-situation-derived-functor} において、
\begin{enumerate}
\item $RF$ が $X$ で定義され、標準写像 $F(X) \to RF(X)$ が同型であるとき、
$\mathcal{D}$ の対象 $X$ は $RF$ を {\it 計算する}という。
\item $LF$ が $X$ で定義され、標準写像 $LF(X) \to F(X)$ が同型であるとき、
$\mathcal{D}$ の対象 $X$ は $LF$ を {\it 計算する}という。
\end{enumerate}
\end{definition}

\begin{lemma}
\label{derived-lemma-computes-shift}
Situation \ref{derived-situation-derived-functor} の仮定と記法を用いる。
$X$ を $\mathcal{D}$ の対象、$n \in \mathbf{Z}$ とする。
\begin{enumerate}
\item $X$ が $RF$ を計算するための必要十分条件は $X[n]$ が $RF$ を計算することである。
\item $X$ が $LF$ を計算するための必要十分条件は $X[n]$ が $LF$ を計算することである。
\end{enumerate}
\end{lemma}

\begin{proof}
省略する。
\end{proof}

\begin{lemma}
\label{derived-lemma-2-out-of-3-computes}
Situation \ref{derived-situation-derived-functor} の仮定と記法を用いる。
$(X, Y, Z, f, g, h)$ を $\mathcal{D}$ の完全三角とする。
$X, Y$ が $RF$ を計算するならば、$Z$ もこれを計算する。
$LF$ についても同様である。
\end{lemma}

\begin{proof}
補題 \ref{derived-lemma-2-out-of-3-defined} により、$RF$ は $Z$ で定義され、$RF$ を
この三角に適用すると完全三角が得られる。完全三角の射
$$
\xymatrix{
(F(X), F(Y), F(Z), F(f), F(g), F(h)) \ar[d] \\
(RF(X), RF(Y), RF(Z), RF(f), RF(g), RF(h))
}
$$
を考える。三つの写像のうち二つが同型なので、残る一つも同型である。
\end{proof}

\begin{lemma}
\label{derived-lemma-direct-sum-computes}
Situation \ref{derived-situation-derived-functor} の仮定と記法を用いる。
$X, Y$ を $\mathcal{D}$ の対象とする。$X \oplus Y$ が $RF$ を計算するならば、
$X$ と $Y$ も $RF$ を計算する。$LF$ についても同様である。
\end{lemma}

\begin{proof}
$X \oplus Y$ が $RF$ を計算するならば、
$RF(X \oplus Y) = F(X) \oplus F(Y)$ である。補題
\ref{derived-lemma-direct-sum-defined} の証明で、関手
$X/S \times Y/S \to (X \oplus Y)/S$, $(s, s') \mapsto s \oplus s'$ が
余終であることを見た。したがって Categories, Lemma
\ref{categories-lemma-cofinal-essentially-constant} と Categories, Lemma
\ref{categories-lemma-characterize-essentially-constant-ind} の特徴づけ (4) により、
$\mathcal{D}'$ の任意の対象 $W$ に対して写像
$$
\Hom_{\mathcal{D}'}(F(X \oplus Y), W)
\longrightarrow
\colim_{s : X \to X', s' : Y \to Y'} \Hom_{\mathcal{D}'}(F(X' \oplus Y'), W)
$$
は全単射である。この矢印は両辺の直和分解と明らかに両立するので、写像
$$
\Hom_{\mathcal{D}'}(F(X), W)
\longrightarrow
\colim_{s : X \to X'} \Hom_{\mathcal{D}'}(F(X'), W)
$$
も全単射である（細部は省略する）。したがって Categories, Lemma
\ref{categories-lemma-characterize-essentially-constant-ind} により、$RF$ は $X$ で
定義され、その値は $F(X)$ である。$Y$ についても同様である。
\end{proof}

\begin{lemma}
\label{derived-lemma-existence-computes}
Situation \ref{derived-situation-derived-functor} の仮定と記法を用いる。
\begin{enumerate}
\item すべての対象 $X \in \Ob(\mathcal{D})$ に対して、$X'$ が $RF$ を計算するような
$S$ の矢印 $s : X \to X'$ が存在するならば、$RF$ は至る所で定義される。
\item すべての対象 $X \in \Ob(\mathcal{D})$ に対して、$X'$ が $LF$ を計算するような
$S$ の矢印 $s : X' \to X$ が存在するならば、$LF$ は至る所で定義される。
\end{enumerate}
\end{lemma}

\begin{proof}
定義から明らかである。
\end{proof}

\begin{lemma}
\label{derived-lemma-find-existence-computes}
Situation \ref{derived-situation-derived-functor} の仮定と記法を用いる。部分集合
$\mathcal{I} \subset \Ob(\mathcal{D})$ であって、
\begin{enumerate}
\item すべての $X \in \Ob(\mathcal{D})$ に対して、$X' \in \mathcal{I}$ を満たす
$S$ の射 $s : X \to X'$ が存在し、
\item $X, X' \in \mathcal{I}$ を満たす $S$ の任意の矢印 $s : X \to X'$ に対し、
写像 $F(s) : F(X) \to F(X')$ が同型である
\end{enumerate}
ものが存在するとする。このとき $RF$ は至る所で定義され、すべての
$X \in \mathcal{I}$ は $RF$ を計算する。双対的に、部分集合
$\mathcal{P} \subset \Ob(\mathcal{D})$ であって、
\begin{enumerate}
\item すべての $X \in \Ob(\mathcal{D})$ に対して、$X' \in \mathcal{P}$ を満たす
$S$ の射 $s : X' \to X$ が存在し、
\item $X, X' \in \mathcal{P}$ を満たす $S$ の任意の矢印 $s : X \to X'$ に対し、
写像 $F(s) : F(X) \to F(X')$ が同型である
\end{enumerate}
ものが存在するとする。このとき $LF$ は至る所で定義され、すべての
$X \in \mathcal{P}$ は $LF$ を計算する。
\end{lemma}

\begin{proof}
$X$ を $\mathcal{D}$ の対象とする。仮定 (1) により、$X' \in \mathcal{I}$ を満たす
$S$ の矢印 $s : X \to X'$ は圏 $X/S$ で余終である。仮定 (2) により、
$F$ はこの余終部分圏上で定値である。したがって、$X' \in \mathcal{I}$ を満たす
$S$ の任意の $s : X \to X'$ に対し、
$F : (X/S) \to \mathcal{D}'$ は値 $F(X')$ を持つ本質的に定値な関手である。
\end{proof}

\begin{lemma}
\label{derived-lemma-compose-derived-functors-general}
$\mathcal{A}, \mathcal{B}, \mathcal{C}$ を三角圏とする。$S$、それぞれ $S'$ を、
$\mathcal{A}$、それぞれ $\mathcal{B}$ の三角構造と両立する飽和乗法系とする。
$F : \mathcal{A} \to \mathcal{B}$ および $G : \mathcal{B} \to \mathcal{C}$ を
完全関手とする。$F$ と局所化関手との合成を
$F' : \mathcal{A} \to (S')^{-1}\mathcal{B}$ と書く。
\begin{enumerate}
\item $RF'$、$RG$、$R(G \circ F)$ が至る所で定義されるならば、関手の標準自然変換
$t : R(G \circ F) \longrightarrow RG \circ RF'$ が存在する。
\item $LF'$、$LG$、$L(G \circ F)$ が至る所で定義されるならば、関手の標準自然変換
$t : LG \circ LF' \to L(G \circ F)$ が存在する。
\end{enumerate}
\end{lemma}

\begin{proof}
この証明では慎重に議論する。そこで導来関手を関手
$$
RF' : S^{-1}\mathcal{A} \to (S')^{-1}\mathcal{B}, \quad
R(G \circ F) : S^{-1}\mathcal{A} \to \mathcal{C}, \quad
RG : (S')^{-1}\mathcal{B} \to \mathcal{C}.
$$
と考える。局所化関手を
$Q_A : \mathcal{A} \to S^{-1}\mathcal{A}$ および
$Q_B : \mathcal{B} \to (S')^{-1}\mathcal{B}$ と書く。このとき
$F' = Q_B \circ F$ である。$\mathcal{B}$ の任意の対象 $Y$ に対して標準写像
$$
G(Y) \longrightarrow RG(Q_B(Y))
$$
があることに注意する。言い換えれば、関手の自然変換
$t' : G \to RG \circ Q_B$ がある。$X$ を $\mathcal{A}$ の対象とする。このとき
\begin{align*}
R(G \circ F)(Q_A(X))
& = \colim_{s : X \to X' \in S} G(F(X')) \\
& \xrightarrow{t'} \colim_{s : X \to X' \in S} RG(Q_B(F(X'))) \\
& = \colim_{s : X \to X' \in S} RG(F'(X')) \\
& = RG(\colim_{s : X \to X' \in S} F'(X')) \\
& = RG(RF'(X)).
\end{align*}
系 $F'(X')$ は圏 $(S')^{-1}\mathcal{B}$ において本質的に定値である。したがって、
上の図式の第三の等式で余極限を関手 $RG$ の内側に移せる。Categories, Lemma
\ref{categories-lemma-image-essentially-constant} とその証明を参照せよ。これが関手の
自然変換を定めることの証明は省略する。左導来関手の場合も同様である。
\end{proof}




\section{導来圏上の導来関手}
\label{derived-section-derived-functors-classical}

\noindent
実際には、導来関手はアーベル圏の間の加法的関手が与えられたときに
現れることが最も多い。

\begin{situation}
\label{derived-situation-classical}
ここで $F : \mathcal{A} \to \mathcal{B}$ はアーベル圏の間の加法的関手である。
これは完全関手
$$
F : K(\mathcal{A}) \to K(\mathcal{B}), \quad
K^{+}(\mathcal{A}) \to K^{+}(\mathcal{B}), \quad
K^{-}(\mathcal{A}) \to K^{-}(\mathcal{B}).
$$
を誘導する。補題 \ref{derived-lemma-additive-exact-homotopy-category} を参照せよ。
また、$F$ と局所化関手 $K(\mathcal{B}) \to D(\mathcal{B})$ などとの合成
$K(\mathcal{A}) \to D(\mathcal{B})$、
$K^{+}(\mathcal{A}) \to D^{+}(\mathcal{B})$、および
$K^{-}(\mathcal{A}) \to D^-(\mathcal{B})$ も $F$ と書く。この状況から、
以下で考察する四つの導来関手が得られる。
\begin{enumerate}
\item 乗法系 $\text{Qis}(\mathcal{A})$ に関する
$F : K(\mathcal{A}) \to D(\mathcal{B})$
の右導来関手。
\item 乗法系 $\text{Qis}^{+}(\mathcal{A})$ に関する
$F : K^{+}(\mathcal{A}) \to D^{+}(\mathcal{B})$
の右導来関手。
\item 乗法系 $\text{Qis}(\mathcal{A})$ に関する
$F : K(\mathcal{A}) \to D(\mathcal{B})$
の左導来関手。
\item 乗法系 $\text{Qis}^-(\mathcal{A})$ に関する
$F : K^{-}(\mathcal{A}) \to D^{-}(\mathcal{B})$
の左導来関手。
\end{enumerate}
これらはいずれも Situation \ref{derived-situation-derived-functor} の例である。
\end{situation}

\noindent
生じ得る曖昧さの一部は、次の結果によって解消される。

\begin{lemma}
\label{derived-lemma-irrelevant}
Situation \ref{derived-situation-classical} の仮定のもとで、次が成り立つ。
\begin{enumerate}
\item $X$ を $K^{+}(\mathcal{A})$ の対象とする。
$K(\mathcal{A}) \to D(\mathcal{B})$ の右導来関手が $X$ で定義されるための
必要十分条件は、$K^{+}(\mathcal{A}) \to D^{+}(\mathcal{B})$ の右導来関手が
$X$ で定義されることである。さらに、その値は標準的に同型である。
\item $X$ を $K^{+}(\mathcal{A})$ の対象とする。このとき $X$ が
$K(\mathcal{A}) \to D(\mathcal{B})$ の右導来関手を計算するための
必要十分条件は、$X$ が
$K^{+}(\mathcal{A}) \to D^{+}(\mathcal{B})$ の右導来関手を計算することである。
\item $X$ を $K^{-}(\mathcal{A})$ の対象とする。
$K(\mathcal{A}) \to D(\mathcal{B})$ の左導来関手が $X$ で定義されるための
必要十分条件は、$K^{-}(\mathcal{A}) \to D^{-}(\mathcal{B})$ の左導来関手が
$X$ で定義されることである。さらに、その値は標準的に同型である。
\item $X$ を $K^{-}(\mathcal{A})$ の対象とする。このとき $X$ が
$K(\mathcal{A}) \to D(\mathcal{B})$ の左導来関手を計算するための
必要十分条件は、$X$ が
$K^{-}(\mathcal{A}) \to D^{-}(\mathcal{B})$ の左導来関手を計算することである。
\end{enumerate}
\end{lemma}

\begin{proof}
$X$ を $K^{+}(\mathcal{A})$ の対象とする。$K(\mathcal{A})$ における擬同型
$s : X \to X'$ を考える。補題 \ref{derived-lemma-complex-cohomology-bounded} により、
$X''$ が下に有界となる擬同型 $X' \to X''$ が存在する。したがって
$X/\text{Qis}^+(\mathcal{A})$ は $X/\text{Qis}(\mathcal{A})$ で余終である。
これより (1) は明らかである。(2) は (1) から直ちに従う。
(3) と (4) はそれぞれ (1) と (2) の双対である。
\end{proof}

\noindent
アーベル圏 $\mathcal{A}$ の対象 $A$ が与えられると、$A$ を次数ゼロに置いた複体
$$
A[0] = ( \ldots \to 0 \to A \to 0 \to \ldots )
$$
が得られる。したがって関手
$\mathcal{A} \to K(\mathcal{A})$, $A \mapsto A[0]$ が得られる。
しばらくの間、部分圏上で定義された関手を部分関手と呼ぶことにする。

\begin{definition}
\label{derived-definition-derived-functor}
Situation \ref{derived-situation-classical} の仮定のもとで、次のように定める。
\begin{enumerate}
\item {\it $F$ の右導来関手}とは、Situation \ref{derived-situation-classical} の
(1) と (2) に対応する部分関手 $RF$ のことである。
\item {\it $F$ の左導来関手}とは、Situation \ref{derived-situation-classical} の
(3) と (4) に対応する部分関手 $LF$ のことである。
\item $\mathcal{A}$ の対象 $A$ が {\it $F$ に関して右非輪状}、または
{\it $RF$ に関して非輪状}であるとは、$A[0]$ が $RF$ を計算することをいう。
\item $\mathcal{A}$ の対象 $A$ が {\it $F$ に関して左非輪状}、または
{\it $LF$ に関して非輪状}であるとは、$A[0]$ が $LF$ を計算することをいう。
\end{enumerate}
\end{definition}

\noindent
次のいくつかの補題は、十分多くの非輪状対象が存在するための判定条件を与える。

\begin{lemma}
\label{derived-lemma-subcategory-left-resolution}
$\mathcal{A}$ をアーベル圏とする。$\mathcal{P} \subset \Ob(\mathcal{A})$ を
$0$ を含む部分集合で、$\mathcal{A}$ の任意の対象が $\mathcal{P}$ の元の
商となるものとする。$a \in \mathbf{Z}$ とする。
\begin{enumerate}
\item $K^n = 0$ が $n > a$ に対して成り立つ $K^\bullet$ が与えられると、
擬同型 $P^\bullet \to K^\bullet$ であって、すべての $n$ に対して
$P^n \in \mathcal{P}$ かつ $P^n \to K^n$ が全射であり、
$n > a$ に対して $P^n = 0$ となるものが存在する。
\item $H^n(K^\bullet) = 0$ が $n > a$ に対して成り立つ $K^\bullet$ が
与えられると、擬同型 $P^\bullet \to K^\bullet$ であって、
すべての $n$ に対して $P^n \in \mathcal{P}$ であり、
$n > a$ に対して $P^n = 0$ となるものが存在する。
\end{enumerate}
\end{lemma}

\begin{proof}
(1) を証明する。次の帰納法の仮定 $IH_n$ を考える。
$P^j \in \mathcal{P}$, $j \geq n$ で、$j > a$ に対して $P^j = 0$ となるもの、
$j \geq n$ に対する写像 $d^j : P^j \to P^{j + 1}$、および
$j \geq n$ に対する全射 $\alpha^j : P^j \to K^j$ が存在し、図式
$$
\xymatrix{
& &
P^n \ar[d]^\alpha \ar[r] &
P^{n + 1} \ar[d]^\alpha \ar[r] &
P^{n + 2} \ar[d]^\alpha \ar[r] &
\ldots \\
\ldots \ar[r] &
K^{n - 1} \ar[r] &
K^n \ar[r] &
K^{n + 1} \ar[r] &
K^{n + 2} \ar[r] &
\ldots
}
$$
が可換であり、$j \geq n$ に対して $d^{j + 1} \circ d^j = 0$ であり、
$j > n$ に対して $\alpha$ が同型
$\Ker(d^j)/\Im(d^{j - 1}) \to H^j(K^\bullet)$ を誘導し、さらに
$\alpha : \Ker(d^n) \to \Ker(d_K^n)$ が全射である、と仮定する。
このとき $P^{n - 1}$ が $\mathcal{P}$ に属する全射
$$
P^{n - 1}
\longrightarrow
K^{n - 1} \times_{K^n} \Ker(d^n) =
K^{n - 1} \times_{\Ker(d_K^n)} \Ker(d^n)
$$
を選ぶ。これにより、上の図式を
$$
\xymatrix{
& P^{n - 1} \ar[d]^\alpha \ar[r] &
P^n \ar[d]^\alpha \ar[r] &
P^{n + 1} \ar[d]^\alpha \ar[r] &
P^{n + 2} \ar[d]^\alpha \ar[r] &
\ldots \\
\ldots \ar[r] &
K^{n - 1} \ar[r] &
K^n \ar[r] &
K^{n + 1} \ar[r] &
K^{n + 2} \ar[r] &
\ldots
}
$$
へ拡張できる。この選択によって $IH_{n - 1}$ が成り立つことは容易に確かめられる。

\medskip\noindent
次のようにして (1) の証明を終える。$j > a$ に対して $P^j = 0$ と取ればよいので、
$n = a + 1$ に対して $IH_n$ が成り立つ。したがって降下帰納法を進めると、
$\mathcal{P}$ の対象からなり、$n > a$ に対して $P^n = 0$ となる複体
$P^\bullet$ と、所望の項ごとに全射な擬同型
$\alpha : P^\bullet \to K^\bullet$ が得られる。

\medskip\noindent
(2) を証明する。仮定より射
$\tau_{\leq a}K^\bullet \to K^\bullet$
（Homology, Section \ref{homology-section-truncations}）は擬同型である。
(1) を適用して $P^\bullet \to \tau_{\leq a}K^\bullet$ を取る。
合成 $P^\bullet \to K^\bullet$ が所望の擬同型である。
\end{proof}

\begin{lemma}
\label{derived-lemma-subcategory-right-resolution}
$\mathcal{A}$ をアーベル圏とする。$\mathcal{I} \subset \Ob(\mathcal{A})$ を
$0$ を含む部分集合で、$\mathcal{A}$ の任意の対象が $\mathcal{I}$ の元の
部分対象となるものとする。$a \in \mathbf{Z}$ とする。
\begin{enumerate}
\item $K^n = 0$ が $n < a$ に対して成り立つ $K^\bullet$ が与えられると、
擬同型 $K^\bullet \to I^\bullet$ であって、すべての $n$ に対して
$K^n \to I^n$ が単射かつ $I^n \in \mathcal{I}$ であり、
$n < a$ に対して $I^n = 0$ となるものが存在する。
\item $H^n(K^\bullet) = 0$ が $n < a$ に対して成り立つ $K^\bullet$ が
与えられると、擬同型 $K^\bullet \to I^\bullet$ であって、
$I^n \in \mathcal{I}$ かつ $n < a$ に対して $I^n = 0$ となるものが存在する。
\end{enumerate}
\end{lemma}

\begin{proof}
本補題は補題 \ref{derived-lemma-subcategory-left-resolution} の双対である。
\end{proof}

\begin{lemma}
\label{derived-lemma-subcategory-right-acyclics}
Situation \ref{derived-situation-classical} の仮定のもとで、
$\mathcal{I} \subset \Ob(\mathcal{A})$ を次の性質を持つ部分集合とする。
\begin{enumerate}
\item $\mathcal{A}$ の任意の対象は $\mathcal{I}$ の元の部分対象である。
\item $\mathcal{A}$ の任意の短完全列 $0 \to P \to Q \to R \to 0$ に対し、
$P, Q \in \mathcal{I}$ ならば $R \in \mathcal{I}$ であり、
$0 \to F(P) \to F(Q) \to F(R) \to 0$ は完全である。
\end{enumerate}
このとき $\mathcal{I}$ のすべての対象は $RF$ に関して非輪状である。
\end{lemma}

\begin{proof}
$A \in \mathcal{I}$ を取る。$K^\bullet$ が下に有界である擬同型
$A[0] \to K^\bullet$ を取る。このとき $I^\bullet$ が下に有界で、各
$I^n \in \mathcal{I}$ となる擬同型 $K^\bullet \to I^\bullet$ を取れる。
補題 \ref{derived-lemma-subcategory-right-resolution}\footnote{(1) より
$\mathcal{I}$ は空でない。$P$ を $\mathcal{I}$ の元とする。このとき短完全列
$0 \to P \to P \to 0 \to 0$ と仮定 (2) により $0$ は
$\mathcal{I}$ に属する。したがって補題を適用できる。}を参照せよ。
したがって、これらの分解は圏
$A[0]/\text{Qis}^{+}(\mathcal{A})$ で余終である。ゆえに証明を終えるには、
$I^\bullet$ が下に有界かつ $I^n \in \mathcal{I}$ となる任意の擬同型
$A[0] \to I^\bullet$ に対し、$F(A)[0] \to F(I^\bullet)$ が擬同型であることを
示せば十分である。これを見るため、$n < n_0$ に対して $I^n = 0$ と仮定する。
もちろん $n_0 < 0$ としてよい。$n = n_0$ から始め、性質 (2) と完全列
$$
0 \to \Ker(d^n) \to I^n \to \Im(d^n) \to 0.
$$
を用いる帰納法により、$\Im(d^{n - 1}) = \Ker(d^n)$ と $\Im(d^{-1})$ が
$\mathcal{I}$ の元であることを示す。さらに性質 (2) により、複体
$$
0 \to F(I^{n_0}) \to F(I^{n_0 + 1}) \to \ldots \to F(I^{-1}) \to
F(\Im(d^{-1})) \to 0
$$
も完全である。完全列
$0 \to \Im(d^{-1}) \to I^0 \to I^0/\Im(d^{-1}) \to 0$
より、$I^0/\Im(d^{-1})$ は $\mathcal{I}$ の元である。
% P02-E0034: frozen source has “is an elements”; Japanese supplies grammatical agreement.
次に完全列 $0 \to A \to I^0/\Im(d^{-1}) \to \Im(d^0) \to 0$ より、
$\Im(d^0) = \Ker(d^1)$ は $\mathcal{I}$ の元である。その後も同様にして、
すべての $n > 0$ に対して $\Im(d^{n - 1}) = \Ker(d^n)$ が
$\mathcal{I}$ の元であることを示せる。上で述べた各短完全列に $F$ を適用し、
(2) を用いると、次が分かる。
% P02-E0033: frozen source says “is an isomorphism”; the formal source claim is preserved.
$F(A)[0] \to F(I^\bullet)$ は所望の同型である。
\end{proof}

\begin{lemma}
\label{derived-lemma-subcategory-left-acyclics}
Situation \ref{derived-situation-classical} の仮定のもとで、
$\mathcal{P} \subset \Ob(\mathcal{A})$ を次の性質を持つ部分集合とする。
\begin{enumerate}
\item $\mathcal{A}$ の任意の対象は $\mathcal{P}$ の元の商である。
\item $\mathcal{A}$ の任意の短完全列 $0 \to P \to Q \to R \to 0$ に対し、
$Q, R \in \mathcal{P}$ ならば $P \in \mathcal{P}$ であり、
$0 \to F(P) \to F(Q) \to F(R) \to 0$ は完全である。
\end{enumerate}
このとき $\mathcal{P}$ のすべての対象は $LF$ に関して非輪状である。
\end{lemma}

\begin{proof}
補題 \ref{derived-lemma-subcategory-right-acyclics} の証明の双対である。
\end{proof}







\section{高次導来関手}
\label{derived-section-higher-derived}

\noindent
次の簡単な補題は、右導来関手が「右へ移る」ことを示す。

\begin{lemma}
\label{derived-lemma-negative-vanishing}
$F : \mathcal{A} \to \mathcal{B}$ をアーベル圏の間の加法的関手とする。
$K^\bullet$ を $\mathcal{A}$ の複体とし、$a \in \mathbf{Z}$ とする。
\begin{enumerate}
\item すべての $i < a$ に対して $H^i(K^\bullet) = 0$ であり、$RF$ が
$K^\bullet$ で定義されるならば、すべての $i < a$ に対して
$H^i(RF(K^\bullet)) = 0$ である。
\item $RF$ が $K^\bullet$ と $\tau_{\leq a}K^\bullet$ で定義されるならば、
すべての $i \leq a$ に対して
$H^i(RF(\tau_{\leq a}K^\bullet)) = H^i(RF(K^\bullet))$ である。
\end{enumerate}
\end{lemma}

\begin{proof}
$K^\bullet$ が (1) の仮定を満たすとする。任意の擬同型
$s : K^\bullet \to L^\bullet$ を取る。$K^\bullet$ に関する仮定により、
$K^\bullet \to \tau_{\geq a}L^\bullet$ も擬同型である。したがって圏
$K^\bullet/\text{Qis}^{+}(\mathcal{A})$ において、$n < a$ に対して
$L^n = 0$ となる擬同型 $s : K^\bullet \to L^\bullet$ は余終である。
Categories, Lemma
\ref{categories-lemma-cofinal-essentially-constant}
により、$RF$ はこれらの $s$ に対する本質的に定値な ind-対象
$F(L^\bullet)$ の値である。これは、ある複体 $L^\bullet$ で
$n < a$ に対して $L^n = 0$ となるものに対し、
$\text{id} : RF(K^\bullet) \to RF(K^\bullet)$ が
$F(L^\bullet)$ を通じて分解することを意味する。
したがって $i < a$ に対して $H^i(RF(K^\bullet)) = 0$ である。

\medskip\noindent
(2) を証明するため、Remark \ref{derived-remark-truncation-distinguished-triangle} の完全三角
$$
\tau_{\leq a}K^\bullet \to K^\bullet \to \tau_{\geq a + 1}K^\bullet
\to (\tau_{\leq a}K^\bullet)[1]
$$
を用いる。補題 \ref{derived-lemma-2-out-of-3-defined} により、$RF$ は
$\tau_{\geq a + 1}K^\bullet$ でも定義され、$D(\mathcal{B})$ に完全三角
$$
RF(\tau_{\leq a}K^\bullet) \to RF(K^\bullet) \to RF(\tau_{\geq a + 1}K^\bullet)
\to RF(\tau_{\leq a}K^\bullet)[1]
$$
が得られる。(1) により $RF(\tau_{\geq a + 1}K^\bullet)$ の
次数 $< a + 1$ のコホモロジーは消える。この完全三角のコホモロジー長完全列から
所望の結論が従う。
\end{proof}

\begin{definition}
\label{derived-definition-higher-derived-functors}
$F : \mathcal{A} \to \mathcal{B}$ をアーベル圏の間の加法的関手とする。
$RF : D^{+}(\mathcal{A}) \to D^{+}(\mathcal{B})$ が至る所で定義されると仮定し、
$i \in \mathbf{Z}$ とする。{\it $F$ の第 $i$ 右導来関手 $R^iF$} とは関手
$$
R^iF = H^i \circ RF :
\mathcal{A}
\longrightarrow
\mathcal{B}
$$
のことである。
\end{definition}

\noindent
次の補題は、関手が左完全でない限り、その右導来関手を取ることには
実際あまり意味がないことを示す。

\begin{lemma}
\label{derived-lemma-left-exact-higher-derived}
$F : \mathcal{A} \to \mathcal{B}$ をアーベル圏の間の加法的関手とし、
$RF : D^{+}(\mathcal{A}) \to D^{+}(\mathcal{B})$ が至る所で定義されると仮定する。
\begin{enumerate}
\item $i < 0$ に対して $R^iF = 0$ である。
\item $R^0F$ は左完全である。
\item 写像 $F \to R^0F$ が同型であるための必要十分条件は、$F$ が左完全なことである。
\end{enumerate}
\end{lemma}

\begin{proof}
$A$ を $\mathcal{A}$ の対象とする。補題 \ref{derived-lemma-negative-vanishing} により、
% P02-E0032: frozen source has an unmatched opening parenthesis in this formula.
$i < 0$ に対して $H^i(RF(A[0]) = 0$ である。これで (1) が示された。

\medskip\noindent
$0 \to A \to B \to C \to 0$ を $\mathcal{A}$ の短完全列とする。
補題 \ref{derived-lemma-derived-canonical-delta-functor} により、
$D^{+}(\mathcal{A})$ に完全三角
$(A[0], B[0], C[0], a, b, c)$ が得られる。コホモロジー長完全列と、
上で示した $i < 0$ における消滅から、
$0 \to R^0F(A) \to R^0F(B) \to R^0F(C)$ は完全である。
したがって $R^0F$ は左完全である。またこれにより、$F \to R^0F$ が同型ならば
$F$ が左完全であることも示される。

\medskip\noindent
$F$ が左完全であると仮定する。$RF(A[0])$ は、擬同型
$s : A[0] \to K^\bullet$ に対する本質的に定値な系 $F(K^\bullet)$ の値であることを
思い出そう。したがって $R^0F(A)$ は、擬同型 $s : A[0] \to K^\bullet$ に対する
本質的に定値な系 $H^0(F(K^\bullet))$ の値である。Categories, Lemma
\ref{categories-lemma-image-essentially-constant} を参照せよ。
しかし $s : A[0] \to K^\bullet$ が擬同型ならば、
$A[0] \to \tau_{\geq 0}K^\bullet$ は擬同型である。したがって圏
$A[0]/\text{Qis}^{+}(\mathcal{A})$ において、$n < 0$ に対して
$K^n = 0$ となる擬同型 $s : A[0] \to K^\bullet$ は余終である。
Categories, Lemma \ref{categories-lemma-cofinal-essentially-constant}
により、そのような $s$ に制限してよい。さらに、そのような $s$ に対して列
$$
0 \to A \to K^0 \to K^1
$$
は完全である。$F$ は左完全なので
$0 \to F(A) \to F(K^0) \to F(K^1)$ も完全である。
したがって $F(A) \to H^0(F(K^\bullet))$ は同型であり、この系は実際に
$F(A)$ を値とする定値系である。ゆえに所望のとおり $R^0F = F$ である。
\end{proof}

\begin{lemma}
\label{derived-lemma-F-acyclic}
$F : \mathcal{A} \to \mathcal{B}$ をアーベル圏の間の加法的関手とし、
$RF : D^{+}(\mathcal{A}) \to D^{+}(\mathcal{B})$ が至る所で定義されると仮定する。
$A$ を $\mathcal{A}$ の対象とする。
\begin{enumerate}
\item $A$ が $F$ に関して右非輪状であるための必要十分条件は、
$F(A) \to R^0F(A)$ が同型であり、すべての $i > 0$ に対して
$R^iF(A) = 0$ となることである。
\item $F$ が左完全ならば、$A$ が $F$ に関して右非輪状であるための必要十分条件は、
すべての $i > 0$ に対して $R^iF(A) = 0$ となることである。
\end{enumerate}
\end{lemma}

\begin{proof}
$A$ が $F$ に関して右非輪状ならば $RF(A[0]) = F(A)[0]$ であり、特に
$F(A) \to R^0F(A)$ は同型で、$i \not = 0$ に対して $R^iF(A) = 0$ である。
逆に、$F(A) \to R^0F(A)$ が同型で、すべての $i > 0$ に対して
$R^iF(A) = 0$ ならば、補題 \ref{derived-lemma-left-exact-higher-derived} の (1) により
$F(A[0]) \to RF(A[0])$ は擬同型であり、したがって $A$ は非輪状である。
$F$ が左完全ならば、補題 \ref{derived-lemma-left-exact-higher-derived} により
$F = R^0F$ である。
\end{proof}

\begin{lemma}
\label{derived-lemma-F-acyclic-ses}
$F : \mathcal{A} \to \mathcal{B}$ をアーベル圏の間の左完全関手とし、
$RF : D^{+}(\mathcal{A}) \to D^{+}(\mathcal{B})$ が至る所で定義されると仮定する。
$0 \to A \to B \to C \to 0$ を $\mathcal{A}$ の短完全列とする。
\begin{enumerate}
\item $A$ と $C$ が $F$ に関して右非輪状ならば、$B$ もそうである。
\item $A$ と $B$ が $F$ に関して右非輪状ならば、$C$ もそうである。
\item $B$ と $C$ が $F$ に関して右非輪状であり、$F(B) \to F(C)$ が
全射ならば、$A$ は $F$ に関して右非輪状である。
\end{enumerate}
三つのいずれの場合にも
$$
0 \to F(A) \to F(B) \to F(C) \to 0
$$
は $\mathcal{B}$ の短完全列である。
\end{lemma}

\begin{proof}
補題 \ref{derived-lemma-derived-canonical-delta-functor} により、
$D^{+}(\mathcal{A})$ に完全三角
$(A[0], B[0], C[0], a, b, c)$ が得られる。$RF$ は完全関手であり、
$i < 0$ に対して $R^iF = 0$ かつ $R^0F = F$ なので
（補題 \ref{derived-lemma-left-exact-higher-derived}）、アーベル圏 $\mathcal{B}$ に完全な
コホモロジー列
$$
0 \to F(A) \to F(B) \to F(C) \to R^1F(A) \to \ldots
$$
が得られる。したがって補題 \ref{derived-lemma-F-acyclic} の非輪状対象の特徴づけから
結論が従う。
\end{proof}

\begin{lemma}
\label{derived-lemma-right-derived-delta-functor}
$F : \mathcal{A} \to \mathcal{B}$ をアーベル圏の間の加法的関手とし、
$RF : D^{+}(\mathcal{A}) \to D^{+}(\mathcal{B})$ が至る所で定義されると仮定する。
\begin{enumerate}
\item 関手 $R^iF$, $i \geq 0$ には、$\mathcal{A} \to \mathcal{B}$ の
$\delta$-関手の標準構造が備わる。Homology, Definition
\ref{homology-definition-cohomological-delta-functor} を参照せよ。
\item $\mathcal{A}$ の任意の対象が $F$ に関して右非輪状な対象の部分対象ならば、
$\{R^iF, \delta\}_{i \geq 0}$ は普遍 $\delta$-関手である。Homology, Definition
\ref{homology-definition-universal-delta-functor} を参照せよ。
\end{enumerate}
\end{lemma}

\begin{proof}
関手 $\mathcal{A} \to \text{Comp}^{+}(\mathcal{A})$,
$A \mapsto A[0]$ は完全である。関手
$\text{Comp}^{+}(\mathcal{A}) \to D^{+}(\mathcal{A})$
は $\delta$-関手である。補題 \ref{derived-lemma-derived-canonical-delta-functor} を参照せよ。
関手 $RF : D^{+}(\mathcal{A}) \to D^{+}(\mathcal{B})$ は完全である。
最後に、関手 $H^0 : D^{+}(\mathcal{B}) \to \mathcal{B}$ はホモロジー的関手である。
定義 \ref{derived-definition-unbounded-derived-category} を参照せよ。
したがって補題 \ref{derived-lemma-compose-delta-functor-homological} と
補題 \ref{derived-lemma-exact-compose-delta-functor} により、$\delta$-関手の構造が得られる。
(2) は Homology, Lemma \ref{homology-lemma-efface-implies-universal} と
補題 \ref{derived-lemma-F-acyclic} による非輪状対象の記述から従う。
\end{proof}

\begin{lemma}[Leray の非輪状性補題]
\label{derived-lemma-leray-acyclicity}
$F : \mathcal{A} \to \mathcal{B}$ をアーベル圏の間の加法的関手とする。
$A^\bullet$ を右 $F$-非輪状対象からなる下に有界な複体で、$RF$ が
$A^\bullet$ で定義されるものとする\footnote{たとえば
$RF : D^{+}(\mathcal{A}) \to D^{+}(\mathcal{B})$ が至る所で定義されればよい。}。
標準写像
$$
F(A^\bullet) \longrightarrow RF(A^\bullet)
$$
は $D^{+}(\mathcal{B})$ で同型である。すなわち、$A^\bullet$ は $RF$ を計算する。
\end{lemma}

\begin{proof}
$A^\bullet$ を右 $F$-非輪状対象からなる有界複体とする。$RF$ は
$A^\bullet$ で定義され、$F(A^\bullet) \to RF(A^\bullet)$ は
$D^+(\mathcal{B})$ で同型であることを示す。実際、非零な右 $F$-非輪状対象を
高々一つしか持たない複体については、定義 \ref{derived-definition-derived-functor} により
成り立つ。次に $n \not \in [a, b]$ に対して $A^n = 0$ と仮定する。
「素朴な」切断を用いると、複体の項ごとに分裂する短完全列
$$
0 \to \sigma_{\geq a + 1} A^\bullet \to A^\bullet \to
\sigma_{\leq a} A^\bullet \to 0
$$
が得られる。Homology, Section \ref{homology-section-truncations} を参照せよ。
したがって完全三角
$(\sigma_{\geq a + 1} A^\bullet, A^\bullet, \sigma_{\leq a} A^\bullet)$
が得られる。帰納法の仮定により、$RF$ は両端の二つの複体で定義され、それらの
複体は $RF$ を計算する。このとき補題 \ref{derived-lemma-2-out-of-3-computes} により、
中央の複体についても同じことが成り立つ。

\medskip\noindent
$A^\bullet$ を非輪状対象からなる下に有界な複体で、$RF$ が
$A^\bullet$ で定義されるものとする。
$F(A^\bullet) \to RF(A^\bullet)$ が $D^{+}(\mathcal{B})$ で同型であることを
示すには、すべての $i$ に対して
$H^i(F(A^\bullet)) \to H^i(RF(A^\bullet))$ が同型であることを示せば十分である。
$i$ を取る。複体の項ごとに分裂する短完全列
$$
0 \to \sigma_{\geq i + 2} A^\bullet \to A^\bullet \to
\sigma_{\leq i + 1} A^\bullet \to 0.
$$
を考える。これは項ごとに分裂する短完全列
$$
0 \to \sigma_{\geq i + 2} F(A^\bullet) \to F(A^\bullet) \to
\sigma_{\leq i + 1} F(A^\bullet) \to 0.
$$
を誘導することに注意する。したがって完全三角
$$
(\sigma_{\geq i + 2} A^\bullet, A^\bullet,
\sigma_{\leq i + 1} A^\bullet)
\quad\text{and}\quad
(\sigma_{\geq i + 2} F(A^\bullet), F(A^\bullet),
\sigma_{\leq i + 1} F(A^\bullet))
$$
が得られる。$RF$ は $A^\bullet$ で仮定により定義され、
$\sigma_{\leq i + 1}A^\bullet$ では第一段落により定義されるので、
$RF$ は $\sigma_{\geq i + 2}A^\bullet$ でも定義され、完全三角
$$
(RF(\sigma_{\geq i + 2} A^\bullet), RF(A^\bullet),
RF(\sigma_{\leq i + 1} A^\bullet))
$$
が得られる。補題 \ref{derived-lemma-2-out-of-3-defined} を参照せよ。
これらの完全三角を用いると、完全列の射
$$
\xymatrix{
H^i(\sigma_{\geq i + 2} F(A^\bullet)) \ar[r] \ar[d] &
H^i(F(A^\bullet)) \ar[r] \ar[d]^\alpha &
H^i(\sigma_{\leq i + 1} F(A^\bullet)) \ar[r] \ar[d]^\beta &
H^{i + 1}(\sigma_{\geq i + 2} F(A^\bullet)) \ar[d] \\
H^i(RF(\sigma_{\geq i + 2} A^\bullet)) \ar[r] &
H^i(RF(A^\bullet)) \ar[r] &
H^i(RF(\sigma_{\leq i + 1} A^\bullet)) \ar[r] &
H^{i + 1}(RF(\sigma_{\geq i + 2} A^\bullet))
}
$$
が得られる。第一段落の結果により、写像 $\beta$ は同型である。
上段左端と上段右端の対象は、見れば分かるように零である。したがって証明を
終えるには、
$H^i(RF(\sigma_{\geq i + 2} A^\bullet)) = 0$ および
$H^{i + 1}(RF(\sigma_{\geq i + 2} A^\bullet)) = 0$
を示せば十分である。これは補題 \ref{derived-lemma-negative-vanishing} から直ちに従う。
\end{proof}

\begin{proposition}
\label{derived-proposition-enough-acyclics}
\begin{slogan}
アーベル圏上の関手は、その関手に関して非輪状な複体で分解することにより、
下または上に有界な導来圏へ拡張される。
\end{slogan}
$F : \mathcal{A} \to \mathcal{B}$ をアーベル圏の間の加法的関手とする。
\begin{enumerate}
\item $\mathcal{A}$ の任意の対象が $RF$ に関して非輪状な対象へ単射で写るならば、
$RF$ は $K^{+}(\mathcal{A})$ のすべてで定義され、完全関手
$$
RF : D^{+}(\mathcal{A}) \longrightarrow D^{+}(\mathcal{B})
$$
が得られる。(\ref{derived-equation-everywhere}) を参照せよ。さらに、各項が
$RF$ に関して非輪状な任意の下に有界な複体 $A^\bullet$ は $RF$ を計算する。
\item
% P02-E0035: frozen source omits the article in “is quotient of”.
$\mathcal{A}$ の任意の対象が $LF$ に関して非輪状な対象の商ならば、
$LF$ は $K^{-}(\mathcal{A})$ のすべてで定義され、完全関手
$$
LF : D^{-}(\mathcal{A}) \longrightarrow D^{-}(\mathcal{B})
$$
が得られる。(\ref{derived-equation-everywhere}) を参照せよ。さらに、各項が
$LF$ に関して非輪状な任意の上に有界な複体 $A^\bullet$ は $LF$ を計算する。
\end{enumerate}
\end{proposition}

\begin{proof}
$\mathcal{A}$ の任意の対象が $RF$ に関して非輪状な対象へ単射で写ると仮定する。
$\mathcal{I}$ を $RF$ に関して非輪状な対象の集合とする。
$K^\bullet$ を $\mathcal{A}$ の下に有界な複体とする。補題
\ref{derived-lemma-subcategory-right-resolution} により、$I^\bullet$ が下に有界で
$I^n \in \mathcal{I}$ となる擬同型
$\alpha : K^\bullet \to I^\bullet$ が存在する。したがって (1) を証明するには、
$s : I^\bullet \to (I')^\bullet$ が $\mathcal{I}$ の対象からなる下に有界な
複体の擬同型であるとき、$F(I^\bullet) \to F((I')^\bullet)$ が擬同型であることを
示せば十分である。補題 \ref{derived-lemma-find-existence-computes} を参照せよ。
錐 $C(s)^\bullet$ は、すべての項が $\mathcal{I}$ に属する非輪状かつ下に有界な
複体であることに注意する。したがって、すべての項が $\mathcal{I}$ に属する
非輪状かつ下に有界な複体 $I^\bullet$ が与えられたとき、
$F(I^\bullet)$ が非輪状であることを示せば十分である。

\medskip\noindent
$n < n_0$ に対して $I^n = 0$ とする。$J^n = \Im(d^n)$ と置くことにより、
$n \geq n_0$ に対する短完全列
$0 \to J^n \to I^{n + 1} \to J^{n + 1} \to 0$
へ $I^\bullet$ を分解する。補題 \ref{derived-lemma-derived-canonical-delta-functor} により、
これらの列は $D^+(\mathcal{A})$ に完全三角
$(J^n, I^{n + 1}, J^{n + 1})$ を誘導する。各 $k \in \mathbf{Z}$ に対して
$H_k$ を「すべての $n \leq k$ に対して対象 $J^n$ が $\mathcal{I}$ に属する」
という主張とする。$k < n_0$ ならば $H_k$ は自明に成り立つ。$H_n$ が成り立つと、
補題 \ref{derived-lemma-2-out-of-3-computes} により $J^{n + 1}$ は
$\mathcal{I}$ に属し、$H_{n + 1}$ が成り立つ。命題
\ref{derived-proposition-derived-functor} により完全三角
$(RF(J^n), RF(I^{n + 1}), RF(J^{n + 1}))$ が得られる。
$J^n, I^{n + 1}, J^{n + 1}$ は $\mathcal{I}$ に属するので、この完全三角に付随する
コホモロジー長完全列
(\ref{derived-equation-long-exact-cohomology-sequence-D})
は完全列
$$
0 \to F(J^n) \to F(I^{n + 1}) \to F(J^{n + 1}) \to 0
$$
に退化する。これから $F(I^\bullet)$ が完全であることが従う。

\medskip\noindent
$LF$ の場合の証明は双対的である。
\end{proof}

\begin{lemma}
\label{derived-lemma-right-derived-exact-functor}
$F : \mathcal{A} \to \mathcal{B}$ をアーベル圏の間の完全関手とする。このとき
\begin{enumerate}
\item $\mathcal{A}$ のすべての対象は $F$ に関して右非輪状である。
\item $RF : D^{+}(\mathcal{A}) \to D^{+}(\mathcal{B})$ は至る所で定義される。
\item $RF : D(\mathcal{A}) \to D(\mathcal{B})$ は至る所で定義される。
\item すべての複体は $RF$ を計算する。言い換えれば、すべての複体に対して標準写像
$F(K^\bullet) \to RF(K^\bullet)$ は同型である。
\item $i \not = 0$ に対して $R^iF = 0$ である。
\end{enumerate}
\end{lemma}

\begin{proof}
$F$ は非輪状複体を非輪状複体へ、擬同型を擬同型へ写すので、これは成り立つ。
詳細は省略する。
\end{proof}


\section{導来圏の三角部分圏}
\label{derived-section-triangulated-sub}

\noindent
$\mathcal{A}$ をアーベル圏とする。本節では、ある厳密充満飽和三角部分圏
$\mathcal{D}' \subset D(\mathcal{A})$ を考察する。

\medskip\noindent
$\mathcal{B} \subset \mathcal{A}$ を弱 Serre 部分圏とする。Homology, Definition
\ref{homology-definition-serre-subcategory} および
Lemma \ref{homology-lemma-characterize-weak-serre-subcategory} を参照せよ。
% P02-E0036: frozen source omits “be” in the definition of the full subcategory.
$D_\mathcal{B}(\mathcal{A})$ を、対象が
$$
\Ob(D_\mathcal{B}(\mathcal{A}))
=
\{X \in \Ob(D(\mathcal{A})) \mid
H^n(X) \text{ is an object of }\mathcal{B}\text{ for all }n\}
$$
である $D(\mathcal{A})$ の充満部分圏とする。また
$D^{+}_\mathcal{B}(\mathcal{A}) =
D^{+}(\mathcal{A}) \cap D_\mathcal{B}(\mathcal{A})$
と定め、他の有界版についても同様に定める。

\begin{lemma}
\label{derived-lemma-cohomology-in-serre-subcategory}
$\mathcal{A}$ をアーベル圏とし、
$\mathcal{B} \subset \mathcal{A}$ を弱 Serre 部分圏とする。
圏 $D_\mathcal{B}(\mathcal{A})$ は $D(\mathcal{A})$ の
厳密充満飽和三角部分圏である。有界版についても同様である。
\end{lemma}

\begin{proof}
$D_\mathcal{B}(\mathcal{A})$ が、平行移動関手で保たれる加法的部分圏であることは
明らかである。$X \oplus Y$ が $D_\mathcal{B}(\mathcal{A})$ に属するならば、
% P02-E0037: frozen source repeats “maps” in “kernels of maps between maps of objects”.
$H^n(X \oplus Y) = H^n(X) \oplus H^n(Y)$ なので、$H^n(X)$ と $H^n(Y)$ はともに
$\mathcal{B}$ の対象間の射の核である。したがって $X$ と $Y$ はともに
$D_\mathcal{B}(\mathcal{A})$ に属する。ゆえに補題
\ref{derived-lemma-triangulated-subcategory} により、$X$ と $Y$ が
$D_\mathcal{B}(\mathcal{A})$ に属する完全三角
$(X, Y, Z, f, g, h)$ が与えられたとき、$Z$ が
$D_\mathcal{B}(\mathcal{A})$ の対象であることを示せば十分である。
コホモロジー長完全列 (\ref{derived-equation-long-exact-cohomology-sequence-D}) と
弱 Serre 部分圏の定義（Homology, Definition
\ref{homology-definition-serre-subcategory} を参照）により、すべての $n$ に対して
$H^n(Z)$ は $\mathcal{B}$ の対象である。したがって $Z$ は
$D_\mathcal{B}(\mathcal{A})$ の対象である。
\end{proof}

\noindent
引き続き、$\mathcal{B}$ をアーベル圏 $\mathcal{A}$ の弱 Serre 部分圏と仮定する。
このとき $\mathcal{B}$ はアーベル圏であり、包含関手
$\mathcal{B} \to \mathcal{A}$ は完全である。したがって補題
\ref{derived-lemma-right-derived-exact-functor} により導来関手
$D(\mathcal{B}) \to D(\mathcal{A})$ が得られる。関手
$D(\mathcal{B}) \to D(\mathcal{A})$ が標準完全関手
\begin{equation}
\label{derived-equation-compare}
D(\mathcal{B}) \longrightarrow D_\mathcal{B}(\mathcal{A})
\end{equation}
を経由することは明らかである。実際、$\mathcal{B}$ の対象からなる複体は
確かに $D_\mathcal{B}(\mathcal{A})$ の対象を与える。また、
$D_\mathcal{B}(\mathcal{A})$ の完全三角は、その頂点が
$D_\mathcal{B}(\mathcal{A})$ に属する $D(\mathcal{A})$ の完全三角にほかならない。
したがって $D(\mathcal{B}) \to D(\mathcal{A})$ が完全であることから、
この関手も完全である。同様に有界版について、関手
$D^+(\mathcal{B}) \to D^+_\mathcal{B}(\mathcal{A})$、
$D^-(\mathcal{B}) \to D^-_\mathcal{B}(\mathcal{A})$、および
$D^b(\mathcal{B}) \to D^b_\mathcal{B}(\mathcal{A})$
が得られる。多くの場合、表示された関手が圏同値であるかどうかが重要な問題となる。

\medskip\noindent
ここで $\mathcal{B}$ が $\mathcal{A}$ の Serre 部分圏であると仮定する。
この場合、商関手 $\mathcal{A} \to \mathcal{A}/\mathcal{B}$ がある。
Homology, Lemma \ref{homology-lemma-serre-subcategory-is-kernel} を参照せよ。
この場合 $D_\mathcal{B}(\mathcal{A})$ は関手
$D(\mathcal{A}) \to D(\mathcal{A}/\mathcal{B})$ の核である。
したがって補題 \ref{derived-lemma-universal-property-quotient} により標準関手
$$
D(\mathcal{A})/D_\mathcal{B}(\mathcal{A})
\longrightarrow
D(\mathcal{A}/\mathcal{B})
$$
が得られる。有界版についても同様である。

\begin{lemma}
\label{derived-lemma-derived-of-quotient}
$\mathcal{A}$ をアーベル圏とし、
$\mathcal{B} \subset \mathcal{A}$ を Serre 部分圏とする。このとき
$D(\mathcal{A}) \to D(\mathcal{A}/\mathcal{B})$ は本質的全射である。
\end{lemma}

\begin{proof}
Homology, Lemma \ref{homology-lemma-serre-subcategory-is-kernel} の証明における
圏 $\mathcal{A}/\mathcal{B}$ の記述を用いる。
$(X^\bullet, d^\bullet)$ を $\mathcal{A}/\mathcal{B}$ の複体とする。
これは、$X^i$ が $\mathcal{A}$ の対象であり、
$d^i : X^i \to X^{i + 1}$ が $\mathcal{A}/\mathcal{B}$ の射で、
$\mathcal{A}/\mathcal{B}$ において $d^i \circ d^{i - 1} = 0$ となることを意味する。

\medskip\noindent
$i \geq 0$ に対して $d^i = (s^i, f^i)$ と書ける。ここで
$s^i : Y^i \to X^i$ は、その核と余核が $\mathcal{B}$ に属する
$\mathcal{A}$ の射であり（同値なことに、$s^i$ は商圏で同型となる）、
$f^i : Y^i \to X^{i + 1}$ は $\mathcal{A}$ の射である。
帰納法により可換図式
$$
\xymatrix{
& (X')^1 \ar@{..>}[r] & (X')^2 \ar@{..>}[r] & \ldots \\
X^0 \ar@{..>}[ru] &
X^1 \ar@{..>}[u] &
X^2 \ar@{..>}[u] &
\ldots \\
Y^0 \ar[u]_{s^0} \ar[ru]_{f^0} &
Y^1 \ar[u]_{s^1} \ar[ru]_{f^1} &
Y^2 \ar[u]_{s^2} \ar[ru]_{f^2} &
\ldots
}
$$
で、垂直な矢印 $X^i \to (X')^i$ が商圏で同型となるものを構成する。
実際、まず
$(X')^1 = \Coker(Y^0 \to X^0 \oplus X^1)$
（より正確には、矢印 $s^0$ と $f^0$ からなる図式の押し出し）と置けば、
最初の可換図式が得られる。次に
$(X')^2 = \Coker(Y^1 \to (X')^1 \oplus X^2)$ と置き、以下同様に進める。
さらに $n \leq 0$ に対して $(X')^n = X^n$ と置くと、写像
$(X^\bullet, d^\bullet) \to ((X')^\bullet, (d')^\bullet)$ は
$\mathcal{A}/\mathcal{B}$ における複体の同型であることが分かる。
したがって $n \geq 0$ に対して、$d^n : X^n \to X^{n + 1}$ が
$\mathcal{A}$ の写像 $X^n \to X^{n + 1}$ によって与えられると仮定してよい。

\medskip\noindent
双対的に、$i < 0$ に対して $d^i = (g^i, t^{i + 1})$ と書ける。ここで
$t^{i + 1} : X^{i + 1} \to Z^{i + 1}$ は商圏における同型であり、
$g^i : X^i \to Z^{i + 1}$ は射である。帰納法により可換図式
% P02-E0040: frozen source uses subscripts for these cochain components.
$$
\xymatrix{
\ldots &
Z^{-2} &
Z^{-1} &
Z^0 \\
\ldots &
X^{-2} \ar[u]_{t_{-2}} \ar[ru]_{g_{-2}} &
X^{-1} \ar[u]_{t_{-1}} \ar[ru]_{g_{-1}} &
X^0 \ar[u]_{t^0} \\
\ldots &
(X')^{-2} \ar@{..>}[u] \ar@{..>}[r] &
(X')^{-1} \ar@{..>}[u] \ar@{..>}[ru]
}
$$
で、垂直な矢印 $(X')^i \to X^i$ が商圏で同型となるものを構成する。
実際、$(X')^{-1} = X^{-1} \times_{Z^0} X^0$ と置き、次に
$(X')^{-2} = X^{-2} \times_{Z^{-1}} (X')^{-1}$ と置き、以下同様に進める。
さらに $n \geq 0$ に対して $(X')^n = X^n$ と置くと、写像
$((X')^\bullet, (d')^\bullet) \to (X^\bullet, d^\bullet)$ は
$\mathcal{A}/\mathcal{B}$ における複体の同型であることが分かる。
したがってすべての $n \in \mathbf{Z}$ に対して、
$d^n : X^n \to X^{n + 1}$ が $\mathcal{A}$ の写像
$d^n : X^n \to X^{n + 1}$ によって与えられると仮定してよい。

\medskip\noindent
この場合、合成 $d^n \circ d^{n - 1}$ が
$\mathcal{A}/\mathcal{B}$ で零であることが分かっている。$n > 0$ に対して
$X^n$ を
$$
(X')^n = X^n/\sum\nolimits_{0 < k \leq n} \Im(\Im(X^{k - 2} \to X^k) \to X^n)
$$
で置き換えると、$n \geq 0$ に対して合成
$d^n \circ d^{n - 1}$ は零となる。（上の第二段落と同様に、複体の同型
$(X^\bullet, d^\bullet) \to ((X')^\bullet, (d')^\bullet)$ が得られる。）
最後に、$n < 0$ に対して $X^n$ を
$$
(X')^n = \bigcap\nolimits_{n \leq k < 0}
(X^n \to X^k)^{-1}\Ker(X^k \to X^{k + 2})
$$
で置き換え、同様に議論すると、その $\mathcal{A}/\mathcal{B}$ における像が
与えられた複体と同型であるような $\mathcal{A}$ の複体が得られる。
\end{proof}

\begin{lemma}
\label{derived-lemma-quotient-by-serre-easy}
$\mathcal{A}$ をアーベル圏とし、
$\mathcal{B} \subset \mathcal{A}$ を Serre 部分圏とする。
関手 $v : \mathcal{A} \to \mathcal{A}/\mathcal{B}$ が左随伴
$u : \mathcal{A}/\mathcal{B} \to \mathcal{A}$ を持ち、
$vu \cong \text{id}$ であると仮定する。このとき
$$
D(\mathcal{A})/D_\mathcal{B}(\mathcal{A}) = D(\mathcal{A}/\mathcal{B})
$$
であり、有界版についても同様である。
\end{lemma}

\begin{proof}
補題 \ref{derived-lemma-derived-of-quotient} により、関手
$D(v) : D(\mathcal{A}) \to D(\mathcal{A}/\mathcal{B})$ は本質的全射である。
$D(\mathcal{A})$ の対象 $X$ に対する随伴写像
$c_X : uvX \to X$ は $D(\mathcal{A}/\mathcal{B})$ で同型へ写る。
実際、$vu \cong \text{id}$ という仮定により $vuv \cong v$ である。
したがって完全三角 $(uvX, X, Z, c_X, g, h)$ において、コホモロジー長完全列を
見れば、$Z$ は $D_\mathcal{B}(\mathcal{A})$ の対象であることが分かる。
ゆえに $c_X$ は商圏 $D(\mathcal{A})/D_\mathcal{B}(\mathcal{A})$ を定義するための
乗法系の元である。したがって
$D(\mathcal{A})/D_\mathcal{B}(\mathcal{A})$ において $uvX \cong X$ である。
% P02-E0038: frozen source has an extra closing parenthesis in Ob(A)).
$X, Y \in \Ob(\mathcal{A}))$ に対して写像
$$
\Hom_{D(\mathcal{A})/D_\mathcal{B}(\mathcal{A})}(X, Y)
\longrightarrow
\Hom_{D(\mathcal{A}/\mathcal{B})}(vX, vY)
$$
は全単射である。上の注意により、$u$ がその逆を与えるからである。
\end{proof}

\noindent
ある Serre 部分圏 $\mathcal{B} \subset \mathcal{A}$ に対しては、関手
$D(\mathcal{B}) \to D_\mathcal{B}(\mathcal{A})$
が充満忠実であることを示せる。

\begin{lemma}
\label{derived-lemma-fully-faithful-embedding}
$\mathcal{A}$ をアーベル圏とし、
$\mathcal{B} \subset \mathcal{A}$ を Serre 部分圏とする。
$X \in \Ob(\mathcal{A})$ かつ $Y \in \Ob(\mathcal{B})$ である任意の全射
$X \to Y$ に対して、$Y$ へ全射となる $X' \subset X$,
$X' \in \Ob(\mathcal{B})$ が存在すると仮定する。このとき
(\ref{derived-equation-compare}) の関手
$D^-(\mathcal{B}) \to D^-_\mathcal{B}(\mathcal{A})$ は圏同値である。
\end{lemma}

\begin{proof}
$X^\bullet$ を $\mathcal{A}$ の上に有界な複体で、すべての
$i \in \mathbf{Z}$ に対して $H^i(X^\bullet) \in \Ob(\mathcal{B})$ となるものとする。
さらに、すべての $i \in \mathbf{Z}$ に対して
$B^i \subset X^i$, $B^i \in \Ob(\mathcal{B})$ が与えられているとする。
主張：次を満たす部分複体 $Y^\bullet \subset X^\bullet$ が存在する。
\begin{enumerate}
\item $Y^\bullet \to X^\bullet$ は擬同型である。
\item すべての $i \in \mathbf{Z}$ に対して $Y^i \in \Ob(\mathcal{B})$ である。
\item すべての $i \in \mathbf{Z}$ に対して $B^i \subset Y^i$ である。
\end{enumerate}
この主張を示すため、まず補題の仮定を用いて、
$H^i(X^\bullet)$ へ全射となる
$C^i \subset \Ker(d^i : X^i \to X^{i + 1})$,
$C^i \in \Ob(\mathcal{B})$ を選ぶ。
$D^i = C^i + d^{i - 1}(B^{i - 1}) + B^i$ と置くと、(2) と (3) を満たす
部分複体 $D^\bullet$ で、すべての $i \in \mathbf{Z}$ に対して
$H^i(D^\bullet) \to H^i(X^\bullet)$ が全射となるものが得られる。
$E^i \subset X^i$, $E^i \in \Ob(\mathcal{B})$ で
$d^i(E^i) \subset D^{i + 1} + E^{i + 1}$ を満たす任意の選択に対して、
$Y^i = D^i + E^i$ と置けば、その各項が $\mathcal{B}$ に属し、その
コホモロジーが $X^\bullet$ のコホモロジーへ全射となる部分複体が得られる。
明らかに、
$d^i(E^i) = (D^{i + 1} + E^{i + 1}) \cap \Im(d^i)$
ならば、コホモロジー上の写像は単射でもある。$n \gg 0$ に対しては
$E^n$ を $0$ と取れる。降下帰納法により、所望の性質を持つ $E^i$ を
すべての $i$ に対して選べる。実際、
$E^{i + 1}, E^{i + 2}, \ldots$ が与えられたとき、
$$
d^i(E^i) = (D^{i + 1} + E^{i + 1}) \cap \Im(d^i)
$$
となる $E^i \subset X^i$ を選ぶ。これは補題の仮定と、
$\mathcal{B}$ が $\mathcal{A}$ の Serre 部分圏であるため
$(D^{i + 1} + E^{i + 1}) \cap \Im(d^i)$ が $\mathcal{B}$ に属するという事実から
可能である。

\medskip\noindent
上の主張から補題が従う。本質的全射性は主張から直ちに従う。
忠実性を証明しよう。$\mathcal{B}$ の上に有界な複体の射
$f : U^\bullet \to V^\bullet$ で、その $D(\mathcal{A})$ における像が零となるものを
取る。このとき $\mathcal{A}$ の上に有界な複体への擬同型
$s : V^\bullet \to X^\bullet$ で、$s \circ f$ が零とホモトピックとなるものが
存在する。$0$ と $s \circ f$ の間のホモトピー
$h^i : U^i \to X^{i - 1}$ を選ぶ。上の主張を
$B^i = h^{i + 1}(U^{i + 1}) + s^i(V^i)$ に適用する。
得られる写像 $s' : V^\bullet \to Y^\bullet$ も擬同型であり、
$h^i$ が $Y^{i - 1}$ を経由することから明らかなように、
$s' \circ f$ は零とホモトピックである。これで忠実性が示された。
% P02-E0039: frozen heading says “Fully faithfulness”; Japanese uses the standard noun.
充満忠実性もまったく同じ方法で示される。
\end{proof}



\section{入射分解}
\label{derived-section-injective-resolutions}

\noindent
本節では、十分多くの入射対象を持つアーベル圏における入射分解の存在に関する
いくつかの補題を証明する。

\begin{definition}
\label{derived-definition-injective-resolution}
$\mathcal{A}$ をアーベル圏とし、$A \in \Ob(\mathcal{A})$ とする。
{\it $A$ の入射分解}とは、複体 $I^\bullet$ と写像 $A \to I^0$ であって、
次を満たすものをいう。
\begin{enumerate}
\item $n < 0$ に対して $I^n = 0$ である。
\item 各 $I^n$ は $\mathcal{A}$ の入射対象である。
\item 写像 $A \to I^0$ は $\Ker(d^0)$ への同型である。
\item $i > 0$ に対して $H^i(I^\bullet) = 0$ である。
\end{enumerate}
したがって $A[0] \to I^\bullet$ は擬同型である。言い換えれば、複体
$$
\ldots \to 0 \to A \to I^0 \to I^1 \to \ldots
$$
は非輪状である。$K^\bullet$ を $\mathcal{A}$ の複体とする。
{\it $K^\bullet$ の入射分解}とは、複体 $I^\bullet$ と複体の写像
$\alpha : K^\bullet \to I^\bullet$ であって、次を満たすものをいう。
\begin{enumerate}
\item $n \ll 0$ に対して $I^n = 0$、すなわち $I^\bullet$ は下に有界である。
\item 各 $I^n$ は $\mathcal{A}$ の入射対象である。
\item 写像 $\alpha : K^\bullet \to I^\bullet$ は擬同型である。
\end{enumerate}
\end{definition}

\noindent
言い換えれば、入射分解 $K^\bullet \to I^\bullet$ は図式
$$
\xymatrix{
\ldots \ar[r] & K^{n - 1} \ar[d] \ar[r] & K^n \ar[d] \ar[r] &
K^{n + 1} \ar[d] \ar[r] & \ldots \\
\ldots \ar[r] & I^{n - 1} \ar[r] & I^n \ar[r] & I^{n + 1} \ar[r] & \ldots
}
$$
を与え、これは各次数のコホモロジー対象上で同型を誘導する。
$\mathcal{A}$ の対象 $A$ の入射分解は、複体 $A[0]$ の入射分解とほぼ同じものである。

\begin{lemma}
\label{derived-lemma-cohomology-bounded-below}
$\mathcal{A}$ をアーベル圏とし、$K^\bullet$ を $\mathcal{A}$ の複体とする。
\begin{enumerate}
\item $K^\bullet$ が入射分解を持つならば、$n \ll 0$ に対して
$H^n(K^\bullet) = 0$ である。
\item すべての $n \ll 0$ に対して $H^n(K^\bullet) = 0$ ならば、
$L^\bullet$ が下に有界である擬同型 $K^\bullet \to L^\bullet$ が存在する。
\end{enumerate}
\end{lemma}

\begin{proof}
省略する。第二の主張には、ある $n \ll 0$ に対して
$L^\bullet = \tau_{\geq n}K^\bullet$ と取ればよい。
切断 $\tau_{\geq n}$ の定義については Homology, Section
\ref{homology-section-truncations} を参照せよ。
\end{proof}

\begin{lemma}
\label{derived-lemma-injective-resolutions-exist}
$\mathcal{A}$ をアーベル圏とする。$\mathcal{A}$ は十分多くの入射対象を持つと仮定する。
\begin{enumerate}
\item $\mathcal{A}$ の任意の対象は入射分解を持つ。
\item すべての $n \ll 0$ に対して $H^n(K^\bullet) = 0$ ならば、
$K^\bullet$ は入射分解を持つ。
\item $K^\bullet$ が $n < a$ に対して $K^n = 0$ となる複体ならば、
各 $\alpha^n : K^n \to I^n$ が単射であり、$n < a$ に対して
$I^n = 0$ となる入射分解 $\alpha : K^\bullet \to I^\bullet$ が存在する。
\end{enumerate}
\end{lemma}

\begin{proof}
(1) を証明する。まず $A$ から $\mathcal{A}$ の入射対象への単射
$A \to I^0$ を選ぶ。次に、
% P02-E0041: frozen source uses I_0 although the constructed object is I^0.
$\mathcal{A}$ の入射対象への単射 $I_0/A \to I^1$ を選ぶ。
誘導される写像 $I^0 \to I^1$ を $d^0$ と書く。次に、
$\mathcal{A}$ の入射対象への単射 $I^1/\Im(d^0) \to I^2$ を選び、
誘導される写像 $I^1 \to I^2$ を $d^1$ と書く。以下同様に進める。
補題 \ref{derived-lemma-cohomology-bounded-below} により、(2) は (3) から従う。
(3) は補題 \ref{derived-lemma-subcategory-right-resolution} の特別な場合である。
\end{proof}

\begin{lemma}
\label{derived-lemma-acyclic-is-zero}
$\mathcal{A}$ をアーベル圏とする。$K^\bullet$ を非輪状複体とし、
$I^\bullet$ を入射対象からなる下に有界な複体とする。
任意の射 $K^\bullet \to I^\bullet$ は零とホモトピックである。
\end{lemma}

\begin{proof}
$\alpha : K^\bullet \to I^\bullet$ を複体の射とする。$I^\bullet$ は
下に有界なので、$j \ll 0$ に対して $\alpha^j = 0$ であることに注意する。
特に、$j \leq n$ に対する $h^j :  K^j \to I^{j - 1}$ で、
$j < n$ に対して
$\alpha^j = d^{j - 1} \circ h^j + h^{j + 1} \circ d^j$
となるものが存在するような $n$ を取れる。
$\alpha^n = d^{n - 1} \circ h^n + h^{n + 1} \circ d^n$
となる射 $h^{n + 1} : K^{n + 1} \to I^n$ が存在することを示す。
% P02-E0042: frozen source has several incorrect degree components in this block.
\begin{align*}
(\alpha^n - d^{n - 1} \circ h^{n - 1}) \circ d^{n - 1}
& =
\alpha^{n - 1} \circ d^{n - 1} - d^{n - 1} \circ h^{n - 1} \circ d^{n - 1} \\
& =
d^{n - 1} \circ \alpha^{n - 1} - d^{n - 1} \circ h^{n - 1} \circ d^{n - 1} \\
& =
d^{n - 1} \circ (d^{n - 2} \circ h^{n - 1} + h^n \circ d^{n - 1})
- d^{n - 1} \circ h^{n - 1} \circ d^{n - 1} \\
& = 0
\end{align*}
$K^\bullet$ は非輪状なので
$d^{n - 1}(K^{n - 1}) = \Ker(K^n \to K^{n + 1})$ である。
したがって $\alpha^n -  d^{n - 1} \circ h^{n - 1}$ を、
$K^{n + 1}$ の部分対象 $\Im(K^n \to K^{n + 1})$ 上で定義された
$I^n$ への写像とみなせる。対象 $I^n$ の入射性により、これを写像
$h^{n + 1} : K^{n + 1} \to I^n$ へ拡張できる。この選択に対して実際に
$\alpha^n = d^{n - 1} \circ h^n + h^{n + 1} \circ d^n$
となることは容易に確かめられる。

\medskip\noindent
$n$ に関する帰納法により、所望の $\alpha$ と $0$ の間のホモトピーをなす
$h  = (h^j)_{j \in \mathbf{Z}}$ が得られる。
\end{proof}

\begin{remark}
\label{derived-remark-easier-proofs}
$\mathcal{A}$ をアーベル圏とする。$K(\mathcal{A})$ が三角圏であることと
補題 \ref{derived-lemma-acyclic-is-zero} を用いれば、通常は図式追跡で証明される
以下の補題のいくつかについて別の証明が得られる。
実際、$\alpha : K^\bullet \to L^\bullet$ を複体の擬同型とする。このとき
$$
(K^\bullet, L^\bullet, C(\alpha)^\bullet, \alpha, i, -p)
$$
は $K(\mathcal{A})$ の完全三角であり
（補題 \ref{derived-lemma-the-same-up-to-isomorphisms}）、
$C(\alpha)^\bullet$ は非輪状複体である（補題 \ref{derived-lemma-acyclic}）。
次に、$I^\bullet$ を入射対象からなる下に有界な複体とする。このとき
$$
\xymatrix{
\Hom_{K(\mathcal{A})}(C(\alpha)^\bullet, I^\bullet) \ar[r] &
\Hom_{K(\mathcal{A})}(L^\bullet, I^\bullet) \ar[r] &
\Hom_{K(\mathcal{A})}(K^\bullet, I^\bullet) \ar[lld] \\
\Hom_{K(\mathcal{A})}(C(\alpha)^\bullet[-1], I^\bullet)
}
$$
はアーベル群の完全列である。補題 \ref{derived-lemma-representable-homological} を
参照せよ。ここで補題 \ref{derived-lemma-acyclic-is-zero} により両端の二群は零であり、
したがって
$\Hom_{K(\mathcal{A})}(L^\bullet, I^\bullet) =
\Hom_{K(\mathcal{A})}(K^\bullet, I^\bullet)$ である。
\end{remark}

\begin{lemma}
\label{derived-lemma-morphisms-lift}
$\mathcal{A}$ をアーベル圏とする。実線からなる図式
$$
\xymatrix{
K^\bullet \ar[r]_\alpha \ar[d]_\gamma & L^\bullet \ar@{-->}[dl]^\beta \\
I^\bullet
}
$$
を考える。ここで $I^\bullet$ は入射対象からなる下に有界な複体であり、
$\alpha$ は擬同型である。
\begin{enumerate}
\item ホモトピーを除いて図式を可換にする複体の写像 $\beta$ が存在する。
\item $\alpha$ が各次数で単射ならば、図式を可換にする $\beta$ を取れる。
\end{enumerate}
\end{lemma}

\begin{proof}
(1) の「正しい」証明は Remark \ref{derived-remark-easier-proofs} で説明した。
ここでは直接的な証明も与える。

\medskip\noindent
まず (2) から (1) が従うことを示す。補題 \ref{derived-lemma-make-injective} と同じく、
% P02-E0043: frozen source writes K instead of the complex K^\bullet here.
$\tilde \alpha : K \to \tilde L^\bullet$, $\pi$, $s$ を取る。
$\tilde \alpha$ は単射なので、(2) により
$\gamma = \tilde \beta \circ \tilde \alpha$ となる射
$\tilde \beta : \tilde L^\bullet \to I^\bullet$ が存在する。
$\beta = \tilde \beta \circ s$ と置く。このとき所望のとおり
$$
\beta \circ \alpha
=
\tilde \beta \circ s \circ \pi \circ \tilde \alpha
\sim
\tilde \beta \circ \tilde \alpha
=
\gamma
$$
である。

\medskip\noindent
$\alpha : K^\bullet \to L^\bullet$ が単射であると仮定する。
すべての次数 $\leq n - 1$ で $\beta$ が既に定義され、微分と両立し、
すべての $j \leq n - 1$ に対して
$\gamma^j = \beta^j \circ \alpha^j$ となると仮定する。
実線からなる可換図式
$$
\xymatrix{
K^{n - 1} \ar[r] \ar@/_2pc/[dd]_\gamma \ar[d]^\alpha &
K^n \ar@/^2pc/[dd]^\gamma \ar[d]^\alpha \\
L^{n - 1} \ar[r] \ar[d]^\beta &
L^n \ar@{-->}[d] \\
I^{n - 1} \ar[r] &
I^n
}
$$
を考える。点線の矢印は部分対象 $\alpha(K^n)$ と
$d^{n - 1}(L^{n - 1})$ 上で指定されることが分かる。さらに、この二つの矢印は
$\alpha(d^{n - 1}(K^{n - 1}))$ 上で一致する。したがって
\begin{equation}
\label{derived-equation-qis}
\alpha(d^{n - 1}(K^{n - 1}))
=
\alpha(K^n) \cap d^{n - 1}(L^{n - 1})
\end{equation}
ならば、これらの射は貼り合わされて射
$\alpha(K^n) + d^{n - 1}(L^{n - 1}) \to I^n$ となり、
$I^n$ の入射性を用いて $L^n$ 全体から $I^n$ への射へ拡張できる。
この後、帰納法によりすべての $n$ に対する射 $\beta$ が同時に得られる
（$I^\bullet$ は下に有界なので、すべての $n \ll 0$ に対して
$\beta^n = 0$ と置ける。このようにして帰納法を開始する）。

\medskip\noindent
残るのは等式 (\ref{derived-equation-qis}) の証明である。適切な図式追跡によって
各自で示すことを勧めるが、以下に議論を与える。
包含
$\alpha(d^{n - 1}(K^{n - 1})) \subset \alpha(K^n) \cap d^{n - 1}(L^{n - 1})$
は明らかである。$\mathcal{A}$ の対象 $T$ と、その像が部分対象
$\alpha(K^n) \cap d^{n - 1}(L^{n - 1})$ に含まれる射
$x : T \to L^n$ を取る。$\alpha$ は単射なので、ある
$x' : T \to K^n$ に対して $x = \alpha \circ x'$ である。
さらに $x$ は $d^{n - 1}(L^{n - 1})$ に属するので
$d^n \circ x = 0$ である。再び $\alpha$ の単射性により
$d^n \circ x' = 0$ である。したがって $x'$ は射
$[x'] : T \to H^n(K^\bullet)$ を与える。一方、$x$ が誘導する対応する写像
$[x] : T \to H^n(L^\bullet)$ は仮定により零である。
$\alpha$ は擬同型なので、$[x'] = 0$ と結論できる。これはまさに $x'$ の像が
$d^{n - 1}(K^{n - 1})$ に含まれることを意味し、証明が終わる。
\end{proof}

\begin{lemma}
\label{derived-lemma-morphisms-equal-up-to-homotopy}
$\mathcal{A}$ をアーベル圏とする。実線からなる図式
$$
\xymatrix{
K^\bullet \ar[r]_\alpha \ar[d]_\gamma & L^\bullet \ar@{-->}[dl]^{\beta_i} \\
I^\bullet
}
$$
を考える。ここで $I^\bullet$ は入射対象からなる下に有界な複体であり、
$\alpha$ は擬同型である。ホモトピーを除いて図式を可換にする任意の二射
$\beta_1, \beta_2$ は互いにホモトピックである。
\end{lemma}

\begin{proof}
これは Remark \ref{derived-remark-easier-proofs} から従う。
ここでも直接的な議論を与える。

\medskip\noindent
補題 \ref{derived-lemma-make-injective} と同じく、
% P02-E0043: second frozen occurrence writes K instead of K^\bullet.
$\tilde \alpha : K \to \tilde L^\bullet$, $\pi$, $s$ を取る。
$\beta_1 \circ\pi$ が $\beta_2 \circ \pi$ とホモトピックであることを示せれば、
$\pi \circ s$ は恒等射なので $\beta_1 \sim \beta_2$ が従う。
したがって、すべての $n$ に対して
$\alpha^n : K^n \to L^n$ が直和因子の包含であると仮定してよい。
このとき複体の短完全列
$$
0 \to K^\bullet \to L^\bullet \to M^\bullet \to 0
$$
で、項ごとに分裂し、$M^\bullet$ が非輪状であるものが得られる。
分裂 $L^n = K^n \oplus M^n$ を選ぶ。すると
$\beta_i^n : K^n \oplus M^n \to I^n$ および
$\gamma^n : K^n \to I^n$ がある。この場合、$\beta_i$ に関する条件は、
射 $h_i^n : K^n \to I^{n - 1}$ で
$$
\gamma^n - \beta_i^n|_{K^n} = d \circ h_i^n + h_i^{n + 1} \circ d
$$
となるものが存在することである。したがって
$$
\beta_1^n|_{K^n} - \beta_2^n|_{K^n}
=
d \circ (h_1^n - h_2^n) + (h_1^{n + 1} - h_2^{n + 1}) \circ d
$$
である。第一因子上で $h_1^n - h_2^n$ に等しく、第二因子上で零となる写像
$h^n : K^n \oplus M^n \to I^{n - 1}$ を考える。このとき
% P02-E0044: frozen source composes d with the parenthesized whole sum.
$$
\beta_1^n - \beta_2^n
-
(d \circ h^n + h^{n + 1}) \circ d
$$
は複体の射 $L^\bullet \to I^\bullet$ であり、部分複体 $K^\bullet$ 上で
恒等的に零である。したがってこれは
$L^\bullet \to M^\bullet \to I^\bullet$ と分解する。
ゆえに補題 \ref{derived-lemma-acyclic-is-zero} から本補題の結論が従う。
\end{proof}

\begin{lemma}
\label{derived-lemma-morphisms-into-injective-complex}
$\mathcal{A}$ をアーベル圏とする。$I^\bullet$ を入射対象からなる下に有界な
複体とし、$L^\bullet \in K(\mathcal{A})$ とする。このとき
$$
\Mor_{K(\mathcal{A})}(L^\bullet, I^\bullet)
=
\Mor_{D(\mathcal{A})}(L^\bullet, I^\bullet).
$$
\end{lemma}

\begin{proof}
$a$ を右辺の元とする。ここで
$\alpha : K^\bullet \to L^\bullet$ は擬同型、
$\gamma : K^\bullet \to I^\bullet$ は複体の写像として、
$a = \gamma\alpha^{-1}$ と表示できる。補題
\ref{derived-lemma-morphisms-lift} により、$\beta \circ \alpha$ が $\gamma$ と
ホモトピックになるような射 $\beta : L^\bullet \to I^\bullet$ を取れる。
これで写像が全射であることが示された。$b$ を左辺の元で、右辺で零へ写るものとする。
このとき $b$ は射 $\beta : L^\bullet \to I^\bullet$ のホモトピー類であり、
$\beta  \circ \alpha$ が零とホモトピックとなるような擬同型
$\alpha : K^\bullet \to L^\bullet$ が存在する。補題
\ref{derived-lemma-morphisms-equal-up-to-homotopy} により $\beta$ 自身も零と
ホモトピックである。すなわち $b = 0$ である。
\end{proof}

\begin{lemma}
\label{derived-lemma-injective-resolution-ses}
$\mathcal{A}$ をアーベル圏とする。$\mathcal{A}$ は十分多くの入射対象を持つと仮定する。
$\text{Comp}^{+}(\mathcal{A})$ の任意の短完全列
$0 \to A^\bullet \to B^\bullet \to C^\bullet \to 0$ に対して、
$\text{Comp}^{+}(\mathcal{A})$ に可換図式
$$
\xymatrix{
0 \ar[r] &
A^\bullet \ar[r] \ar[d] &
B^\bullet \ar[r] \ar[d] &
C^\bullet \ar[r] \ar[d] &
0 \\
0 \ar[r] &
I_1^\bullet \ar[r] &
I_2^\bullet \ar[r] &
I_3^\bullet \ar[r] &
0
}
$$
であって、垂直な矢印が入射分解で、各行が複体の短完全列となるものが存在する。
さらに、次が成り立つ。
\begin{enumerate}
\item 任意の入射分解 $A^\bullet \to I^\bullet$ が与えられたとき、
$I_1^\bullet = I^\bullet$ と仮定してよい。
\item $n < 0$ に対して $A^n = B^n = C^n = 0$ ならば、
$n < 0$ に対して $I_j^n = 0$ と仮定してよい。
\item さらに $n < 0$ に対して $I^n = 0$ ならば、(1) と (2) を同時に満たせる。
% P02-E0045: frozen source contains the unfinished editorial placeholder.
\item ここにさらに追加する。
\end{enumerate}
\end{lemma}

\begin{proof}
ステップ 1。入射分解 $A^\bullet \to I^\bullet$ を選ぶ
（補題 \ref{derived-lemma-injective-resolutions-exist} を参照）か、与えられたものを用いる。
$\text{Comp}^{+}(\mathcal{A})$ はアーベル圏であることを思い出そう。
Homology, Lemma \ref{homology-lemma-cat-cochain-abelian} を参照せよ。
したがって写像 $A^\bullet \to I^\bullet$ に沿った押し出しを作ると、
$$
\xymatrix{
0 \ar[r] &
A^\bullet \ar[r] \ar[d] &
B^\bullet \ar[r] \ar[d] &
C^\bullet \ar[r] \ar[d] &
0 \\
0 \ar[r] &
I^\bullet \ar[r] &
E^\bullet \ar[r] &
C^\bullet \ar[r] &
0
}
$$
が得られる。$5$-補題と Homology, Lemma
\ref{homology-lemma-long-exact-sequence-cochain} の最後の主張により、
写像 $B^\bullet \to E^\bullet$ は擬同型である。下の短完全列は項ごとに
分裂することに注意する。Homology, Lemma
\ref{homology-lemma-characterize-injectives} を参照せよ。したがって
$0 \to A^\bullet \to B^\bullet \to C^\bullet \to 0$ が項ごとに分裂する場合に
補題を証明すれば十分である。

\medskip\noindent
ステップ 2。分裂を選ぶ。すなわち $B^n = A^n \oplus C^n$ と書く。
Homology, Lemma \ref{homology-lemma-ses-termwise-split-cochain} と同じく、
射 $\delta : C^\bullet \to A^\bullet[1]$ を書く。
入射分解 $f_1 : A^\bullet \to I_1^\bullet$ と
$f_3 : C^\bullet \to I_3^\bullet$ を選ぶ。
（$A^\bullet$ が入射対象からなる複体ならば
$I_1^\bullet = A^\bullet$ を用いる。$n < 0$ に対して
$A^n = C^n = 0$ ならば、補題 \ref{derived-lemma-injective-resolutions-exist} により
$n < 0$ に対して $I_1^n = I_3^n = 0$ となるように選ぶ。）
$f_3$ は各次数で単射であると仮定してよい。補題
\ref{derived-lemma-morphisms-lift} により、複体の射の等式
$\delta' \circ f_3 = f_1[1] \circ \delta$ を満たす射
$\delta' : I_3^\bullet \to I_1^\bullet[1]$ を取れる。
$I_2^n = I_1^n \oplus I_3^n$ と置き、
$$
d_{I_2}^n =
\left(
\begin{matrix}
d_{I_1}^n & (\delta')^n \\
0 & d_{I_3}^n
\end{matrix}
\right)
$$
と定める。また、写像 $B^n \to I_2^n$ を、写像
$A^n \to I_1^n$ と $C^n \to I_3^n$ の和によって定める。
すべて明らかである。
\end{proof}




\section{射影分解}
\label{derived-section-projective-resolutions}

\noindent
本節は Section \ref{derived-section-injective-resolutions} の双対である。
定義と結果の主張を与えるが、補題は再証明しない。

\begin{definition}
\label{derived-definition-projective-resolution}
$\mathcal{A}$ をアーベル圏とし、$A \in \Ob(\mathcal{A})$ とする。
{\it $A$ の射影分解}とは、複体 $P^\bullet$ と写像 $P^0 \to A$ であって、
次を満たすものをいう。
\begin{enumerate}
\item $n > 0$ に対して $P^n = 0$ である。
\item 各 $P^n$ は $\mathcal{A}$ の射影対象である。
\item 写像 $P^0 \to A$ は同型 $\Coker(d^{-1}) \to A$ を誘導する。
\item $i < 0$ に対して $H^i(P^\bullet) = 0$ である。
\end{enumerate}
したがって $P^\bullet \to A[0]$ は擬同型である。言い換えれば、複体
$$
\ldots \to P^{-1} \to P^0 \to A \to 0 \to \ldots
$$
は非輪状である。$K^\bullet$ を $\mathcal{A}$ の複体とする。
{\it $K^\bullet$ の射影分解}とは、複体 $P^\bullet$ と複体の写像
$\alpha : P^\bullet \to K^\bullet$ であって、次を満たすものをいう。
\begin{enumerate}
\item $n \gg 0$ に対して $P^n = 0$、すなわち $P^\bullet$ は上に有界である。
\item 各 $P^n$ は $\mathcal{A}$ の射影対象である。
\item 写像 $\alpha : P^\bullet \to K^\bullet$ は擬同型である。
\end{enumerate}
\end{definition}

\begin{lemma}
\label{derived-lemma-cohomology-bounded-above}
$\mathcal{A}$ をアーベル圏とし、$K^\bullet$ を $\mathcal{A}$ の複体とする。
\begin{enumerate}
\item $K^\bullet$ が射影分解を持つならば、$n \gg 0$ に対して
$H^n(K^\bullet) = 0$ である。
\item $n \gg 0$ に対して $H^n(K^\bullet) = 0$ ならば、
$L^\bullet$ が上に有界である擬同型 $L^\bullet \to K^\bullet$ が存在する。
\end{enumerate}
\end{lemma}

\begin{proof}
補題 \ref{derived-lemma-cohomology-bounded-below} の双対である。
\end{proof}

\begin{lemma}
\label{derived-lemma-projective-resolutions-exist}
$\mathcal{A}$ をアーベル圏とする。$\mathcal{A}$ は十分多くの射影対象を持つと仮定する。
\begin{enumerate}
\item $\mathcal{A}$ の任意の対象は射影分解を持つ。
\item すべての $n \gg 0$ に対して $H^n(K^\bullet) = 0$ ならば、
$K^\bullet$ は射影分解を持つ。
\item $K^\bullet$ が $n > a$ に対して $K^n = 0$ となる複体ならば、
各 $\alpha^n : P^n \to K^n$ が全射であり、$n > a$ に対して
$P^n = 0$ となる射影分解 $\alpha : P^\bullet \to K^\bullet$ が存在する。
\end{enumerate}
\end{lemma}

\begin{proof}
補題 \ref{derived-lemma-injective-resolutions-exist} の双対である。
\end{proof}

\begin{lemma}
\label{derived-lemma-projective-into-acyclic-is-zero}
$\mathcal{A}$ をアーベル圏とする。$K^\bullet$ を非輪状複体とし、
$P^\bullet$ を射影対象からなる上に有界な複体とする。
任意の射 $P^\bullet \to K^\bullet$ は零とホモトピックである。
\end{lemma}

\begin{proof}
補題 \ref{derived-lemma-acyclic-is-zero} の双対である。
\end{proof}

\begin{remark}
\label{derived-remark-easier-projective}
$\mathcal{A}$ をアーベル圏とする。$\alpha : K^\bullet \to L^\bullet$ を
複体の擬同型とし、$P^\bullet$ を射影対象からなる上に有界な複体とする。このとき
$$
\Hom_{K(\mathcal{A})}(P^\bullet, K^\bullet)
\longrightarrow
\Hom_{K(\mathcal{A})}(P^\bullet, L^\bullet)
$$
は同型である。これは Remark \ref{derived-remark-easier-proofs} の双対である。
\end{remark}

\begin{lemma}
\label{derived-lemma-morphisms-lift-projective}
$\mathcal{A}$ をアーベル圏とする。実線からなる図式
$$
\xymatrix{
K^\bullet & L^\bullet \ar[l]^\alpha \\
P^\bullet \ar[u] \ar@{-->}[ru]_\beta
}
$$
を考える。ここで $P^\bullet$ は射影対象からなる上に有界な複体であり、
$\alpha$ は擬同型である。
\begin{enumerate}
\item ホモトピーを除いて図式を可換にする複体の写像 $\beta$ が存在する。
\item $\alpha$ が各次数で全射ならば、図式を可換にする $\beta$ を取れる。
\end{enumerate}
\end{lemma}

\begin{proof}
補題 \ref{derived-lemma-morphisms-lift} の双対である。
\end{proof}

\begin{lemma}
\label{derived-lemma-morphisms-equal-up-to-homotopy-projective}
$\mathcal{A}$ をアーベル圏とする。実線からなる図式
$$
\xymatrix{
K^\bullet & L^\bullet \ar[l]^\alpha \\
P^\bullet \ar[u] \ar@{-->}[ru]_{\beta_i}
}
$$
を考える。ここで $P^\bullet$ は射影対象からなる上に有界な複体であり、
$\alpha$ は擬同型である。ホモトピーを除いて図式を可換にする任意の二射
$\beta_1, \beta_2$ は互いにホモトピックである。
\end{lemma}

\begin{proof}
補題 \ref{derived-lemma-morphisms-equal-up-to-homotopy} の双対である。
\end{proof}

\begin{lemma}
\label{derived-lemma-morphisms-from-projective-complex}
$\mathcal{A}$ をアーベル圏とする。$P^\bullet$ を射影対象からなる上に有界な
複体とし、$L^\bullet \in K(\mathcal{A})$ とする。このとき
$$
\Mor_{K(\mathcal{A})}(P^\bullet, L^\bullet)
=
\Mor_{D(\mathcal{A})}(P^\bullet, L^\bullet).
$$
\end{lemma}

\begin{proof}
補題 \ref{derived-lemma-morphisms-into-injective-complex} の双対である。
\end{proof}

\begin{lemma}
\label{derived-lemma-projective-resolution-ses}
$\mathcal{A}$ をアーベル圏とする。$\mathcal{A}$ は十分多くの射影対象を持つと仮定する。
% P02-E0046: frozen projective dual uses Comp^+ where the dual should use Comp^-.
$\text{Comp}^{+}(\mathcal{A})$ の任意の短完全列
$0 \to A^\bullet \to B^\bullet \to C^\bullet \to 0$ に対して、
$\text{Comp}^{+}(\mathcal{A})$ に可換図式
$$
\xymatrix{
0 \ar[r] &
P_1^\bullet \ar[r] \ar[d] &
P_2^\bullet \ar[r] \ar[d] &
P_3^\bullet \ar[r] \ar[d] &
0 \\
0 \ar[r] &
A^\bullet \ar[r] &
B^\bullet \ar[r] &
C^\bullet \ar[r] &
0
}
$$
であって、垂直な矢印が射影分解で、各行が複体の短完全列となるものが存在する。
実際、任意の射影分解 $P^\bullet \to C^\bullet$ が与えられたとき、
$P_3^\bullet = P^\bullet$ と仮定してよい。
\end{lemma}

\begin{proof}
補題 \ref{derived-lemma-injective-resolution-ses} の双対である。
\end{proof}

\begin{lemma}
\label{derived-lemma-precise-vanishing}
$\mathcal{A}$ をアーベル圏とする。$P^\bullet$, $K^\bullet$ を複体とし、
$n \in \mathbf{Z}$ とする。次を仮定する。
\begin{enumerate}
\item $P^\bullet$ は射影対象からなる有界複体である。
\item $i < n$ に対して $P^i = 0$ である。
\item $i \geq n$ に対して $H^i(K^\bullet) = 0$ である。
\end{enumerate}
このとき
$\Hom_{K(\mathcal{A})}(P^\bullet, K^\bullet) =
\Hom_{D(\mathcal{A})}(P^\bullet, K^\bullet) = 0$ である。
\end{lemma}

\begin{proof}
第一の等式は補題 \ref{derived-lemma-morphisms-from-projective-complex} から従う。
Remark \ref{derived-remark-truncation-distinguished-triangle} により完全三角
$$
(\tau_{\leq n - 1}K^\bullet, K^\bullet, \tau_{\geq n}K^\bullet, f, g, h)
$$
があることに注意する。したがって補題 \ref{derived-lemma-representable-homological} により、
$\Hom_{K(\mathcal{A})}(P^\bullet, \tau_{\leq n - 1}K^\bullet) = 0$ および
$\Hom_{K(\mathcal{A})}(P^\bullet, \tau_{\geq n} K^\bullet) = 0$
を示せば十分である。第一の消滅は自明であり、第二は補題
\ref{derived-lemma-projective-into-acyclic-is-zero} である。
\end{proof}

\begin{lemma}
\label{derived-lemma-lift-map}
$\mathcal{A}$ をアーベル圏とする。$\beta : P^\bullet \to L^\bullet$ と
$\alpha : E^\bullet \to L^\bullet$ を複体の写像とし、
$n \in \mathbf{Z}$ とする。次を仮定する。
\begin{enumerate}
\item $P^\bullet$ は射影対象からなる有界複体であり、
$i < n$ に対して $P^i = 0$ である。
\item $H^i(\alpha)$ は $i > n$ に対して同型であり、$i = n$ に対して全射である。
\end{enumerate}
このとき、$\alpha \circ \gamma$ と $\beta$ がホモトピックになる複体の写像
$\gamma : P^\bullet \to E^\bullet$ が存在する。
\end{lemma}

\begin{proof}
錐 $C^\bullet = C(\alpha)^\bullet$ と写像
$i : L^\bullet \to C^\bullet$ を考える。補題
\ref{derived-lemma-precise-vanishing} により $i \circ \beta$ は零である。
したがって補題 \ref{derived-lemma-representable-homological} により
$\beta$ を $E^\bullet$ へ持ち上げられる。
\end{proof}





















\section{右導来関手と入射分解}
\label{derived-section-right-derived-functor}

\noindent
ここまでの結果を用いると、十分多くの入射対象を持つアーベル圏から一般の
アーベル圏への加法的関手の右導来関手を定義できる。

\begin{lemma}
\label{derived-lemma-injective-acyclic}
$\mathcal{A}$ をアーベル圏とする。$I \in \Ob(\mathcal{A})$ を入射対象とし、
$I^\bullet$ を $\mathcal{A}$ の入射対象からなる下に有界な複体とする。
\begin{enumerate}
\item 任意の三角圏 $\mathcal{D}$ への任意の完全関手
$F : K^{+}(\mathcal{A}) \to \mathcal{D}$ に対して、
$I^\bullet$ は $\text{Qis}^{+}(\mathcal{A})$ に関する $RF$ を計算する。
\item 任意のアーベル圏 $\mathcal{B}$ への任意の加法的関手
$F : \mathcal{A} \to \mathcal{B}$ に対して、$I$ は右非輪状である。
\end{enumerate}
\end{lemma}

\begin{proof}
% P02-E0047: frozen source says “a direct consequences”.
(2) は (1) と定義 \ref{derived-definition-derived-functor} の直接の帰結である。
(1) を証明するため、複体への擬同型
$\alpha : I^\bullet \to K^\bullet$ を取る。補題
\ref{derived-lemma-morphisms-lift} により、$\alpha$ は左逆を持つ。したがって圏
$I^\bullet/\text{Qis}^{+}(\mathcal{A})$ は
$\text{id} : I^\bullet \to I^\bullet$ を値として本質的に定値である。
ゆえに ind-対象
$$
I^\bullet/\text{Qis}^{+}(\mathcal{A}) \longrightarrow \mathcal{D}, \quad
(I^\bullet \to K^\bullet) \longmapsto F(K^\bullet)
$$
も $F(I^\bullet)$ を値として本質的に定値である。これで (1) が示された。
定義 \ref{derived-definition-right-derived-functor-defined} および
\ref{derived-definition-computes} を参照せよ。
\end{proof}

\begin{lemma}
\label{derived-lemma-enough-injectives-right-derived}
$\mathcal{A}$ を十分多くの入射対象を持つアーベル圏とする。
\begin{enumerate}
\item 任意の三角圏 $\mathcal{D}$ への任意の完全関手
$F : K^{+}(\mathcal{A}) \to \mathcal{D}$ に対して、右導来関手
$$
RF : D^{+}(\mathcal{A}) \longrightarrow \mathcal{D}
$$
は至る所で定義される。
\item 任意のアーベル圏 $\mathcal{B}$ への任意の加法的関手
$F : \mathcal{A} \to \mathcal{B}$ に対して、右導来関手
$$
RF : D^{+}(\mathcal{A}) \longrightarrow D^{+}(\mathcal{B})
$$
は至る所で定義される。
\end{enumerate}
\end{lemma}

\begin{proof}
第二の主張には補題 \ref{derived-lemma-injective-acyclic} と命題
\ref{derived-proposition-enough-acyclics} を組み合わせる。第一の主張には補題
\ref{derived-lemma-injective-resolutions-exist}、補題
\ref{derived-lemma-injective-acyclic}、補題
\ref{derived-lemma-existence-computes}、および式
(\ref{derived-equation-everywhere}) を組み合わせる。
\end{proof}

\begin{lemma}
\label{derived-lemma-right-derived-properties}
$\mathcal{A}$ を十分多くの入射対象を持つアーベル圏とし、
$F : \mathcal{A} \to \mathcal{B}$ を加法的関手とする。
\begin{enumerate}
\item 関手 $RF$ は完全関手
$D^{+}(\mathcal{A}) \to D^{+}(\mathcal{B})$ である。
\item 関手 $RF$ は完全関手
$K^{+}(\mathcal{A}) \to D^{+}(\mathcal{B})$ を誘導する。
\item 関手 $RF$ は $\delta$-関手
$\text{Comp}^{+}(\mathcal{A}) \to D^{+}(\mathcal{B})$ を誘導する。
\item 関手 $RF$ は $\delta$-関手
$\mathcal{A} \to D^{+}(\mathcal{B})$ を誘導する。
\end{enumerate}
\end{lemma}

\begin{proof}
本補題は、これまでに得た結果のいくつかをまとめたものにすぎない。
補題 \ref{derived-lemma-enough-injectives-right-derived} により $RF$ は
至る所で定義されることに注意する。参照先を挙げる。
\begin{enumerate}
\item 導来関手が完全であることは、補題
\ref{derived-lemma-2-out-of-3-defined} に帰着する。
\item $K^{+}(\mathcal{A}) \to D^{+}(\mathcal{A})$ は完全であり、
完全関手の合成は完全なので成り立つ。
\item $\text{Comp}^{+}(\mathcal{A}) \to D^{+}(\mathcal{A})$ は
$\delta$-関手なので成り立つ。補題
\ref{derived-lemma-derived-canonical-delta-functor} を参照せよ。
\item $\mathcal{A} \to \text{Comp}^{+}(\mathcal{A})$ は完全であり、
$\delta$-関手に完全関手を前合成すると $\delta$-関手が得られるので成り立つ。
\end{enumerate}
\end{proof}

\begin{lemma}
\label{derived-lemma-higher-derived-functors}
$\mathcal{A}$ を十分多くの入射対象を持つアーベル圏とし、
$F : \mathcal{A} \to \mathcal{B}$ を左完全関手とする。
\begin{enumerate}
\item $\text{Comp}^{+}(\mathcal{A})$ の複体の任意の短完全列
$0 \to A^\bullet \to B^\bullet \to C^\bullet \to 0$ に対して、
付随する長完全列
$$
\ldots \to
H^i(RF(A^\bullet)) \to
H^i(RF(B^\bullet)) \to
H^i(RF(C^\bullet)) \to
H^{i + 1}(RF(A^\bullet)) \to \ldots
$$
がある。
\item 関手 $R^iF : \mathcal{A} \to \mathcal{B}$ は $i < 0$ に対して零である。
また $R^0F = F : \mathcal{A} \to \mathcal{B}$ である。
\item $i > 0$ かつ $I$ が入射的ならば $R^iF(I) = 0$ である。
\item 列 $(R^iF, \delta)$ は $\mathcal{A}$ から $\mathcal{B}$ への
普遍 $\delta$-関手をなす（Homology, Definition
\ref{homology-definition-universal-delta-functor} を参照）。
\end{enumerate}
\end{lemma}

\begin{proof}
本補題は、これまでに得た結果のいくつかをまとめたものにすぎない。
補題 \ref{derived-lemma-enough-injectives-right-derived} により $RF$ は
至る所で定義されることに注意する。参照先を挙げる。
\begin{enumerate}
\item 補題 \ref{derived-lemma-right-derived-properties} の (3) と、
$D^{+}(\mathcal{B})$ に対するコホモロジー長完全列
(\ref{derived-equation-long-exact-cohomology-sequence-D}) から従う。
\item これは補題 \ref{derived-lemma-left-exact-higher-derived} である。
\item これは入射対象が非輪状であるという事実である。
\item これは補題 \ref{derived-lemma-right-derived-delta-functor} である。
\end{enumerate}
\end{proof}



\section{Cartan--Eilenberg 分解}
\label{derived-section-cartan-eilenberg}

\noindent
本節は拡張の余地がある。ここでの内容は一般化でき、より多くの場合に適用できる。
分解に入射対象を用いる必要はなく、状況によっては非有界の場合にもこの方法は働く。

\begin{definition}
\label{derived-definition-cartan-eilenberg}
$\mathcal{A}$ をアーベル圏とし、$K^\bullet$ を下に有界な複体とする。
$K^\bullet$ の {\it Cartan--Eilenberg 分解}とは、二重複体
$I^{\bullet, \bullet}$ と複体の射
$\epsilon : K^\bullet \to I^{\bullet, 0}$ であって、次の性質を持つものをいう。
\begin{enumerate}
% P02-E0048: frozen source uses the wrong article in “a i”.
\item ある $i \ll 0$ が存在し、すべての $p < i$ とすべての $q$ に対して
$I^{p, q} = 0$ である。
\item $q < 0$ ならば $I^{p, q} = 0$ である。
\item 複体 $I^{p, \bullet}$ は $K^p$ の入射分解である。
\item 複体 $\Ker(d_1^{p, \bullet})$ は $\Ker(d_K^p)$ の入射分解である。
\item 複体 $\Im(d_1^{p, \bullet})$ は $\Im(d_K^p)$ の入射分解である。
\item 複体 $H^p_I(I^{\bullet, \bullet})$ は $H^p(K^\bullet)$ の入射分解である。
\end{enumerate}
\end{definition}

\begin{lemma}
\label{derived-lemma-cartan-eilenberg}
$\mathcal{A}$ を十分多くの入射対象を持つアーベル圏とし、
$K^\bullet$ を下に有界な複体とする。このとき $K^\bullet$ の
Cartan--Eilenberg 分解が存在する。
\end{lemma}

\begin{proof}
$p < n$ に対して $K^p = 0$ と仮定する。$K^\bullet$ を次のように短完全列へ
分解する。$Z^p = \Ker(d^p)$, $B^p = \Im(d^{p - 1})$,
$H^p = Z^p/B^p$ と置き、
$$
\begin{matrix}
0 \to Z^n \to K^n \to B^{n + 1} \to 0 \\
0 \to B^{n + 1} \to Z^{n + 1} \to H^{n + 1} \to 0 \\
0 \to Z^{n + 1} \to K^{n + 1} \to B^{n + 2} \to 0 \\
0 \to B^{n + 2} \to Z^{n + 2} \to H^{n + 2} \to 0 \\
\ldots
\end{matrix}
$$
を考える。$p < n$ に対して $I^{p, q} = 0$ と置く。
帰納的に、次のように入射分解を選ぶ。
\begin{enumerate}
\item $m < 0$ に対して $J_Z^{n, m} = 0$ となる入射分解
$Z^n \to J_Z^{n, \bullet}$ を選ぶ。補題
\ref{derived-lemma-injective-resolutions-exist} を参照せよ。
\item 補題 \ref{derived-lemma-injective-resolution-ses} を用いて、入射分解
$K^n \to I^{n, \bullet}$ と
$B^{n + 1} \to J_B^{n + 1, \bullet}$ で、$m < 0$ に対して
$I^{n, m} = 0$、かつ $m < 0$ に対して $J_B^{n + 1, m} = 0$ となるもの、
および短完全列
$0 \to Z^n \to K^n \to B^{n + 1} \to 0$ と両立する複体の完全列
$0 \to J_Z^{n, \bullet} \to I^{n, \bullet} \to J_B^{n + 1, \bullet} \to 0$
を選ぶ。
\item 補題 \ref{derived-lemma-injective-resolution-ses} を用いて、入射分解
$Z^{n + 1} \to J_Z^{n + 1, \bullet}$ と
$H^{n + 1} \to J_H^{n + 1, \bullet}$ で、$m < 0$ に対して
$J_Z^{n + 1, m} = 0$、かつ $m < 0$ に対して $J_H^{n + 1, m} = 0$
となるもの、および
短完全列 $0 \to B^{n + 1} \to Z^{n + 1} \to H^{n + 1} \to 0$ と
両立する複体の完全列
$0 \to J_B^{n + 1, \bullet} \to J_Z^{n + 1, \bullet}
\to J_H^{n + 1, \bullet} \to 0$
を選ぶ。
\item 以下同様である。
\end{enumerate}
写像 $d_1^\bullet : I^{p, \bullet} \to I^{p + 1, \bullet}$ を合成
$I^{p, \bullet} \to J_B^{p + 1, \bullet} \to
J_Z^{p + 1, \bullet} \to I^{p + 1, \bullet}$ とすれば、すべて明らかである。
\end{proof}

\begin{lemma}
\label{derived-lemma-two-ss-complex-functor}
$F : \mathcal{A} \to \mathcal{B}$ をアーベル圏の間の左完全関手とする。
$K^\bullet$ を $\mathcal{A}$ の下に有界な複体とし、
$I^{\bullet, \bullet}$ を $K^\bullet$ の Cartan--Eilenberg 分解とする。
二重複体 $F(I^{\bullet, \bullet})$ に付随するスペクトル系列
$({}'E_r, {}'d_r)_{r \geq 0}$ と $({}''E_r, {}''d_r)_{r \geq 0}$ は関係式
$$
{}'E_1^{p, q} = R^qF(K^p)
\quad
\text{and}
\quad
{}''E_2^{p, q} = R^pF(H^q(K^\bullet))
$$
を満たす。さらに、これらのスペクトル系列は有界で、
$H^*(RF(K^\bullet))$ に収束し、$H^n(RF(K^\bullet))$ 上に誘導される
付随フィルトレーションは有限である。
\end{lemma}

\begin{proof}
以下の注意を断りなく用いる。
\begin{enumerate}
\item $I^{p, \bullet}$ は $K^p$ の入射分解なので、$RF$ は
$K^p[0]$ で定義され、その値は $F(I^{p, \bullet})$ である。
\item
% P02-E0049: frozen proof uses p where the vertical homology degree is q.
$H^p_I(I^{\bullet, \bullet})$ は $H^p(K^\bullet)$ の入射分解なので、
導来関手 $RF$ は $H^p(K^\bullet)[0]$ で定義され、その値は
$F(H^p_I(I^{\bullet, \bullet}))$ である。
\item Homology, Lemma \ref{homology-lemma-double-complex-gives-resolution}
により、全複体 $\text{Tot}(I^{\bullet, \bullet})$ は $K^\bullet$ の入射分解である。
したがって $RF$ は $K^\bullet$ で定義され、その値は
$F(\text{Tot}(I^{\bullet, \bullet}))$ である。
\end{enumerate}
二重複体 $L^{\bullet, \bullet} = F(I^{\bullet, \bullet})$ に付随する二つの
スペクトル系列を考える。Homology, Lemma
\ref{homology-lemma-ss-double-complex} を参照せよ。これらはいずれも有界で、
$H^*(\text{Tot}(L^{\bullet, \bullet}))$ に収束し、
$H^n(\text{Tot}(L^{\bullet, \bullet}))$ 上に有限フィルトレーションを誘導する。
Homology, Lemma \ref{homology-lemma-first-quadrant-ss} を参照せよ。
$$
\text{Tot}(L^{\bullet, \bullet}) =
\text{Tot}(F(I^{\bullet, \bullet})) =
F(\text{Tot}(I^{\bullet, \bullet}))
$$
は $RF(K^\bullet)$ を計算するので、補題の最後の主張が成り立つ。

\medskip\noindent
第一スペクトル系列を計算する。
${}'E_1^{p, q} = H^q(L^{p, \bullet})$、言い換えれば所望のとおり
$$
{}'E_1^{p, q} = H^q(F(I^{p, \bullet})) = R^qF(K^p)
$$
である。後で用いるため、写像
${}'d_1^{p, q} : {}'E_1^{p, q} \to {}'E_1^{p + 1, q}$ は、
$K^p \to K^{p + 1}$ と $R^qF$ が関手であることから誘導される写像
$R^qF(K^p) \to R^qF(K^{p + 1})$ であることに注意する。

\medskip\noindent
第二スペクトル系列を計算する。
${}''E_1^{p, q} = H^q(L^{\bullet, p}) = H^q(F(I^{\bullet, p}))$ である。
複体 $I^{\bullet, p}$ は下に有界で、入射対象からなり、さらに微分の各核、像、
およびコホモロジー群は $\mathcal{A}$ の入射対象であることに注意する。
したがって微分を分裂できる。すなわち各微分は、直和因子への分裂全射である。
$F$ を適用した後にも同じことが成り立つ。ゆえに
${}''E_1^{p, q} = F(H^q(I^{\bullet, p})) = F(H^q_I(I^{\bullet, p}))$ である。
% P02-E0049: second frozen occurrence again uses p instead of q.
この上の微分は、$H^p(K^\bullet)$ の入射分解である複体
$H^p_I(I^{\bullet, \bullet})$ の微分に $F$ を適用したものの
$(-1)^q$ 倍である。これで $E_2$ 項の記述が得られる。
\end{proof}

\begin{remark}
\label{derived-remark-functorial-ss}
補題 \ref{derived-lemma-two-ss-complex-functor} のスペクトル系列は、複体
$K^\bullet$ に関して関手的である。これは Cartan--Eilenberg 分解の
関手性から従う。一方、これらはいずれも、$\mathcal{A}$ の
フィルトレーション付き複体に付随させられる、より一般のスペクトル系列の例である。
その関手性は構成から従う。フィルトレーション付き導来圏を扱う節でこの点に戻る。
Remark \ref{derived-remark-final-functorial} を参照せよ。
\end{remark}









\section{右導来関手の合成}
\label{derived-section-composition-right-derived-functors}

\noindent
合成の右導来関手を計算できる場合がある。
% P02-E0050: frozen source says “Suppose ... be abelian categories”.
$\mathcal{A}, \mathcal{B}, \mathcal{C}$ をアーベル圏とする。
$F : \mathcal{A} \to \mathcal{B}$ および
$G : \mathcal{B} \to \mathcal{C}$ を左完全関手とする。右導来関手
$RF : D^{+}(\mathcal{A}) \to D^{+}(\mathcal{B})$,
$RG : D^{+}(\mathcal{B}) \to D^{+}(\mathcal{C})$, および
$R(G \circ F) : D^{+}(\mathcal{A}) \to D^{+}(\mathcal{C})$
が至る所で定義されると仮定する。このとき
$D^{+}(\mathcal{A})$ から $D^{+}(\mathcal{C})$ への関手の間の
標準自然変換
$$
t : R(G \circ F) \longrightarrow RG \circ RF
$$
が存在する。補題 \ref{derived-lemma-compose-derived-functors-general} を参照せよ。
この変換は常に同型であるとは限らない。

\begin{lemma}
\label{derived-lemma-compose-derived-functors}
$\mathcal{A}, \mathcal{B}, \mathcal{C}$ をアーベル圏とする。
$F : \mathcal{A} \to \mathcal{B}$ および
$G : \mathcal{B} \to \mathcal{C}$ を左完全関手とする。
$\mathcal{A}$, $\mathcal{B}$ は十分多くの入射対象を持つと仮定する。
次の条件は同値である。
\begin{enumerate}
\item $\mathcal{A}$ の各入射対象 $I$ に対して、$F(I)$ は
$G$ に関して右非輪状である。
\item 標準写像
% P02-E0051: frozen source has a premature display period and omits the article in “is isomorphism”.
$$
t : R(G \circ F) \longrightarrow RG \circ RF.
$$
は $D^{+}(\mathcal{A})$ から $D^{+}(\mathcal{C})$ への関手の同型である。
\end{enumerate}
\end{lemma}

\begin{proof}
(2) が成り立つなら、同型 $t$ を $RF(I) = F(I)$ において評価すれば
(1) が従う。逆に (1) が成り立つと仮定する。
$A^\bullet$ を $\mathcal{A}$ の下に有界な複体とし、
入射分解 $A^\bullet \to I^\bullet$ を選ぶ。写像 $t$ は
（補題 \ref{derived-lemma-compose-derived-functors-general} の証明を参照）
次の写像によって与えられる。
% P02-E0052: frozen source has an extra closing parenthesis in G(F(I^\bullet))).
$$
R(G \circ F)(A^\bullet) =
(G \circ F)(I^\bullet) =
G(F(I^\bullet))) \to
RG(F(I^\bullet)) =
RG(RF(A^\bullet))
$$
ここで矢印は補題 \ref{derived-lemma-leray-acyclicity} により同型である。
\end{proof}

\begin{lemma}[Grothendieck スペクトル系列]
\label{derived-lemma-grothendieck-spectral-sequence}
% P02-E0053: frozen opening is a sentence fragment (“With assumptions ... and assuming ...”).
補題 \ref{derived-lemma-compose-derived-functors} と同じ仮定のもとで、
その同値な条件 (1), (2) が成り立つと仮定する。
$X$ を $D^{+}(\mathcal{A})$ の対象とする。
$\mathcal{C}$ の二重次数付き対象 $E_r$ からなるスペクトル系列
$(E_r, d_r)_{r \geq 0}$ で、$d_r$ の双次数が
$(r, - r + 1)$ であり、かつ
$$
E_2^{p, q} = R^pG(H^q(RF(X)))
$$
となるものが存在する。さらに、このスペクトル系列は有界で、
$H^*(R(G \circ F)(X))$ に収束し、各
$H^n(R(G \circ F)(X))$ 上に有限フィルトレーションを誘導する。
\end{lemma}

\noindent
$\mathcal{A}$ の対象 $A$ に対しては
$E_2^{p, q} = R^pG(R^qF(A))$ が
$R^{p + q}(G \circ F)(A)$ に収束する。

\begin{proof}
$X$ を下に有界な複体 $A^\bullet$ で表してよい。
入射分解 $A^\bullet \to I^\bullet$ を選ぶ。
補題 \ref{derived-lemma-cartan-eilenberg} を用いて Cartan--Eilenberg 分解
$F(I^\bullet) \to I^{\bullet, \bullet}$ を選ぶ。
補題 \ref{derived-lemma-two-ss-complex-functor} の第二スペクトル系列を適用する。
\end{proof}







\section{分解関手}
\label{derived-section-derived-category}

\noindent
$\mathcal{A}$ を十分多くの入射対象を持つアーベル圏とする。
$\mathcal{I}$ で、対象が $\mathcal{A}$ の入射対象である
$\mathcal{A}$ の充満加法部分圏を表す。この場合
$K^{+}(\mathcal{I})$ と $D^{+}(\mathcal{A})$ は同値であることが分かる
（命題 \ref{derived-proposition-derived-category} を参照）。
したがって多くの目的において、この場合の $D^{+}(\mathcal{A})$ を
（より把握しやすい）圏 $K^{+}(\mathcal{I})$ と考えることが有用である。

\begin{proposition}
\label{derived-proposition-derived-category}
$\mathcal{A}$ をアーベル圏とする。
$\mathcal{A}$ は十分多くの入射対象を持つと仮定する。
$\mathcal{I} \subset \mathcal{A}$ で、対象が $\mathcal{A}$ の
入射対象である厳密充満加法部分圏を表す。関手
$$
K^{+}(\mathcal{I}) \longrightarrow D^{+}(\mathcal{A})
$$
は完全かつ充満忠実で本質的全射である。すなわち三角圏の同値である。
\end{proposition}

\begin{proof}
この関手が完全であることは明らかである。
補題 \ref{derived-lemma-injective-resolutions-exist} により本質的全射である。
充満忠実性は補題 \ref{derived-lemma-morphisms-into-injective-complex} から従う。
\end{proof}

\noindent
命題 \ref{derived-proposition-derived-category} により、分解関手を見つけることができる。
実際、アーベル圏 $\mathcal{A}$ が「大きい」圏、すなわち対象が類をなす場合にも、
分解関手の存在を証明できることがある。

\begin{definition}
\label{derived-definition-localization-functor}
$\mathcal{A}$ を十分多くの入射対象を持つアーベル圏とする。
$\mathcal{A}$ の{it 分解関手}\footnote{これはおそらく標準的な用語ではない。}
とは、次のデータのことである。
\begin{enumerate}
\item すべての $K^\bullet \in \Ob(K^{+}(\mathcal{A}))$ に対する、
入射対象からなる下に有界な複体 $j(K^\bullet)$。
\item すべての $K^\bullet \in \Ob(K^{+}(\mathcal{A}))$ に対する、
擬同型 $i_{K^\bullet} : K^\bullet \to j(K^\bullet)$。
\end{enumerate}
\end{definition}

\begin{lemma}
\label{derived-lemma-resolution-functor}
$\mathcal{A}$ を十分多くの入射対象を持つアーベル圏とする。
分解関手 $(j, i)$ が与えられたとき、$j$ を関手にし、
$i$ を $2$-同型にして、$2$-可換図
$$
\xymatrix{
K^{+}(\mathcal{A}) \ar[rd] \ar[rr]_j & & K^{+}(\mathcal{I}) \ar[ld] \\
& D^{+}(\mathcal{A})
}
$$
を生じさせる仕方が一意に存在する。ここで $\mathcal{I}$ は、
入射対象からなる $\mathcal{A}$ の充満加法部分圏である。
\end{lemma}

\begin{proof}
$K^{+}(\mathcal{A})$ の各射 $\alpha : K^\bullet \to L^\bullet$
に対して、$K^{+}(\mathcal{I})$ における一意な射
$j(\alpha) : j(K^\bullet) \to j(L^\bullet)$ で、
$$
\xymatrix{
K^\bullet \ar[r]_\alpha \ar[d]_{i_{K^\bullet}} &
L^\bullet \ar[d]^{i_{L^\bullet}} \\
j(K^\bullet) \ar[r]^{j(\alpha)} & j(L^\bullet)
}
$$
が $K^{+}(\mathcal{A})$ において可換となるものが存在する。
これを見るには補題 \ref{derived-lemma-morphisms-lift} と
\ref{derived-lemma-morphisms-equal-up-to-homotopy}、または同値な補題
\ref{derived-lemma-morphisms-into-injective-complex} を用いればよい。
一意性から $j$ は関手であり、図の可換性から $i$ は、
補題の主張にある圏の図の $2$-可換性を立証する $2$-射を与える。
\end{proof}

\begin{lemma}
\label{derived-lemma-into-derived-category}
$\mathcal{A}$ をアーベル圏とする。
$\mathcal{A}$ は十分多くの入射対象を持つと仮定する。
このとき分解関手 $j$ が存在し、関手の一意な同型を除いて一意である。
\end{lemma}

\begin{proof}
$K^{+}(\mathcal{A})$ のすべての対象 $K^\bullet$ からなる集合を考える。
（別途言及しない限り、どの圏も対象の集合を持つという規約を思い出そう。）
補題 \ref{derived-lemma-injective-resolutions-exist} により、各対象は入射分解を持つ。
選択公理により、各 $K^\bullet$ に対して入射分解
$i_{K^\bullet} : K^\bullet \to j(K^\bullet)$ を選ぶことができる。
\end{proof}

\begin{lemma}
\label{derived-lemma-j-is-exact}
$\mathcal{A}$ を十分多くの入射対象を持つアーベル圏とする。
任意の分解関手
$j : K^{+}(\mathcal{A}) \to K^{+}(\mathcal{I})$
は完全である。
\end{lemma}

\begin{proof}
定義 \ref{derived-definition-localization-functor} の標準写像を
$i_{K^\bullet} : K^\bullet \to j(K^\bullet)$ と表す。
まず、関手的同型 $j(K^\bullet[1]) \to j(K^\bullet)[1]$ の存在を論じる。
図
$$
\xymatrix{
K^\bullet[1] \ar[d]^{i_{K^\bullet[1]}} \ar@{=}[rr] & &
K^\bullet[1] \ar[d]^{i_{K^\bullet}[1]} \\
j(K^\bullet[1]) \ar@{..>}[rr]^{\xi_{K^\bullet}} & & j(K^\bullet)[1]
}
$$
を考える。補題 \ref{derived-lemma-morphisms-lift} と
\ref{derived-lemma-morphisms-equal-up-to-homotopy} により、
$K^{+}(\mathcal{I})$ において一意な点線矢印 $\xi_{K^\bullet}$ で、
図を $K^{+}(\mathcal{A})$ において可換にするものが存在する。
これが関手的同型を与えることの確認は省略する。
（ヒント：補題 \ref{derived-lemma-morphisms-equal-up-to-homotopy} を再び用いよ。）

\medskip\noindent
$(K^\bullet, L^\bullet, M^\bullet, f, g, h)$ を
$K^{+}(\mathcal{A})$ の完全三角とする。
$(j(K^\bullet), j(L^\bullet), j(M^\bullet), j(f), j(g),
\xi_{K^\bullet} \circ j(h))$ が
$K^{+}(\mathcal{I})$ の完全三角であることを示さなければならない。
$K^{+}(\mathcal{A})$ において、縦矢印が擬同型 $i_K, i_L, i_M$ である
次の可換図を得る。
$$
\xymatrix{
K^\bullet \ar[r]_f \ar[d] &
L^\bullet \ar[r]_g \ar[d] &
M^\bullet \ar[rr]_h \ar[d] & &
K^\bullet[1] \ar[d] \\
j(K^\bullet) \ar[r]^{j(f)} &
j(L^\bullet) \ar[r]^{j(g)} &
j(M^\bullet) \ar[rr]^{\xi_{K^\bullet} \circ j(h)} & &
j(K^\bullet)[1]
}
$$
したがって
$(j(K^\bullet), j(L^\bullet), j(M^\bullet), j(f), j(g),
\xi_{K^\bullet} \circ j(h))$ の $D^{+}(\mathcal{A})$ における像は
完全三角と同型であり、従って TR1 により完全三角である。
ゆえに補題 \ref{derived-lemma-exact-equivalence} から
$(j(K^\bullet), j(L^\bullet), j(M^\bullet), j(f), j(g),
\xi_{K^\bullet} \circ j(h))$ は $K^{+}(\mathcal{I})$ における
完全三角である。
\end{proof}

\begin{lemma}
\label{derived-lemma-resolution-functor-quasi-inverse}
$\mathcal{A}$ を十分多くの入射対象を持つアーベル圏とする。
$j$ を分解関手とする。自然な関手を
$Q : K^{+}(\mathcal{A}) \to D^{+}(\mathcal{A})$ と書く。
このとき $j = j' \circ Q$ であり、標準関手
$K^{+}(\mathcal{I}) \to D^{+}(\mathcal{A})$ の擬逆となる一意な関手
$j' : D^{+}(\mathcal{A}) \to K^{+}(\mathcal{I})$ が存在する。
\end{lemma}

\begin{proof}
補題 \ref{derived-lemma-bounded-derived} により、関手 $Q$ は局所化である。
$j'$ の存在を証明するには、補題
\ref{derived-lemma-universal-property-localization} により、
$\text{Qis}^{+}(\mathcal{A})$ の任意の元が関手 $j$ によって
同型に写されることを示せば十分である。
補題 \ref{derived-lemma-j-is-exact} の証明にある可換な正方形を考える。
この正方形において $\alpha$ が擬同型なら、
$i_{L^\bullet}\circ\alpha=j(\alpha)\circ i_{K^\bullet}$
も擬同型であり、従って $j(\alpha)$ も擬同型である。
ゆえに命題 \ref{derived-proposition-derived-category} により、射
$j(\alpha)$ は $K^+(\mathcal{I})$ における同型である。
$j'$ が $K^{+}(\mathcal{I}) \to D^{+}(\mathcal{A})$ の擬逆であることの
確認は省略する。
\end{proof}

\begin{remark}
\label{derived-remark-big-localization}
$\mathcal{A}$ が、アーベル群の圏のような十分多くの入射対象を持つ
「大きい」アーベル圏であると仮定する。この場合、類の上を走る量化子に
選択公理を用いることはできないため、分解関手の構成には少し注意が必要である。
しかし、この補題の証明が任意の二つの局所化関手は標準的に同型であることを
示す点に注意する。すなわち、下に有界な複体 $K^\bullet$ から、
入射対象からなる下に有界な複体への擬同型
$i : K^\bullet \to I^\bullet$ および
$i' : K^\bullet \to J^\bullet$ が与えられると、
$K^{+}(\mathcal{I})$ において一意の（！）射
$a : I^\bullet \to J^\bullet$ で、
% P02-E0054: frozen source has the ill-typed composition i' = i \circ a; it should compose a after i.
$K^{+}(\mathcal{I})$ における射として $i' = i \circ a$ を満たすものが存在する。
したがって唯一の問題は存在性であり、次節でこれをどのように扱うかを見る。
\end{remark}









\section{関手的入射埋め込みと分解関手}
\label{derived-section-functorial-injective-resolutions}

\noindent
本節では、圏 $\mathcal{A}$ が関手的入射埋め込みを持つ場合に、
分解関手 $K^{+}(\mathcal{A}) \to K^{+}(\mathcal{I})$ の構成を
やり直す。理由は二つある。(1) 証明がより容易であり、(2)
$\mathcal{A}$ が「大きい」アーベル圏である場合にも構成が機能するからである。
下の Remark \ref{derived-remark-big-abelian-category} を参照せよ。

\medskip\noindent
$\mathcal{A}$ をアーベル圏とする。以前と同様に $\mathcal{I}$ で、
入射対象からなる $\mathcal{A}$ の充満加法部分圏を表す。
矢印 $\alpha : K^\bullet \to I^\bullet$ からなる圏
$\text{InjRes}(\mathcal{A})$ を考える。ここで $K^\bullet$ は
$\mathcal{A}$ の下に有界な複体、$I^\bullet$ は
$\mathcal{A}$ の入射対象からなる下に有界な複体であり、
$\alpha$ は擬同型である。言い換えると、$\alpha$ は入射分解であり、
$K^\bullet$ は下に有界である。明らかな関手
$$
s : \text{InjRes}(\mathcal{A}) \longrightarrow \text{Comp}^{+}(\mathcal{A})
$$
があり、$(\alpha : K^\bullet \to I^\bullet) \mapsto K^\bullet$
によって定義される。また、関手
$$
t : \text{InjRes}(\mathcal{A}) \longrightarrow K^{+}(\mathcal{I})
$$
があり、$(\alpha : K^\bullet \to I^\bullet) \mapsto I^\bullet$
によって定義される。

\begin{lemma}
\label{derived-lemma-functorial-injective-resolutions}
$\mathcal{A}$ をアーベル圏とする。$\mathcal{A}$ は関手的入射埋め込みを
持つと仮定する。Homology, Definition
\ref{homology-definition-functorial-injective-embedding} を参照せよ。
\begin{enumerate}
\item $s \circ inj = \text{id}$ を満たす関手
$inj : \text{Comp}^{+}(\mathcal{A}) \to \text{InjRes}(\mathcal{A})$
が存在する。
\item $s \circ inj = \text{id}$ を満たす任意の関手
$inj : \text{Comp}^{+}(\mathcal{A}) \to \text{InjRes}(\mathcal{A})$
から、分解関手が得られる。定義 \ref{derived-definition-localization-functor}
を参照せよ。
\end{enumerate}
\end{lemma}

\begin{proof}
$A \mapsto (A \to J(A))$ を関手的入射埋め込みとする。
Homology, Definition \ref{homology-definition-functorial-injective-embedding}
を参照せよ。まず $J(0) = 0$ と仮定してよいことに注意する。
実際、そうでなければ任意の対象 $A$ に対して $0 \to A \to 0$ があり、
これは直和分解 $J(A) = J(0) \oplus \Ker(J(A) \to J(0))$ を与える。
関手的な射 $A \to J(A)$ は第二直和因子へ写らなければならないことに
注意する。したがって必要なら、関手を
$J'(A) = \Ker(J(A) \to J(0))$ で置き換えることができる。

\medskip\noindent
$K^\bullet$ を $\mathcal{A}$ の下に有界な複体とする。
$p < B$ ならば $K^p = 0$ であるとする。
入射対象からなる二重複体 $I^{\bullet, \bullet}$ と写像
$\alpha : K^\bullet \to I^{\bullet, 0}$ で、$\alpha$ が
$K^\bullet$ と $I^{\bullet, \bullet}$ に付随する全複体との間の
擬同型を誘導するものを構成する。
まず $q < 0$ のときは常に $I^{p, q} = 0$ と置く。
次に $I^{p, 0} = J(K^p)$ と置き、
$\alpha^p : K^p \to I^{p, 0}$ を関手的埋め込みとする。
$J$ は関手なので $I^{\bullet, 0}$ は複体であり、$\alpha$ は
複体の射である。各 $\alpha^p$ は単射である。また $J(0) = 0$
なので、$p < B$ に対して $I^{p, 0} = 0$ である。
次に $I^{p, 1} = J(\Coker(K^p \to I^{p, 0}))$ と置く。
再び関手性から $I^{\bullet, 1}$ は複体である。また
$p < B$ に対して $I^{p, 1} = 0$ である。
$K^p$ が $\Ker(I^{p, 0} \to I^{p, 1})$ の上へ同型に写ることも明らかである。
第三段階として
$I^{p, 2} = J(\Coker(I^{p, 0} \to I^{p, 1}))$ を取り、以下同様に続ける。

\medskip\noindent
ここで Homology, Lemma
\ref{homology-lemma-double-complex-gives-resolution} を適用して、写像
$$
\alpha : K^\bullet \longrightarrow \text{Tot}(I^{\bullet, \bullet})
$$
が擬同型であることを得る。
% P02-E0056: frozen source says “it rests to show”; “it remains to show” is intended.
関手 $inj$ が得られることを証明するには、上の構成が関手的であることを
示すことが残る。この確認は省略する。

\medskip\noindent
$s \circ inj = \text{id}$ を満たす関手 $inj$ があると仮定する。
$\text{Comp}^{+}(\mathcal{A})$ の各対象 $K^\bullet$ に対して
$$
inj(K^\bullet) = (i_{K^\bullet} : K^\bullet \to j(K^\bullet))
$$
と書くことができる。これは定義
\ref{derived-definition-localization-functor} の分解関手を与える。
\end{proof}

\begin{remark}
\label{derived-remark-match}
$inj$ が、補題 \ref{derived-lemma-functorial-injective-resolutions} の (2) にある
$s \circ inj = \text{id}$ を満たす関手であると仮定する。
同補題の証明と同様に
$inj(K^\bullet) = (i_{K^\bullet} : K^\bullet \to j(K^\bullet))$
と書く。$\alpha : K^\bullet \to L^\bullet$ を下に有界な複体の写像とする。
圏 $\text{InjRes}(\mathcal{A})$ における写像 $inj(\alpha)$ を考える。
これは複体の射の可換図
% P02-E0057: frozen diagram writes j(K)^\bullet and j(L)^\bullet, inconsistent with j(K^\bullet), j(L^\bullet).
$$
\xymatrix{
K^\bullet
\ar[rr]^-{\alpha}
\ar[d]_{i_K} & &
L^\bullet \ar[d]^{i_L} \\
j(K)^\bullet
\ar[rr]^-{inj(\alpha)}
& &
j(L)^\bullet
}
$$
を誘導する。したがって補題 \ref{derived-lemma-resolution-functor} の証明を見ると、
関手 $j : K^{+}(\mathcal{A}) \to K^{+}(\mathcal{I})$ は規則
$$
j(\alpha\text{ up to homotopy}) = inj(\alpha)\text{ up to homotopy}\in
\Hom_{K^{+}(\mathcal{I})}(j(K^\bullet), j(L^\bullet))
$$
によって与えられることが分かる。ゆえにこの場合 $j$ は
$t \circ inj$ と一致する。すなわち図
$$
\xymatrix{
\text{Comp}^{+}(\mathcal{A}) \ar[rr]_{t \circ inj} \ar[rd] & &
K^{+}(\mathcal{I}) \\
& K^{+}(\mathcal{A}) \ar[ru]_j
}
$$
は可換である。
\end{remark}

\begin{remark}
\label{derived-remark-big-abelian-category}
$\textit{Mod}(\mathcal{O}_X)$ を、環付き空間
$(X, \mathcal{O}_X)$（またはより一般に環付きサイト）上の
$\mathcal{O}_X$-加群の圏とする。後で
$\textit{Mod}(\mathcal{O}_X)$ が十分多くの入射対象を持ち、
実際に関手的入射埋め込みを持つことを見る。Injectives, Theorem
\ref{injectives-theorem-sheaves-modules-injectives} を参照せよ。
補題 \ref{derived-lemma-into-derived-category} の証明は
$\textit{Mod}(\mathcal{O}_X)$ に適用できないことに注意する。
しかし補題 \ref{derived-lemma-functorial-injective-resolutions} の証明は
$\textit{Mod}(\mathcal{O}_X)$ に適用できる。したがって
$$
j : K^{+}(\textit{Mod}(\mathcal{O}_X))
\longrightarrow
K^{+}(\mathcal{I})
$$
を得る。これは分解関手であり、ここで $\mathcal{I}$ は入射的
$\mathcal{O}_X$-加群の加法圏である。この議論は次の場合にも機能する。
\begin{enumerate}
\item 環 $R$ 上の $R$-加群の圏 $\text{Mod}_R$。
\item 環の前層を備えたサイト上の $\mathcal{O}$-加群の前層の圏
$\textit{PMod}(\mathcal{O})$。
\item 環付きサイト上の $\mathcal{O}$-加群の層の圏
$\textit{Mod}(\mathcal{O})$。
% P02-E0058: frozen source contains the unfinished placeholder “Add more here as needed.”
\item 必要に応じてここに追加する。
\end{enumerate}
\end{remark}









\section{分解関手による右導来関手}
\label{derived-section-right-derived-functor-via-resolutions}

\noindent
次の補題の内容は、分解関手 $j$ が与えられているならば、単に
$RF(K^\bullet) = F(j(K^\bullet))$ と定義できるということである。

\begin{lemma}
\label{derived-lemma-right-derived-functor}
% P02-E0059: frozen source lacks punctuation after “enough injectives”.
$\mathcal{A}$ を十分多くの入射対象を持つアーベル圏とする。
$F : \mathcal{A} \to \mathcal{B}$ をアーベル圏への加法関手とする。
$(i, j)$ を分解関手とする。定義
\ref{derived-definition-localization-functor} を参照せよ。
$F$ の右導来関手 $RF$ は次の $2$-可換図に組み込まれる。
$$
\xymatrix{
D^{+}(\mathcal{A}) \ar[rd]_{RF} \ar[rr]^{j'} & &
K^{+}(\mathcal{I}) \ar[ld]^F \\
& D^{+}(\mathcal{B})
}
$$
ここで $j'$ は補題
\ref{derived-lemma-resolution-functor-quasi-inverse} の関手である。
\end{lemma}

\begin{proof}
補題 \ref{derived-lemma-injective-acyclic} により
$RF(K^\bullet) = F(j(K^\bullet))$ である。
\end{proof}

\begin{remark}
\label{derived-remark-right-derived-functor}
補題 \ref{derived-lemma-right-derived-functor} の状況では、実際に右導来関手を
完全関手 $F \circ j' : D^{+}(\mathcal{A}) \to K^{+}(\mathcal{B})$
へ持ち上げている。このような分解を用いることが有用な場合がある。
\end{remark}









\section{フィルトレーション付き導来圏と入射分解}
\label{derived-section-filtered-derived}

\noindent
$\mathcal{A}$ をアーベル圏とする。本節では $\mathcal{A}$ が
十分多くの入射対象を持つならば、ある意味で圏
$\text{Fil}^f(\mathcal{A})$ もそうであることを示す。
この観察を用いて $\mathcal{A}$ のフィルトレーション付き導来圏で
計算することができる。

\medskip\noindent
圏 $\text{Fil}^f(\mathcal{A})$ は完全圏の一例である。
Injectives, Remark \ref{injectives-remark-embed-exact-category} を参照せよ。
厳密射が特別な役割を果たす。Homology, Definition
\ref{homology-definition-strict} を参照せよ。すなわち
$\Coim(f) = \Im(f)$ を満たす射 $f$ である。
$\text{Fil}^f(\mathcal{A})$ の複体 $A \to B \to C$ について、
列 $\text{gr}(A) \to \text{gr}(B) \to \text{gr}(C)$ が
$\mathcal{A}$ において完全ならば、その複体を{it 完全}という。
これは $A \to B$ と $B \to C$ が厳密射であることを意味する。
Homology, Lemma \ref{homology-lemma-filtered-acyclic} を参照せよ。

\begin{definition}
\label{derived-definition-filtered-complexes-notation}
$\mathcal{A}$ をアーベル圏とする。
$\text{Fil}^f(\mathcal{A})$ の対象 $I$ について、各
$\text{gr}^p(I)$ が $\mathcal{A}$ の入射対象であるとき、
この対象を{it フィルトレーション付き入射対象}という。
\end{definition}

\begin{lemma}
\label{derived-lemma-filtered-injective}
$\mathcal{A}$ をアーベル圏とする。
$\text{Fil}^f(\mathcal{A})$ の対象 $I$ が
フィルトレーション付き入射対象であることと、$a \leq b$、
$\mathcal{A}$ の入射対象 $I_n$（$a \leq n \leq b$）、および同型
$I \cong \bigoplus_{a \leq n \leq b} I_n$ で
$F^pI = \bigoplus_{n \geq p} I_n$ を満たすものが存在することとは同値である。
\end{lemma}

\begin{proof}
$J$ が入射対象ならば $\mathcal{A}$ の任意の単射 $J \to M$ が
分裂するという事実から従う。詳細は省略する。
\end{proof}

\begin{lemma}
\label{derived-lemma-split-strict-monomorphism}
$\mathcal{A}$ をアーベル圏とする。
$\text{Fil}^f(\mathcal{A})$ の任意の厳密単射
$u : I \to A$ は、$I$ がフィルトレーション付き入射対象ならば分裂単射である。
\end{lemma}

\begin{proof}
$p$ を $F^pI \not = 0$ を満たす最大の整数とする。
特に $\text{gr}^p(I) = F^pI$ である。
$I'$ を次のような $\text{Fil}^f(\mathcal{A})$ の対象とする。
その $\mathcal{A}$ における基礎対象は $F^pI$ であり、
フィルトレーションは $n > p$ に対して $F^nI' = 0$、
$n \leq p$ に対して $F^nI' = I' = F^pI$ で与えられる。
$I' \to I$ も厳密単射であることに注意する。
$u$ が厳密単射であることから
$F^pI \to A/F^{p + 1}(A)$ は単射である。Homology, Lemma
\ref{homology-lemma-characterize-strict} を参照せよ。
$\mathcal{A}$ において分裂 $s : A/F^{p + 1}A \to F^pI$ を選ぶ。
誘導される射 $s' : A \to I'$ は、合成 $I' \to I \to A$ を
分裂させるフィルトレーション付き対象の厳密射である。
したがって $A = I' \oplus \Ker(s')$ および
$I = I' \oplus \Ker(s'|_I)$ と書ける。
$\Ker(s'|_I) \to \Ker(s')$ は厳密単射であり、
$\Ker(s'|_I)$ はフィルトレーション付き入射対象であることに注意する。
$I$ のフィルトレーションの長さに関する帰納法により、写像
$\Ker(s'|_I) \to \Ker(s')$ は分裂単射である。これで証明は終わる。
\end{proof}

\begin{lemma}
\label{derived-lemma-injective-property-filtered-injective}
$\mathcal{A}$ をアーベル圏とする。
$u : A \to B$ を $\text{Fil}^f(\mathcal{A})$ の厳密単射とし、
$f : A \to I$ を $A$ から $\text{Fil}^f(\mathcal{A})$ の
フィルトレーション付き入射対象への射とする。このとき
$f = g \circ u$ を満たす射 $g : B \to I$ が存在する。
\end{lemma}

\begin{proof}
$f$ を $u$ に沿って押し出して得られる
$f' : I \to I \amalg_A B$ は厳密単射である。Homology, Lemma
\ref{homology-lemma-pushout-filtered} を参照せよ。
したがって結論は補題 \ref{derived-lemma-split-strict-monomorphism} から
形式的に従う。
\end{proof}

\begin{lemma}
\label{derived-lemma-strict-monomorphism-into-filtered-injective}
$\mathcal{A}$ を十分多くの入射対象を持つアーベル圏とする。
$\text{Fil}^f(\mathcal{A})$ の任意の対象 $A$ に対して、
$I$ がフィルトレーション付き入射対象であるような厳密単射
$A \to I$ が存在する。
\end{lemma}

\begin{proof}
$p \in \{a, a + 1, \ldots, b\}$ でない限り
$\text{gr}^p(A) = 0$ となるように $a \leq b$ を選ぶ。
各 $n \in \{a, a + 1, \ldots, b\}$ に対し、
$I_n$ が入射対象であるような単射
$u_n : A/F^{n + 1}A \to I_n$ を選ぶ。
$I = \bigoplus_{a \leq n \leq b} I_n$ と置き、そのフィルトレーションを
$F^pI = \bigoplus_{n \geq p} I_n$ とし、$u : A \to I$ を
写像 $u_n$ の直和とする。
\end{proof}

\begin{lemma}
\label{derived-lemma-filtered-injective-right-resolution-single-object}
$\mathcal{A}$ を十分多くの入射対象を持つアーベル圏とする。
$\text{Fil}^f(\mathcal{A})$ の任意の対象 $A$ に対して、
フィルトレーション付き擬同型 $A[0] \to I^\bullet$ で、
$I^\bullet$ がフィルトレーション付き入射対象の複体であり、
$n < 0$ に対して $I^n = 0$ であるものが存在する。
\end{lemma}

\begin{proof}
まず、$A$ からフィルトレーション付き入射対象への厳密単射
$u_0 : A \to I^0$ を選ぶ。補題
\ref{derived-lemma-strict-monomorphism-into-filtered-injective} を参照せよ。
次に、$\mathcal{A}$ のフィルトレーション付き入射対象への厳密単射
% P02-E0060: filtered injectives are objects of Fil^f(A), not of A as frozen prose states here and below.
$u_1 : \Coker(u_0) \to I^1$ を選ぶ。$d^0$ で誘導される写像
$I^0 \to I^1$ を表す。次に、$\mathcal{A}$ の
フィルトレーション付き入射対象への厳密単射
$u_2 : \Coker(u_1) \to I^2$ を選ぶ。$d^1$ で誘導される写像
$I^1 \to I^2$ を表す。以下同様である。これが機能するのは、各列
$$
0 \to \Coker(u_n) \to I^{n + 1} \to \Coker(u_{n + 1}) \to 0
$$
が短完全、すなわち $\text{gr}$ を適用すると短完全列を誘導するからである。
これを見るには Homology, Lemma
\ref{homology-lemma-characterize-strict} を用いる。
\end{proof}

\begin{lemma}
\label{derived-lemma-filtered-injective-right-resolution-map}
$\mathcal{A}$ を十分多くの入射対象を持つアーベル圏とする。
$f : A \to B$ を $\text{Fil}^f(\mathcal{A})$ の射とする。
フィルトレーション付き擬同型 $A[0] \to I^\bullet$ および
$B[0] \to J^\bullet$ が与えられ、$I^\bullet, J^\bullet$ が
フィルトレーション付き入射対象の複体で
$n < 0$ に対して $I^n = J^n = 0$ ならば、可換図
$$
\xymatrix{
A[0] \ar[r] \ar[d] &
B[0] \ar[d] \\
I^\bullet \ar[r] &
J^\bullet
}
$$
が存在する。
\end{lemma}

\begin{proof}
% P02-E0061: the second object is B[0] in the lemma, but the frozen proof says C[0].
$A[0] \to I^\bullet$ と $C[0] \to J^\bullet$ は
フィルトレーション付き擬同型なので、$a : A \to I^0$,
$b : B \to J^0$ およびすべての射 $d_I^n$, $d_J^n$ は厳密である。
Homology, Lemma \ref{homology-lemma-filtered-acyclic} を参照せよ。
次の可換図の写像 $f^n$ を帰納的に構成する。
$$
\xymatrix{
A \ar[r]_a \ar[d]_f &
I^0 \ar[r] \ar[d]^{f^0} &
I^1 \ar[r] \ar[d]^{f^1} &
I^2 \ar[r] \ar[d]^{f^2} &
\ldots \\
B \ar[r]^b &
J^0 \ar[r] &
J^1 \ar[r] &
J^2 \ar[r] &
\ldots
}
$$
$A \to I^0$ は厳密単射で $J^0$ はフィルトレーション付き入射対象なので、
補題 \ref{derived-lemma-injective-property-filtered-injective} により
$f^0 \circ a = b \circ f$ を満たす射 $f^0 : I^0 \to J^0$ を
見つけることができる。合成 $d_J^0 \circ b \circ f$ は零なので
$d_J^0 \circ f^0 \circ a = 0$、従って $d_J^0 \circ f^0$ は
一意な射
$$
\Coker(a) = \Coim(d_I^0) = \Im(d_I^0) \longrightarrow J^1.
$$
を経由する。$\Im(d_I^0) \to I^1$ は厳密単射なので、
補題 \ref{derived-lemma-injective-property-filtered-injective} を再び用いて、
表示された矢印を射 $f^1 : I^1 \to J^1$ に延長できる。
以下同様である。
\end{proof}

\begin{lemma}
\label{derived-lemma-filtered-injective-right-resolution-ses}
$\mathcal{A}$ を十分多くの入射対象を持つアーベル圏とする。
$0 \to A \to B \to C \to 0$ を
$\text{Fil}^f(\mathcal{A})$ の短完全列とする。
フィルトレーション付き擬同型 $A[0] \to I^\bullet$ および
$C[0] \to J^\bullet$ が与えられ、$I^\bullet, J^\bullet$ が
フィルトレーション付き入射対象の複体で
$n < 0$ に対して $I^n = J^n = 0$ ならば、可換図
$$
\xymatrix{
0 \ar[r] &
A[0] \ar[r] \ar[d] &
B[0] \ar[r] \ar[d] &
C[0] \ar[r] \ar[d] &
0 \\
0 \ar[r] &
I^\bullet \ar[r] &
M^\bullet \ar[r] &
J^\bullet \ar[r] &
0
}
$$
で、下段が複体の項ごとに分裂する列であるものが存在する。
\end{lemma}

\begin{proof}
$A[0] \to I^\bullet$ と $C[0] \to J^\bullet$ は
フィルトレーション付き擬同型なので、$a : A \to I^0$,
$c : C \to J^0$ およびすべての射 $d_I^n$, $d_J^n$ は厳密である。
Homology, Lemma
% P02-E0062: the frozen cross-reference omits the homology- namespace used by the actual label.
\ref{derived-lemma-filtered-acyclic} を参照せよ。
次の可換図の南東向きおよび南向きの矢印を段階的に構成する。
$$
\xymatrix{
B \ar[r]_\beta \ar[rd]^b &
C \ar[r]_c \ar[rd]^{\overline{b}} &
J^0 \ar[d]^{\delta^0} \ar[r] &
J^1 \ar[d]^{\delta^1} \ar[r] & \ldots \\
A \ar[u]^\alpha \ar[r]^a &
I^0 \ar[r] &
I^1 \ar[r] &
I^2 \ar[r] & \ldots
}
$$
$A \to B$ は厳密単射なので、補題
\ref{derived-lemma-injective-property-filtered-injective} により
$b \circ \alpha = a$ を満たす射 $b : B \to I^0$ を見つけられる。
$A$ は厳密射 $I^0 \to I^1$ の核であり、
$\beta = \Coker(\alpha)$ なので、図に組み込まれる一意な射
$\overline{b} : C \to I^1$ を得る。
$c$ は厳密単射で $I^1$ はフィルトレーション付き入射対象なので、
補題 \ref{derived-lemma-injective-property-filtered-injective} により
$\delta^0 : J^0 \to I^1$ を見つけられる。
$B \to C$ は厳密全射で、$B \to I^0 \to I^1 \to I^2$ は零なので、
$C \to I^1 \to I^2$ は零である。従って
$d_I^1 \circ \delta^0$ は $C \cong \Im(c)$ 上で零である。
ゆえに $d_I^1 \circ \delta^0$ は一意な射
$$
\Coker(c) = \Coim(d_J^0) = \Im(d_J^0) \longrightarrow I^2.
$$
を経由する。$I^2$ はフィルトレーション付き入射対象で、
$\Im(d_J^0) \to J^1$ は厳密単射なので、補題
\ref{derived-lemma-injective-property-filtered-injective} により、表示された射を
$\delta^1 : J^1 \to I^2$ に延長できる。以下同様である。
$M^\bullet = I^\bullet \oplus J^\bullet$ と置き、微分を
$$
d_M^n =
\left(
\begin{matrix}
d_I^n & (-1)^{n + 1}\delta^n \\
0 & d_J^n
\end{matrix}
\right)
$$
とする。最後に、写像 $B[0] \to M^\bullet$ は
% P02-E0063: the frozen formula gives domain M although the map just defined has domain B.
$b \oplus c \circ \beta : M \to I^0 \oplus J^0$ によって与えられる。
\end{proof}

\begin{lemma}
\label{derived-lemma-right-resolution-by-filtered-injectives}
$\mathcal{A}$ を十分多くの入射対象を持つアーベル圏とする。
すべての $K^\bullet \in K^{+}(\text{Fil}^f(\mathcal{A}))$ に対して、
フィルトレーション付き擬同型 $K^\bullet \to I^\bullet$ で、
$I^\bullet$ が下に有界、各 $I^n$ がフィルトレーション付き入射対象、
各 $K^n \to I^n$ が厳密単射であるものが存在する。
\end{lemma}

\begin{proof}
$K^\bullet$ をシフトで置き換えた後（証明には影響しない）、
$n < 0$ に対して $K^n = 0$ と仮定してよい。完全圏
$\text{Fil}^f(\mathcal{A})$ の短完全列
$$
\begin{matrix}
0 \to \Ker(d_K^0) \to K^0 \to \Coim(d_K^0) \to 0 \\
0 \to \Ker(d_K^1) \to K^1 \to \Coim(d_K^1) \to 0 \\
0 \to \Ker(d_K^2) \to K^2 \to \Coim(d_K^2) \to 0 \\
\ldots
\end{matrix}
$$
および写像 $u_i : \Coim(d_K^i) \to \Ker(d_K^{i + 1})$ を考える。
各 $i \geq 0$ に対して、フィルトレーション付き擬同型
$$
\begin{matrix}
\Ker(d_K^i)[0] \to I_{ker, i}^\bullet \\
\Coim(d_K^i)[0] \to I_{coim, i}^\bullet
\end{matrix}
$$
で、$I_{ker, i}^n, I_{coim, i}^n$ がフィルトレーション付き入射対象であり、
$n < 0$ に対して零であるものを選べる。補題
\ref{derived-lemma-filtered-injective-right-resolution-single-object} を参照せよ。
補題 \ref{derived-lemma-filtered-injective-right-resolution-map} により、
$u_i$ を複体の射
$u_i^\bullet : I_{coim, i}^\bullet \to I_{ker, i + 1}^\bullet$
へ持ち上げることができる。最後に、各 $i \geq 0$ に対して、図
$$
\xymatrix{
0 \ar[r] &
\Ker(d_K^i)[0] \ar[r] \ar[d] &
K^i[0] \ar[r] \ar[d] &
\Coim(d_K^i)[0] \ar[r] \ar[d] &
0 \\
0 \ar[r] &
I_{ker, i}^\bullet \ar[r]^{\alpha_i} &
I_i^\bullet \ar[r]^{\beta_i} &
I_{coim, i}^\bullet \ar[r] &
0
}
$$
を、下段が項ごとに分裂する完全列となるように完成できる。
補題 \ref{derived-lemma-filtered-injective-right-resolution-ses} を参照せよ。
$i \geq 0$ に対して $d_i : I_i^\bullet \to I_{i + 1}^\bullet$ を
$d_i =  \alpha_{i + 1} \circ u_i^\bullet \circ \beta_i$ と置く。
$\beta_i \circ \alpha_i = 0$ なので
$d_i \circ d_{i - 1} = 0$ であることに注意する。
従って可換図
$$
\xymatrix{
I_0^\bullet \ar[r] &
I_1^\bullet \ar[r] &
I_2^\bullet \ar[r] & \ldots \\
K^0[0] \ar[r] \ar[u] &
K^1[0] \ar[r] \ar[u] &
K^2[0] \ar[r] \ar[u] &
\ldots
}
$$
を構成した。ここで縦矢印はフィルトレーション付き擬同型である。
上段は複体の複体であり、各複体はフィルトレーション付き入射対象からなり、
次数 $< 0$ に非零対象を持たない。従って
$I^{a, b} = I_a^b$ と置き、
$$
d_1^{a, b} : I^{a, b} = I_a^b \to I_{a + 1}^b = I^{a + 1, b}
$$
には写像 $d_a^b$ を用い、
$$
d_2^{a, b} : I^{a, b} = I_a^b \to I_a^{b + 1} = I^{a, b + 1}
$$
には写像 $d_{I_a}^b$ を用いることで二重複体を得る。
この二重複体に付随する全複体を
$\text{Tot}(I^{\bullet, \bullet})$ と表す。Homology, Definition
\ref{homology-definition-associated-simple-complex} を参照せよ。
写像 $K^n[0] \to I_n^\bullet$ は写像
$K^n \to I^{n, 0}$ から来ており、複体の写像
$$
K^\bullet \longrightarrow \text{Tot}(I^{\bullet, \bullet})
$$
を生じさせることに注意する。これはフィルトレーション付き擬同型であると主張する。
$\text{gr}(-)$ は加法関手なので
$\text{gr}(\text{Tot}(I^{\bullet, \bullet})) =
\text{Tot}(\text{gr}(I^{\bullet, \bullet}))$ である。
従って Homology, Lemma
\ref{homology-lemma-double-complex-gives-resolution} を用いて、
望みどおり
$\text{gr}(K^\bullet) \to \text{gr}(\text{Tot}(I^{\bullet, \bullet}))$
が擬同型であると結論できる。
\end{proof}

\begin{lemma}
\label{derived-lemma-filtered-acyclic-is-zero}
$\mathcal{A}$ をアーベル圏とする。
$K^\bullet, I^\bullet \in K(\text{Fil}^f(\mathcal{A}))$ とする。
$K^\bullet$ はフィルトレーション付き非輪状で、
$I^\bullet$ は下に有界かつフィルトレーション付き入射対象からなると仮定する。
任意の射 $K^\bullet \to I^\bullet$ は零とホモトピックである。
$\Hom_{K(\text{Fil}^f(\mathcal{A}))}(K^\bullet, I^\bullet) = 0$。
\end{lemma}

\begin{proof}
$\alpha : K^\bullet \to I^\bullet$ を複体の射とする。
$j < n$ に対して $\alpha^j = 0$ と仮定する。
$\alpha^n = h \circ d$ を満たす射
$h : K^{n + 1} \to I^n$ が存在することを示す。
すると $\alpha$ は、次で定義される複体の射 $\beta$ とホモトピックになる。
$$
\beta^j =
\left\{
\begin{matrix}
0 & \text{if} & j \leq n \\
\alpha^{n + 1} - d \circ h & \text{if} & j = n + 1 \\
\alpha^j & \text{if} & j > n + 1
\end{matrix}
\right.
$$
これで帰納法により補題が証明されることは明らかである。
$h$ の存在を示すため、$\alpha^{n - 1} = 0$ なので
$\alpha^n \circ d_K^{n - 1} = 0$ であることに注意する。
$K^\bullet$ はフィルトレーション付き非輪状なので、
$d_K^{n - 1}$ と $d_K^n$ は厳密であり、
$$
0 \to \Im(d_K^{n - 1}) \to K^n \to \Im(d_K^n) \to 0
$$
は完全圏 $\text{Fil}^f(\mathcal{A})$ の完全列である。
Homology, Lemma \ref{homology-lemma-filtered-acyclic} を参照せよ。
従って $\alpha^n$ を $\Im(d_K^n)$ 上で定義された $I^n$ への写像と
考えることができる。$\Im(d_K^n) \to K^{n + 1}$ が厳密単射で、
$I^n$ がフィルトレーション付き入射対象であることを用いれば、
この写像を望みの写像 $h : K^{n + 1} \to I^n$ に持ち上げられる。
補題 \ref{derived-lemma-injective-property-filtered-injective} を参照せよ。
\end{proof}

\begin{lemma}
\label{derived-lemma-morphisms-into-filtered-injective-complex}
$\mathcal{A}$ をアーベル圏とする。
$I^\bullet \in K(\text{Fil}^f(\mathcal{A}))$ を、
フィルトレーション付き入射対象からなる下に有界な複体とする。
\begin{enumerate}
\item $\alpha : K^\bullet \to L^\bullet$ を
$K(\text{Fil}^f(\mathcal{A}))$ におけるフィルトレーション付き擬同型とする。
このとき写像
$$
\Hom_{K(\text{Fil}^f(\mathcal{A}))}(L^\bullet, I^\bullet)
\to
\Hom_{K(\text{Fil}^f(\mathcal{A}))}(K^\bullet, I^\bullet)
$$
は全単射である。
\item $L^\bullet \in K(\text{Fil}^f(\mathcal{A}))$ とする。このとき
$$
\Hom_{K(\text{Fil}^f(\mathcal{A}))}(L^\bullet, I^\bullet)
=
\Hom_{DF(\mathcal{A})}(L^\bullet, I^\bullet).
$$
\end{enumerate}
\end{lemma}

\begin{proof}
(1) の証明。補題 \ref{derived-lemma-the-same-up-to-isomorphisms} により
$$
(K^\bullet, L^\bullet, C(\alpha)^\bullet, \alpha, i, -p)
$$
は $K(\text{Fil}^f(\mathcal{A}))$ における完全三角であり、
補題 \ref{derived-lemma-filtered-acyclic} により
$C(\alpha)^\bullet$ はフィルトレーション付き非輪状複体である。
従って
$$
\xymatrix{
\Hom_{K(\text{Fil}^f(\mathcal{A}))}(C(\alpha)^\bullet, I^\bullet) \ar[r] &
\Hom_{K(\text{Fil}^f(\mathcal{A}))}(L^\bullet, I^\bullet) \ar[r] &
\Hom_{K(\text{Fil}^f(\mathcal{A}))}(K^\bullet, I^\bullet) \ar[lld] \\
\Hom_{K(\text{Fil}^f(\mathcal{A}))}(C(\alpha)^\bullet[-1], I^\bullet)
}
$$
はアーベル群の完全列である。補題
\ref{derived-lemma-representable-homological} を参照せよ。
ここで補題 \ref{derived-lemma-filtered-acyclic-is-zero} により外側の二つの群は
零であり、従って
% P02-E0064: the frozen conclusion drops Fil^f(A) from both homotopy categories.
$\Hom_{K(\mathcal{A})}(L^\bullet, I^\bullet) =
\Hom_{K(\mathcal{A})}(K^\bullet, I^\bullet)$
である。

\medskip\noindent
(2) の証明。$a$ を右辺の元とする。
$a = \gamma\alpha^{-1}$ と表すことができる。ここで
$\alpha : K^\bullet \to L^\bullet$ はフィルトレーション付き擬同型で、
$\gamma : K^\bullet \to I^\bullet$ は複体の写像である。
(1) により、$\beta \circ \alpha$ が $\gamma$ とホモトピックとなる射
$\beta : L^\bullet \to I^\bullet$ を見つけることができる。
これで写像が全射であることが示された。
$b$ を左辺の元で、右辺における像が零であるものとする。
すると $b$ は射 $\beta : L^\bullet \to I^\bullet$ のホモトピー類であり、
$\beta  \circ \alpha$ が零とホモトピックとなる
フィルトレーション付き擬同型
$\alpha : K^\bullet \to L^\bullet$ が存在する。
このとき (1) により $\beta$ も零とホモトピック、すなわち $b = 0$ である。
\end{proof}

\begin{lemma}
\label{derived-lemma-filtered-localization-functor}
$\mathcal{A}$ を十分多くの入射対象を持つアーベル圏とする。
$\mathcal{I}^f \subset \text{Fil}^f(\mathcal{A})$ で、
対象がフィルトレーション付き入射対象である厳密充満加法部分圏を表す。
標準関手
$$
K^{+}(\mathcal{I}^f)
\longrightarrow
DF^{+}(\mathcal{A})
$$
は完全かつ充満忠実で本質的全射である。すなわち三角圏の同値である。
さらに図
$$
\xymatrix{
K^{+}(\mathcal{I}^f) \ar[d]_{\text{gr}^p} \ar[r] &
DF^{+}(\mathcal{A}) \ar[d]_{\text{gr}^p} \\
K^{+}(\mathcal{I}) \ar[r] &
D^{+}(\mathcal{A})
}
\quad
\xymatrix{
K^{+}(\mathcal{I}^f) \ar[d]^{\text{forget }F} \ar[r] &
DF^{+}(\mathcal{A}) \ar[d]^{\text{forget }F} \\
K^{+}(\mathcal{I}) \ar[r] &
D^{+}(\mathcal{A})
}
$$
は可換である。ここで $\mathcal{I} \subset \mathcal{A}$ は、
対象が入射対象である厳密充満加法部分圏である。
\end{lemma}

\begin{proof}
補題 \ref{derived-lemma-right-resolution-by-filtered-injectives} により、関手
$K^{+}(\mathcal{I}^f) \to DF^{+}(\mathcal{A})$ は本質的全射である。
補題 \ref{derived-lemma-morphisms-into-filtered-injective-complex} により
充満忠実である。完全三角に関する定義により、これは完全関手である。
正方形の可換性は明らかである。
\end{proof}

\begin{remark}
\label{derived-remark-filtered-localization-big}
補題の矢印を逆にできるのは、$\mathcal{A}$ がわれわれの意味での圏、
すなわち対象の集合を持つ場合に限られる。しかし例えば
$\textit{Mod}(\mathcal{O}_X)$ のように、十分多くの入射対象を持つ
大きなアーベル圏 $\mathcal{A}$ が与えられたと仮定しよう。
このとき、任意に与えられた対象の集合 $\{A_i\}_{i\in I}$ に対し、
それらをすべて含み、十分多くの入射対象を持つアーベル部分圏
$\mathcal{A}' \subset \mathcal{A}$ が存在する。Sets, Lemma
\ref{sets-lemma-abelian-injectives} を参照せよ。
従って上の補題を $\mathcal{A}'$ に用いることができる。
これは本質的に、集合個分の図などを用いる限り、補題を使って問題に
遭遇することは決してないということである。
\end{remark}

\noindent
$\mathcal{A}, \mathcal{B}$ をアーベル圏とする。
$T : \mathcal{A} \to \mathcal{B}$ を左完全関手とする。
（関手に文字 $F$ を使うことはできない。フィルトレーションを表す文字
$F$ の用法とあまりにも衝突するからである。）
$T$ は規則 $T(A, F) = (T(A), F)$ により加法関手
$$
T : \text{Fil}^f(\mathcal{A}) \to \text{Fil}^f(\mathcal{B})
$$
を誘導することに注意する。ここで $F^pT(A) = T(F^pA)$ であり、
$T$ は左完全なのでこれは意味を持つ。
（警告：$\text{gr}(T(A)) = T(\text{gr}(A))$ とは限らない。）
これは三角圏の関手
\begin{equation}
\label{derived-equation-induced-T-filtered}
T :
K^{+}(\text{Fil}^f(\mathcal{A}))
\longrightarrow
K^{+}(\text{Fil}^f(\mathcal{B}))
\end{equation}
を誘導する。$T$ のフィルトレーション付き右導来関手とは、この完全関手と
完全関手 $K^{+}(\text{Fil}^f(\mathcal{B})) \to DF^{+}(\mathcal{B})$
との合成および乗法系 $\text{FQis}^{+}(\mathcal{A})$ に対する、
定義 \ref{derived-definition-right-derived-functor-defined} の右導来関手である。
$\mathcal{A}$ は十分多くの入射対象を持つと仮定する。
ここで Section \ref{derived-section-right-derived-functor} の議論を繰り返して、
関手 $T$ の{\it フィルトレーション付き右導来関手}
\begin{equation}
\label{derived-equation-filtered-derived-functor}
RT : DF^{+}(\mathcal{A}) \longrightarrow DF^{+}(\mathcal{B})
\end{equation}
を定義できる。

\medskip\noindent
しかし代わりに Section
\ref{derived-section-right-derived-functor-via-resolutions} と同様に進める。
すると $T$ が単に加法的である場合にも $RT$ を定義できることが分かる。
まず、補題 \ref{derived-lemma-filtered-localization-functor} の同値の擬逆
$j' : DF^{+}(\mathcal{A}) \to K^{+}(\mathcal{I}^f)$ を選ぶ。
補題 \ref{derived-lemma-exact-equivalence} により、$j'$ は三角圏の完全関手である。
次に、フィルトレーション付き入射対象 $I$ に対して、
補題 \ref{derived-lemma-filtered-injective} により（非標準的な）分解
\begin{equation}
\label{derived-equation-decompose}
I \cong \bigoplus\nolimits_{p \in \mathbf{Z}} I_p,
\quad\text{with}\quad
F^pI = \bigoplus\nolimits_{q \geq p} I_q
\end{equation}
があることに注意する。従って $T$ が任意の加法関手
$T : \mathcal{A} \to \mathcal{B}$ ならば、
$T_{ext}(I) = \bigoplus T(I_p)$ および
$F^pT_{ext}(I) = \bigoplus_{q \geq p} T(I_q)$ と置くことで加法関手
\begin{equation}
\label{derived-equation-extend-T}
T_{ext} : \mathcal{I}^f \to \text{Fil}^f(\mathcal{B})
\end{equation}
を得る。構成により
$\text{gr}(T_{ext}(I)) = T(\text{gr}(I))$ という性質を持つことに注意する。
従って $\text{gr}$ と可換する関手
\begin{equation}
\label{derived-equation-extend-T-complexes}
T_{ext} : K^{+}(\mathcal{I}^f) \to K^{+}(\text{Fil}^f(\mathcal{B}))
\end{equation}
を得る。そして (\ref{derived-equation-filtered-derived-functor}) を合成
\begin{equation}
\label{derived-equation-definition-filtered-derived-functor}
RT = T_{ext} \circ j'.
\end{equation}
により定義する。$RT : D^{+}(\mathcal{A}) \to D^{+}(\mathcal{B})$ も
入射分解によって計算されるので（補題 \ref{derived-lemma-injective-acyclic}
を参照）、$T$ と $\text{gr}$ の可換性および補題
\ref{derived-lemma-filtered-localization-functor} の可換図から
\begin{equation}
\label{derived-equation-commute-gr}
\text{gr}^p \circ RT \cong RT \circ \text{gr}^p
\end{equation}
および
\begin{equation}
\label{derived-equation-commute-forget}
(\text{forget }F) \circ RT \cong RT \circ (\text{forget }F)
\end{equation}
を $DF^{+}(\mathcal{A}) \to D^{+}(\mathcal{B})$ の関手として得る。

\medskip\noindent
フィルトレーション付き導来関手 $RT$
（\ref{derived-equation-filtered-derived-functor}）は関手
$$
\begin{matrix}
RT : \text{Fil}^f(\mathcal{A}) \to DF^{+}(\mathcal{B}), \\
RT : \text{Comp}^{+}(\text{Fil}^f(\mathcal{A})) \to DF^{+}(\mathcal{B}), \\
RT : KF^{+}(\mathcal{A}) \to DF^{+}(\mathcal{B}).
\end{matrix}
$$
を誘導する。$\text{Fil}^f(\mathcal{A})$ および
$\text{Comp}^{+}(\text{Fil}^f(\mathcal{A}))$ はもはやアーベル圏ではないので、
$RT$ がそれらの上で $\delta$-関手に制限されると言うことは意味を持たない。
（これらの圏を完全圏と考え、完全圏から三角圏への $\delta$-関手の概念を
定式化すれば修復できる。）
しかし $RT$ が三角圏 $KF^{+}(\mathcal{A})$ 上の完全関手であると言うことには
意味があり、構成により実際にそうである。

\begin{lemma}
\label{derived-lemma-ss-filtered-derived}
$\mathcal{A}, \mathcal{B}$ をアーベル圏とする。
$T : \mathcal{A} \to \mathcal{B}$ を左完全関手とする。
$\mathcal{A}$ は十分多くの入射対象を持つと仮定する。
$(K^\bullet, F)$ を
$\text{Comp}^{+}(\text{Fil}^f(\mathcal{A}))$ の対象とする。
$\mathcal{B}$ の二重次数付き対象 $E_r$ と双次数 $(r, - r + 1)$ の
$d_r$ からなるスペクトル系列 $(E_r, d_r)_{r\geq 0}$ で、
$$
E_1^{p, q} = R^{p + q}T(\text{gr}^p(K^\bullet))
$$
を満たすものが存在する。さらに、このスペクトル系列は有界で、
$R^*T(K^\bullet)$ に収束し、各 $R^nT(K^\bullet)$ 上に
有限フィルトレーションを誘導する。このスペクトル系列の構成は
$\text{Comp}^{+}(\text{Fil}^f(\mathcal{A}))$ の対象
$K^\bullet$ に関して関手的であり、$r \geq 1$ に対する項
$(E_r, d_r)$ はいかなる選択にも依存しない。
\end{lemma}

\begin{proof}
フィルトレーション付き擬同型 $K^\bullet \to I^\bullet$ で、
$I^\bullet$ がフィルトレーション付き入射対象からなる下に有界な複体であるものを
選ぶ。補題 \ref{derived-lemma-right-resolution-by-filtered-injectives} を参照せよ。
複体 $RT(K^\bullet) = T_{ext}(I^\bullet)$ を考える。
(\ref{derived-equation-definition-filtered-derived-functor}) を参照せよ。
これを $\mathcal{B}$ のフィルトレーション付き複体と見て、付随する
スペクトル系列 $(E_r, d_r)_{r \geq 0}$ を考えることができる。
Homology, Section \ref{homology-section-filtered-complex} を参照せよ。
Homology, Lemma
\ref{homology-lemma-spectral-sequence-filtered-complex} により
$E_1^{p, q} = H^{p + q}(\text{gr}^p(T(I^\bullet)))$ である。
Equation (\ref{derived-equation-decompose}) により
$E_1^{p, q} = H^{p + q}(T(\text{gr}^p(I^\bullet)))$ であり、
フィルトレーション付き入射分解の定義により、写像
$\text{gr}^p(K^\bullet) \to \text{gr}^p(I^\bullet)$ は入射分解である。
従って $E_1^{p, q} = R^{p + q}T(\text{gr}^p(K^\bullet))$ である。

\medskip\noindent
一方、各 $I^n$ は有限フィルトレーションを持つので、
各 $T(I^n)$ も有限フィルトレーションを持つ。
従って Homology, Lemma
\ref{homology-lemma-biregular-ss-converges} を適用して、
スペクトル系列が有界で $H^n(T(I^\bullet)) = R^nT(K^\bullet)$ に収束し、
さらに各項に有限フィルトレーションを誘導すると結論できる。

\medskip\noindent
$K^\bullet \to L^\bullet$ を
$\text{Comp}^{+}(\text{Fil}^f(\mathcal{A}))$ の射と仮定する。
フィルトレーション付き擬同型 $L^\bullet \to J^\bullet$ で、
$J^\bullet$ がフィルトレーション付き入射対象からなる下に有界な複体であるものを
選ぶ。補題 \ref{derived-lemma-right-resolution-by-filtered-injectives} を参照せよ。
上の結果、例えば補題
\ref{derived-lemma-morphisms-into-filtered-injective-complex} により、図
$$
\xymatrix{
K^\bullet \ar[r] \ar[d] & L^\bullet \ar[d] \\
I^\bullet \ar[r] & J^\bullet
}
$$
でホモトピーを除いて可換なものが存在する。従って
フィルトレーション付き複体の射
$T(I^\bullet) \to T(J^\bullet)$ を得て、これはスペクトル系列の射を生じさせる。
Homology, Lemma
\ref{homology-lemma-spectral-sequence-filtered-complex-functorial} を参照せよ。
最後の主張はこれから従う。
\end{proof}

\begin{remark}
\label{derived-remark-final-functorial}
Remark \ref{derived-remark-functorial-ss} で約束したとおり、上の補題と
Cartan--Eilenberg 分解を用いる構成との関係を論じる。
すなわち $T : \mathcal{A} \to \mathcal{B}$ をアーベル圏の間の
左完全関手とし、$\mathcal{A}$ は十分多くの入射対象を持つと仮定し、
$K^\bullet$ を $\mathcal{A}$ の下に有界な複体とする。
補題 \ref{derived-lemma-two-ss-complex-functor} のスペクトル系列
${}'E$ と ${}''E$ の別の構成を与える。

\medskip\noindent
第一スペクトル系列。$F^p(K^\bullet) = \sigma_{\geq p}(K^\bullet)$
と置くことで得られる $K^\bullet$ 上の「素朴な」フィルトレーションを考える。
Homology, Section \ref{homology-section-truncations} を参照せよ。
% P02-E0065: frozen prose reads “Note that this stupid” and omits “is”.
$d(F^p(K^\bullet)) \subset F^{p + 1}(K^\bullet)$ という意味で
これが「素朴」であることに注意する。Homology, Lemma
\ref{homology-lemma-spectral-sequence-filtered-complex-d1} と比較せよ。
このフィルトレーションについて
$\text{gr}^p(K^\bullet) = K^p[-p]$ であることに注意する。
補題 \ref{derived-lemma-ss-filtered-derived} により、$E_1$ 項が
$$
E_1^{p, q} = R^{p + q}T(K^p[-p]) = R^qT(K^p)
$$
となるスペクトル系列が存在し、これはスペクトル系列 ${}'E_r$ と同様である。
さらに、微分 $E_1^{p, q} \to E_1^{p + 1, q}$ は $'{}E_1$ における
微分と一致することに注意する。Homology, Lemma
\ref{homology-lemma-spectral-sequence-filtered-complex-d1} の (2) および
補題 \ref{derived-lemma-two-ss-complex-functor} の証明における
${}'d_1$ の記述を参照せよ。

\medskip\noindent
第二スペクトル系列。$F^p(K^\bullet) = \tau_{\leq -p}(K^\bullet)$
と置くことで得られる複体 $K^\bullet$ 上のフィルトレーションを考える。
Homology, Section \ref{homology-section-truncations} を参照せよ。
減少フィルトレーションを得るために負号が必要である。
このフィルトレーションについて $\text{gr}^p(K^\bullet)$ は
$H^{-p}(K^\bullet)[p]$ と擬同型であることに注意する。
補題 \ref{derived-lemma-ss-filtered-derived} により、$E_1$ 項が
$$
E_1^{p, q} = R^{p + q}T(H^{-p}(K^\bullet)[p])
= R^{2p + q}T(H^{-p}(K^\bullet)) = {}''E_2^{i, j}
$$
となるスペクトル系列が存在する。ここで $i = 2p + q$ および $j = -p$ である。
（これは不自然に見えるが、フィルトレーション付き複体の理論全体を
増加フィルトレーションを用いて展開することもでき、その場合はこちらが自然に見え、
他方が不自然に見えることに注意する。）
微分が一致することの確認は読者に委ねる。

\medskip\noindent
実際、Cartan--Eilenberg 分解
$K^\bullet \to I^{\bullet, \bullet}$ が与えられると、付随する全複体への
誘導射 $K^\bullet \to \text{Tot}(I^{\bullet, \bullet})$ は、
$\text{Tot}(I^{\bullet, \bullet})$ 上に適切なフィルトレーションを用いることで、
いずれのフィルトレーションについてもフィルトレーション付き入射分解になる。
これを用いてスペクトル系列を正確に一致させることができる。
\end{remark}







\section{Ext 群}
\label{derived-section-ext}

\noindent
本節ではアーベル圏の対象の Ext 群の記述を始める。
まず次の非常に一般的な定義がある。

\begin{definition}
\label{derived-definition-ext}
$\mathcal{A}$ をアーベル圏とする。$i \in \mathbf{Z}$ とする。
$X, Y$ を $D(\mathcal{A})$ の対象とする。
$Y$ による $X$ の{\it 第 $i$ Ext 群}とは、群
$$
\Ext^i_\mathcal{A}(X, Y) =
\Hom_{D(\mathcal{A})}(X, Y[i]) =
\Hom_{D(\mathcal{A})}(X[-i], Y).
$$
のことである。$A, B \in \Ob(\mathcal{A})$ ならば
$\Ext^i_\mathcal{A}(A, B) = \text{Ext}^i_\mathcal{A}(A[0], B[0])$ と置く。
\end{definition}

\noindent
$\Hom_{D(\mathcal{A})}(X, -)$ はホモロジー的関手、
$\Hom_{D(\mathcal{A})}(-, Y)$ はコホモロジー的関手である。
補題 \ref{derived-lemma-representable-homological} を参照せよ。
従って完全三角 $(Y, Y', Y'')$ は長完全列
$$
\ldots \to
\Ext^i_\mathcal{A}(X, Y) \to
\Ext^i_\mathcal{A}(X, Y') \to
\Ext^i_\mathcal{A}(X, Y'') \to
\Ext^{i + 1}_\mathcal{A}(X, Y) \to \ldots
$$
を、完全三角 $(X, X', X'')$ はそれぞれ長完全列
$$
\ldots \to
\Ext^i_\mathcal{A}(X'', Y) \to
\Ext^i_\mathcal{A}(X', Y) \to
\Ext^i_\mathcal{A}(X, Y) \to
\Ext^{i + 1}_\mathcal{A}(X'', Y) \to \ldots
$$
を与える。$D^+(\mathcal{A})$, $D^-(\mathcal{A})$, $D^b(\mathcal{A})$
は充満部分圏なので、$X$, $Y$ がそれらに含まれるならば、
これらの圏における Hom 群によって Ext 群を計算してよいことに注意する。

\medskip\noindent
圏 $\mathcal{A}$ が十分多くの入射対象または十分多くの射影対象を持つ場合、
入射分解または射影分解を用いて Ext 群を計算できる。
混同を避けるため、入射分解（それぞれ射影分解）を持つことは、
すべての低い（それぞれ高い）次数におけるホモロジーの消滅を意味することを
思い出そう。補題 \ref{derived-lemma-cohomology-bounded-below} および
\ref{derived-lemma-cohomology-bounded-above} を参照せよ。

\begin{lemma}
\label{derived-lemma-compute-ext-resolutions}
$\mathcal{A}$ をアーベル圏とする。
$X^\bullet, Y^\bullet \in \Ob(K(\mathcal{A}))$ とする。
\begin{enumerate}
\item $Y^\bullet \to I^\bullet$ を入射分解とする
（定義 \ref{derived-definition-injective-resolution}）。このとき
$$
\Ext^i_\mathcal{A}(X^\bullet, Y^\bullet) =
\Hom_{K(\mathcal{A})}(X^\bullet, I^\bullet[i]).
$$
\item $P^\bullet \to X^\bullet$ を射影分解とする
（定義 \ref{derived-definition-projective-resolution}）。このとき
$$
\Ext^i_\mathcal{A}(X^\bullet, Y^\bullet) =
\Hom_{K(\mathcal{A})}(P^\bullet[-i], Y^\bullet).
$$
\end{enumerate}
\end{lemma}

\begin{proof}
補題 \ref{derived-lemma-morphisms-into-injective-complex} および
補題 \ref{derived-lemma-morphisms-from-projective-complex} から直ちに従う。
\end{proof}

\noindent
本節の残りではアーベル圏そのものの対象の拡大を論じる。
まず次を観察する。

\begin{lemma}
\label{derived-lemma-negative-exts}
$\mathcal{A}$ をアーベル圏とする。
\begin{enumerate}
\item $X$, $Y$ を $D(\mathcal{A})$ の対象とする。
$i > a$ に対して $H^i(X) = 0$、$j < b$ に対して $H^j(Y) = 0$
となる $a, b \in \mathbf{Z}$ が与えられたとする。
このとき $n < b - a$ に対して $\Ext^n_\mathcal{A}(X, Y) = 0$ であり、
$$
\Ext^{b - a}_\mathcal{A}(X, Y) = \Hom_\mathcal{A}(H^a(X), H^b(Y))
$$
である。
\item $A, B \in \Ob(\mathcal{A})$ とする。
$i < 0$ に対して $\Ext^i_\mathcal{A}(B, A) = 0$ である。
$\Ext^0_\mathcal{A}(B, A) = \Hom_\mathcal{A}(B, A)$ である。
\end{enumerate}
\end{lemma}

\begin{proof}
$X$ と $Y$ を表す複体 $X^\bullet$ と $Y^\bullet$ を選ぶ。
$Y^\bullet \to \tau_{\geq b}Y^\bullet$ は擬同型なので、
$j < b$ に対して $Y^j = 0$ と仮定してよい。
$L^\bullet \to X^\bullet$ を任意の擬同型とする。このとき
$\tau_{\leq a}L^\bullet \to X^\bullet$ は擬同型である。
従って $D(\mathcal{A})$ の射 $X \to Y[n]$ は
$fs^{-1}$ と表せる。ここで $s : L^\bullet \to X^\bullet$ は擬同型、
$f : L^\bullet \to Y^\bullet[n]$ は射であり、
% P02-E0067: after tau_{<=a}, the frozen proof should say L^i=0 for i>a, not i<a.
$i < a$ に対して $L^i = 0$ である。
$f$ は $L^i$ を $Y^{i + n}$ に写すことに注意する。
従って $n < b - a$ ならば、常に $L^i$ または $Y^{i + n}$ の一方が
零なので $f = 0$ である。$n = b - a$ ならば、$f$ はちょうど射
$H^a(X) \to H^b(Y)$ に対応する。(2) は (1) の特別な場合である。
\end{proof}

\noindent
$\mathcal{A}$ をアーベル圏とする。
$0 \to A \to A' \to A'' \to 0$ を
$\mathcal{A}$ の対象の短完全列と仮定する。このとき
$0 \to A[0] \to A'[0] \to A''[0] \to 0$ は
$D(\mathcal{A})$ における完全三角を生じさせ
（補題 \ref{derived-lemma-derived-canonical-delta-functor} を参照）、
従って Ext 群の長完全列
$$
0 \to \Ext^0_\mathcal{A}(B, A) \to
\Ext^0_\mathcal{A}(B, A') \to
\Ext^0_\mathcal{A}(B, A'') \to
\Ext^1_\mathcal{A}(B, A) \to \ldots
$$
を生じさせる。同様に、短完全列 $0 \to B \to B' \to B'' \to 0$
が与えられると、Ext 群の長完全列
$$
0 \to \Ext^0_\mathcal{A}(B'', A) \to
\Ext^0_\mathcal{A}(B', A) \to
\Ext^0_\mathcal{A}(B, A) \to
\Ext^1_\mathcal{A}(B'', A) \to \ldots
$$
を得る。これらの Ext 群を導来圏の構成の応用と見ることができる。
これは、十分多くの入射対象または射影対象の存在を必要とせずに Ext 群を定義し、
Ext 群の長完全列を構成できることを示す。
Yoneda による Ext 群の別の構成は導来圏の使用を避ける。
\cite{Yoneda} を参照せよ。

\begin{definition}
\label{derived-definition-yoneda-extension}
$\mathcal{A}$ をアーベル圏とする。$A, B \in \Ob(\mathcal{A})$ とする。
$i \geq 1$ に対して、$A$ による $B$ の次数 $i$ の
{\it Yoneda 拡大}とは、$\mathcal{A}$ における完全列
$$
E : 0 \to A \to Z_{i - 1} \to Z_{i - 2} \to \ldots \to Z_0 \to B \to 0
$$
のことである。同じ次数の二つの Yoneda 拡大 $E$ と $E'$ について、
可換図
$$
\xymatrix{
0 \ar[r] & A \ar[r] & Z_{i - 1} \ar[r] & \ldots \ar[r] &
Z_0 \ar[r] & B \ar[r] & 0 \\
0 \ar[r] &
A \ar[r] \ar[u]^{\text{id}} \ar[d]_{\text{id}} &
Z''_{i - 1} \ar[r] \ar[u] \ar[d] &
\ldots \ar[r] &
Z''_0 \ar[r] \ar[u] \ar[d] &
B \ar[r] \ar[u]_{\text{id}} \ar[d]^{\text{id}} & 0 \\
0 \ar[r] & A \ar[r] & Z'_{i - 1} \ar[r] & \ldots \ar[r] &
Z'_0 \ar[r] & B \ar[r] & 0
}
$$
で中段も Yoneda 拡大であるものが存在するとき、両者は{\it 同値}であるという。
\end{definition}

\noindent
定義の同値性が同値関係であることは直ちには明らかでない。
これを直接証明することは有益だが、下の補題
\ref{derived-lemma-yoneda-extension} からも従う。

\medskip\noindent
$\mathcal{A}$ を対象 $A$, $B$ を持つアーベル圏とする。
Yoneda 拡大
$E : 0 \to A \to Z_{i - 1} \to Z_{i - 2} \to \ldots \to Z_0 \to B \to 0$
が与えられたとき、付随する元
$\delta(E) \in \Ext^i_\mathcal{A}(B, A)$ を射
$\delta(E) = fs^{-1} : B[0] \to A[i]$ として定義する。
ここで $s$ は擬同型
$$
(\ldots \to 0 \to A \to Z_{i - 1} \to \ldots \to Z_0 \to 0 \to \ldots)
\longrightarrow
B[0]
$$
であり、$f$ は複体の射
$$
(\ldots \to 0 \to A \to Z_{i - 1} \to \ldots \to Z_0 \to 0 \to \ldots)
\longrightarrow
A[i]
$$
である。$\delta(E) = fs^{-1}$ を Yoneda 拡大の{\it 類}という。
この類は Yoneda 拡大の同値類を特徴づけることが分かる。

\begin{lemma}
\label{derived-lemma-yoneda-extension}
$\mathcal{A}$ を対象 $A$, $B$ を持つアーベル圏とし、$i \geq 1$ とする。
$\Ext^i_\mathcal{A}(B, A)$ の任意の元は、$A$ による $B$ の次数 $i$ の
ある Yoneda 拡大に対する $\delta(E)$ である。
同じ次数の二つの Yoneda 拡大 $E$, $E'$ が与えられたとき、
$E$ と $E'$ が同値であることと
$\delta(E) = \delta(E')$ であることとは同値である。
\end{lemma}

\begin{proof}
$\xi : B[0] \to A[i]$ を $\Ext^i_\mathcal{A}(B, A)$ の元とする。
ある擬同型 $s : L^\bullet \to B[0]$ と写像
$f : L^\bullet \to A[i]$ を用いて $\xi = f s^{-1}$ と書ける。
$L^\bullet$ を $\tau_{\leq 0}L^\bullet$ で置き換えた後、
$j > 0$ に対して $L^j = 0$ と仮定してよい。図示すると
$$
\xymatrix{
L^{- i - 1} \ar[r] & L^{-i} \ar[r] \ar[d] & \ldots \ar[r] &
L^0 \ar[r] & B \ar[r] & 0 \\
& A
}
$$
である。$Z_{i - 1} = (L^{- i + 1} \oplus A)/L^{-i}$ と置き、
$j = i - 2, \ldots, 0$ に対して $Z_j = L^{-j}$ と置けば、
$A$ による $B$ の次数 $i$ の拡大 $E$ で、
その類 $\delta(E)$ が $\xi$ に等しいものを得る。

\medskip\noindent
同値な Yoneda 拡大が同じ類を持つことは定義から直ちに分かる。
$E : 0 \to A \to Z_{i - 1} \to Z_{i - 2} \to \ldots \to Z_0 \to B \to 0$
および
$E' : 0 \to A \to Z'_{i - 1} \to Z'_{i - 2} \to \ldots \to Z'_0 \to B \to 0$
を同じ類を持つ Yoneda 拡大と仮定する。
$D(\mathcal{A})$ を $K(\mathcal{A})$ の擬同型の集合における局所化として
構成したことにより、これは複体 $L^\bullet$ と擬同型
$$
t : L^\bullet \to
(\ldots \to 0 \to A \to Z_{i - 1} \to \ldots \to Z_0 \to 0 \to \ldots)
$$
および
$$
t' : L^\bullet \to
(\ldots \to 0 \to A \to Z'_{i - 1} \to \ldots \to Z'_0 \to 0 \to \ldots)
$$
で、$s \circ t = s' \circ t'$ および
$f \circ t = f' \circ t'$ を満たすものが存在することを意味する。
Categories, Section \ref{categories-section-localization} を参照せよ。
$E''$ を、証明の第一段落で一対の写像
$L^\bullet \to B[0]$ と $L^\bullet \to A[i]$ から構成した、
$A$ による $B$ の次数 $i$ の拡大とする。
このとき同値な Yoneda 拡大の定義にあるような、次数 $i$ の
Yoneda 拡大の「射」$E'' \to E$ および $E'' \to E'$ が存在することは
読者に容易に分かる（詳細は省略する）。これで証明が完了する。
\end{proof}

\begin{lemma}
\label{derived-lemma-ext-1}
$\mathcal{A}$ をアーベル圏とする。$A$, $B$ を $\mathcal{A}$ の対象とする。
このとき $\Ext^1_\mathcal{A}(B, A)$ は、Homology, Definition
\ref{homology-definition-ext-group} で構成した群
$\Ext_\mathcal{A}(B, A)$ である。
\end{lemma}

\begin{proof}
これは補題 \ref{derived-lemma-yoneda-extension} の $i = 1$ の場合である。
\end{proof}

\noindent
アーベル圏 $\mathcal{A}$ と $D(\mathcal{A})$ の対象 $X, Y, Z$
が与えられたとき、双線形かつ結合的な合成則
$$
\Ext^j_\mathcal{A}(Y, Z) \times \Ext^i_\mathcal{A}(X, Y)
\longrightarrow
\Ext^{i + j}_\mathcal{A}(X, Z),\quad
(\eta, \xi) \longmapsto \eta \circ \xi
$$
がある。すなわち $\xi : X \to Y[i]$ および
$\eta : Y \to Z[j]$ ならば、$\eta \circ \xi$ を
$\xi$ と $\eta[i] : Y[i] \to Z[i + j]$ の合成として定義する。
$A, B, C$ が $\mathcal{A}$ の対象で $i, j \geq 1$ ならば、この合成則
$$
\Ext^j_\mathcal{A}(B, C) \times \Ext^i_\mathcal{A}(A, B)
\longrightarrow
\Ext^{i + j}_\mathcal{A}(A, C)
$$
は Yoneda 拡大を用いて次のように記述できる。次の拡大
$$
0 \to A \to Z_{i - 1} \to Z_{i - 2} \to \ldots \to Z_0 \to B \to 0
$$
と
$$
0 \to B \to Z'_{j - 1} \to Z'_{j - 2} \to \ldots \to Z'_0 \to C \to 0
$$
の合成は Yoneda 拡大
$$
0 \to A \to Z_{i - 1} \to Z_{i - 2} \to \ldots \to Z_0 \to 
Z'_{j - 1} \to Z'_{j - 2} \to \ldots \to Z'_0 \to C \to 0
$$
である。

\begin{lemma}
\label{derived-lemma-cup-ext-1-zero}
$\mathcal{A}$ をアーベル圏とする。
$0 \to A \to Z \to B \to 0$ および
$0 \to B \to Z' \to C \to 0$ を $\mathcal{A}$ の短完全列とする。
それらの類を $[Z] \in \Ext^1_\mathcal{A}(B, A)$ および
$[Z'] \in \Ext^1_\mathcal{A}(C, B)$ と表す。
このとき $[Z] \circ [Z'] \in \Ext^2_\mathcal{A}(C, A)$ が $0$ であることと、
可換図
$$
\xymatrix{
&
&
0 \ar[d] &
0 \ar[d]
\\
0 \ar[r] &
A \ar[r] \ar[d]^1 &
Z \ar[r] \ar[d] &
B \ar[r] \ar[d] &
0 \\
0 \ar[r] &
A \ar[r]  &
W \ar[r] \ar[d] &
Z' \ar[r] \ar[d] &
0 \\
&
&
C \ar[r]^1  \ar[d]&
C \ar[d]\\
&
&
0 &
0 
}
$$
で行と列が $\mathcal{A}$ において完全であるものが存在することとは同値である。
\end{lemma}

\begin{proof}
省略する。ヒント：補題 \ref{derived-lemma-yoneda-extension} の結果を用い、
$2$-拡大の類が零であることの意味を具体的に調べればよい。
または、$[Z] \circ [Z'] \in \Ext^2_\mathcal{A}(C, A)$ が零ならば、
$\Ext$ の長完全コホモロジー列により
$[Z] \in \Ext^1_\mathcal{A}(B, A)$ が
$\Ext^1_\mathcal{A}(Z', A)$ のある元の像であることを用いてもよい。
\end{proof}

\begin{lemma}
\label{derived-lemma-higher-ext-zero}
$\mathcal{A}$ をアーベル圏とし、$p \geq 0$ とする。
$\mathcal{A}$ の任意の対象の対 $A$, $B$ に対して
$\Ext^p_\mathcal{A}(B, A) = 0$ ならば、$i \geq p$ および
$\mathcal{A}$ の任意の対象の対 $A$, $B$ に対して
$\Ext^i_\mathcal{A}(B, A) = 0$ である。
\end{lemma}

\begin{proof}
$i > p$ に対して、任意の類 $\xi$ を $\delta(E)$ と書く。
ここで $E$ は Yoneda 拡大
$$
E : 0 \to A \to Z_{i - 1} \to Z_{i - 2} \to \ldots \to Z_0 \to B \to 0
$$
である。これは補題 \ref{derived-lemma-yoneda-extension} により可能である。
$C = \Ker(Z_{p - 1} \to Z_{p - 2}) = \Im(Z_p \to Z_{p - 1})$ と置く。
このとき $\delta(E)$ は $\delta(E')$ と $\delta(E'')$ の合成である。
ここで
$$
E' : 0 \to C \to Z_{p - 1} \to \ldots \to Z_0 \to B \to 0
$$
および
$$
E'' : 0 \to A \to Z_{i - 1} \to Z_{i - 2} \to \ldots \to Z_p \to C \to 0
$$
である。$\delta(E') \in \Ext^p_\mathcal{A}(B, C) = 0$ なので結論を得る。
\end{proof}

\begin{lemma}
\label{derived-lemma-ext-2-zero-pre}
$\mathcal{A}$ をアーベル圏とし、$K$ を $D^b(\mathcal{A})$ の対象で、
すべての $p \geq 2$ と $i > j$ に対して
$\Ext^p_\mathcal{A}(H^i(K), H^j(K)) = 0$ を満たすものとする。
このとき $K$ はそのコホモロジーの直和と同型である。
$K \cong \bigoplus H^i(K)[-i]$。
\end{lemma}

\begin{proof}
$i \not \in [a, b]$ に対して $H^i(K) = 0$ となる $a, b$ を選ぶ。
$b - a$ に関する帰納法で補題を証明する。
$b - a \leq 0$ ならば結論は明らかである。
$b - a > 0$ ならば、切断の完全三角
$$
\tau_{\leq b - 1}K \to K \to H^b(K)[-b] \to (\tau_{\leq b - 1}K)[1]
$$
を見る。Remark \ref{derived-remark-truncation-distinguished-triangle} を参照せよ。
最後の矢印が零ならば、補題 \ref{derived-lemma-split} により
$K \cong \tau_{\leq b - 1}K \oplus H^b(K)[-b]$ であり、
帰納法により結論を得る。再び帰納法を用いると
$$
\Hom_{D(\mathcal{A})}(H^b(K)[-b], (\tau_{\leq b - 1}K)[1]) =
\bigoplus\nolimits_{i < b} \Ext_\mathcal{A}^{b - i + 1}(H^b(K), H^i(K))
$$
である。仮定により直和は零なので、証明が完了する。
\end{proof}

\begin{lemma}
\label{derived-lemma-ext-2-zero}
$\mathcal{A}$ をアーベル圏とする。
$\mathcal{A}$ の任意の対象の対 $A$, $B$ に対して
$\Ext^2_\mathcal{A}(B, A) = 0$ と仮定する。
このとき $D^b(\mathcal{A})$ の任意の対象 $K$ はそのコホモロジーの直和と
同型である。$K \cong \bigoplus H^i(K)[-i]$。
\end{lemma}

\begin{proof}
補題 \ref{derived-lemma-higher-ext-zero} により、仮定から $i \geq 2$ および
$\mathcal{A}$ の任意の対象の対 $A, B$ に対して
$\Ext^i_\mathcal{A}(B, A) = 0$ が従う。
従ってこの補題は補題 \ref{derived-lemma-ext-2-zero-pre} の特別な場合である。
\end{proof}







\section{K 群}
\label{derived-section-K-groups}

\noindent
三角圏の $K_0$ について少しだけ述べる。

\begin{definition}
\label{derived-definition-K-zero}
$\mathcal{D}$ を三角圏とする。$K_0(\mathcal{D})$ で
{\it $\mathcal{D}$ の第零 $K$ 群}を表す。
これは次のように構成されるアーベル群である。
% P02-E0068: frozen prose says “objects on D” rather than “objects of D”.
$\mathcal{D}$ の対象上の自由アーベル群を取り、各完全三角
$X \to Y \to Z$ に対して関係 $[Y] - [X] - [Z] = 0$ を課す。
\end{definition}

\noindent
完全三角 $(X, 0, X[1], 0, 0, -\text{id}[1])$ があるので、
これは $[X[n]] = (-1)^n[X]$ を意味することに注意する。

\begin{lemma}
\label{derived-lemma-K-bounded-derived}
$\mathcal{A}$ をアーベル圏とする。このとき第零 $K$ 群の標準的な同一視
$K_0(D^b(\mathcal{A})) = K_0(\mathcal{A})$ がある。
\end{lemma}

\begin{proof}
$\mathcal{A}$ の対象 $A$ が与えられたとき、$A[0]$ で次数 $0$ に置かれた
複体と見た対象 $A$ を表す。
$0 \to A \to A' \to A'' \to 0$ が短完全列ならば、完全三角
$A[0] \to A'[0] \to A''[0] \to A[1]$ を得る。
Section \ref{derived-section-canonical-delta-functor} を参照せよ。
従って、ひどい記法で申し訳ないが、$[A]$ を $[A[0]]$ に送ることで写像
$K_0(\mathcal{A}) \to K_0(D^b(\mathcal{A}))$ を得る。

\medskip\noindent
一方、$D^b(\mathcal{A})$ の対象 $X$ が与えられたとき、元
$$
c(X) = \sum (-1)^i[H^i(X)] \in K_0(\mathcal{A})
$$
を考えることができる。完全三角 $X \to Y \to Z$ が与えられると、
コホモロジーの長完全列
(\ref{derived-equation-long-exact-cohomology-sequence-D}) と
$K_0(\mathcal{A})$ の関係から $c(Y) = c(X) + c(Z)$ が従う。
従って $c$ は写像 $c : K_0(D^b(\mathcal{A})) \to K_0(\mathcal{A})$
を経由する。

\medskip\noindent
上の二つの写像が互いに逆であることを示したい。
合成 $K_0(\mathcal{A}) \to
K_0(D^b(\mathcal{A})) \to K_0(\mathcal{A})$ が恒等写像であることは明らかである。
$X^\bullet$ を $\mathcal{A}$ の有界複体と仮定する。
「素朴な切断」の完全三角（Homology, Section
\ref{homology-section-truncations} を参照）
$$
\sigma_{\geq n}X^\bullet \to \sigma_{\geq n - 1}X^\bullet \to
X^{n - 1}[-n + 1] \to (\sigma_{\geq n}X^\bullet)[1]
$$
の存在と帰納法により、$K_0(D^b(\mathcal{A}))$ において
$$
[X^\bullet] = \sum (-1)^i[X^i[0]]
$$
であることが分かる（記法について再びお詫びする）。
従って写像 $K_0(\mathcal{A}) \to K_0(D^b(\mathcal{A}))$ は全射であり、
証明が完了する。
\end{proof}

\begin{lemma}
\label{derived-lemma-map-K}
$F : \mathcal{D} \to \mathcal{D}'$ を三角圏の完全関手とする。
このとき $F$ は群準同型
$K_0(\mathcal{D}) \to K_0(\mathcal{D}')$ を誘導する。
\end{lemma}

\begin{proof}
省略する。
\end{proof}

\begin{lemma}
\label{derived-lemma-homological-map-K}
$H : \mathcal{D} \to \mathcal{A}$ を三角圏からアーベル圏への
ホモロジー的関手とする。$\mathcal{D}$ の任意の $X$ に対して、
$H(X[i])$ のうち $\mathcal{A}$ において非零であるものは有限個だけだと
仮定する。このとき $H$ は、$[X]$ を
$\sum (-1)^i[H(X[i])]$ に送る群準同型
$K_0(\mathcal{D}) \to K_0(\mathcal{A})$ を誘導する。
\end{lemma}

\begin{proof}
省略する。
\end{proof}

\begin{lemma}
\label{derived-lemma-DBA-map-K}
$\mathcal{B}$ をアーベル圏 $\mathcal{A}$ の弱 Serre 部分圏とする。
標準同型
$$
K_0(\mathcal{B}) \longrightarrow
K_0(D^b_\mathcal{B}(\mathcal{A})),\quad
[B] \longmapsto [B[0]]
$$
が存在する。その逆は $X$ の類 $[X]$ を元
$\sum (-1)^i[H^i(X)]$ に送る。
\end{lemma}

\begin{proof}
逆写像の規則が良定義な写像
$K_0(D^b_\mathcal{B}(\mathcal{A})) \to K_0(\mathcal{B})$
を与えることの確認は省略する。合成
$K_0(\mathcal{B}) \to
K_0(D^b_\mathcal{B}(\mathcal{A})) \to
K_0(\mathcal{B})$ が恒等写像であることは直ちに分かる。
一方、Remark \ref{derived-remark-truncation-distinguished-triangle} の完全三角と
帰納法を用いて、補題の主張に表示された矢印が全射であることを
読者は示せる（詳細は省略する）。補題が従う。
\end{proof}

\begin{lemma}
\label{derived-lemma-bilinear-map-K}
$\mathcal{D}$, $\mathcal{D}'$, $\mathcal{D}''$ を三角圏とする。
関手
$$
\otimes : \mathcal{D} \times \mathcal{D}' \longrightarrow \mathcal{D}''
$$
について、$\mathcal{D}$ の $X$ を固定した関手
$X \otimes - : \mathcal{D}' \to \mathcal{D}''$ が完全関手であり、
$\mathcal{D}'$ の $X'$ を固定した関手
$- \otimes X' : \mathcal{D} \to \mathcal{D}''$ が完全関手であるとする。
このとき $\otimes$ は、$([X], [X'])$ を $[X \otimes X']$ に送る双線形写像
$K_0(\mathcal{D}) \times K_0(\mathcal{D}') \to K_0(\mathcal{D}'')$
を誘導する。
\end{lemma}

\begin{proof}
省略する。
\end{proof}

\section{非有界複体}
\label{derived-section-unbounded}

\noindent
本節の内容については \cite{Spaltenstein} を参照されたい。
次の補題は、非有界複体の「良い」左分解を見つけるのに有用である。

\begin{lemma}
\label{derived-lemma-special-direct-system}
$\mathcal{A}$ をアーベル圏とし、
$\mathcal{P} \subset \Ob(\mathcal{A})$ を部分集合とする。
$\mathcal{P}$ は $0$ を含み、（有限）直和について閉じており、
$\mathcal{A}$ の任意の対象は $\mathcal{P}$ の元の商であると仮定する。
$K^\bullet$ を複体とする。このとき、複体の圏における可換図式
$$
\xymatrix{
P_1^\bullet \ar[d] \ar[r] & P_2^\bullet \ar[d] \ar[r] & \ldots \\
\tau_{\leq 1}K^\bullet \ar[r] & \tau_{\leq 2}K^\bullet \ar[r] & \ldots
}
$$
であって、次を満たすものが存在する。
\begin{enumerate}
\item 縦の射は擬同型であり、項ごとに全射である。
\item $P_n^\bullet$ は、各項が $\mathcal{P}$ に属する上に有界な複体である。
\item 射 $P_n^\bullet \to P_{n + 1}^\bullet$ は項ごとに分裂する単射であり、
各余核 $P^i_{n + 1}/P^i_n$ は $\mathcal{P}$ の元である。
\end{enumerate}
\end{lemma}

\begin{proof}
ホモトピー圏 $K(\mathcal{A})$ が三角圏であることを用いる。命題
\ref{derived-proposition-homotopy-category-triangulated} を参照されたい。
補題 \ref{derived-lemma-subcategory-left-resolution} により、項ごとに全射な複体の写像
$P_1^\bullet \to \tau_{\leq 1}K^\bullet$ であって擬同型であり、しかも
$P_1^\bullet$ の各項が $\mathcal{P}$ に属するものを見つけられる。
帰納法により、$P_1^\bullet, \ldots, P_n^\bullet$ が与えられたときに
$P_{n + 1}^\bullet$ と写像
$P_n^\bullet \to P_{n + 1}^\bullet$ および
$P_{n + 1}^\bullet \to \tau_{\leq n + 1}K^\bullet$
を構成すれば十分である。

\medskip\noindent
$K(\mathcal{A})$ において完全三角
$P_n^\bullet \to \tau_{\leq n + 1}K^\bullet \to C^\bullet \to P_n^\bullet[1]$
を選ぶ。補題 \ref{derived-lemma-subcategory-left-resolution} を適用して、
擬同型である複体の写像 $Q^\bullet \to C^\bullet$ であって、
$Q^\bullet$ の各項が $\mathcal{P}$ に属するものを選ぶ。
三角圏の公理により、合成
$Q^\bullet \to C^\bullet \to P_n^\bullet[1]$ を
$K(\mathcal{A})$ における完全三角
$P_n^\bullet \to P_{n + 1}^\bullet \to Q^\bullet \to P_n^\bullet[1]$
に組み込める。
補題 \ref{derived-lemma-improve-distinguished-triangle-homotopy} により、
$$
0 \to P_n^\bullet \to P_{n + 1}^\bullet \to Q^\bullet \to 0
$$
が項ごとに分裂する短完全列であると仮定してよいし、実際そうする。
これより、$P_{n + 1}^\bullet$ の各項は $\mathcal{P}$ に属し、
$P_n^\bullet \to P_{n + 1}^\bullet$ は項ごとに分裂する単射で、
その余核は $\mathcal{P}$ に属する。
三角圏の公理により、三角圏 $K(\mathcal{A})$ における完全三角の射
$$
\xymatrix{
P_n^\bullet \ar[r] \ar[d] &
P_{n + 1}^\bullet \ar[r] \ar[d] &
Q^\bullet \ar[r] \ar[d] &
P_n^\bullet[1] \ar[d] \\
P_n^\bullet \ar[r] &
\tau_{\leq n + 1}K^\bullet \ar[r] &
C^\bullet \ar[r] &
P_n^\bullet[1]
}
$$
を得る。実際の複体の射
$f : P_{n + 1}^\bullet \to \tau_{\leq n + 1}K^\bullet$ を選ぶ。
上図の左の四角形はホモトピー可換であるが、
$P_n^\bullet \to P_{n + 1}^\bullet$ は項ごとに分裂する単射なので、
このホモトピーを持ち上げ、$f$ の選び方を修正して四角形を可換にできる。
最後に、$P_n^\bullet \to P_n^\bullet$ と
$Q^\bullet \to C^\bullet$ がともに擬同型であるため、$f$ も擬同型である。

\medskip\noindent
ここまでで所望の性質はすべて得られたが、写像
$f : P_{n + 1}^\bullet \to \tau_{\leq n + 1}K^\bullet$
が項ごとに全射であるかどうかだけはまだ分からない。複体の可換図式
$$
\xymatrix{
P_n^\bullet \ar[d] \ar[r] & P_{n + 1}^\bullet \ar[d] \\
\tau_{\leq n}K^\bullet \ar[r] & \tau_{\leq n + 1}K^\bullet
}
$$
があるので、帰納法の仮定により、次数 $n$ と $n + 1$ を除く
すべての次数で $f$ は項の上で全射である。対象
$P \in \mathcal{P}$ と全射 $q : P \to K^n$ を選ぶ。
写像
$$
g :
P^\bullet = (\ldots \to 0 \to P \xrightarrow{1} P \to 0 \to \ldots)
\longrightarrow
\tau_{\leq n + 1}K^\bullet
$$
を考える。ここで最初の $P$ は次数 $n$ にあり、写像は次数 $n$ では
$q$、次数 $n + 1$ では $d_K \circ q$ で与えられる。
これは次数 $n$ で全射であり、次数 $n + 1$ における余核は
$H^{n + 1}(\tau_{\leq n + 1}K^\bullet)$ である。実際、
$\tau_{\leq n + 1}K^\bullet$ の次数 $n + 1$ の項は
$\Ker(d_K^{n + 1})$ であることを思い出せばよい。
しかし $f$ は擬同型なので、$H^{n + 1}(f)$ が全射であることが分かる。
したがって
$f : P_{n + 1}^\bullet \to \tau_{\leq n + 1}K^\bullet$
を
$f \oplus g : P_{n + 1}^\bullet \oplus P^\bullet \to \tau_{\leq n + 1}K^\bullet$
で置き換えれば、目的が達せられる。
\end{proof}

\noindent
場合によっては、上の補題を用いて左導来関手が至る所で定義されることを
示せる。

\begin{proposition}
\label{derived-proposition-left-derived-exists}
$F : \mathcal{A} \to \mathcal{B}$ をアーベル圏の間の右完全関手とする。
$\mathcal{P} \subset \Ob(\mathcal{A})$ を部分集合とし、次を仮定する。
\begin{enumerate}
\item $\mathcal{P}$ は $0$ を含み、（有限）直和について閉じており、
$\mathcal{A}$ の任意の対象は $\mathcal{P}$ の元の商である。
\item $\mathcal{A}$ の任意の上に有界な非輪状複体 $P^\bullet$ について、
すべての $n$ で $P^n \in \mathcal{P}$ ならば、複体
$F(P^\bullet)$ は完全である。
\item $\mathcal{A}$ と $\mathcal{B}$ は
$\mathbf{N}$ 上の系の余極限をもつ。
\item $\mathbf{N}$ 上の余極限は $\mathcal{A}$ と
$\mathcal{B}$ の両方において完全である。
\item $F$ は $\mathbf{N}$ 上の余極限と可換である。
\end{enumerate}
このとき $LF$ は $D(\mathcal{A})$ の全体で定義される。
\end{proposition}

\begin{proof}
(1) と補題 \ref{derived-lemma-subcategory-left-resolution} により、任意の
上に有界な複体 $K^\bullet$ に対して擬同型
$P^\bullet \to K^\bullet$ であって、$P^\bullet$ は上に有界で、
すべての $n$ について $P^n \in \mathcal{P}$ となるものが存在する。
$s : P^\bullet \to (P')^\bullet$ を、$\mathcal{P}$ の対象からなる
上に有界な複体の間の擬同型とする。このとき仮定 (2) により
$F(C(s)^\bullet)$ は非輪状なので、
$F(P^\bullet) \to F((P')^\bullet)$ は擬同型である。
これだけで $LF$ が $D^{-}(\mathcal{A})$ 上で定義され、
$\mathcal{P}$ の対象からなる上に有界な複体が $LF$ を計算することが分かる。
補題 \ref{derived-lemma-find-existence-computes} を参照されたい。

\medskip\noindent
次に、$K^\bullet$ を $\mathcal{A}$ の任意の複体とする。
補題 \ref{derived-lemma-special-direct-system} にある図式
$$
\xymatrix{
P_1^\bullet \ar[d] \ar[r] & P_2^\bullet \ar[d] \ar[r] & \ldots \\
\tau_{\leq 1}K^\bullet \ar[r] & \tau_{\leq 2}K^\bullet \ar[r] & \ldots
}
$$
を選ぶ。$\mathcal{A}$ における $\mathbf{N}$ 上の余極限は完全であり、
$n > i$ ならば $H^i(P_n^\bullet) = H^i(K^\bullet)$ なので、写像
$\colim P_n^\bullet \to K^\bullet$ は擬同型であることに注意する。
ここで
$$
F(\colim P_n^\bullet) = \colim F(P_n^\bullet)
$$
（項ごとの余極限）が $LF(K^\bullet)$ である、すなわち
$\colim P_n^\bullet$ が $LF$ を計算すると主張する。
これを示すには、補題 \ref{derived-lemma-find-existence-computes} により、
次の主張を証明すれば十分である。
$$
\colim Q_n^\bullet = Q^\bullet
\xrightarrow{\ \alpha\ }
P^\bullet = \colim P_n^\bullet
$$
が複体の擬同型であり、各 $P_n^\bullet$, $Q_n^\bullet$ は
各項が $\mathcal{P}$ に属する上に有界な複体で、写像
$P_n^\bullet \to \tau_{\leq n}P^\bullet$ と
$Q_n^\bullet \to \tau_{\leq n}Q^\bullet$ は擬同型であると仮定する。
主張は、$F(\alpha)$ が擬同型だということである。

\medskip\noindent
問題は、$\alpha$ が複体 $P_n^\bullet$ と $Q_n^\bullet$ の間の
写像の余極限として与えられるとは仮定していないことである。
しかし各 $n$ に対し、次の図式の実線の射が擬同型であることは分かっている。
% P02-E0070: alpha is declared Q -> P, while the entire downstream construction consistently uses P -> Q.
$$
\xymatrix{
& R^\bullet \ar@{..>}[d] \\
P_n^\bullet \ar[d] &
L^\bullet \ar@{..>}[l] \ar@{..>}[r] &
Q_n^\bullet \ar[d] \\
\tau_{\leq n}P^\bullet \ar[rr]^{\tau_{\leq n}\alpha} & &
\tau_{\leq n}Q^\bullet
}
$$
擬同型は $K(\mathcal{A})$ において乗法系をなすので
（補題 \ref{derived-lemma-acyclic} を参照）、擬同型
$L^\bullet \to P_n^\bullet$ と複体の写像 $L^\bullet \to Q_n^\bullet$
であって、上図がホモトピー可換となるものを見つけられる。
すると $\tau_{\leq n}L^\bullet \to L^\bullet$ は擬同型である。
したがって（証明の第一部分により）、各項が $\mathcal{P}$ に属する
上に有界な複体 $R^\bullet$ と擬同型 $R^\bullet \to L^\bullet$
を見つけられる（図中に示した通りである）。
証明の第一段落の結果を用いると、
$F(R^\bullet) \to F(P_n^\bullet)$ と $F(R^\bullet) \to F(Q_n^\bullet)$
が擬同型であることが分かる。したがって、可換図式
$$
\xymatrix{
H^i(F(P_n^\bullet)) \ar[r] \ar[d] &
H^i(F(Q_n^\bullet)) \ar[d] \\
H^i(F(P^\bullet)) \ar[r] &
H^i(F(Q^\bullet))
}
$$
に組み込まれる同型
$H^i(F(P_n^\bullet)) \to H^i(F(Q_n^\bullet))$ を得る。
まったく同じ議論により、$n$ が変化するときにもこれらの写像は両立する。
(4) と (5) により
$$
H^i(F(P^\bullet)) =
H^i(F(\colim P_n^\bullet)) =
H^i(\colim F(P_n^\bullet)) = \colim H^i(F(P_n^\bullet))
$$
であり、$Q^\bullet$ についても同様なので、
% P02-E0071: the frozen displayed inline map has unmatched parentheses and omits F(alpha); direction follows P02-E0070.
$H^i(\alpha) : H^i(F(P^\bullet) \to H^i(F(Q^\bullet)$ は同型であると結論でき、
主張が従う。
\end{proof}

\begin{lemma}
\label{derived-lemma-special-inverse-system}
$\mathcal{A}$ をアーベル圏とし、
$\mathcal{I} \subset \Ob(\mathcal{A})$ を部分集合とする。
$\mathcal{I}$ は $0$ を含み、（有限）積について閉じており、
$\mathcal{A}$ の任意の対象は $\mathcal{I}$ の元の部分対象であると仮定する。
$K^\bullet$ を複体とする。このとき、複体の圏における可換図式
$$
\xymatrix{
\ldots \ar[r] &
\tau_{\geq -2}K^\bullet \ar[r] \ar[d] &
\tau_{\geq -1}K^\bullet \ar[d] \\
\ldots \ar[r] & I_2^\bullet \ar[r] & I_1^\bullet
}
$$
であって、次を満たすものが存在する。
\begin{enumerate}
\item 縦の射は擬同型であり、項ごとに単射である。
\item $I_n^\bullet$ は、各項が $\mathcal{I}$ に属する下に有界な複体である。
\item 射 $I_{n + 1}^\bullet \to I_n^\bullet$ は項ごとに分裂する全射であり、
$\Ker(I^i_{n + 1} \to I^i_n)$ は $\mathcal{I}$ の元である。
\end{enumerate}
\end{lemma}

\begin{proof}
本補題は補題 \ref{derived-lemma-special-direct-system} の双対である。
\end{proof}

\section{随伴関手の導来}
\label{derived-section-deriving-adjoints}

\noindent
$F : \mathcal{D} \to \mathcal{D}'$ と
$G : \mathcal{D}' \to \mathcal{D}$ を三角圏の完全関手とする。
$S$、それぞれ $S'$ を、$\mathcal{D}$、それぞれ $\mathcal{D}'$ の
三角構造と両立する乗法系とする。
$Q : \mathcal{D} \to S^{-1}\mathcal{D}$ と
$Q' : \mathcal{D}' \to (S')^{-1}\mathcal{D}'$ で局所化関手を表す。
この状況では記法を濫用して、
$Q' \circ F : \mathcal{D} \to (S')^{-1}\mathcal{D}'$ と乗法系 $S$
に対応する部分的に定義された右導来関手をしばしば $RF$ と表す。
同様に、$Q \circ G : \mathcal{D}' \to S^{-1}\mathcal{D}$ と乗法系
$S'$ に対応する部分的に定義された左導来関手を $LG$ と表す。図式は
$$
\vcenter{
\xymatrix{
\mathcal{D} \ar[r]_F \ar[d]_Q & \mathcal{D}' \ar[d]^{Q'} \\
S^{-1}\mathcal{D} \ar@{..>}[r]^{RF} &
(S')^{-1}\mathcal{D}'
}
}
\quad\text{and}\quad
\vcenter{
\xymatrix{
\mathcal{D}' \ar[r]_G \ar[d]_{Q'} & \mathcal{D} \ar[d]^Q \\
(S')^{-1}\mathcal{D}' \ar@{..>}[r]^{LG} &
S^{-1}\mathcal{D}
}
}
$$

\begin{lemma}
\label{derived-lemma-pre-derived-adjoint-functors-general}
上の状況で $F$ は $G$ の右随伴であると仮定する。
$K \in \Ob(\mathcal{D})$ および $M \in \Ob(\mathcal{D}')$ とする。
$RF$ が $K$ において定義され、$LG$ が $M$ において定義されるならば、
標準同型
$$
\Hom_{(S')^{-1}\mathcal{D}'}(M, RF(K)) =
\Hom_{S^{-1}\mathcal{D}}(LG(M), K)
$$
が存在する。この同型は、$RF$ と $LG$ が定義される
$S^{-1}\mathcal{D}$ と $(S')^{-1}\mathcal{D}'$ の三角部分圏上で、
両方の変数について関手的である。
\end{lemma}

\begin{proof}
$RF$ は $K$ において定義されるので、$S$ の
$s : K \to I$ に対象 $F(I)$ を対応させる規則は、
$(S')^{-1}\mathcal{D}'$ の ind-対象として本質的に定値であり、
その値は $RF(K)$ である。同様に、$S'$ の
$t : P \to M$ に対象 $G(P)$ を対応させる規則は、
$S^{-1}\mathcal{D}$ の pro-対象として本質的に定値であり、
その値は $LG(M)$ である。したがって
\begin{align*}
\Hom_{(S')^{-1}\mathcal{D}'}(M, RF(K))
& =
\colim_{s : K \to I} \Hom_{(S')^{-1}\mathcal{D}'}(M, F(I)) \\
& =
\colim_{s : K \to I} \colim_{t : P \to M} \Hom_{\mathcal{D}'}(P, F(I)) \\
& =
\colim_{t : P \to M} \colim_{s : K \to I} \Hom_{\mathcal{D}'}(P, F(I)) \\
& =
\colim_{t : P \to M} \colim_{s : K \to I} \Hom_{\mathcal{D}}(G(P), I) \\
& =
\colim_{t : P \to M} \Hom_{S^{-1}\mathcal{D}}(G(P), K) \\
& =
\Hom_{S^{-1}\mathcal{D}}(LG(M), K)
\end{align*}
第一の等号は Categories, Lemma
\ref{categories-lemma-characterize-essentially-constant-ind} による。
% P02-E0072: the generalized proof refers to D(B), although no abelian category B is present here.
第二の等号は $D(\mathcal{B})$ における射の定義による。Categories, Remark
\ref{categories-remark-right-localization-morphisms-colimit} を参照されたい。
第三の等号は Categories, Lemma
\ref{categories-lemma-colimits-commute} による。
第四の等号は $F$ と $G$ が随伴であることによる。
% P02-E0073: similarly, the generalized proof refers to D(A), although only S^{-1}D is defined.
第五の等号は $D(\mathcal{A})$ における射の定義による。Categories, Remark
\ref{categories-remark-left-localization-morphisms-colimit} を参照されたい。
第六の等号は Categories, Lemma
\ref{categories-lemma-characterize-essentially-constant-pro} による。
関手性の証明は省略する。
\end{proof}

\begin{lemma}
\label{derived-lemma-pre-derived-adjoint-functors}
$F : \mathcal{A} \to \mathcal{B}$ と $G : \mathcal{B} \to \mathcal{A}$
をアーベル圏の関手とし、$F$ は $G$ の右随伴であるとする。
$K^\bullet$ を $\mathcal{A}$ の複体、$M^\bullet$ を
$\mathcal{B}$ の複体とする。$RF$ が $K^\bullet$ において定義され、
$LG$ が $M^\bullet$ において定義されるならば、標準同型
$$
\Hom_{D(\mathcal{B})}(M^\bullet, RF(K^\bullet)) =
\Hom_{D(\mathcal{A})}(LG(M^\bullet), K^\bullet)
$$
が存在する。この同型は、$RF$ と $LG$ が定義される
$D(\mathcal{A})$ と $D(\mathcal{B})$ の三角部分圏上で、
両方の変数について関手的である。
\end{lemma}

\begin{proof}
これは非常に一般的な補題
\ref{derived-lemma-pre-derived-adjoint-functors-general} の特別な場合である。
\end{proof}

\noindent
次の補題は、非有界導来圏を用いる方が容易である理由の一例である。
すなわち、非有界な導来関手がなければ、この補題は述べることすらできない。

\begin{lemma}
\label{derived-lemma-derived-adjoint-functors}
$F : \mathcal{A} \to \mathcal{B}$ と $G : \mathcal{B} \to \mathcal{A}$
をアーベル圏の関手とし、$F$ は $G$ の右随伴であるとする。
導来関手 $RF : D(\mathcal{A}) \to D(\mathcal{B})$ と
$LG : D(\mathcal{B}) \to D(\mathcal{A})$ が存在するならば、
$RF$ は $LG$ の右随伴である。
\end{lemma}

\begin{proof}
補題 \ref{derived-lemma-pre-derived-adjoint-functors} から直ちに従う。
\end{proof}

\section{K-入射的複体}
\label{derived-section-K-injective}

\noindent
次の種類の複体は、非有界導来圏上の右導来関手を計算するために使える。

\begin{definition}
\label{derived-definition-K-injective}
$\mathcal{A}$ をアーベル圏とする。複体 $I^\bullet$ が
{\it K-入射的}であるとは、任意の非輪状複体 $M^\bullet$ に対して
$\Hom_{K(\mathcal{A})}(M^\bullet, I^\bullet) = 0$ となることをいう。
\end{definition}

\noindent
この定義の状況では、非輪状複体の平行移動も非輪状なので、
実際、すべての $i$ について
$\Hom_{K(\mathcal{A})}(M^\bullet[i], I^\bullet) = 0$ である。

\begin{lemma}
\label{derived-lemma-K-injective}
$\mathcal{A}$ をアーベル圏とし、$I^\bullet$ を複体とする。
次は同値である。
\begin{enumerate}
\item $I^\bullet$ は K-入射的である。
\item 任意の擬同型 $M^\bullet \to N^\bullet$ に対して写像
$$
\Hom_{K(\mathcal{A})}(N^\bullet, I^\bullet)
\to \Hom_{K(\mathcal{A})}(M^\bullet, I^\bullet)
$$
は全単射である。
\item 任意の複体 $N^\bullet$ に対して写像
$$
\Hom_{K(\mathcal{A})}(N^\bullet, I^\bullet)
\to \Hom_{D(\mathcal{A})}(N^\bullet, I^\bullet)
$$
は同型である。
\end{enumerate}
\end{lemma}

\begin{proof}
(1) を仮定する。このとき関手
$\Hom_{K(\mathcal{A})}( - , I^\bullet)$ はコホモロジー的であり、
擬同型の錐は非輪状なので、(2) が成り立つ。

\medskip\noindent
(2) を仮定する。$D(\mathcal{A})$ における射
$N^\bullet \to I^\bullet$ は
$fs^{-1} : N^\bullet \to I^\bullet$ の形である。ここで
$s : M^\bullet \to N^\bullet$ は擬同型で、
$f : M^\bullet \to I^\bullet$ は写像である。(2) により、これは
$K(\mathcal{A})$ における一意な射 $N^\bullet \to I^\bullet$
に対応する。すなわち (3) が成り立つ。

\medskip\noindent
(3) を仮定する。$M^\bullet$ が非輪状ならば、$M^\bullet$ は
$D(\mathcal{A})$ において零複体と同型なので
$\Hom_{D(\mathcal{A})}(M^\bullet, I^\bullet) = 0$ である。
従って (3) により
$\Hom_{K(\mathcal{A})}(M^\bullet, I^\bullet) = 0$ となる。
すなわち (1) が成り立つ。
\end{proof}

\begin{lemma}
\label{derived-lemma-triangle-K-injective}
$\mathcal{A}$ をアーベル圏とする。$(K, L, M, f, g, h)$ を
$K(\mathcal{A})$ の完全三角とする。$K$, $L$, $M$ のうち二つが
K-入射的複体ならば、残る一つも K-入射的複体である。
\end{lemma}

\begin{proof}
定義、補題 \ref{derived-lemma-representable-homological}、および
$K(\mathcal{A})$ が三角圏であることから従う
（命題 \ref{derived-proposition-homotopy-category-triangulated}）。
\end{proof}

\begin{lemma}
\label{derived-lemma-bounded-below-injectives-K-injective}
$\mathcal{A}$ をアーベル圏とする。入射対象からなる下に有界な複体は
K-入射的である。
\end{lemma}

\begin{proof}
補題 \ref{derived-lemma-K-injective} と
\ref{derived-lemma-morphisms-into-injective-complex} から従う。
\end{proof}

\begin{lemma}
\label{derived-lemma-product-K-injective}
$\mathcal{A}$ をアーベル圏とする。$T$ を集合とし、各
$t \in T$ に対して $I_t^\bullet$ を K-入射的複体とする。
すべての $n$ に対して $I^n = \prod_t I_t^n$ が存在するならば、
$I^\bullet$ は K-入射的複体である。さらに、$I^\bullet$ は
$D(\mathcal{A})$ において対象 $I_t^\bullet$ の積を表現する。
\end{lemma}

\begin{proof}
% P02-E0074: frozen prose reads “an complex” rather than “a complex”.
$K^\bullet$ を複体とする。複体
$$
C :
\prod\nolimits_b \Hom(K^{-b}, I^{b - 1}) \to
\prod\nolimits_b \Hom(K^{-b}, I^b) \to
\prod\nolimits_b \Hom(K^{-b}, I^{b + 1})
$$
の中央のコホモロジーは
$\Hom_{K(\mathcal{A})}(K^\bullet, I^\bullet)$ であることに注意する。
同様に、複体
$$
C_t :
\prod\nolimits_b \Hom(K^{-b}, I_t^{b - 1}) \to
\prod\nolimits_b \Hom(K^{-b}, I_t^b) \to
\prod\nolimits_b \Hom(K^{-b}, I_t^{b + 1})
$$
は $\Hom_{K(\mathcal{A})}(K^\bullet, I_t^\bullet)$ を計算する。
次に、$I$ の選び方により、アーベル群の複体として
$$
C = \prod\nolimits_{t \in T} C_t
$$
であることに注意する。積を取ることはアーベル群の圏上の完全関手である。
従って $K^\bullet$ が非輪状ならば
$\Hom_{K(\mathcal{A})}(K^\bullet, I_t^\bullet) = 0$ なので
$C_t$ は非輪状であり、従って $C$ も非輪状である。ゆえに
$\Hom_{K(\mathcal{A})}(K^\bullet, I^\bullet) = 0$ を得る。
従って $I^\bullet$ は K-入射的である。
以上を踏まえ、補題 \ref{derived-lemma-K-injective} を用いると
$$
\Hom_{D(\mathcal{A})}(K^\bullet, I^\bullet)
=
\prod\nolimits_{t \in T} \Hom_{D(\mathcal{A})}(K^\bullet, I_t^\bullet)
$$
と結論でき、実際 $I^\bullet$ は導来圏における積を表現する。
\end{proof}

\begin{lemma}
\label{derived-lemma-K-injective-defined}
$\mathcal{A}$ をアーベル圏とする。
$F : K(\mathcal{A}) \to \mathcal{D}'$ を三角圏の完全関手とする。
このとき $RF$ は、$K(\mathcal{A})$ のうち K-入射的複体と
擬同型な任意の複体において定義される。実際、任意の K-入射的複体は
$RF$ を計算する。
\end{lemma}

\begin{proof}
補題 \ref{derived-lemma-derived-inverts} により、$RF$ が K-入射的複体で
定義されること、すなわち K-入射的複体 $I^\bullet$ が $RF$ を
計算することを示せば十分である。擬同型
$I^\bullet \to N^\bullet$ を考える。補題
\ref{derived-lemma-K-injective} により、これは左逆をもつ。
従って $I^\bullet \to I^\bullet$ は
$I^\bullet/\text{Qis}(\mathcal{A})$ の終対象であり、目的が達せられる。
\end{proof}

\begin{lemma}
\label{derived-lemma-enough-K-injectives-implies}
$\mathcal{A}$ をアーベル圏とする。任意の複体から K-入射的複体への
擬同型が存在すると仮定する。このとき、三角圏の任意の完全関手
$F : K(\mathcal{A}) \to \mathcal{D}'$ は右導来関手
$$
RF : D(\mathcal{A}) \longrightarrow \mathcal{D}'
$$
をもち、K-入射的複体 $I^\bullet$ に対して
$RF(I^\bullet) = F(I^\bullet)$ である。
\end{lemma}

\begin{proof}
これを見るため、K-入射的複体の集まりを $\mathcal{I}$ として
補題 \ref{derived-lemma-find-existence-computes} を適用する。
(1) は仮定により成り立つので、$I^\bullet \to J^\bullet$ が
K-入射的複体の間の擬同型ならば
$F(I^\bullet) \to F(J^\bullet)$ が同型であることを示せば十分である。
これは明らかである。実際、補題 \ref{derived-lemma-K-injective} により
$I^\bullet \to J^\bullet$ はホモトピー同値、すなわち
$K(\mathcal{A})$ における同型である。
\end{proof}

\noindent
次の補題は、より大きい順序数上の極限へ一般化できる。

\begin{lemma}
\label{derived-lemma-limit-K-injectives}
\begin{slogan}
K-入射的複体の「分裂した」塔の極限は K-入射的である。
\end{slogan}
$\mathcal{A}$ をアーベル圏とし、
$$
\ldots \to I_3^\bullet \to I_2^\bullet \to I_1^\bullet
$$
を複体の逆系とする。次を仮定する。
\begin{enumerate}
\item 各 $I_n^\bullet$ は $K$-入射的である。
\item 各写像 $I_{n + 1}^m \to I_n^m$ は分裂全射である。
\item 極限 $I^m = \lim I_n^m$ が存在する。
\end{enumerate}
このとき複体 $I^\bullet$ は K-入射的である。
\end{lemma}

\begin{proof}
読者にはこの補題の証明を飛ばすことを強く勧める。
$M^\bullet$ を非輪状複体とし、
$H_n(a, b) = \Hom_\mathcal{A}(M^a, I_n^b)$ と略記する。この記法では
$\Hom_{K(\mathcal{A})}(M^\bullet, I^\bullet)$ は複体
$$
\prod_m \lim\limits_n H_n(m, m - 2)
\to
\prod_m \lim\limits_n H_n(m, m - 1)
\to
\prod_m \lim\limits_n H_n(m, m)
\to
\prod_m \lim\limits_n H_n(m, m + 1)
$$
の左から三番目におけるコホモロジーである。
$\prod$ と $\lim$ の順序を交換でき、複体
$$
\prod_m H_n(m, m - 2)
\to
\prod_m H_n(m, m - 1)
\to
\prod_m H_n(m, m)
\to
\prod_m H_n(m, m + 1)
$$
はそれぞれ仮定 (1) により完全である。仮定 (2) により、系
$$
\ldots \to
\prod_m H_3(m, m - 2) \to
\prod_m H_2(m, m - 2) \to
\prod_m H_1(m, m - 2)
$$
の写像は全射である。従って補題は Homology, Lemma
\ref{homology-lemma-apply-Mittag-Leffler} から従う。
\end{proof}

\noindent
補題 \ref{derived-lemma-special-inverse-system}、
\ref{derived-lemma-bounded-below-injectives-K-injective}、および
\ref{derived-lemma-limit-K-injectives} を組み合わせると、十分多くの入射対象と
可算積をもつ任意のアーベル圏に対して「十分多くの K-入射的複体」が
得られるように見える。実際には、これはうまくいかないかもしれない。
説明については補題 \ref{derived-lemma-difficulty-K-injectives} を参照されたい。

\begin{lemma}
\label{derived-lemma-adjoint-preserve-K-injectives}
$\mathcal{A}$ と $\mathcal{B}$ をアーベル圏とする。
$u : \mathcal{A} \to \mathcal{B}$ と
$v : \mathcal{B} \to \mathcal{A}$ を加法的関手とし、次を仮定する。
\begin{enumerate}
\item $u$ は $v$ の右随伴である。
\item $v$ は完全である。
\end{enumerate}
このとき $u$ は K-入射的複体を K-入射的複体へ移す。
\end{lemma}

\begin{proof}
$I^\bullet$ を $\mathcal{A}$ の K-入射的複体とする。
% P02-E0075: frozen prose reads “a acyclic complex”.
$M^\bullet$ を $\mathcal{B}$ の非輪状複体とする。
$v$ は完全なので、$v(M^\bullet)$ は非輪状複体である。
随伴性により
$$
0 = \Hom_{K(\mathcal{A})}(v(M^\bullet), I^\bullet) =
\Hom_{K(\mathcal{B})}(M^\bullet, u(I^\bullet))
$$
を得る。従って補題が従う。
\end{proof}

\section{有界コホモロジー次元}
\label{derived-section-bounded}

\noindent
非有界導来関手が存在するもう一つの場合がある。
すなわち、関手が有界コホモロジー次元をもつ場合である。

\begin{lemma}
\label{derived-lemma-replace-resolution}
$\mathcal{A}$ をアーベル圏とする。
$d : \Ob(\mathcal{A}) \to \{0, 1, 2, \ldots, \infty\}$ を関数とし、
次を仮定する。
\begin{enumerate}
\item $\mathcal{A}$ の任意の対象は、$d(A) = 0$ を満たす対象
$A$ の部分対象である。
\item $A, B \in \mathcal{A}$ に対して
$d(A \oplus B) \leq \max \{d(A), d(B)\}$ である。
\item $0 \to A \to B \to C \to 0$ が短完全ならば
$d(C) \leq \max\{d(A) - 1, d(B)\}$ である。
\end{enumerate}
$K^\bullet$ を、$n \to -\infty$ のとき $n + d(K^n)$ が
$-\infty$ に収束する複体とする。このとき擬同型
$K^\bullet \to L^\bullet$ であって、すべての
$n \in \mathbf{Z}$ に対して $d(L^n) = 0$ となるものが存在する。
\end{lemma}

\begin{proof}
補題 \ref{derived-lemma-subcategory-right-resolution} により、擬同型
$\sigma_{\geq 0}K^\bullet \to M^\bullet$ であって、
$n < 0$ ならば $M^n = 0$、$n \geq 0$ ならば $d(M^n) = 0$
となるものを見つけられる。このとき $K^\bullet$ は複体
$$
\ldots \to K^{-2} \to K^{-1} \to M^0 \to M^1 \to \ldots
$$
と擬同型である。従って $n \gg 0$ のとき
$d(K^n) = 0$ であると仮定してよい。
$n \to -\infty$ のとき $n + d(K^n) \to -\infty$ という条件は、
この置換によって損なわれないことに注意する。

\medskip\noindent
（無限）列をなす初等置換により $K^\bullet$ を改良する。
{\it 初等置換}とは次のものである。$d(K^n) > 0$ を満たす添字
$n$ を選ぶ。$d(M) = 0$ となるような単射 $K^n \to M$ を選ぶ。
$M' = \Coker(K^n \to M \oplus K^{n + 1})$ と置く。複体の写像
$$
\xymatrix{
K^\bullet : \ar[d] &
K^{n - 1} \ar[d] \ar[r] &
K^n \ar[d] \ar[r] &
K^{n + 1} \ar[d] \ar[r] &
K^{n + 2} \ar[d] \\
(K')^\bullet : &
K^{n - 1} \ar[r] &
M \ar[r] &
M' \ar[r] &
K^{n + 2}
}
$$
を考える。$K^\bullet \to (K')^\bullet$ が擬同型であることは明らかである。
さらに $d((K')^n) = 0$ であることも明らかであり、
$$
d((K')^{n + 1}) \leq \max\{d(K^n) - 1, d(M \oplus K^{n + 1})\} \leq
\max\{d(K^n) - 1, d(K^{n + 1})\}
$$
である。他の値は変わらない。

\medskip\noindent
証明を完了するため、任意の整数 $m$ に対して複体
$\sigma_{\geq m}K^\bullet$ が有限回しか変更されないように、
初等置換を行う順序を注意深く選ぶ。そのため
$$
\xi(K^\bullet) = \max \{n + d(K^n) \mid d(K^n) > 0\}
$$
および
$$
I = \{n \in \mathbf{Z} \mid \xi(K^\bullet) = n + d(K^n)
\text{ and }
 d(K^n) > 0\}
$$
と置く。$n \to -\infty$ のとき $n + d(K^n)$ が
$-\infty$ に収束するという仮定と、$n >> 0$ のとき
$d(K^n) = 0$ であるという事実から、
$\xi(K^\bullet) < +\infty$ であり、$I$ は有限集合である。
上のような初等変換に対して
$\xi((K')^\bullet) \leq \xi(K^\bullet)$ であることは明らかである。
一回の初等変換は次数 $\leq \xi(K^\bullet) + 1$ において複体を変える。
% P02-E0076: frozen prose says “find finite sequence” and omits the article.
従って、$\xi(K^\bullet)$ を減少させる有限列の初等変換を
見つけられれば、目的が達せられる。しかし、最小の元
$n \in I$ から始めて初等変換を行えば、$I$ の大きさが減少するか、
$\min I$ が増加するかのいずれかである。$I$ の任意の元は
$\leq \xi(K^\bullet)$ なので、有限回の操作で目的が達せられる。
\end{proof}

\begin{lemma}
\label{derived-lemma-unbounded-right-derived}
$F : \mathcal{A} \to \mathcal{B}$ をアーベル圏の間の左完全関手とする。
次を仮定する。
\begin{enumerate}
\item $\mathcal{A}$ の任意の対象は、$F$ に関して右非輪状な対象の
部分対象である。
\item $R^nF = 0$ を満たす整数 $n \geq 0$ が存在する。
\end{enumerate}
このとき次が成り立つ。
\begin{enumerate}
\item $RF : D(\mathcal{A}) \to D(\mathcal{B})$ が存在する。
\item $F$ に関して右非輪状な対象からなる任意の複体は $RF$ を計算する。
\item 任意の複体から、$F$ に関して右非輪状な対象からなる複体への
擬同型が存在する。
\item $E \in D(\mathcal{A})$ に対して次が成り立つ。
\begin{enumerate}
% P02-E0077: the frozen map omits a closing parenthesis after RF(tau_{<=a}E).
\item $i \leq a$ ならば
$H^i(RF(\tau_{\leq a}E) \to H^i(RF(E))$ は同型である。
\item $i \geq b$ ならば
$H^i(RF(E)) \to H^i(RF(\tau_{\geq b - n + 1}E))$ は同型である。
\item ある $-\infty \leq a \leq b \leq \infty$ に対して
$i \not \in [a, b]$ ならば $H^i(E) = 0$ であるならば、
$i \not \in [a, b + n - 1]$ ならば $H^i(RF(E)) = 0$ である。
\end{enumerate}
\end{enumerate}
\end{lemma}

\begin{proof}
第一の仮定から
$RF : D^+(\mathcal{A}) \to D^+(\mathcal{B})$ が存在することに注意する。
命題 \ref{derived-proposition-enough-acyclics} を参照されたい。
$A$ を $\mathcal{A}$ の対象とする。$A'$ が非輪状となるような単射
$A \to A'$ を選ぶ。長完全コホモロジー列により
$R^{n + 1}F(A) = R^nF(A'/A) = 0$ であることが分かる。
従って $R^{n + 1}F = 0$ と結論する。このように帰納法を続けると、
すべての $m \geq n$ に対して $R^mF = 0$ となる。

\medskip\noindent
関数 $d : \Ob(\mathcal{A}) \to \{0, 1, 2, \ldots \}$ を
% P02-E0078: frozen d(A) is malformed as max {0} union {i | ...}, rather than max of the union.
$d(A) = \max \{0\} \cup \{i \mid R^iF(A) \not = 0\}$
で定め、補題 \ref{derived-lemma-replace-resolution} を適用する。
補題 \ref{derived-lemma-replace-resolution} の第一の仮定は、ここでの仮定 (1) である。
補題 \ref{derived-lemma-replace-resolution} の第二の仮定は
$RF(A \oplus B) = RF(A) \oplus RF(B)$ から従う。
補題 \ref{derived-lemma-replace-resolution} の第三の仮定は
長完全コホモロジー列から従う。
従って任意の複体 $K^\bullet$ に対して、
$F$ に関して右非輪状な対象からなる複体への擬同型
$K^\bullet \to L^\bullet$ が存在する。これで主張 (3) が示された。

\medskip\noindent
$L^\bullet \to M^\bullet$ が、$F$ に関して右非輪状な対象からなる
複体の間の擬同型ならば、
$F(L^\bullet) \to F(M^\bullet)$ は擬同型であると主張する。
この主張を証明すれば、補題 \ref{derived-lemma-find-existence-computes} により
補題の主張 (1) と (2) が得られる。主張を証明するため、整数
$i \in \mathbf{Z}$ を選ぶ。完全三角
$$
\sigma_{\geq i - n - 1}L^\bullet \to
\sigma_{\geq i - n - 1}M^\bullet \to Q^\bullet,
$$
を考える。すなわち、$Q^\bullet$ を第一の写像の錐とする。
$Q^\bullet$ は下に有界であり、その $H^j(Q^\bullet)$ は、
$j = i - n - 1$ または $j = i - n - 2$ の場合を除いて零である。
$Q^\bullet$ に $RF$ を適用できる。
補題 \ref{derived-lemma-two-ss-complex-functor} の第二スペクトル系列と、
仮定 (2) のコホモロジー消滅を用いると、
$j \in \{i - n - 2, \ldots, i - 1\}$ の場合を除いて
$H^j(RF(Q^\bullet))$ が零であると結論できる。従って
$RF(\sigma_{\geq i - n - 1}L^\bullet) \to RF(\sigma_{\geq i - n - 1}M^\bullet)$
は次数 $\geq i$ のコホモロジー対象上で同型を誘導する。
命題 \ref{derived-proposition-enough-acyclics} により
$RF(\sigma_{\geq i - n - 1}L^\bullet) = \sigma_{\geq i - n - 1}F(L^\bullet)$
および
$RF(\sigma_{\geq i - n - 1}M^\bullet) = \sigma_{\geq i - n - 1}F(M^\bullet)$
である。従って所望の通り
$F(L^\bullet) \to F(M^\bullet)$ は次数 $i$ で同型である。

\medskip\noindent
(4)(a) は補題 \ref{derived-lemma-negative-vanishing} から従う。

\medskip\noindent
(4)(b) について、$E$ を $F$ に関して右非輪状な対象からなる複体
$L^\bullet$ で表す。(2) により $RF(E)$ は複体
$F(L^\bullet)$ で表され、
% P02-E0079: frozen proof introduces sigma_{>=c} although c is not defined in this argument.
$RF(\sigma_{\geq c}L^\bullet)$ は
$\sigma_{\geq c}F(L^\bullet)$ で表される。Remark
\ref{derived-remark-truncation-distinguished-triangle} の完全三角
$$
H^{b - n}(L^\bullet)[n - b] \to
\tau_{\geq b - n}L^\bullet \to
\tau_{\geq b - n + 1}L^\bullet
$$
を考える。上で確立した消滅性により、$i \geq b$ に対して
$H^i(RF(\tau_{\geq b - n}L^\bullet))$ は
$H^i(RF(\tau_{\geq b - n + 1}L^\bullet))$ と一致する。
複体の短完全列
$$
0 \to
\Im(L^{b - n - 1} \to L^{b - n})[n - b] \to
\sigma_{\geq b - n}L^\bullet \to
\tau_{\geq b - n}L^\bullet \to 0
$$
を考える。これに付随する完全三角
（Section \ref{derived-section-canonical-delta-functor} を参照）と先ほどと同じ
消滅性を用いると、$i \geq b$ に対して
$H^i(RF(\tau_{\geq b - n}L^\bullet))$ は
$H^i(RF(\sigma_{\geq b - n}L^\bullet))$ と一致すると結論できる。
写像 $RF(\sigma_{\geq b - n}L^\bullet) \to RF(L^\bullet)$ は
$\sigma_{\geq b - n}F(L^\bullet) \to F(L^\bullet)$ で表されるので、
これはさらに、所望の通り $i \geq b$ に対して
$H^i(RF(L^\bullet))$ と一致すると結論できる。

\medskip\noindent
(4)(c) を証明する。$E$ に関する仮定の下で
$\tau_{\leq a - 1}E = 0$ であるから、(4)(a) により
$i \leq a - 1$ に対して $H^i(RF(E))$ は零である。
同様に $\tau_{\geq b + 1}E = 0$ であるから、(4)(b) により
$i \geq b + n$ に対して $H^i(RF(E))$ は零である。
\end{proof}

\begin{lemma}
\label{derived-lemma-unbounded-left-derived}
$F : \mathcal{A} \to \mathcal{B}$ をアーベル圏の間の右完全関手とする。
次を仮定する。
\begin{enumerate}
\item $\mathcal{A}$ の任意の対象は、$F$ に関して左非輪状な対象の商である。
\item $L^nF = 0$ を満たす整数 $n \geq 0$ が存在する。
\end{enumerate}
このとき次が成り立つ。
\begin{enumerate}
\item $LF : D(\mathcal{A}) \to D(\mathcal{B})$ が存在する。
\item $F$ に関して左非輪状な対象からなる任意の複体は $LF$ を計算する。
\item 任意の複体へ、$F$ に関して左非輪状な対象からなる複体からの
擬同型が存在する。
\item $E \in D(\mathcal{A})$ に対して次が成り立つ。
\begin{enumerate}
% P02-E0080: the frozen map omits a closing parenthesis after LF(tau_{<=a+n-1}E).
\item $i \leq a$ ならば
$H^i(LF(\tau_{\leq a + n - 1}E) \to H^i(LF(E))$ は同型である。
\item $i \geq b$ ならば
$H^i(LF(E)) \to H^i(LF(\tau_{\geq b}E))$ は同型である。
\item ある $-\infty \leq a \leq b \leq \infty$ に対して
$i \not \in [a, b]$ ならば $H^i(E) = 0$ であるならば、
$i \not \in [a - n + 1, b]$ ならば $H^i(LF(E)) = 0$ である。
\end{enumerate}
\end{enumerate}
\end{lemma}

\begin{proof}
これは補題 \ref{derived-lemma-unbounded-right-derived} の双対である。
\end{proof}

\section{導来余極限}
\label{derived-section-derived-colimit}

\noindent
三角圏には導来余極限という概念がある。

\begin{definition}
\label{derived-definition-derived-colimit}
$\mathcal{D}$ を三角圏とする。
$(K_n, f_n)$ を $\mathcal{D}$ の対象の系とする。
直和 $\bigoplus K_n$ が存在し、完全三角
$$
\bigoplus K_n \to \bigoplus K_n \to K \to \bigoplus K_n[1]
$$
であって、写像 $\bigoplus K_n \to \bigoplus K_n$ が次数 $n$ で
$1 - f_n$ によって与えられるものが存在するとき、対象 $K$ を
系 $(K_n)$ の{\it 導来余極限}または{\it ホモトピー余極限}という。
この場合、これを記法 $K = \text{hocolim} K_n$ で表すことがある。
\end{definition}

\noindent
TR3 により、導来余極限は、存在するならば（非一意的な）同型を除いて
一意である。さらに TR1 により、$\bigoplus K_n$ が存在しさえすれば
$K_n$ の導来余極限は存在する。アーベル群の圏の導来圏
$D(\textit{Ab})$ は、すべてのホモトピー余極限が存在する三角圏の例である。

\medskip\noindent
非一意性のため、導来余極限を正確に指定することは難しい。
More on Algebra, Lemma \ref{more-algebra-lemma-map-from-hocolim} には、
通常の $\Hom$ を用いてホモトピー余極限から出る $\Hom$ を記述する完全列
$$
0 \to R^1\lim \Hom(K_n, L[-1]) \to \Hom(\text{hocolim} K_n, L)
\to \lim \Hom(K_n, L) \to 0
$$
がある。

\begin{remark}
\label{derived-remark-uniqueness-derived-colimit}
$\mathcal{D}$ を三角圏とし、
$(K_n, f_n)$ を $\mathcal{D}$ の対象の系とする。
導来余極限を、$\mathcal{D}$ の対象 $K$ と射
$i_n : K_n \to K$ の組であって、$i_{n + 1} \circ f_n = i_n$ を満たし、
さらに
$$
\bigoplus K_n \xrightarrow{1 - f_n} \bigoplus K_n \xrightarrow{i_n}
K \xrightarrow{c} \bigoplus K_n[1]
$$
が完全三角となるような射 $c : K \to \bigoplus K_n[1]$ が存在するものと
考えてよい。$(K', i'_n, c')$ を第二の導来余極限とすると、
$\varphi \circ i_n = i'_n$ と $c' \circ \varphi = c$ を満たす同型
$\varphi : K \to K'$ が存在する。$\varphi$ の存在は TR3 により、
$\varphi$ が同型であることは補題
\ref{derived-lemma-third-isomorphism-triangle} による。
\end{remark}

\begin{remark}
\label{derived-remark-functoriality-derived-colimit}
$\mathcal{D}$ を三角圏とする。
$(a_n) : (K_n, f_n) \to (L_n, g_n)$ を $\mathcal{D}$ の対象の系の射とする。
$(K, i_n, c)$ を第一の系の導来余極限、$(L, j_n, d)$ を第二の系の
導来余極限とし、記法は Remark
\ref{derived-remark-uniqueness-derived-colimit} のものとする。
このとき、$a \circ i_n = j_n$ および
$d \circ a = (a_n[1]) \circ c$ を満たす射 $a : K \to L$ が存在する。
これは定義に現れる完全三角に TR3 を適用すれば従う。
\end{remark}

\begin{lemma}
\label{derived-lemma-hocolim-subsequence}
$\mathcal{D}$ を三角圏とし、
$(K_n, f_n)$ を $\mathcal{D}$ の対象の系とする。
$n_1 < n_2 < n_3 < \ldots$ を整数の列とする。
$\bigoplus K_n$ と $\bigoplus K_{n_i}$ が存在すると仮定する。
このとき同型
$\text{hocolim} K_{n_i} \to \text{hocolim} K_n$
であって、すべての $i$ について図式
$$
\xymatrix{
K_{n_i} \ar[r] \ar[d]_{\text{id}} & \text{hocolim} K_{n_i} \ar[d] \\
K_{n_i} \ar[r] & \text{hocolim} K_n
}
$$
が可換となるものが存在する。
\end{lemma}

\begin{proof}
$g_i : K_{n_i} \to K_{n_{i + 1}}$ を合成
$f_{n_{i + 1} - 1} \circ \ldots \circ f_{n_i}$ とする。
次のように可換図式を構成する。
$$
\vcenter{
\xymatrix{
\bigoplus\nolimits_i K_{n_i} \ar[r]_{1 - g_i} \ar[d]_b &
\bigoplus\nolimits_i K_{n_i} \ar[d]^a \\
\bigoplus\nolimits_n K_n \ar[r]^{1 - f_n} &
\bigoplus\nolimits_n K_n
}
}
\quad\text{and}\quad
\vcenter{
\xymatrix{
\bigoplus\nolimits_n K_n \ar[r]_{1 - f_n} \ar[d]_d &
\bigoplus\nolimits_n K_n \ar[d]^c \\
\bigoplus\nolimits_i K_{n_i} \ar[r]^{1 - g_i} &
\bigoplus\nolimits_i K_{n_i}
}
}
$$
$a_i = a|_{K_{n_i}}$ を $K_{n_i}$ から直和への包含射とする。
言い換えれば、$a$ は自然な包含射である。
$b_i = b|_{K_{n_i}}$ を写像
$$
K_{n_i}
\xrightarrow{1,\ f_{n_i},\ f_{n_i + 1} \circ f_{n_i},
\ \ldots,\ f_{n_{i + 1} - 2} \circ \ldots \circ f_{n_i}}
K_{n_i} \oplus K_{n_i + 1} \oplus \ldots \oplus K_{n_{i + 1} - 1}
$$
とする。$n_{i - 1} < j \leq n_i$ ならば、
$c_j = c|_{K_j}$ を写像
$$
K_j \xrightarrow{f_{n_i - 1} \circ \ldots \circ f_j} K_{n_i}
$$
とする。$d_j = d|_{K_j}$ は、どの $i$ に対しても
$j \not = n_i$ ならば零とし、$d_{n_i}$ は $K_{n_i}$ から直和への
自然な包含射とする。言い換えれば、$d$ は自然な射影である。
TR3 により、これらの図式は射
$$
\varphi : \text{hocolim} K_{n_i} \to \text{hocolim} K_n
\quad\text{and}\quad
\psi : \text{hocolim} K_n \to \text{hocolim} K_{n_i}
$$
を定める。$c \circ a$ と $d \circ b$ は恒等写像なので、補題
\ref{derived-lemma-third-isomorphism-triangle} により
$\varphi \circ \psi$ は同型であることが分かる。
逆の順序では、射 $a \circ c$ と $b \circ d$ を得る。規則として、
$n_{i - 1} < j < n_i$ に対して
$$
h_j : K_j
\xrightarrow{1,\ f_j,\ f_{j + 1} \circ f_j,
\ \ldots,\ f_{n_i - 1} \circ \ldots \circ f_j}
K_j \oplus \ldots \oplus K_{n_i}
$$
と置くことにより与えられる射
$h = (h_j) : \bigoplus K_n \to \bigoplus K_n$ を考える。
すると
$(1 - f) \circ h = \text{id} - a \circ c$ および
$h \circ (1 - f) = \text{id} - b \circ d$ であることを読者は確認できる。
これは、補題 \ref{derived-lemma-third-map-square-zero} により
$\text{id} - \psi \circ \varphi$ の平方が零であることを意味する
（短い議論を省略する）。言い換えれば、
$\psi \circ \varphi$ と恒等写像との差は冪零自己準同型なので、これは同型である。
従って所望の通り $\varphi$ と $\psi$ は同型である。
\end{proof}

\begin{lemma}
\label{derived-lemma-direct-sums}
$\mathcal{A}$ をアーベル圏とする。$\mathcal{A}$ が完全な可算直和をもつならば、
$D(\mathcal{A})$ は可算直和をもつ。実際、可算な添字集合 $I$ で添字付けられた
複体の集まり $K_i^\bullet$ が与えられたとき、項ごとの直和
$\bigoplus K_i^\bullet$ は $D(\mathcal{A})$ における
$K_i^\bullet$ の直和である。
\end{lemma}

\begin{proof}
$L^\bullet$ を複体とする。$D(\mathcal{A})$ における写像
$\alpha_i : K_i^\bullet \to L^\bullet$ が与えられたと仮定する。
これは、複体の擬同型 $s_i : M_i^\bullet \to K_i^\bullet$ と
複体の写像 $f_i : M_i^\bullet \to L^\bullet$ であって
$\alpha_i = f_is_i^{-1}$ となるものが存在することを意味する。
仮定により、複体の写像
$$
s : \bigoplus M_i^\bullet \longrightarrow \bigoplus K_i^\bullet
$$
は擬同型である。従って $f = \bigoplus f_i$ と置けば、
$\alpha = fs^{-1}$ は $D(\mathcal{A})$ における写像であり、
標準包含射 $K_i^\bullet \to \bigoplus K_i^\bullet$ との合成は
$\alpha_i$ である。$\alpha$ の一意性の確認は省略する。
\end{proof}

\begin{lemma}
\label{derived-lemma-compute-colimit}
$\mathcal{A}$ をアーベル圏とする。$\mathbf{N}$ 上の余極限が存在し、
完全であると仮定する。このとき可算直和が存在し、完全である。
さらに、$(A_n, f_n)$ が $\mathbf{N}$ 上の系ならば、短完全列
$$
0 \to \bigoplus A_n \to \bigoplus A_n \to \colim A_n \to 0
$$
であって、最初の写像が次数 $n$ で $1 - f_n$ によって与えられるものが存在する。
\end{lemma}

\begin{proof}
第一の主張は
$\bigoplus A_n = \colim (A_1 \oplus \ldots \oplus A_n)$
から従う。第二の主張については、各 $n$ に対して短完全列
$$
0 \to
A_1 \oplus \ldots \oplus A_{n - 1} \to
A_1 \oplus \ldots \oplus A_n \to A_n \to 0
$$
があることに注意する。ここで第一の写像は写像 $1 - f_i$ によって
与えられ、第二の写像は遷移写像の和である。
余極限を取れば補題の列を得る。
\end{proof}

\begin{lemma}
\label{derived-lemma-colim-hocolim}
$\mathcal{A}$ をアーベル圏とする。$L_n^\bullet$ を
$\mathcal{A}$ の複体の系とする。$\mathcal{A}$ において
$\mathbf{N}$ 上の余極限が存在し、完全であると仮定する。
このとき項ごとの余極限 $L^\bullet = \colim L_n^\bullet$ は、
$D(\mathcal{A})$ におけるこの系のホモトピー余極限である。
\end{lemma}

\begin{proof}
補題 \ref{derived-lemma-compute-colimit} により複体の完全列
$$
0 \to \bigoplus L_n^\bullet \to \bigoplus L_n^\bullet \to L^\bullet \to 0
$$
がある。補題 \ref{derived-lemma-direct-sums} により、これらの直和は
$D(\mathcal{A})$ における直和である。従って結果は、Definition
\ref{derived-definition-derived-colimit} における導来余極限の定義と、
複体の短完全列が完全三角を与えるという事実
（補題 \ref{derived-lemma-derived-canonical-delta-functor}）から従う。
\end{proof}

\begin{lemma}
\label{derived-lemma-cohomology-of-hocolim}
$\mathcal{D}$ を可算直和をもつ三角圏とする。
$\mathcal{A}$ を $\mathbf{N}$ 上の完全な余極限をもつアーベル圏とする。
$H : \mathcal{D} \to \mathcal{A}$ を可算直和と可換するホモロジー的関手とする。
このとき $\mathcal{D}$ の任意の対象の系に対して
$H(\text{hocolim} K_n) = \colim H(K_n)$ である。
\end{lemma}

\begin{proof}
$K = \text{hocolim} K_n$ と書く。定義に現れる完全三角に $H$ を適用して
$$
\bigoplus H(K_n) \to \bigoplus H(K_n)
\to H(K) \to
\bigoplus H(K_n[1]) \to \bigoplus H(K_n[1])
$$
を得る。ここで第一の写像は $1 - H(f_n)$ によって与えられ、
最後の写像は $1 - H(f_n[1])$ によって与えられる。
補題 \ref{derived-lemma-compute-colimit} を適用すれば、この補題が示される。
\end{proof}

\noindent
次の補題は、コンパクト対象（後で定義する）から出る写像を取ることが
導来余極限と可換することを述べる。

\begin{lemma}
\label{derived-lemma-commutes-with-countable-sums}
$\mathcal{D}$ を可算直和をもつ三角圏とする。
$K \in \mathcal{D}$ を、任意の可算な対象の集合
$E_n \in \mathcal{D}$ に対して標準写像
$$
\bigoplus \Hom_\mathcal{D}(K, E_n)
\longrightarrow
\Hom_\mathcal{D}(K, \bigoplus E_n)
$$
が全単射となる対象とする。このとき、$\mathcal{D}$ の任意の
$\mathbf{N}$ 上の系 $L_n$ で、導来余極限
$L = \text{hocolim} L_n$ が存在するものに対して
$$
\colim \Hom_\mathcal{D}(K, L_n) \longrightarrow \Hom_\mathcal{D}(K, L)
$$
は全単射である。
\end{lemma}

\begin{proof}
補題 \ref{derived-lemma-representable-homological} により、これは補題
\ref{derived-lemma-cohomology-of-hocolim} の特別な場合である。
\end{proof}

\section{導来極限}
\label{derived-section-derived-limit}

\noindent
三角圏には導来極限という概念がある。

\begin{definition}
\label{derived-definition-derived-limit}
$\mathcal{D}$ を三角圏とする。
$(K_n, f_n)$ を $\mathcal{D}$ の対象の逆系とする。
積 $\prod K_n$ が存在し、完全三角
$$
K \to \prod K_n \to \prod K_n \to K[1]
$$
であって、写像 $\prod K_n \to \prod K_n$ が
$(k_n) \mapsto (k_n - f_{n + 1}(k_{n + 1}))$ によって与えられるものが
存在するとき、対象 $K$ を系 $(K_n)$ の{\it 導来極限}または
{\it ホモトピー極限}という。この場合、記法 $K = R\lim K_n$ で表すことがある。
\end{definition}

\noindent
TR3 により、導来極限は、存在するならば（非一意的な）同型を除いて一意である。
さらに TR1 により、$\prod K_n$ が存在しさえすれば導来極限
$R\lim K_n$ は存在する。アーベル群の圏の導来圏
$D(\textit{Ab})$ は、すべての導来極限が存在する三角圏の例である。

\medskip\noindent
非一意性のため、導来極限を正確に指定することは難しい。
More on Algebra, Lemma \ref{more-algebra-lemma-map-into-Rlim} には、
通常の $\Hom$ を用いて導来極限への $\Hom$ を記述する完全列
$$
0 \to R^1\lim \Hom(L, K_n[-1]) \to \Hom(L, R\lim K_n)
\to \lim \Hom(L, K_n) \to 0
$$
がある。

\begin{lemma}
\label{derived-lemma-products}
$\mathcal{A}$ を完全な可算積をもつアーベル圏とする。このとき
\begin{enumerate}
\item $D(\mathcal{A})$ は可算積をもつ。
\item $D(\mathcal{A})$ における可算積 $\prod K_i$ は、
$K_i$ を表現する任意の複体を項ごとに積に取ることで得られる。
\item $H^p(\prod K_i) = \prod H^p(K_i)$ である。
\end{enumerate}
\end{lemma}

\begin{proof}
$K_i^\bullet$ を $D(\mathcal{A})$ において $K_i$ を表現する複体とする。
$L^\bullet$ を複体とする。$D(\mathcal{A})$ における写像
$\alpha_i : L^\bullet \to K_i^\bullet$ が与えられたと仮定する。
これは、複体の擬同型 $s_i : K_i^\bullet \to M_i^\bullet$ と
複体の写像 $f_i : L^\bullet \to M_i^\bullet$ であって
$\alpha_i = s_i^{-1}f_i$ となるものが存在することを意味する。
仮定により、複体の写像
$$
s : \prod K_i^\bullet \longrightarrow \prod M_i^\bullet
$$
は擬同型である。従って $f = \prod f_i$ と置けば、
$\alpha = s^{-1}f$ は $D(\mathcal{A})$ における写像であり、
射影 $\prod K_i^\bullet \to K_i^\bullet$ との合成は $\alpha_i$ である。
$\alpha$ の一意性の確認は省略する。
\end{proof}

\noindent
補題 \ref{derived-lemma-compute-colimit}、
\ref{derived-lemma-colim-hocolim}、および
\ref{derived-lemma-commutes-with-countable-sums} の双対をここで述べ、証明すべきである。
しかし現在のところ、これらの補題の応用を何も知らない。

\begin{lemma}
\label{derived-lemma-inverse-limit-bounded-below}
$\mathcal{A}$ を可算積と十分多くの入射対象をもつアーベル圏とする。
$(K_n)$ を $D^+(\mathcal{A})$ の逆系とする。このとき
$R\lim K_n$ は存在する。
\end{lemma}

\begin{proof}
$\prod K_n$ が $D(\mathcal{A})$ において存在することを示せば十分である。
各 $n$ に対して、$K_n$ を入射対象からなる下に有界な複体
$I_n^\bullet$ で表現できる
（補題 \ref{derived-lemma-injective-resolutions-exist}）。
このとき補題 \ref{derived-lemma-product-K-injective} により、
$\prod K_n$ は $\prod I_n^\bullet$ で表現される。
\end{proof}

\begin{remark}
\label{derived-remark-functoriality-derived-limit}
$\mathcal{D}$ を三角圏とする。$(K_n)$ と $(L_n)$ を
$\mathcal{D}$ の対象の逆系とし、その導来極限を $K$ と $L$ とする。
pro-対象の射 $a : (K_n) \to (L_n)$ があると仮定する。Categories, Example
\ref{categories-example-pro-morphism-inverse-systems} を参照されたい。
これは写像 $a_n : K_{m(n)} \to L_n$ が与えられることを意味する。
$m(1) < m(2) < m(3) < \ldots$ であり、写像 $a_n$ が両系の遷移写像と
両立すると仮定してよい。このとき、規則
$a'((k_n)) = (a_n(k_{m(n)}))$ および
$$
a''((k_n)) =
(a_n(k_{m(n)} + f(k_{m(n) + 1}) + \ldots + f(k_{m(n + 1) - 1})))
$$
によって定義される写像
$$
a', a'' : \prod K_n \longrightarrow \prod L_n
$$
を考えられる。
% P02-E0081: frozen source misspells “occurrence” as “occurence”.
ここで現れる各 $f$ は、逆系 $(K_n)$ の適切な遷移写像を表す。
すると実線からなる図式
$$
\xymatrix{
K \ar[r] \ar@{..>}[d] & \prod K_n \ar[r] \ar[d]^{a'} &
\prod K_n \ar[d]^{a''} \\
L \ar[r] & \prod L_n \ar[r] & \prod L_n
}
$$
は可換であり、TR3 により点線の射を得て、完全三角の射が生じる。
写像 $K \to L$ は一意ではないことを注意しておく。後で、$a$ が
pro-同型ならば $K \to L$ が同型であることを見る。More on Algebra, Lemma
\ref{more-algebra-lemma-pro-isomorphism-Rlim} を参照されたい。
\end{remark}

\begin{remark}
\label{derived-remark-map-into-derived-limit-truncations}
$\mathcal{A}$ をアーベル圏とし、$K^\bullet$ を
$\mathcal{A}$ の複体とする。このとき
% P02-E0082: frozen source says “is an inverse system of complexes and which in particular determines”.
$\tau_{\geq -n}K^\bullet$ は複体の逆系であり、特に
$D(\mathcal{A})$ の逆系を定める。
$R\lim \tau_{\geq -n}K^\bullet$ が存在すると仮定する。
標準写像 $c_n : K^\bullet \to \tau_{\geq -n}K^\bullet$ は
この逆系の遷移写像と両立する。Definition
\ref{derived-definition-derived-limit} の定義に現れる完全三角と
補題 \ref{derived-lemma-representable-homological} により、$D(\mathcal{A})$ の射
$$
c : K^\bullet \longrightarrow R\lim \tau_{\geq -n}K^\bullet
$$
であって、$c$ と射影
$R\lim \tau_{\geq -n}K^\bullet \to \tau_{\geq -m}K^\bullet$
との合成が $c_m$ に等しいものが存在すると結論できる。
射 $c$ は一意でないかもしれないが、$c$ が同型であるかどうかは、
$c$ の選択（およびホモトピー極限の選択）によらないと主張する。
実際、$i \in \mathbf{Z}$ および $m > -i$ に対して、合成
$$
H^i(K^\bullet) \xrightarrow{H^i(c)} H^i(R\lim \tau_{\geq -n}K^\bullet)
\to H^i(\tau_{\geq -m}K^\bullet) = H^i(K^\bullet)
$$
は恒等写像である。従って $H^i(c)$ が同型であることと
第二の写像が同型であることは同値である。これは $c$ によらず、
ホモトピー極限の選択にもよらない（任意の二つの選択は同型だからである）。
\end{remark}

\begin{lemma}
\label{derived-lemma-difficulty-K-injectives}
$\mathcal{A}$ を可算積と十分多くの入射対象をもつアーベル圏とする。
$K^\bullet$ を複体とする。$I_n^\bullet$ を、補題
\ref{derived-lemma-special-inverse-system} により得られる入射対象からなる
下に有界な複体の逆系とする。このとき
$I^\bullet = \lim I_n^\bullet$ は存在し、K-入射的であり、
$D(\mathcal{A})$ において $R\lim \tau_{\geq -n}K^\bullet$ を表現する。
さらに次は同値である。
\begin{enumerate}
\item 写像 $K^\bullet \to I^\bullet$（証明を参照）は擬同型である。
\item Remark \ref{derived-remark-map-into-derived-limit-truncations} の写像
$K^\bullet \to R\lim \tau_{\geq -n}K^\bullet$ は
$D(\mathcal{A})$ における同型である。
\end{enumerate}
\end{lemma}

\begin{proof}
$R\lim \tau_{\geq -n}K^\bullet$ は補題
\ref{derived-lemma-inverse-limit-bounded-below} により存在するので、
補題の主張には意味がある。各複体 $I_n^\bullet$ は補題
\ref{derived-lemma-bounded-below-injectives-K-injective} により K-入射的である。
すべての $n \geq 1$ に対して直和分解
$I_{n + 1}^p = C_{n + 1}^p \oplus I_n^p$ を選び、$C_1^p = I_1^p$ と置く。
複体 $I^\bullet = \lim I_n^\bullet$ は存在する。実際、
$I^p = \prod_{n \geq 1} C_n^p$ と取ればよい。
$p \in \mathbf{Z}$ を固定する。$\mathcal{A}$ の対象の分裂短完全列
$$
0 \to I^p \to \prod I_n^p \to \prod I_n^p \to 0
$$
が存在すると主張する。ここで第一の写像は射影
$I^p \to I_n^p$ によって、第二の写像は
$(x_n) \mapsto (x_n - f^p_{n + 1}(x_{n + 1}))$ によって与えられ、
$f^p_n : I_n^p \to I_{n - 1}^p$ は遷移写像である。分裂は写像
$\prod I_n^p \to \prod C_n^p = I^p$ から得られる。
複体の項ごとに分裂する短完全列
$$
0 \to I^\bullet \to \prod I_n^\bullet \to \prod I_n^\bullet \to 0
$$
を得る。従って $K(\mathcal{A})$ と $D(\mathcal{A})$ において
対応する完全三角を得る。補題 \ref{derived-lemma-product-K-injective} により、
これらの積は K-入射的であり、$D(\mathcal{A})$ における対応する積を表現する。
従って $I^\bullet$ は $R\lim I_n^\bullet$ を表現する
（Definition \ref{derived-definition-derived-limit}）。
導来極限は導来圏のレベルで定義されるので
$R\lim I_n^\bullet \cong R\lim \tau_{\geq -n}K^\bullet$ であり、
$I^\bullet$ が $R\lim \tau_{\geq -n}K^\bullet$ を表現することが分かる。
さらに複体 $I^\bullet$ は補題
\ref{derived-lemma-triangle-K-injective} により K-入射的である。
補題 \ref{derived-lemma-special-inverse-system} の可換図式と、
$n \gg 0$ のとき $K^i = (\tau_{\geq -n}K^\bullet)^i$ であることにより、
図式
$$
\xymatrix{
K^\bullet \ar[r] \ar[d]_\gamma & \tau_{\geq -n} K^\bullet \ar[d] \\
I^\bullet \ar[r] & I_n^\bullet
}
$$
を可換にする一意な写像 $\gamma : K^\bullet \to I^\bullet$ を得る。
これより $\gamma$ は複体の写像であり、$D(\mathcal{A})$ において
Remark \ref{derived-remark-map-into-derived-limit-truncations} の写像
$c : K^\bullet \to R\lim \tau_{\geq -n}K^\bullet$ を表現する。
言い換えれば、図式
$$
\xymatrix{
K^\bullet \ar[r]_-c \ar[d]_\gamma & R\lim \tau_{\geq -n} K^\bullet
\ar[d]^{\cong} \\
I^\bullet \ar[r]^-{\cong} & R\lim I_n^\bullet
}
$$
は $D(\mathcal{A})$ において可換である。補題が従う。
\end{proof}

\begin{lemma}
\label{derived-lemma-enough-K-injectives-Ab4-star}
$\mathcal{A}$ を十分多くの入射対象と完全な可算積をもつアーベル圏とする。
このとき任意の複体から K-入射的複体への擬同型が存在する。
\end{lemma}

\begin{proof}
補題 \ref{derived-lemma-difficulty-K-injectives} により、$D(\mathcal{A})$ の
すべての $K$ に対して $K \to R\lim\tau_{\geq -n}K$ が同型であることを
示せば十分である。定義に現れる完全三角
$$
R\lim\tau_{\geq -n}K \to
\prod \tau_{\geq -n}K \to
\prod \tau_{\geq -n}K \to
(R\lim\tau_{\geq -n}K)[1]
$$
を考える。補題 \ref{derived-lemma-products} により
$$
H^p(\prod \tau_{\geq -n}K) = \prod\nolimits_{p \geq -n} H^p(K)
$$
である。表示した完全三角の長完全コホモロジー列から直ちに
$H^p(R\lim \tau_{\geq -n}K) = H^p(K)$ が従う。
\end{proof}

\section{充満部分圏上の演算}
\label{derived-section-operate-on-full}

\noindent
$\mathcal{T}$ を三角圏とする。$\mathcal{T}$ の充満部分圏を
$\Ob(\mathcal{T})$ の部分集合と同一視する。充満部分圏
$\mathcal{A}, \mathcal{B}, \ldots$ が与えられたとき、次のように定める。
\begin{enumerate}
\item $-\infty \leq a \leq b \leq \infty$ に対して
$\mathcal{A}[a, b]$ を、$i \in [a, b] \cap \mathbf{Z}$ および
$A \in \Ob(\mathcal{A})$ を満たす $A[-i]$ の形のすべての対象からなる
$\mathcal{T}$ の充満部分圏とする（負号に注意せよ）。
\item $smd(\mathcal{A})$ を、$\mathcal{A}$ の対象の直和因子と
同型なすべての対象からなる $\mathcal{T}$ の充満部分圏とする。
\item $add(\mathcal{A})$ を、$\mathcal{A}$ の対象の有限直和と
同型なすべての対象からなる $\mathcal{T}$ の充満部分圏とする。
\item $\mathcal{A} \star \mathcal{B}$ を、
$A \in \Ob(\mathcal{A})$ と $B \in \Ob(\mathcal{B})$ に対する完全三角
$A \to X \to B$ に組み込める $\mathcal{T}$ のすべての対象 $X$ からなる
$\mathcal{T}$ の充満部分圏とする。
\item $\mathcal{A}^{\star n} = \mathcal{A} \star \ldots \star \mathcal{A}$
とし、因子は $n \geq 1$ 個とする（下で $\star$ が結合的であることを見る）。
\item $smd(add(\mathcal{A})^{\star n}) =
smd(add(\mathcal{A}) \star \ldots \star add(\mathcal{A}))$
とし、因子は $n \geq 1$ 個とする。
\end{enumerate}
$E$ が $\mathcal{T}$ の対象ならば、$E$ を、唯一の対象が $E$ である
$\mathcal{T}$ の充満部分圏とみなすこともある。すると
$add(E[-1, 2])$ などを考えられる。この記法は一般に受け入れられている
わけではないことを注意しておく。

\begin{lemma}
\label{derived-lemma-associativity-star}
$\mathcal{T}$ を三角圏とする。充満部分圏
$\mathcal{A}$, $\mathcal{B}$, $\mathcal{C}$ が与えられたとき
$(\mathcal{A} \star \mathcal{B}) \star \mathcal{C} =
\mathcal{A} \star (\mathcal{B} \star \mathcal{C})$ である。
\end{lemma}

\begin{proof}
完全三角 $A \to X \to B$ と $X \to Y \to C$ があるならば、
公理 TR4 により完全三角 $A \to Y \to Z$ と $B \to Z \to C$ がある。
\end{proof}

\begin{lemma}
\label{derived-lemma-smd-star}
$\mathcal{T}$ を三角圏とする。充満部分圏 $\mathcal{A}$, $\mathcal{B}$
が与えられたとき
$smd(\mathcal{A}) \star smd(\mathcal{B}) \subset
smd(\mathcal{A} \star \mathcal{B})$ および
$smd(smd(\mathcal{A}) \star smd(\mathcal{B})) =
smd(\mathcal{A} \star \mathcal{B})$ である。
\end{lemma}

\begin{proof}
$A_1 \oplus A_2 \in \Ob(\mathcal{A})$ および
$B_1 \oplus B_2 \in \Ob(\mathcal{B})$ を満たす完全三角
$A_1 \to X \to B_1$ があると仮定する。このとき完全三角
$A_1 \oplus A_2 \to A_2 \oplus X \oplus B_2 \to B_1 \oplus B_2$
を得る。これにより $X$ が $smd(\mathcal{A} \star \mathcal{B})$
に属することが示される。これで包含関係が示された。
等式はこの包含関係から直ちに従う。
\end{proof}

\begin{lemma}
\label{derived-lemma-add-star}
$\mathcal{T}$ を三角圏とする。充満部分圏
$\mathcal{A}$, $\mathcal{B}$ が与えられたとき、充満部分圏
$add(\mathcal{A}) \star add(\mathcal{B})$ および
$smd(add(\mathcal{A}))$ は直和について閉じている。
\end{lemma}

\begin{proof}
実際、$A \to X \to B$ と $A' \to X' \to B'$ が完全三角で、
$A, A' \in add(\mathcal{A})$ および $B, B' \in add(\mathcal{B})$
ならば、$A \oplus A' \to X \oplus X' \to B \oplus B'$ は完全三角で、
$A \oplus A' \in add(\mathcal{A})$ および
$B \oplus B' \in add(\mathcal{B})$ である。
$smd(add(\mathcal{A}))$ に関する結果は明らかである。
\end{proof}

\begin{lemma}
\label{derived-lemma-cone-n}
$\mathcal{T}$ を三角圏とする。充満部分圏 $\mathcal{A}$ が与えられたとき、
$n \geq 1$ に対して上で定義した部分圏
$$
\mathcal{C}_n = smd(add(\mathcal{A})^{\star n}) =
smd(add(\mathcal{A}) \star \ldots \star add(\mathcal{A}))
$$
は、直和と直和因子について閉じた $\mathcal{T}$ の厳密充満部分圏であり、
すべての $n, m \geq 1$ に対して
$\mathcal{C}_{n + m} = smd(\mathcal{C}_n \star \mathcal{C}_m)$ である。
\end{lemma}

\begin{proof}
補題 \ref{derived-lemma-associativity-star}、\ref{derived-lemma-smd-star}、および
\ref{derived-lemma-add-star} から直ちに従う。
\end{proof}

\begin{remark}
\label{derived-remark-operations-functor}
$F : \mathcal{T} \to \mathcal{T}'$ を三角圏の完全関手とする。
$\mathcal{T}$ の充満部分圏 $\mathcal{A}$ が与えられたとき、
$F(\mathcal{A})$ で、$A \in \Ob(\mathcal{A})$ に対する対象
$F(A)$ すべてからなる $\mathcal{T}'$ の充満部分圏を表す。
% P02-E0083: frozen source has subject--verb disagreement in “whose objects consists”.
次が成り立つ。
$$
F(\mathcal{A}[a, b]) = F(\mathcal{A})[a, b]
$$
$$
F(smd(\mathcal{A})) \subset smd(F(\mathcal{A})),
$$
$$
F(add(\mathcal{A})) \subset add(F(\mathcal{A})),
$$
$$
F(\mathcal{A} \star \mathcal{B}) \subset F(\mathcal{A}) \star F(\mathcal{B}),
$$
$$
F(\mathcal{A}^{\star n}) \subset F(\mathcal{A})^{\star n}.
$$
明らかな確認は省略する。
\end{remark}

\begin{remark}
\label{derived-remark-operations-unions}
$\mathcal{T}$ を三角圏とする。$\mathcal{T}$ の充満部分圏
$\mathcal{A}_1 \subset \mathcal{A}_2 \subset \mathcal{A}_3 \subset \ldots$
および $\mathcal{B}$ が与えられたとき、次が成り立つ。
$$
\left(\bigcup \mathcal{A}_i\right)[a, b] = \bigcup \mathcal{A}_i[a, b]
$$
$$
smd\left(\bigcup \mathcal{A}_i\right) = \bigcup smd(\mathcal{A}_i),
$$
$$
add\left(\bigcup \mathcal{A}_i\right) = \bigcup add(\mathcal{A}_i),
$$
$$
\left(\bigcup \mathcal{A}_i\right) \star \mathcal{B} =
\bigcup \mathcal{A}_i \star \mathcal{B},
$$
$$
\mathcal{B} \star \left(\bigcup \mathcal{A}_i\right) =
\bigcup \mathcal{B} \star \mathcal{A}_i,
$$
$$
\left(\bigcup \mathcal{A}_i\right)^{\star n} =
\bigcup \mathcal{A}_i^{\star n}.
$$
明らかな確認は省略する。
\end{remark}

\begin{lemma}
\label{derived-lemma-in-cone-n}
$\mathcal{A}$ をアーベル圏とし、$\mathcal{D} = D(\mathcal{A})$ と置く。
$\mathcal{E} \subset \Ob(\mathcal{A})$ を部分集合とし、これを
$\Ob(\mathcal{D})$ の部分集合ともみなす。$K$ を $\mathcal{D}$ の対象とする。
\begin{enumerate}
\item $b \geq a$ とし、$i \not \in [a, b]$ ならば $H^i(K)$ は零であり、
$i \in [a, b]$ ならば $H^i(K) \in \mathcal{E}$ であると仮定する。
このとき $K$ は
$smd(add(\mathcal{E}[a, b])^{\star (b - a + 1)})$ に属する。
\item $b \geq a$ とし、$i \not \in [a, b]$ ならば $H^i(K)$ は零であり、
$i \in [a, b]$ ならば $H^i(K) \in smd(add(\mathcal{E}))$ であると仮定する。
このとき $K$ は
$smd(add(\mathcal{E}[a, b])^{\star (b - a + 1)})$ に属する。
\item $b \geq a$ とし、$K$ は、$i \not \in [a, b]$ ならば
$K^i = 0$、$i \in [a, b]$ ならば $K^i \in \mathcal{E}$ となる複体
$K^\bullet$ で表現できると仮定する。このとき $K$ は
$smd(add(\mathcal{E}[a, b])^{\star (b - a + 1)})$ に属する。
\item $b \geq a$ とし、$K$ は、$i \not \in [a, b]$ ならば
$K^i = 0$、$i \in [a, b]$ ならば $K^i \in smd(add(\mathcal{E}))$
となる複体 $K^\bullet$ で表現できると仮定する。このとき $K$ は
$smd(add(\mathcal{E}[a, b])^{\star (b - a + 1)})$ に属する。
\end{enumerate}
\end{lemma}

\begin{proof}
補題 \ref{derived-lemma-cone-n} を断りなく用いる。
(1) を直ちに導く (2) を証明する。$b - a$ に関する帰納法を用いる。
$b - a = 0$ ならば、$K$ は $\mathcal{D}$ において
% P02-E0086: frozen source uses H^i(K) although the one surviving degree is a.
$H^i(K)[-a]$ と同型であり、結果は直ちに従う。
$b - a > 0$ ならば、完全三角
$$
\tau_{\leq b - 1}K^\bullet \to K^\bullet \to K^b[-b]
$$
を考え、$b - a$ に関する帰納法で結論する。(3) と (4) の証明は省略する。
\end{proof}

\begin{lemma}
\label{derived-lemma-forward-cone-n}
$\mathcal{T}$ を三角圏とする。
$H : \mathcal{T} \to \mathcal{A}$ をアーベル圏 $\mathcal{A}$ への
ホモロジー的関手とする。$a \leq b$ とし、
$\mathcal{E} \subset \Ob(\mathcal{T})$ を、$E \in \mathcal{E}$ および
$i \not \in [a, b]$ ならば $H^i(E) = 0$ となる部分集合とする。
このとき $X \in smd(add(\mathcal{E}[-m, m])^{\star n})$ に対して、
$i \not \in [-m + na, m + nb]$ ならば $H^i(X) = 0$ である。
\end{lemma}

\begin{proof}
省略する。定義に関する楽しい演習である。
\end{proof}

\section{三角圏の生成元}
\label{derived-section-generators}

\noindent
本節では、三角圏の生成元に関するいくつかの異なる概念を簡潔に導入する。
用語は \cite{BvdB} に従う（ただし、彼らが「\'epaisse」と呼ぶものを
ここでは「飽和」と呼ぶ。Definition \ref{derived-definition-saturated} を参照。
また、$add(\mathcal{A})$ の定義も異なる）。

\medskip\noindent
$\mathcal{D}$ を三角圏とし、$E$ を $\mathcal{D}$ の対象とする。
$\langle E \rangle_1$ で、$E$ の平行移動の有限直和
$$
\bigoplus\nolimits_{i = 1, \ldots, r} E[n_i]
$$
の直和因子と同型な $\mathcal{D}$ の対象からなる
$\mathcal{D}$ の厳密充満部分圏を表す。Section
\ref{derived-section-operate-on-full} の記法では、明らかに
$$
\langle E \rangle_1 = smd(add(E[-\infty, \infty]))
$$
である。$n > 1$ に対して、$\langle E \rangle_n$ で、
完全三角
$$
A \to X \to B \to A[1]
$$
に組み込める対象 $X$ の直和因子と同型な $\mathcal{D}$ の対象からなる
$\mathcal{D}$ の充満部分圏を表す。ここで $A$ は
$\langle E \rangle_1$ の対象であり、$B$ は
$\langle E \rangle_{n - 1}$ の対象である。Section
\ref{derived-section-operate-on-full} の記法では
$$
\langle E \rangle_n = smd(\langle E \rangle_1 \star \langle E \rangle_{n - 1})
$$
である。圏 $\langle E \rangle_n$ はそれぞれ、$\mathcal{D}$ の
厳密充満加法的部分圏（補題 \ref{derived-lemma-add-star} による）であり、
平行移動と直和因子を取る操作で保たれる。しかし
$\langle E \rangle_n$ は「錐を取る操作」または「拡大を取る操作」で
閉じているとは限らず、従って三角部分圏であるとは限らない。
以下の補題で示すように、部分圏
$$
\langle E \rangle = \bigcup\nolimits_n \langle E \rangle_n
$$
についてはこれが成り立つ。

\begin{lemma}
\label{derived-lemma-generated-by-E-explicit}
$\mathcal{T}$ を三角圏とし、$E$ を $\mathcal{T}$ の対象とする。
$n \geq 1$ に対して
$$
\langle E \rangle_n =
smd(\langle E \rangle_1 \star \ldots \star \langle E \rangle_1) =
smd({\langle E \rangle_1}^{\star n}) =
\bigcup\nolimits_{m \geq 1} smd(add(E[-m, m])^{\star n})
$$
である。$n, n' \geq 1$ に対して
$\langle E \rangle_{n + n'} =
smd(\langle E \rangle_n \star \langle E \rangle_{n'})$ である。
\end{lemma}

\begin{proof}
表示式の左の等式は、補題 \ref{derived-lemma-associativity-star} と
\ref{derived-lemma-smd-star} および帰納法から従う。
中央の等式は記法上のものである。
% P02-E0084: frozen source has an extra closing bracket in this formula.
$\langle E \rangle_1 = smd(add(E[-\infty, \infty])])$ であり、
$E[-\infty, \infty] = \bigcup_{m \geq 1} E[-m, m]$ なので、Remark
\ref{derived-remark-operations-unions} と補題 \ref{derived-lemma-smd-star} により
右の等式を得る。最後の主張は、この Remark と補題
\ref{derived-lemma-cone-n} の対応する主張から従う。
\end{proof}

\begin{lemma}
\label{derived-lemma-find-smallest-containing-E}
$\mathcal{D}$ を三角圏とし、$E$ を $\mathcal{D}$ の対象とする。部分圏
$$
\langle E \rangle = \bigcup\nolimits_n \langle E \rangle_n
= \bigcup\nolimits_{n, m \geq 1} smd(add(E[-m, m])^{\star n})
$$
は $\mathcal{D}$ の厳密充満飽和三角部分圏であり、対象 $E$ を含む
$\mathcal{D}$ のそのような部分圏のうち最小のものである。
\end{lemma}

\begin{proof}
右の等式は補題 \ref{derived-lemma-generated-by-E-explicit} から従う。
$\langle E \rangle = \bigcup \langle E \rangle_n$ が $E$ を含み、
平行移動、直和、直和因子で保たれることは明らかである。
$A \in \langle E \rangle_a$ および $B \in \langle E \rangle_b$ であり、
$A \to X \to B \to A[1]$ が完全三角ならば、補題
\ref{derived-lemma-generated-by-E-explicit} により
$X \in \langle E \rangle_{a + b}$ である。従って
$\bigcup \langle E \rangle_n$ は拡大についても閉じており、
三角部分圏であることが従う。

\medskip\noindent
最後に、$\mathcal{D}' \subset \mathcal{D}$ を、$E$ を含む
$\mathcal{D}$ の厳密充満飽和三角部分圏とする。このとき
$\mathcal{D}'[-\infty, \infty] \subset \mathcal{D}'$、
% P02-E0085: frozen source uses add(D), although preservation here requires add(D').
$add(\mathcal{D}) \subset \mathcal{D}'$、
$smd(\mathcal{D}') \subset \mathcal{D}'$、および
$\mathcal{D}' \star \mathcal{D}' \subset \mathcal{D}'$ である。
言い換えれば、$E$ から $\langle E \rangle$ を構成するために用いた
演算はすべて $\mathcal{D}'$ を保つ。従って
$\langle E \rangle \subset \mathcal{D}'$ であり、証明が完了する。
\end{proof}

\begin{definition}
\label{derived-definition-generators}
$\mathcal{D}$ を三角圏とし、$E$ を $\mathcal{D}$ の対象とする。
\begin{enumerate}
\item $E$ を含む $\mathcal{D}$ の最小の厳密充満飽和三角部分圏が
$\mathcal{D}$ に等しいとき、すなわち
$\langle E \rangle = \mathcal{D}$ のとき、$E$ を
$\mathcal{D}$ の{\it 古典的生成元}という。
\item ある $n \geq 1$ に対して
$\langle E \rangle_n = \mathcal{D}$ となるとき、$E$ を
$\mathcal{D}$ の{\it 強生成元}という。
\item $\mathcal{D}$ の任意の非零対象 $K$ に対して、整数 $n$ と
非零写像 $E \to K[n]$ が存在するとき、$E$ を $\mathcal{D}$ の
{\it 弱生成元}または{\it 生成元}という。
\end{enumerate}
\end{definition}

\noindent
この定義は対象の族の場合へ一般化できる。

\begin{lemma}
\label{derived-lemma-right-orthogonal}
$\mathcal{D}$ を三角圏とし、$E, K$ を $\mathcal{D}$ の対象とする。
次は同値である。
\begin{enumerate}
\item すべての $i \in \mathbf{Z}$ に対して $\Hom(E, K[i]) = 0$ である。
\item すべての $E' \in \langle E \rangle$ に対して $\Hom(E', K) = 0$ である。
\end{enumerate}
\end{lemma}

\begin{proof}
(2) $\Rightarrow$ (1) は直ちに従う。逆に (1) を仮定する。
$\langle E \rangle_1$ のすべての $X$ に対して
$\Hom(X, K) = 0$ である。$n$ に関する帰納法と補題
\ref{derived-lemma-representable-homological} を用いると、
$\langle E \rangle_n$ のすべての $X$ に対して
$\Hom(X, K) = 0$ であることが分かる。
\end{proof}

\begin{lemma}
\label{derived-lemma-classical-generator-generator}
$\mathcal{D}$ を三角圏とし、$E$ を $\mathcal{D}$ の対象とする。
$E$ が $\mathcal{D}$ の古典的生成元ならば、$E$ は生成元である。
\end{lemma}

\begin{proof}
$E$ が古典的生成元であると仮定する。$K$ を $\mathcal{D}$ の対象とし、
すべての $i \in \mathbf{Z}$ に対して $\Hom(E, K[i]) = 0$ とする。
補題 \ref{derived-lemma-right-orthogonal} により、$\langle E \rangle$ の
すべての $E'$ に対して $\Hom(E', K) = 0$ である。しかし
$\mathcal{D} = \langle E \rangle$ なので
$\text{id}_K = 0$、すなわち $K = 0$ と結論する。
\end{proof}

\begin{lemma}
\label{derived-lemma-classical-generator-strong-generator}
$\mathcal{D}$ を強生成元をもつ三角圏とする。$E$ を $\mathcal{D}$ の
対象とする。$E$ が $\mathcal{D}$ の古典的生成元ならば、
$E$ は強生成元である。
\end{lemma}

\begin{proof}
$E'$ を、$\mathcal{D} = \langle E' \rangle_n$ を満たす
$\mathcal{D}$ の対象とする。$\mathcal{D} = \langle E \rangle$ なので、
補題 \ref{derived-lemma-find-smallest-containing-E} により、ある
$m \geq 1$ に対して $E' \in \langle E \rangle_m$ である。
従って $\langle E' \rangle_1 \subset \langle E \rangle_m$ であり、
$$
\mathcal{D} =
\langle E' \rangle_n = smd(
\langle E' \rangle_1 \star \ldots \star \langle E' \rangle_1)
\subset
smd(
\langle E \rangle_m \star \ldots \star \langle E \rangle_m)
=
\langle E \rangle_{nm}
$$
となるので、所望の結果を得る。ここで補題
\ref{derived-lemma-generated-by-E-explicit} を用いた。
\end{proof}

\begin{remark}
\label{derived-remark-check-on-generator}
$\mathcal{D}$ を三角圏とし、$E$ を $\mathcal{D}$ の対象とする。
$T$ を $\mathcal{D}$ の対象の性質とし、次を仮定する。
\begin{enumerate}
\item $K_i \in \mathcal{D}$、$i = 1, \ldots, r$ であり、すべての
$i = 1, \ldots, r$ に対して $T(K_i)$ ならば、$T(\bigoplus K_i)$ である。
\item $K \to L \to M \to K[1]$ が完全三角であり、二つの対象について
$T$ が成り立つならば、第三の対象についても $T$ が成り立つ。
\item $T(K \oplus L)$ ならば $T(K)$ および $T(L)$ である。
\item すべての $n$ に対して $T(E[n])$ である。
\end{enumerate}
このとき $\langle E \rangle$ のすべての対象について $T$ が成り立つ。
\end{remark}

\section{コンパクト対象}
\label{derived-section-compact}

\noindent
定義は次の通りである。

\begin{definition}
\label{derived-definition-compact-object}
$\mathcal{D}$ を任意直和をもつ加法圏とする。$\mathcal{D}$ の
{\it コンパクト対象}とは、任意の集合 $I$ と、$i \in I$ で
添字付けられた対象 $E_i \in \Ob(\mathcal{D})$ に対して写像
$$
\bigoplus\nolimits_{i \in I} \Hom_{\mathcal{D}}(K, E_i)
\longrightarrow
\Hom_{\mathcal{D}}(K, \bigoplus\nolimits_{i \in I} E_i)
$$
が全単射となる対象 $K$ のことである。
\end{definition}

\noindent
この概念は代数幾何において非常に有用であることが分かる。
これは、対象をまさに何らかの意味で「コンパクト」にする内在的条件である。

\begin{lemma}
\label{derived-lemma-compact-objects-subcategory}
$\mathcal{D}$ を直和をもつ（前）三角圏とする。このとき
$\mathcal{D}$ のコンパクト対象は、$\mathcal{D}$ の Karoubi 的な、
飽和した厳密充満（前）三角部分圏 $\mathcal{D}_c$ の対象全体をなす。
\end{lemma}

\begin{proof}
$(X, Y, Z, f, g, h)$ を $\mathcal{D}$ の完全三角とし、
$X$ と $Y$ はコンパクトであるとする。このとき補題
\ref{derived-lemma-representable-homological} と五項補題
（Homology, Lemma \ref{homology-lemma-five-lemma}）により、
$Z$ もコンパクト対象である。$X \oplus Y$ がコンパクトならば
$X$, $Y$ もコンパクト対象であることは明らかである。
従って $\mathcal{D}_c$ は飽和三角部分圏である。
$\mathcal{D}$ は補題 \ref{derived-lemma-projectors-have-images-triangulated} により
Karoubi 的なので、$\mathcal{D}_c$ も同様であると結論する。
\end{proof}

\begin{lemma}
\label{derived-lemma-write-as-colimit}
$\mathcal{D}$ を直和をもつ三角圏とする。$E_i$、$i \in I$ を
$\mathcal{D}$ のコンパクト対象の族とし、$\bigoplus E_i$ は
$\mathcal{D}$ を生成すると仮定する。このとき $\mathcal{D}$ の
任意の対象 $X$ は
$$
X = \text{hocolim} X_n
$$
と書ける。ここで $X_1$ は $E_i$ の平行移動の直和であり、
各遷移射は完全三角
$Y_n \to X_n \to X_{n + 1} \to Y_n[1]$
に組み込まれ、$Y_n$ は $E_i$ の平行移動の直和である。
\end{lemma}

\begin{proof}
$X_1 = \bigoplus_{(i, m, \varphi)} E_i[m]$ と置く。ここで直和は、
$i \in I$、$m \in \mathbf{Z}$、および
$\varphi : E_i[m] \to X$ を満たすすべての三つ組
$(i, m, \varphi)$ にわたる。このとき $X_1$ には標準射
$X_1 \to X$ が付随する。$X_n \to X$ が与えられたとき、
$Y_n = \bigoplus_{(i, m, \varphi)} E_i[m]$ と置く。ここで直和は、
$i \in I$、$m \in \mathbf{Z}$ を満たし、かつ
$\varphi : E_i[m] \to X_n$ が、$E_i[m] \to X_n \to X$ を
零にする射であるような、すべての三つ組 $(i, m, \varphi)$ にわたる。
完全三角 $Y_n \to X_n \to X_{n + 1} \to Y_n[1]$ を選び、
$X_n \to X_{n + 1} \to X$ が与えられた写像となるような射
$X_{n + 1} \to X$ を一つ選ぶ。そのような射は $Y_n$ の選び方により存在する。
写像 $X_n \to X$ の構成から射 $\text{hocolim} X_n \to X$ を得る。
完全三角
$$
C \to \text{hocolim} X_n \to X \to C[1]
$$
を選ぶ。射 $E_i[m] \to C$ を一つ取る。$E_i$ はコンパクトなので、
補題 \ref{derived-lemma-commutes-with-countable-sums} により、合成
$E_i[m] \to C \to \text{hocolim} X_n$ は、ある $n$ に対して
$X_n$ を経由する。その射を $E_i[m] \to X_n$ とする。
すると $Y_n$ の構成により、合成
$E_i[m] \to X_n \to X_{n + 1}$ は零である。言い換えれば、合成
$E_i[m] \to C \to \text{hocolim} X_n$ は零である。
これは、射 $E_i[m] \to C$ が射 $E_i[m] \to X[-1]$ から来ることを意味する。
$X_1$ の構成により、そのような射は $\text{hocolim} X_n$ へ持ち上がる。
従って射 $E_i[m] \to C$ は零であると結論する。
$\bigoplus E_i$ が $\mathcal{D}$ を生成するという仮定から
$C$ は零であり、証明が完了する。
\end{proof}

\begin{lemma}
\label{derived-lemma-factor-through}
補題 \ref{derived-lemma-write-as-colimit} と同じ仮定および記法を用いる。
% P02-E0087: frozen source sentence “With assumptions and notation ...” lacks a finite verb.
$C$ がコンパクト対象で、$C \to X_n$ が射ならば、分解
$C \to E \to X_n$ であって、ある $i_1, \ldots, i_t \in I$ に対して
$E$ が $\langle E_{i_1} \oplus \ldots \oplus E_{i_t} \rangle$
の対象となるものが存在する。
\end{lemma}

\begin{proof}
$n$ に関する帰納法で証明する。初期の場合 $n = 1$ は明らかである。
$n > 1$ ならば、合成 $C \to X_n \to Y_{n - 1}[1]$ を考える。
これは、$E_i$ の平行移動の有限直和であるある $E'$ に対する
$E'[1] \to Y_{n - 1}[1]$ を経由して分解する。$I' \subset I$ を、
この直和に現れる添字の有限集合とする。従って
$$
\xymatrix{
E' \ar[r] \ar[d] &
C' \ar[r] \ar[d] &
C \ar[r] \ar[d] &
E'[1] \ar[d] \\
Y_{n - 1} \ar[r] &
X_{n - 1} \ar[r] &
X_n \ar[r] &
Y_{n - 1}[1]
}
$$
を得る。帰納法により、射 $C' \to X_{n - 1}$ は
$E'' \to X_{n - 1}$ を経由して分解する。ここで $E''$ は、
ある有限部分集合 $I'' \subset I$ に対する
$\langle \bigoplus_{i \in I''} E_i \rangle$ の対象である。
完全三角
$$
E' \to E'' \to E \to E'[1]
$$
を選ぶと、$E$ は
$\langle \bigoplus_{i \in I' \cup I''} E_i \rangle$ の対象である。
構成と三角圏の公理により、三角の射
$(E', C', C) \to (E', E'', E)$ および
$(E', E'', E) \to (Y_{n - 1}, X_{n - 1}, X_n)$ に組み込まれる射
$C \to E$ と射 $E \to X_n$ を選べる。合成
$C \to E \to X_n$ は与えられた射 $C \to X_n$ と等しいとは限らないが、
$Y_{n - 1}$ への合成は等しい。その差を持ち上げる射
$C \to X_{n - 1}$ を取る。帰納法の仮定により、これは
$E''' \to X_{n - 1}$ を経由して分解する。ただし
% P02-E0088: frozen source says E'' here, although the newly introduced object is E'''.
$E''$ は、ある有限部分集合
% P02-E0089: frozen source repeats I' here, although the new index set is I'''.
$I' \subset I$ に対する $\langle \bigoplus_{i \in I'''} E_i \rangle$
の対象である。従って $E \oplus E''' \to X_n$ を考えれば解を得る。
実際、$E \oplus E'''$ は
$\langle \bigoplus_{i \in I' \cup I'' \cup I'''} E_i \rangle$ の対象である。
\end{proof}

\begin{definition}
\label{derived-definition-compactly-generated}
$\mathcal{D}$ を任意直和をもつ三角圏とする。コンパクト対象の集合
$E_i$、$i \in I$ であって、$\bigoplus E_i$ が $\mathcal{D}$ を
生成するものが存在するとき、$\mathcal{D}$ は{\it コンパクト生成}であるという。
\end{definition}

\noindent
次の命題は、古典的生成元と弱生成元の関係を明らかにする。

\begin{proposition}
\label{derived-proposition-generator-versus-classical-generator}
$\mathcal{D}$ を直和をもつ三角圏とし、$E$ を $\mathcal{D}$ の
コンパクト対象とする。次は同値である。
\begin{enumerate}
\item $E$ は $\mathcal{D}_c$ の古典的生成元であり、
$\mathcal{D}$ はコンパクト生成である。
\item $E$ は $\mathcal{D}$ の生成元である。
\end{enumerate}
\end{proposition}

\begin{proof}
$E$ が $\mathcal{D}_c$ の古典的生成元ならば
$\mathcal{D}_c = \langle E \rangle$ である。$\mathcal{D}$ が
コンパクト生成であるという仮定と補題 \ref{derived-lemma-right-orthogonal} から、
$E$ が $\mathcal{D}$ の生成元であることが形式的に従う。

\medskip\noindent
逆向きの方が興味深い。$E$ が $\mathcal{D}$ の生成元であると仮定する。
$X$ を $\mathcal{D}$ のコンパクト対象とする。補題
\ref{derived-lemma-write-as-colimit} に $I = \{1\}$ および $E_1 = E$ を用いて、
補題の通り
$$
X = \text{hocolim} X_n
$$
と書く。$X$ はコンパクトなので、ある $n$ に対して
$X \to \text{hocolim} X_n$ は $X_n$ を経由して分解する
（補題 \ref{derived-lemma-commutes-with-countable-sums}）。
従って $X$ は $X_n$ の直和因子である。補題
\ref{derived-lemma-factor-through} により、$X$ は $\langle E \rangle$ の対象であり、
% P02-E0091: frozen source closes this proposition proof with “the lemma is proven”.
補題が証明された。
\end{proof}











\section{Brown の表現可能性}
\label{derived-section-brown}

\noindent
本節の内容については \cite{Neeman-Grothendieck} を参照されたい。

\begin{lemma}
\label{derived-lemma-brown}
\begin{reference}
\cite[Theorem 3.1]{Neeman-Grothendieck}.
\end{reference}
直和をもつコンパクト生成三角圏を $\mathcal{D}$ とする。
$H : \mathcal{D} \to \textit{Ab}$ を、直和を直積に移す反変
コホモロジー的関手とする。このとき $H$ は表現可能である。
\end{lemma}

\begin{proof}
$E_i$、$i \in I$ を、$\bigoplus_{i \in I} E_i$ が $\mathcal{D}$ を
生成するようなコンパクト対象の集合とする。対象の集合 $\{E_i\}$ は
平行移動で閉じていると仮定してよく、実際そう仮定する。
$i \in I$ および $a \in H(E_i)$ に対する組 $(i, a)$ を考え、
$$
X_1 = \bigoplus\nolimits_{(i, a)} E_i
$$
と置く。$H(X_1) = \prod_{(i, a)} H(E_i)$ なので、
$(a)_{(i, a)}$ は元 $a_1 \in H(X_1)$ を定める。
$H_1 = \Hom_\mathcal{D}(- , X_1)$ と置く。米田の補題
（Categories, Lemma \ref{categories-lemma-yoneda}）により、元 $a_1$ は
自然変換 $H_1 \to H$ を定める。

\medskip\noindent
$H_n = \Hom_\mathcal{D}(-, X_n)$ として、$X_n$ と変換
$a_n : H_n \to H$ を帰納的に構成する。すなわち、関手
$\Ker(H_n \to H)$ に上の手順を適用して、対象
$$
K_{n + 1} = \bigoplus\nolimits_{(i, k),\ k \in \Ker(H_n(E_i) \to H(E_i))} E_i
$$
と変換 $\Hom_\mathcal{D}(-, K_{n + 1}) \to \Ker(H_n \to H)$ を得る。
米田の補題により、合成 $\Hom_\mathcal{D}(-, K_{n + 1}) \to H_n$ は
射 $K_{n + 1} \to X_n$ を与える。$\mathcal{D}$ における完全三角
$$
K_{n + 1} \to X_n \to X_{n + 1} \to K_{n + 1}[1]
$$
を選ぶ。構成により、元 $a_n \in H(X_n)$ の $H(K_{n + 1})$ における像は
零である。$H$ はコホモロジー的なので、これを元
$a_{n + 1} \in H(X_{n + 1})$ に持ち上げることができる。

\medskip\noindent
$X = \text{hocolim} X_n$ が $H$ を表現すると主張する。定義する完全三角
$$
\bigoplus X_n \to
\bigoplus X_n \to X \to
\bigoplus X_n[1]
$$
に $H$ を適用すると、完全列
$$
\prod H(X_n) \leftarrow 
\prod H(X_n) \leftarrow 
H(X)
$$
を得る。従って、$\prod H(X_n)$ における像が $(a_n)$ である元
$a \in H(X)$ が存在する。従って、
% P02-E0092: frozen source has the sentence fragment “Hence a natural transformation ...”.
$$
\Hom_\mathcal{D}(-, X_1) \to
\Hom_\mathcal{D}(-, X_2) \to
\Hom_\mathcal{D}(-, X_3) \to \ldots \to
\Hom_\mathcal{D}(-, X) \to H
$$
が可換となるような自然変換 $\Hom_\mathcal{D}(- , X) \to H$ が得られる。
各 $i$ に対して、写像 $\Hom_\mathcal{D}(E_i, X) \to H(E_i)$ は、
$X_1$ の構成により全射である。他方、$X_n \to X_{n + 1}$ の構成により、
$\Hom_\mathcal{D}(E_i, X_n) \to H(E_i)$ の核は写像
$\Hom_\mathcal{D}(E_i, X_n) \to \Hom_\mathcal{D}(E_i, X_{n + 1})$
によって零に送られる。補題 \ref{derived-lemma-commutes-with-countable-sums} により
$$
\Hom_\mathcal{D}(E_i, X) = \colim \Hom_\mathcal{D}(E_i, X_n)
$$
なので、$\Hom_\mathcal{D}(E_i, X) \to H(E_i)$ は単射である。

\medskip\noindent
証明を完了するため、部分圏
$$
\mathcal{D}' =
\{Y \in \Ob(\mathcal{D}) \mid \Hom_\mathcal{D}(Y[n], X) \to H(Y[n])
\text{ is an isomorphism for all }n\}
$$
を考える。$\Hom_\mathcal{D}(-, X) \to H$ はコホモロジー的関手の間の
変換なので、部分圏 $\mathcal{D}'$ は $\mathcal{D}$ の厳密充満な飽和三角部分圏
である（詳細は省略する。補題 \ref{derived-lemma-homological-functor-kernel} の証明を
参照されたい）。さらに、$H$ と $\Hom_\mathcal{D}(-, X)$ はともに直和を
直積に移すので、$\mathcal{D}'$ の対象の直和は $\mathcal{D}'$ に属する。
従って、$\mathcal{D}'$ の対象の導来余極限も $\mathcal{D}'$ に属する。
$\{E_i\}$ は平行移動で閉じているので、すべての $i$ に対して $E_i$ は
$\mathcal{D}'$ の対象である。補題 \ref{derived-lemma-write-as-colimit} から
$\mathcal{D}' = \mathcal{D}$ が従い、証明が完了する。
\end{proof}

\begin{proposition}
\label{derived-proposition-brown}
\begin{reference}
\cite[Theorem 4.1]{Neeman-Grothendieck}.
\end{reference}
直和をもつコンパクト生成三角圏を $\mathcal{D}$ とする。
$F : \mathcal{D} \to \mathcal{D}'$ を、直和を直和に移す三角圏の完全関手とする。
このとき $F$ は完全な右随伴をもつ。
\end{proposition}

\begin{proof}
$\mathcal{D}'$ の対象 $Y$ に対し、反変関手
$$
\mathcal{D} \to \textit{Ab},\quad W \mapsto \Hom_{\mathcal{D}'}(F(W), Y)
$$
を考える。$F$ は完全なので、これはコホモロジー的関手である。また、
$F$ は直和を直和に移すので、この関手は直和を直積に移す。従って補題
\ref{derived-lemma-brown} により、
$\Hom_\mathcal{D}(W, X) = \Hom_{\mathcal{D}'}(F(W), Y)$ となる
$\mathcal{D}$ の対象 $X$ が存在する。随伴の存在は Categories, Lemma
\ref{categories-lemma-adjoint-exists} から従い、完全性は補題
\ref{derived-lemma-adjoint-is-exact} から従う。
\end{proof}










\section{Brown の表現可能性、第二形}
\label{derived-section-brown-bis}

\noindent
本節では、適切な生成元の集合をもつ三角圏に対する Brown の表現可能性の
一つの形を説明する。他の形については \cite{Franke}、\cite{Neeman}、
および \cite{Krause} を参照されたい。

\begin{lemma}
\label{derived-lemma-brown-bis}
\begin{reference}
\cite[Theorem A]{Krause} の弱い形。
\end{reference}
直和をもつ三角圏を $\mathcal{D}$ とする。次の条件を満たす
$\mathcal{D}$ の対象の集合 $\mathcal{E}$ が与えられていると仮定する。
\begin{enumerate}
\item $X$ が $\mathcal{D}$ の非零対象ならば、ある
$E \in \mathcal{E}$ と非零写像 $E \to X$ が存在する。
\item $\mathcal{D}$ の対象 $X_n$、$n \in \mathbf{N}$、
$E \in \mathcal{E}$、および $\alpha : E \to \bigoplus X_n$ が
与えられたとき、$E_n \in \mathcal{E}$、$\beta_n : E_n \to X_n$、
および射 $\gamma : E \to \bigoplus E_n$ であって
$\alpha = (\bigoplus \beta_n) \circ \gamma$ となるものが存在する。
\end{enumerate}
$H : \mathcal{D} \to \textit{Ab}$ を、直和を直積に移す反変
コホモロジー的関手とする。このとき $H$ は表現可能である。
\end{lemma}

\begin{proof}
この証明は補題 \ref{derived-lemma-brown} の証明と非常によく似ている。
$\mathcal{E}$ を $\bigcup_{i \in \mathbf{Z}} \mathcal{E}[i]$ で置き換え、
$\mathcal{E}$ が平行移動で閉じていると仮定してよい。
$E \in \mathcal{E}$ および $a \in H(E)$ に対する組 $(E, a)$ を考え、
$$
X_1 = \bigoplus\nolimits_{(E, a)} E
$$
と置く。$H(X_1) = \prod_{(E, a)} H(E)$ なので、
$(a)_{(E, a)}$ は元 $a_1 \in H(X_1)$ を定める。
$H_1 = \Hom_\mathcal{D}(- , X_1)$ と置く。米田の補題
（Categories, Lemma \ref{categories-lemma-yoneda}）により、元 $a_1$ は
自然変換 $H_1 \to H$ を定める。

\medskip\noindent
$H_n = \Hom_\mathcal{D}(-, X_n)$ として、$X_n$ と変換
$a_n : H_n \to H$ を帰納的に構成する。すなわち、関手
$\Ker(H_n \to H)$ に上の手順を適用して、対象
$$
K_{n + 1} = \bigoplus\nolimits_{(E, k),\ k \in \Ker(H_n(E) \to H(E))} E
$$
と変換 $\Hom_\mathcal{D}(-, K_{n + 1}) \to \Ker(H_n \to H)$ を得る。
米田の補題により、合成 $\Hom_\mathcal{D}(-, K_{n + 1}) \to H_n$ は
射 $K_{n + 1} \to X_n$ を与える。$\mathcal{D}$ における完全三角
$$
K_{n + 1} \to X_n \to X_{n + 1} \to K_{n + 1}[1]
$$
を選ぶ。構成により、元 $a_n \in H(X_n)$ の $H(K_{n + 1})$ における像は
零である。$H$ はコホモロジー的なので、これを元
$a_{n + 1} \in H(X_{n + 1})$ に持ち上げることができる。

\medskip\noindent
$X = \text{hocolim} X_n$ と置く。定義する完全三角
$$
\bigoplus X_n \to
\bigoplus X_n \to X \to
\bigoplus X_n[1]
$$
に $H$ を適用すると、完全列
$$
\prod H(X_n) \leftarrow
\prod H(X_n) \leftarrow
H(X)
$$
を得る。従って、$\prod H(X_n)$ における像が $(a_n)$ である元
$a \in H(X)$ が存在する。従って、
% P02-E0093: frozen source repeats the sentence fragment “Hence a natural transformation ...”.
$$
\Hom_\mathcal{D}(-, X_1) \to
\Hom_\mathcal{D}(-, X_2) \to
\Hom_\mathcal{D}(-, X_3) \to \ldots \to
\Hom_\mathcal{D}(-, X) \to H
$$
が可換となるような自然変換 $\Hom_\mathcal{D}(- , X) \to H$ が得られる。
$\Hom_\mathcal{D}(-, X) \to H(-)$ は同型であると主張する。

\medskip\noindent
$E \in \mathcal{E}$ とする。写像
$$
\Hom_\mathcal{D}(E, \bigoplus X_n) \to  \Hom_\mathcal{D}(E, \bigoplus X_n)
$$
が単射であることを示そう。実際、$\alpha : E \to \bigoplus X_n$ とする。
仮定 (2) により分解 $\alpha = (\bigoplus \beta_n) \circ \gamma$ を得る。
構成により $E_n \to X_n \to X_{n + 1}$ は零なので、合成
$\bigoplus E_n \to \bigoplus X_n \to \bigoplus X_n$ は
$\bigoplus \beta_n$ と等しい。従って合成
$E \to \bigoplus X_n \to \bigoplus X_n$ も $\alpha$ と等しい。
これで述べた単射性が証明され、従って写像
$$
\Hom_\mathcal{D}(E, \bigoplus X_n[1]) \to \Hom_\mathcal{D}(E, \bigoplus X_n[1])
$$
も単射である。従って、すべての $E \in \mathcal{E}$ に対して完全列
$$
\Hom_\mathcal{D}(E, \bigoplus X_n) \to
\Hom_\mathcal{D}(E, \bigoplus X_n) \to
\Hom_\mathcal{D}(E, X) \to 0
$$
を得る。

\medskip\noindent
$E \in \mathcal{E}$ とし、$f : E \to X$ を射とする。前段落により、
$f$ を持ち上げる $\alpha : E \to \bigoplus X_n$ を選べる。
仮定 (2) により分解 $\alpha = (\bigoplus \beta_n) \circ \gamma$ を得る。
各 $n$ に対して、$\delta_n$ と $\beta_n$ が $H(E_n)$ の同じ元に写るような
射 $\delta_n : E_n \to X_1$ が存在する。すると、$X_n \to X_{n + 1}$ の
構成により、合成
$$
E_n \to X_n \to X_{n + 1}
\quad\text{and}\quad
E_n \to X_1 \to X_{n + 1}
$$
は等しい。従って、合成
$$
\bigoplus E_n \to \bigoplus X_n \to X
\quad\text{and}\quad
\bigoplus E_n \to \bigoplus X_1 \to X
$$
も等しい。$\bigoplus X_1 \to X$ が
$\bigoplus X_1 \to X_1 \to X$ と分解することに注意すると、写像
$$
\Hom_\mathcal{D}(E, X_1) \to \Hom_\mathcal{D}(E, X)
$$
は全射である。構成により写像 $\Hom_\mathcal{D}(E, X_1) \to H(E)$ は
全射であり、また構成によりこの写像の核は
$\Hom_\mathcal{D}(E, X_1) \to \Hom_\mathcal{D}(E, X)$ によって
零に送られる。従って、すべての $E \in \mathcal{E}$ に対して
$\Hom_\mathcal{D}(E, X) \to H(E)$ は全単射である。

\medskip\noindent
証明を完了するため、部分圏
$$
\mathcal{D}' =
\{Y \in \Ob(\mathcal{D}) \mid \Hom_\mathcal{D}(Y[n], X) \to H(Y[n])
\text{ is an isomorphism for all }n\}
$$
を考える。$\Hom_\mathcal{D}(-, X) \to H$ はコホモロジー的関手の間の
変換なので、部分圏 $\mathcal{D}'$ は $\mathcal{D}$ の厳密充満な飽和三角部分圏
である（詳細は省略する。補題 \ref{derived-lemma-homological-functor-kernel} の証明を
参照されたい）。さらに、$H$ と $\Hom_\mathcal{D}(-, X)$ はともに直和を
直積に移すので、$\mathcal{D}'$ の対象の直和は $\mathcal{D}'$ に属する。
従って、$\mathcal{D}'$ の対象の導来余極限も $\mathcal{D}'$ に属する。
$\mathcal{E}$ は平行移動で閉じているので、前段落の結果から
$\mathcal{E} \subset \Ob(\mathcal{D}')$ と結論できる。証明を完了するには
$\mathcal{D}' = \mathcal{D}$ を示さなければならない。

\medskip\noindent
$Y$ を $\mathcal{D}$ の対象とし、$H(-) = \Hom_\mathcal{D}(-, Y)$ と置く。
すると $H$ は直和を直積に移すコホモロジー的関手である。証明の前半の構成により、
% P02-E0094: frozen source writes colim X_n although X was defined as hocolim X_n.
$\colim X_n = X \to Y$ という射であって、すべての $E \in \mathcal{E}$ に対して
$\Hom_\mathcal{D}(E, X) \to \Hom_\mathcal{D}(E, Y)$ が全単射となるものを得る。
すると仮定 (1) により $X \to Y$ は同型である！ 他方、構成により
$X_1, X_2, \ldots$ は $\mathcal{D}'$ に属し、従って $X$ も属する。
従って $Y \in \mathcal{D}'$ であり、証明が完了する。
\end{proof}

\begin{proposition}
\label{derived-proposition-brown-bis}
直和をもつ三角圏を $\mathcal{D}$ とする。補題 \ref{derived-lemma-brown-bis} の
条件 (1) および (2) を満たす $\mathcal{D}$ の対象の集合 $\mathcal{E}$ が
存在すると仮定する。$F : \mathcal{D} \to \mathcal{D}'$ を、直和を直和に
移す三角圏の完全関手とする。このとき $F$ は完全な右随伴をもつ。
\end{proposition}

\begin{proof}
$\mathcal{D}'$ の対象 $Y$ に対し、反変関手
$$
\mathcal{D} \to \textit{Ab},\quad W \mapsto \Hom_{\mathcal{D}'}(F(W), Y)
$$
を考える。$F$ は完全なので、これはコホモロジー的関手である。また、
$F$ は直和を直和に移すので、この関手は直和を直積に移す。従って補題
\ref{derived-lemma-brown-bis} により、
$\Hom_\mathcal{D}(W, X) = \Hom_{\mathcal{D}'}(F(W), Y)$ となる
$\mathcal{D}$ の対象 $X$ が存在する。随伴の存在は Categories, Lemma
\ref{categories-lemma-adjoint-exists} から従い、完全性は補題
\ref{derived-lemma-adjoint-is-exact} から従う。
\end{proof}

\section{許容部分圏}
\label{derived-section-admissible}

\noindent
本節については \cite[Section 1]{Bondal-Kapranov} を参照されたい。

\begin{definition}
\label{derived-definition-orthogonal}
$\mathcal{D}$ を加法圏とし、$\mathcal{A} \subset \mathcal{D}$ を充満部分圏とする。
$\mathcal{A}$ の{it 右直交} $\mathcal{A}^\perp$ とは、すべての
$A \in \Ob(\mathcal{A})$ に対して $\Hom(A, X) = 0$ となる
$\mathcal{D}$ の対象 $X$ からなる充満部分圏である。
$\mathcal{A}$ の{it 左直交} ${}^\perp\mathcal{A}$ とは、すべての
$A \in \Ob(\mathcal{A})$ に対して $\Hom(X, A) = 0$ となる
$\mathcal{D}$ の対象 $X$ からなる充満部分圏である。
\end{definition}

\begin{lemma}
\label{derived-lemma-pre-prepare-adjoint}
$\mathcal{D}$ を三角圏とし、$\mathcal{A} \subset \mathcal{D}$ を
すべての平行移動で不変な充満部分圏とする。$\mathcal{D}$ の完全三角
$$
X \to Y \to Z \to X[1]
$$
を考える。次は同値である。
\begin{enumerate}
\item $Z$ は $\mathcal{A}^\perp$ に属する。
\item すべての $A \in \Ob(\mathcal{A})$ に対して
$\Hom(A, X) = \Hom(A, Y)$ である。
\end{enumerate}
\end{lemma}

\begin{proof}
補題 \ref{derived-lemma-representable-homological} により、関手
$\Hom(A, -)$ はホモロジー的であり、従って
(\ref{derived-equation-long-exact-cohomology-sequence}) のような長完全列を得る。
(1) を仮定し、$A \in \Ob(\mathcal{A})$ とする。このとき完全列
$$
\Hom(A[1], Z) \to \Hom(A, X) \to \Hom(A, Y) \to \Hom(A, Z)
$$
を考える。$A[1] \in \Ob(\mathcal{A})$ なので、最初と最後の群は零である。
従って (2) を得る。(2) を仮定し、$A \in \Ob(\mathcal{A})$ とする。
このとき完全列
$$
\Hom(A, X) \to \Hom(A, Y) \to \Hom(A, Z) \to \Hom(A[-1], X) \to \Hom(A[-1], Y)
$$
を考えれば、望み通り $\Hom(A, Z) = 0$ と結論できる。
\end{proof}

\begin{lemma}
\label{derived-lemma-pre-prepare-adjoint-dual}
$\mathcal{D}$ を三角圏とし、$\mathcal{B} \subset \mathcal{D}$ を
すべての平行移動で不変な充満部分圏とする。$\mathcal{D}$ の完全三角
$$
X \to Y \to Z \to X[1]
$$
を考える。次は同値である。
\begin{enumerate}
\item $X$ は ${}^\perp\mathcal{B}$ に属する。
\item すべての $B \in \Ob(\mathcal{B})$ に対して
$\Hom(Y, B) = \Hom(Z, B)$ である。
\end{enumerate}
\end{lemma}

\begin{proof}
補題 \ref{derived-lemma-pre-prepare-adjoint} の双対である。
\end{proof}

\begin{lemma}
\label{derived-lemma-orthogonal-triangulated}
$\mathcal{D}$ を三角圏とし、$\mathcal{A} \subset \mathcal{D}$ を
すべての平行移動で不変な充満部分圏とする。このとき $\mathcal{A}$ の
右直交 $\mathcal{A}^\perp$ と左直交 ${}^\perp\mathcal{A}$ はともに、
$\mathcal{D}$ の厳密充満な飽和\footnote{定義 \ref{derived-definition-saturated}。}
三角部分圏である。
% P02-E0095: frozen source misspells “subcategories” as “subcagories”.
\end{lemma}

\begin{proof}
定義から、直交部分圏が平行移動、直和、および直和因子を取る操作で閉じている
ことは直ちに分かる。$\mathcal{D}$ の完全三角
$$
X \to Y \to Z \to X[1]
$$
を考える。補題 \ref{derived-lemma-triangulated-subcategory} により、$X$ と $Y$ が
$\mathcal{A}^\perp$ に属するならば $Z$ も $\mathcal{A}^\perp$ に属することを
示せば十分である。これは補題 \ref{derived-lemma-pre-prepare-adjoint} から直ちに従う。
\end{proof}

\begin{lemma}
\label{derived-lemma-prepare-adjoint}
$\mathcal{D}$ を三角圏とし、$\mathcal{A}$ を $\mathcal{D}$ の
充満三角部分圏とする。$\mathcal{D}$ の対象 $X$ に対して、次の性質 $P(X)$ を
考える。$\mathcal{D}$ における完全三角 $A \to X \to B \to A[1]$ であって、
$A$ は $\mathcal{A}$ に属し、$B$ は $\mathcal{A}^\perp$ に属するものが存在する。
\begin{enumerate}
\item $X_1 \to X_2 \to X_3 \to X_1[1]$ が完全三角で、三つのうち二つに
対して $P$ が成り立つならば、残る一つに対しても成り立つ。
\item $X_1$ と $X_2$ に対して $P$ が成り立つならば、
$X_1 \oplus X_2$ に対しても成り立つ。
\end{enumerate}
\end{lemma}

\begin{proof}
$X_1 \to X_2 \to X_3 \to X_1[1]$ を完全三角とし、$X_1$ と $X_2$ に
対して $P$ が成り立つと仮定する。条件 $P$ にある完全三角
$$
A_1 \to X_1 \to B_1 \to A_1[1]
\quad\text{and}\quad
A_2 \to X_2 \to B_2 \to A_2[1]
$$
を選ぶ。補題 \ref{derived-lemma-pre-prepare-adjoint} により
$\Hom(A_1, A_2) = \Hom(A_1, X_2)$ なので、図式
$$
\xymatrix{
A_1 \ar[d] \ar[r] & X_1 \ar[d] \\
A_2 \ar[r] & X_2
}
$$
が可換となるような射 $A_1 \to A_2$ が一意に存在する。命題
\ref{derived-proposition-9} のように、これを図式
$$
\xymatrix{
A_1 \ar[r] \ar[d] & X_1 \ar[r] \ar[d] & Q_1 \ar[r] \ar[d] & A_1[1] \ar[d] \\
A_2 \ar[r] \ar[d] & X_2 \ar[r] \ar[d] & Q_2 \ar[r] \ar[d] & A_2[1] \ar[d] \\
A_3 \ar[r] \ar[d] & X_3 \ar[r] \ar[d] & Q_3 \ar[r] \ar[d] & A_3[1] \ar[d] \\
A_1[1] \ar[r] & X_1[1] \ar[r] & Q_1[1] \ar[r] & A_1[2]
}
$$
へ拡張するものを選ぶ。TR3 により $Q_1 \cong B_1$ および
$Q_2 \cong B_2$ であり、従って $Q_1, Q_2 \in \Ob(\mathcal{A}^\perp)$ である。
$Q_1 \to Q_2 \to Q_3 \to Q_1[1]$ は完全三角なので、補題
\ref{derived-lemma-orthogonal-triangulated} により
$Q_3 \in \Ob(\mathcal{A}^\perp)$ である。$\mathcal{A}$ は充満三角部分圏なので、
$A_3$ は $\mathcal{A}$ のある対象と同型である。従って $X_3$ は $P$ を満たす。
(1) の他の場合は、平行移動によってこの場合から従う。(2) は補題
\ref{derived-lemma-split} によって (1) の特別な場合である。
\end{proof}

\begin{lemma}
\label{derived-lemma-prepare-adjoint-dual}
$\mathcal{D}$ を三角圏とし、$\mathcal{B}$ を $\mathcal{D}$ の
充満三角部分圏とする。$\mathcal{D}$ の対象 $X$ に対して、次の性質 $P(X)$ を
考える。$\mathcal{D}$ における完全三角 $A \to X \to B \to A[1]$ であって、
$B$ は $\mathcal{B}$ に属し、$A$ は ${}^\perp\mathcal{B}$ に属するものが存在する。
\begin{enumerate}
\item $X_1 \to X_2 \to X_3 \to X_1[1]$ が完全三角で、三つのうち二つに
対して $P$ が成り立つならば、残る一つに対しても成り立つ。
\item $X_1$ と $X_2$ に対して $P$ が成り立つならば、
$X_1 \oplus X_2$ に対しても成り立つ。
\end{enumerate}
\end{lemma}

\begin{proof}
補題 \ref{derived-lemma-prepare-adjoint} の双対である。
\end{proof}

\begin{lemma}
\label{derived-lemma-right-adjoint}
$\mathcal{D}$ を三角圏とし、$\mathcal{A} \subset \mathcal{D}$ を
充満三角部分圏とする。次は同値である。
\begin{enumerate}
\item 包含関手 $\mathcal{A} \to \mathcal{D}$ は右随伴をもつ。
\item $\mathcal{D}$ のすべての $X$ に対して、$\mathcal{D}$ における完全三角
$$
A \to X \to B \to A[1]
$$
であって、$A \in \Ob(\mathcal{A})$ および
$B \in \Ob(\mathcal{A}^\perp)$ となるものが存在する。
\end{enumerate}
これらが成り立つならば、$\mathcal{A}$ は飽和している
（定義 \ref{derived-definition-saturated}）。さらに $\mathcal{A}$ が
$\mathcal{D}$ において厳密充満ならば、
$\mathcal{A} = {}^\perp(\mathcal{A}^\perp)$ である。
\end{lemma}

\begin{proof}
(1) を仮定し、右随伴を $v : \mathcal{D} \to \mathcal{A}$ と書く。
$X \in \Ob(\mathcal{D})$ とし、$A = v(X)$ と置く。随伴射 $A \to X$ を
完全三角 $A \to X \to B \to A[1]$ に延長できる。
$A' \in \Ob(\mathcal{A})$ に対して
$$
\Hom_\mathcal{A}(A', A) =
\Hom_\mathcal{A}(A', v(X)) =
\Hom_\mathcal{D}(A', X)
$$
なので、補題 \ref{derived-lemma-pre-prepare-adjoint} により
$B \in \Ob(\mathcal{A}^\perp)$ と結論する。

\medskip\noindent
(2) を仮定する。随伴 $v$ を明示的に構成する。
$X \in \Ob(\mathcal{D})$ とする。(2) のような
$A \to X \to B \to A[1]$ を選び、$v(X) = A$ と置く。
$f : X \to Y$ を $\mathcal{D}$ の射とする。(2) のような
$A' \to Y \to B' \to A'[1]$ を選ぶ。補題
\ref{derived-lemma-pre-prepare-adjoint} により
$\Hom(A, A') = \Hom(A, Y)$ なので、図式
$$
\xymatrix{
A \ar[d]_{f'} \ar[r] & X \ar[d]^f \\
A' \ar[r] & Y
}
$$
が可換となるような射 $f' : A \to A'$ が一意に存在する。
従って $v(f) = f'$ と置くことで関手を得る。$v$ が包含の右随伴であることを
見るには、補題 \ref{derived-lemma-pre-prepare-adjoint} をもう一度用いればよい。

\medskip\noindent
最後の主張を証明する。$\mathcal{A}$ が飽和していることを証明する際、
$\mathcal{A}$ を、$\mathcal{A}$ と同じ同型類をもつ厳密充満部分圏で置き換えてよい。
詳細は省略する。$\mathcal{A}$ が厳密充満であると仮定する。
$\mathcal{A} = {}^\perp(\mathcal{A}^\perp)$ を示せば、補題
\ref{derived-lemma-orthogonal-triangulated} により $\mathcal{A}$ は飽和する。
包含 $\mathcal{A} \subset {}^\perp(\mathcal{A}^\perp)$ は明らかなので、
% P02-E0096: frozen source misspells “inclusion” as “incusion”.
逆の包含を証明すれば十分である。$X$ を
${}^\perp(\mathcal{A}^\perp)$ の対象とする。(2) のような完全三角
$A \to X \to B \to A[1]$ を選ぶ。仮定により $\Hom(X, B) = 0$ なので、
補題 \ref{derived-lemma-split} により $A \cong X \oplus B[-1]$ である。
$B \in \mathcal{A}^\perp$ なので $\Hom(A, B[-1]) = 0$ であり、
これは $B[-1] = 0$ および $A \cong X$ を含意する。これが望む結論である。
\end{proof}

\begin{lemma}
\label{derived-lemma-left-adjoint}
$\mathcal{D}$ を三角圏とし、$\mathcal{B} \subset \mathcal{D}$ を
充満三角部分圏とする。次は同値である。
\begin{enumerate}
\item 包含関手 $\mathcal{B} \to \mathcal{D}$ は左随伴をもつ。
\item $\mathcal{D}$ のすべての $X$ に対して、$\mathcal{D}$ における完全三角
$$
A \to X \to B \to A[1]
$$
であって、$B \in \Ob(\mathcal{B})$ および
$A \in \Ob({}^\perp\mathcal{B})$ となるものが存在する。
\end{enumerate}
これらが成り立つならば、$\mathcal{B}$ は飽和している
（定義 \ref{derived-definition-saturated}）。さらに $\mathcal{B}$ が
$\mathcal{D}$ において厳密充満ならば、
$\mathcal{B} = ({}^\perp\mathcal{B})^\perp$ である。
\end{lemma}

\begin{proof}
補題 \ref{derived-lemma-right-adjoint} の双対である。
\end{proof}

\begin{definition}
\label{derived-definition-admissible}
$\mathcal{D}$ を三角圏とする。$\mathcal{D}$ の{it 右許容}部分圏とは、
補題 \ref{derived-lemma-right-adjoint} の同値な条件を満たす厳密充満三角部分圏である。
$\mathcal{D}$ の{it 左許容}部分圏とは、補題 \ref{derived-lemma-left-adjoint} の
同値な条件を満たす厳密充満三角部分圏である。右許容かつ左許容である部分圏を
{it 両側許容}部分圏という。
\end{definition}

\noindent
$\mathcal{A}$ を三角圏 $\mathcal{D}$ の右許容部分圏とする。このとき、
$X \in \mathcal{D}$ に対し、$A \in \mathcal{A}$ および
$B \in \mathcal{A}^\perp$ となる完全三角
$$
A \to X \to B \to A[1]
$$
は次の意味で標準的であることに注意する。$A' \in \mathcal{A}$ および
$B' \in \mathcal{A}^\perp$ となる他の任意の完全三角
$A' \to X \to B' \to A'[1]$ に対して、三角の同型
$(\alpha, \text{id}_X, \beta) : (A, X, B) \to (A', X, B')$ が存在する。
次の命題は以上をまとめたものである。

\begin{proposition}
\label{derived-proposition-summarize-admissible}
$\mathcal{D}$ を三角圏とし、$\mathcal{A} \subset \mathcal{D}$ および
$\mathcal{B} \subset \mathcal{D}$ を部分圏とする。次は同値である。
\begin{enumerate}
\item $\mathcal{A}$ は右許容で、$\mathcal{B} = \mathcal{A}^\perp$ である。
\item $\mathcal{B}$ は左許容で、$\mathcal{A} = {}^\perp\mathcal{B}$ である。
\item すべての $A \in \mathcal{A}$ と $B \in \mathcal{B}$ に対して
$\Hom(A, B) = 0$ であり、さらに $\mathcal{D}$ のすべての $X$ に対して、
$A \in \mathcal{A}$ および $B \in \mathcal{B}$ となる $\mathcal{D}$ の完全三角
$A \to X \to B \to A[1]$ が存在する。
\end{enumerate}
これらが成り立つならば、$\mathcal{A} \to \mathcal{D}/\mathcal{B}$ および
$\mathcal{B} \to \mathcal{D}/\mathcal{A}$ は三角圏の同値であり、
包含関手 $\mathcal{A} \to \mathcal{D}$ の右随伴は
$\mathcal{D} \to \mathcal{D}/\mathcal{B} \to \mathcal{A}$、
包含関手 $\mathcal{B} \to \mathcal{D}$ の左随伴は
$\mathcal{D} \to \mathcal{D}/\mathcal{A} \to \mathcal{B}$ である。
\end{proposition}

\begin{proof}
(1)、(2)、(3) の同値性は補題 \ref{derived-lemma-right-adjoint} および
\ref{derived-lemma-left-adjoint} から直ちに従う（細部を一つ省略する）。
$v : \mathcal{D} \to \mathcal{A}$ を包含関手
$i : \mathcal{A} \to \mathcal{D}$ の右随伴と書く。
$\Ker(v) = \mathcal{A}^\perp = \mathcal{B}$ は直ちに分かる。
従って商の普遍性により、$v$ は関手
$\overline{v} : \mathcal{D}/\mathcal{B} \to \mathcal{A}$ を経由して分解する。
Categories, Lemma \ref{categories-lemma-adjoint-fully-faithful} により
$v \circ i = \text{id}_\mathcal{A}$ なので、$\overline{v}$ は
$\overline{i} : \mathcal{A} \to \mathcal{D}/\mathcal{B}$ の左擬逆である。
さらに、合成 $\overline{i} \circ \overline{v}$ は
$\text{id}_{\mathcal{D}/\mathcal{B}}$ と同型であると主張する。
実際、$X$ が (3) のような完全三角 $A \to X \to B \to A[1]$ に
組み込まれていると仮定する。補題 \ref{derived-lemma-right-adjoint} の証明で見たように
$v(X) = A$ である。$X$ を $\mathcal{D}/\mathcal{B}$ の対象とみなすと、
$\overline{i}(\overline{v}(X)) = A$ であり、$\mathcal{D}/\mathcal{B}$ において
関手的同型 $\overline{i}(\overline{v}(X)) = A \to X$ が存在する。
従って $\overline{v} : \mathcal{D}/\mathcal{B} \to \mathcal{A}$ は確かに
同値である。$\mathcal{B} \to \mathcal{D}/\mathcal{A}$ が同値であり、
包含関手 $\mathcal{B} \to \mathcal{D}$ の左随伴が
$\mathcal{D} \to \mathcal{D}/\mathcal{A} \to \mathcal{B}$ であることの証明は、
今述べたことの双対である。
\end{proof}





\section{Postnikov 系}
\label{derived-section-postnikov}

\noindent
本節については \cite{Orlov-K3} を参照されたい。$\mathcal{D}$ を三角圏とし、
$$
X_n \to X_{n - 1} \to \ldots \to X_0
$$
を $\mathcal{D}$ の複体とする。本節では、この複体の「全複体化」を構成する
問題を考える。

\begin{definition}
\label{derived-definition-postnikov-system}
$\mathcal{D}$ を三角圏とし、
$$
X_n \to X_{n - 1} \to \ldots \to X_0
$$
を $\mathcal{D}$ の複体とする。{it Postnikov 系}を次のように帰納的に定義する。
\begin{enumerate}
\item $n = 0$ のとき、これは同型 $Y_0 \to X_0$ である。
\item $n = 1$ のとき、これは同型 $Y_0 \to X_0$ の選択と、完全三角
$$
Y_1 \to X_1 \to Y_0 \to Y_1[1]
$$
の選択からなり、ここで $X_1 \to Y_0$ と $Y_0 \to X_0$ の合成は
与えられた射 $X_1 \to X_0$ である。
\item $n > 1$ のとき、これは $X_{n - 1} \to \ldots \to X_0$ に対する
Postnikov 系の選択と、完全三角
$$
Y_n \to X_n \to Y_{n - 1} \to Y_n[1]
$$
の選択からなり、ここで射 $X_n \to Y_{n - 1}$ と
$Y_{n - 1} \to X_{n - 1}$ の合成は、与えられた射
$X_n \to X_{n - 1}$ である。
\end{enumerate}
$\mathcal{D}$ における同じ長さの複体の間の射
\begin{equation}
\label{derived-equation-map-complexes}
\vcenter{
\xymatrix{
X_n \ar[r] \ar[d] &
X_{n - 1} \ar[r] \ar[d] &
\ldots \ar[r] &
X_0 \ar[d] \\
X'_n \ar[r] &
X'_{n - 1} \ar[r] &
\ldots \ar[r] &
X'_0
}
}
\end{equation}
が与えられたとき、{it Postnikov 系の射}という明らかな概念がある。
\end{definition}

\noindent
重要な例を挙げる。

\begin{example}
\label{derived-example-key-postnikov}
$\mathcal{A}$ をアーベル圏とし、$\ldots \to A_2 \to A_1 \to A_0$ を
$\mathcal{A}$ の鎖複体とする。このとき $D(\mathcal{A})$ の対象
$$
X_n = A_n
\quad\text{and}\quad
Y_n = (A_n \to A_{n - 1} \to \ldots \to A_0)[-n]
$$
を考えることができる。明らかな標準写像 $Y_n \to X_n$ および
$Y_0 \to Y_1[1] \to Y_2[2] \to \ldots$ とともに、完全三角
$Y_n \to X_n \to Y_{n - 1} \to Y_n[1]$ は、定義
\ref{derived-definition-postnikov-system} の意味で
$\ldots \to X_2 \to X_1 \to X_0$ の Postnikov 系を定める。
ここでは $D(\mathcal{A})$ の無限複体に対する Postnikov 系の明らかな拡張を
用いている。最後に、$\mathbf{N}$ 上の余極限が $\mathcal{A}$ に存在して完全ならば、
$$
\text{hocolim} Y_n[n] = (\ldots \to A_2 \to A_1 \to A_0 \to 0 \to \ldots)
$$
が $D(\mathcal{A})$ において成り立つ。これは補題
\ref{derived-lemma-colim-hocolim} から直ちに従う。
\end{example}

\noindent
複体 $X_n \to X_{n - 1} \to \ldots \to X_0$ と定義
\ref{derived-definition-postnikov-system} のような Postnikov 系が与えられたとき、写像
$$
Y_0 \to Y_1[1] \to \ldots \to Y_n[n]
$$
を考えることができる。これらの写像はいくつかの完全三角に組み込まれ、
$X_i$ の間の与えられた写像とも両立する。$n = 3$ の場合の図は次の通りである。
$$
\xymatrix{
Y_0 \ar[rr] & &
Y_1[1] \ar[dl] \ar[rr] & &
Y_2[2] \ar[dl] \ar[rr] & &
Y_3[3] \ar[dl] \\
& X_1[1] \ar[lu]_{+1} & &
X_2[2] \ar[ll]_{+1} \ar[lu]_{+1} & &
X_3[3] \ar[ll]_{+1} \ar[lu]_{+1}
}
$$
$Y_n[n]$ は $X_0, X_1[1], \ldots, X_n[n]$ から得られるものと考えるとよい。
例えば、写像 $X_i \to X_{i - 1}$ が零ならば、
$Y_n[n] = \bigoplus_{i = 0, \ldots, n} X_i[i]$ と取ることができる。
Postnikov 系は常に存在するとは限らない。小さい $n$ に対する簡単な補題を示す。

\begin{lemma}
\label{derived-lemma-postnikov-system-small-cases}
$\mathcal{D}$ を三角圏とする。長さ $n$ の複体に対する Postnikov 系を考える。
\begin{enumerate}
\item $n = 0$ のとき Postnikov 系は常に存在し、複体の任意の射
(\ref{derived-equation-map-complexes}) は Postnikov 系の射に一意に拡張される。
\item $n = 1$ のとき Postnikov 系は常に存在し、複体の任意の射
(\ref{derived-equation-map-complexes}) は Postnikov 系の射に（非一意に）拡張される。
\item $n = 2$ のとき Postnikov 系は常に存在するが、一般には複体の射
(\ref{derived-equation-map-complexes}) は Postnikov 系の射に拡張されない。
\item $n > 2$ のとき Postnikov 系は常に存在するとは限らない。
\end{enumerate}
\end{lemma}

\begin{proof}
$n = 0$ の場合は、同型が可逆なので直ちに従う。$n = 1$ の場合は、
TR1（三角の存在）および TR3（射を三角の射に拡張すること）から直ちに従う。
$n = 2$ の場合は次のように議論する。$Y_0 = X_0$ と置く。
$n = 1$ の場合により、Postnikov 系
$$
Y_1 \to X_1 \to Y_0 \to Y_1[1]
$$
を選べる。合成 $X_2 \to X_1 \to X_0$ は零なので、補題
\ref{derived-lemma-representable-homological} により、$X_2 \to X_1$ を（非一意に）
$X_2 \to Y_1 \to X_1$ と分解できる。そこで、射 $X_2 \to Y_1$ を完全三角
$$
Y_2 \to X_2 \to Y_1 \to Y_2[1]
$$
に組み込めば、$n = 2$ に対する Postnikov 系を得る。
$n > 2$ の場合は同様に議論できない。なぜなら、合成
$X_n \to X_{n - 1} \to Y_{n - 1}$ が $\mathcal{D}$ において零かどうかが
分からないからである。
% P02-E0097: frozen source's last composition is ill-typed for the preceding Postnikov triangle.
\end{proof}

\begin{lemma}
\label{derived-lemma-maps-postnikov-systems-vanishing}
$\mathcal{D}$ を三角圏とする。射 (\ref{derived-equation-map-complexes}) が与えられたとき、
条件
\begin{equation}
\label{derived-equation-P}
\Hom(X_i[i - j - 1], X'_j) = 0 \text{ for }i > j + 1
\end{equation}
を考える。このとき次が成り立つ。
\begin{enumerate}
\item $X'_n \to X'_{n - 1} \to \ldots \to X'_0$ に対する
Postnikov 系があれば、性質 (\ref{derived-equation-P}) は
$$
\Hom(X_i[i - j - 1], Y'_j) = 0 \text{ for }i > j + 1
$$
を含意する。
\item 両方の複体に Postnikov 系が与えられ、(\ref{derived-equation-P}) が成り立つならば、
その射は Postnikov 系の射に（非一意に）拡張される。
\end{enumerate}
\end{lemma}

\begin{proof}
まず $j$ に関する帰納法で (1) を証明する。初期の場合 $j = 0$ には、
$Y'_0 \to X'_0$ が同型なので証明すべきことはない。結論が $j - 1$ に対して
成り立つとする。完全三角
$$
Y'_j \to X'_j \to Y'_{j - 1} \to Y'_j[1]
$$
を考える。補題 \ref{derived-lemma-representable-homological} の長完全列は完全列
$$
\Hom(X_i[i - j - 1], Y'_{j - 1}[-1]) \to
\Hom(X_i[i - j - 1], Y'_j) \to
\Hom(X_i[i - j - 1], X'_j)
$$
を与える。帰納法の仮定と (\ref{derived-equation-P}) により、両端の群は零であり、
結論を得る。

\medskip\noindent
(2) を証明する。$n = 1$ のとき、射の存在は補題
\ref{derived-lemma-postnikov-system-small-cases} で証明済みである。
$n > 1$ のとき、帰納法により長さ $n - 1$ の Postnikov 系の射が
与えられていると仮定してよい。問題は図式
$$
\xymatrix{
X_n \ar[r] \ar[d] & Y_{n - 1} \ar[d] \\
X'_n \ar[r] & Y'_{n - 1}
}
$$
が可換かどうか分からないことである。その差を
$\alpha : X_n \to Y'_{n - 1}$ と書く。しかし、$\alpha$ と
$Y'_{n - 1} \to X'_{n - 1}$ の合成が零であることは分かっている
（長さ $n - 1$ の Postnikov 系の射であることの意味から従う）。完全三角
$Y'_{n - 1} \to X'_{n - 1} \to Y'_{n - 2} \to Y'_{n - 1}[1]$ により、
これは、$\alpha$ がある写像 $\alpha' : X_n \to Y'_{n - 2}[-1]$ と
$Y'_{n - 2}[-1] \to Y'_{n - 1}$ の合成であることを意味する。
すると (\ref{derived-equation-P}) と (1) により $\alpha'$ は零である。
従って $\alpha$ は零である。存在性の証明を完了するには、この可換性と TR3
により、図式に組み込まれる点線の矢印を選べばよい。
$$
\xymatrix{
Y_{n - 1}[-1] \ar[d] \ar[r] &
Y_n \ar[r] \ar@{..>}[d] &
X_n \ar[r] \ar[d] &
Y_{n - 1} \ar[d] \\
Y'_{n - 1}[-1] \ar[r] &
Y'_n \ar[r] &
X'_n \ar[r] &
Y'_{n - 1}
}
$$
\end{proof}

\begin{lemma}
\label{derived-lemma-uniqueness-maps-postnikov-systems}
$\mathcal{D}$ を三角圏とする。射 (\ref{derived-equation-map-complexes}) が与えられ、
両方の複体に Postnikov 系が与えられていると仮定する。次のいずれかが成り立つとする。
\begin{enumerate}
\item $i = 1, \ldots, n$ に対して $\Hom(X_i[i], Y'_n[n]) = 0$ である。
\item $i = 1, \ldots, n$ に対して
$\Hom(Y_n[n], X'_{n - i}[n - i]) = 0$ である。
\item $j \geq i > 0$ に対して $\Hom(X_{j - i}[-i + 1], X'_j) = 0$ および
$\Hom(X_j, X'_{j - i}[-i]) = 0$ である。
\end{enumerate}
このとき、これらの Postnikov 系の間には高々一つの射しか存在しない。
\end{lemma}

\begin{proof}
(1) を証明する。次の図式を考える。
$$
\xymatrix{
Y_0 \ar[r] \ar[d] &
Y_1[1] \ar[r] \ar[ld] &
Y_2[2] \ar[r] \ar[lld] &
\ldots \ar[r] &
Y_n[n] \ar[lllld] \\
Y'_n[n]
}
$$
矢印は射 $Y_n[n] \to Y'_n[n]$ と射 $Y_i[i] \to Y_n[n]$ の合成である。
矢印 $Y_0 \to Y'_n[n]$ は合成
$Y_0 = X_0 \to X'_0 = Y'_0 \to Y'_n[n]$ なので定まっている。
完全三角 $Y_0 \to Y_1[1] \to X_1[1]$ があるので、
$\Hom(X_1[1], Y'_n[n]) = 0$ は二番目の縦の矢印の一意性を保証する。
完全三角 $Y_1[1] \to Y_2[2] \to X_2[2]$ があるので、
$\Hom(X_2[2], Y'_n[n]) = 0$ は三番目の縦の矢印の一意性を保証する。
以下同様である。

\medskip\noindent
(2) を証明する。合成 $Y_n[n] \to Y'_n[n] \to X_n[n]$ は
% P02-E0098: frozen source ends this composition at X_n[n], although the primed system requires X'_n[n].
合成 $Y_n[n] \to X_n[n] \to X'_n[n]$ と同じであり、従って一意である。
完全三角 $Y'_{n - 1}[n - 1] \to Y'_n[n] \to X'_n[n]$ を用いると、
$\Hom(Y_n[n], Y'_{n - 1}[n - 1]) = 0$ を示せば十分であることが分かる。
完全三角
$$
Y'_{n - i - 1}[n - i - 1] \to Y'_{n - i}[n - i] \to X'_{n - i}[n - i]
$$
を用いれば、仮定からこの消滅を得る。細部は省略する。

\medskip\noindent
(3) を証明する。補題 \ref{derived-lemma-maps-postnikov-systems-vanishing} の証明を調べ、
$n$ に関する帰納法で議論すれば、三角の射
$$
\xymatrix{
Y_{n - 1}[-1] \ar[d] \ar[r] &
Y_n \ar[r] \ar@{..>}[d] &
X_n \ar[r] \ar[d] &
Y_{n - 1} \ar[d] \\
Y'_{n - 1}[-1] \ar[r] &
Y'_n \ar[r] &
X'_n \ar[r] &
Y'_{n - 1}
}
$$
における点線の矢印が一意であることを示せば十分である。補題
\ref{derived-lemma-uniqueness-third-arrow} の (5) により、
$\Hom(Y_{n - 1}, X'_n) = 0$ および
$\Hom(X_n, Y'_{n - 1}[-1]) = 0$ を示せば十分である。
最初の消滅を示すには、$i > 0$ に対して完全三角
$Y_{n - i - 1}[-i] \to Y_{n - i}[-(i - 1)] \to X_{n - i}[-(i - 1)]$ を用い、
$i$ に関する帰納法を行えば、仮定した
$\Hom(X_{n - i}[-i + 1], X'_n)$ の消滅で十分だと分かる。
二番目についても同様に、完全三角
$Y'_{n - i - 1}[-i - 1] \to Y'_{n - i}[-i] \to X'_{n - i}[-i]$ を用いれば、
仮定した $\Hom(X_n, X'_{n - i}[-i])$ の消滅で十分だと分かる。
\end{proof}

\begin{lemma}
\label{derived-lemma-existence-postnikov-system}
$\mathcal{D}$ を三角圏とする。
$X_n \to X_{n - 1} \to \ldots \to X_0$ を $\mathcal{D}$ の複体とする。
$$
\Hom(X_i[i - j - 2], X_j) = 0 \text{ for }i > j + 2
$$
ならば、Postnikov 系が存在する。また、
$$
\Hom(X_i[i - j - 1], X_j) = 0 \text{ for }i > j + 1
$$
ならば、任意の二つの Postnikov 系は同型である。
\end{lemma}

\begin{proof}
$n$ に関する帰納法で議論する。$n = 0, 1, 2$ の場合は補題
\ref{derived-lemma-postnikov-system-small-cases} から従う。$n > 2$ と仮定する。
複体 $X_{n - 1} \to X_{n - 2} \to \ldots \to X_0$ に対する
Postnikov 系が与えられているとする。これを長さ $n$ の Postnikov 系に
拡張する際の唯一の障害は、合成 $X_n \to Y_{n - 1} \to X_{n - 1}$ が
与えられた写像 $X_n \to X_{n - 1}$ と等しくなるような射
$X_n \to Y_{n - 1}$ を見つけなければならないことである。完全三角
$$
Y_{n - 1} \to X_{n - 1} \to Y_{n - 2} \to Y_{n - 1}[1]
$$
と、これおよび関手 $\Hom(X_n, -)$ から得られる長完全列
（補題 \ref{derived-lemma-representable-homological} を参照）を考えると、合成
$X_n \to X_{n - 1} \to Y_{n - 2}$ が零であることを示せば十分だと分かる。
$X_n \to X_{n - 1} \to X_{n - 2}$ が零であることは分かっているので、
完全三角
$$
Y_{n - 2} \to X_{n - 2} \to Y_{n - 3} \to Y_{n - 2}[1]
$$
を適用すれば、$\Hom(X_n, Y_{n - 3}[-1]) = 0$ ならば十分だと分かる。
補題 \ref{derived-lemma-maps-postnikov-systems-vanishing} の (1) の証明とまったく同様に
議論すれば、これが補題で述べた条件から従うことは容易に分かる。

\medskip\noindent
同型に関する主張は、補題 \ref{derived-lemma-maps-postnikov-systems-vanishing} の (2) で
証明された、複体上の恒等射を拡張する Postnikov 系の間の射の存在と、
すべての写像が同型であることを示す補題
\ref{derived-lemma-third-isomorphism-triangle} から従う。
\end{proof}









\section{本質的に定値な系}
\label{derived-section-essentially-constant}

\noindent
三角圏における本質的に定値な系について、いくつかの準備的な補題を示す。

\begin{lemma}
\label{derived-lemma-essentially-constant}
$\mathcal{D}$ を三角圏とし、$(A_i)$ を $\mathcal{D}$ の逆系とする。
このとき、$(A_i)$ が本質的に定値であること（Categories, Definition
\ref{categories-definition-essentially-constant-diagram} を参照）と、ある $i$ が
存在し、すべての $j \geq i$ に対して直和分解 $A_j = A \oplus Z_j$ であって、
(a) 写像 $A_{j'} \to A_j$ が直和分解と両立し $A$ 上で恒等であり、
(b) すべての $j \geq i$ に対して、$Z_{j'} \to Z_j$ が零となるような
ある $j' \geq j$ が存在するものがあることとは同値である。
\end{lemma}

\begin{proof}
$(A_i)$ が値 $A$ をもつ本質的に定値な系であると仮定する。このとき
$A = \lim A_i$ であり、ある $i$ と射 $A_i \to A$ が存在して、
(1) 合成 $A \to A_i \to A$ は $A$ 上の恒等写像であり、
(2) すべての $j \geq i$ に対して、ある $j' \geq j$ が存在し、
$A_{j'} \to A_j$ は $A_{j'} \to A_i \to A \to A_j$ と分解する。
(1) により、$j \geq i$ に対して写像 $A \to A_j$ と
$A_j \to A_i \to A$ の合成は $A$ 上の恒等写像である。従って
$A_j \to A$ は核 $Z_j$ をもち、写像 $A \oplus Z_j \to A_j$ は同型である。
補題 \ref{derived-lemma-when-split} および \ref{derived-lemma-split} を参照されたい。
これらの直和分解は明らかに (a) を満たす。(2) により、すべての $j$ に対して、
$Z_{j'} \to Z_j$ が零となるような $j' \geq j$ が存在し、従って (b) が成り立つ。
逆向きの証明は省略する。
\end{proof}

\begin{lemma}
\label{derived-lemma-essentially-constant-2-out-of-3}
$\mathcal{D}$ を三角圏とし、
$$
A_n \to B_n \to C_n \to A_n[1]
$$
を $\mathcal{D}$ における完全三角の逆系とする。$(A_n)$ と $(C_n)$ が
本質的に定値ならば、$(B_n)$ も本質的に定値であり、それらの値は完全三角
$A \to B \to C \to A[1]$ に組み込まれる。さらに、ある $n \geq 1$ に対して、
完全三角の射
$$
\xymatrix{
A_n \ar[d] \ar[r] &
B_n \ar[d] \ar[r] &
C_n \ar[d] \ar[r] &
A_n[1] \ar[d] \\
A \ar[r] &
B \ar[r] &
C \ar[r] &
A[1]
}
$$
であって、同型 $\lim_{n' \geq n} A_{n'} \to A$ を誘導し、$B$ および $C$
についても同様となるものが存在する。
\end{lemma}

\begin{proof}
番号を付け直すことにより、$A_n = A \oplus A'_n$ および
$C_n = C \oplus C'_n$ と仮定してよい。ここで逆系 $(A'_n)$ と $(C'_n)$ は
本質的に零である。補題 \ref{derived-lemma-essentially-constant} を参照されたい。
特に、すべての $n$ に対して射
$$
C \oplus C'_n \to (A \oplus A'_n)[1]
$$
は直和因子 $C$ を直和因子 $A[1]$ に、$n$ に依存しない写像
$\delta : C \to A[1]$ によって写す。完全三角
$$
A \to B \to C \xrightarrow{\delta} A[1]
$$
を選ぶ。次に、完全三角の射
$$
(A_1 \to B_1 \to C_1 \to A_1[1]) \to
(A \to B \to C \to A[1])
$$
を選ぶ。これは TR3 により可能である。$\mathcal{D}$ の任意の対象 $D$ に対して、
これは可換図式
$$
\xymatrix{
\ldots \ar[r] &
\Hom_\mathcal{D}(C, D) \ar[r] \ar[d] &
\Hom_\mathcal{D}(B, D) \ar[r] \ar[d] &
\Hom_\mathcal{D}(A, D) \ar[r] \ar[d] &
\ldots \\
\ldots \ar[r] &
\colim \Hom_\mathcal{D}(C_n, D) \ar[r] &
\colim \Hom_\mathcal{D}(B_n, D) \ar[r] &
\colim \Hom_\mathcal{D}(A_n, D) \ar[r] &
\ldots
}
$$
を誘導する。左と右の縦の矢印は同型であり、それらのさらに左と右にある矢印も
同型である。従って五項補題により、中間の矢印は同型である。
Categories, Remark \ref{categories-remark-pro-category-copresheaves} の議論により、
$(B_n)$ は値 $B$ の定値逆系と同型である。Categories, Remark
\ref{categories-remark-pro-category} により、これは $(B_n)$ が値 $B$ の
本質的に定値な系であることと同値なので、証明が完了する。
\end{proof}

\begin{lemma}
\label{derived-lemma-essentially-constant-cohomology}
$\mathcal{A}$ をアーベル圏とし、$A_n$ を $D(\mathcal{A})$ の対象の逆系とする。
次を仮定する。
\begin{enumerate}
\item 整数 $a \leq b$ が存在し、$i \not \in [a, b]$ に対して
$H^i(A_n) = 0$ である。
\item すべての $i \in \mathbf{Z}$ に対して、$\mathcal{A}$ における逆系
$H^i(A_n)$ は本質的に定値である。
\end{enumerate}
このとき $A_n$ は $D(\mathcal{A})$ の本質的に定値な系であり、その値 $A$ は、
各 $i \in \mathbf{Z}$ に対して $H^i(A)$ が定値系 $H^i(A_n)$ の値となる性質を
満たす。
\end{lemma}

\begin{proof}
Remark \ref{derived-remark-truncation-distinguished-triangle} により、完全三角の逆系
$$
\tau_{\leq a}A_n \to A_n \to \tau_{\geq a + 1}A_n \to (\tau_{\leq a}A_n)[1]
$$
を得る。もちろん $D(\mathcal{A})$ において
$\tau_{\leq a}A_n = H^a(A_n)[-a]$ である。従って仮定により、これらは
本質的に定値な系をなす。$b - a$ に関する帰納法により、逆系
$\tau_{\geq a + 1}A_n$ は本質的に定値であり、その値を $A'$ とする。
補題 \ref{derived-lemma-essentially-constant-2-out-of-3} により、$A_n$ は本質的に
定値な系である。コホモロジーに関する主張の証明は省略する
（ヒント：補題 \ref{derived-lemma-essentially-constant-2-out-of-3} の最後の部分を用いる）。
\end{proof}

\begin{lemma}
\label{derived-lemma-pro-isomorphism}
$\mathcal{D}$ を三角圏とし、
$$
A_n \to B_n \to C_n \to A_n[1]
$$
を完全三角の逆系とする。系 $C_n$ が pro-零（値 $0$ の本質的に定値な系）ならば、
写像 $A_n \to B_n$ は pro-対象 $(A_n)$ と pro-対象 $(B_n)$ の間の
pro-同型を定める。
\end{lemma}

\begin{proof}
$\mathcal{D}$ の任意の対象 $X$ に対して、完全列
$$
\colim \Hom(C_n, X) \to
\colim \Hom(B_n, X) \to
\colim \Hom(A_n, X) \to
\colim \Hom(C_n[-1], X) \to
$$
を考える。完全性は補題 \ref{derived-lemma-representable-homological} と Algebra, Lemma
\ref{algebra-lemma-directed-colimit-exact} から従う。仮定により最初と最後の項は
零である。従って、すべての $X$ に対して写像
$\colim \Hom(B_n, X) \to \colim \Hom(A_n, X)$ は同型である。
本補題はこれと Categories, Remark
\ref{categories-remark-pro-category-copresheaves} から従う。
\end{proof}

\begin{lemma}
\label{derived-lemma-pro-isomorphism-bis}
$\mathcal{A}$ をアーベル圏とする。
$$
A_n \to B_n
$$
% P02-E0100: frozen source omits “Let” before this displayed inverse system.
を $D(\mathcal{A})$ における写像の逆系とする。次を仮定する。
\begin{enumerate}
\item 整数 $a \leq b$ が存在し、$i \not \in [a, b]$ に対して
$H^i(A_n) = 0$ および $H^i(B_n) = 0$ である。
\item すべての $i \in \mathbf{Z}$ に対して、$\mathcal{A}$ における写像の逆系
$H^i(A_n) \to H^i(B_n)$ は $\mathcal{A}$ の pro-対象の同型を定める。
% P02-E0101: frozen source has subject--verb disagreement in “the inverse system ... define”.
\end{enumerate}
このとき写像 $A_n \to B_n$ は pro-対象 $(A_n)$ と pro-対象 $(B_n)$ の間の
pro-同型を定める。
\end{lemma}

\begin{proof}
公理 TR3 により、写像 $A_n \to B_n$ を完全三角の逆系
$A_n \to B_n \to C_n \to A_n[1]$ に帰納的に拡張できる。補題
\ref{derived-lemma-pro-isomorphism} により、$C_n$ が pro-零であることを証明すれば十分である。
補題 \ref{derived-lemma-essentially-constant-cohomology} により、各 $p$ に対して
$H^p(C_n)$ が pro-零であることを証明すれば十分である。これは仮定 (2) と長完全列
$$
H^p(A_n) \xrightarrow{\alpha_n} H^p(B_n)
\xrightarrow{\beta_n}
H^p(C_n) \xrightarrow{\delta_n} H^{p + 1}(A_n)
\xrightarrow{\epsilon_n}
H^{p + 1}(B_n)
$$
から従う。実際、すべての $n$ に対して、$\Im(\beta_m)$ の $H^p(C_n)$ に
おける像が零となるような $m > n$ を見つけることができる。これは
$H^p(B_m) \to H^p(B_n)$ が $\alpha_n : H^p(A_n) \to H^p(B_n)$ を
経由して分解するように $m$ を選べるからである。同様の理由により、次に
$\Im(\delta_k)$ の $H^{p + 1}(A_m)$ における像が零となるような
$k > m$ を選べる。このとき $H^p(C_k) \to H^p(C_n)$ は零である。
なぜなら、$H^p(C_k) \to H^p(C_m)$ は $\Ker(\delta_m)$ に値をもち、
$H^p(C_m) \to H^p(C_n)$ は $\Ker(\delta_m) = \Im(\beta_m)$ を消すからである。
\end{proof}
